\documentclass[11pt]{article}
\usepackage[T1]{fontenc}
\usepackage[utf8]{inputenc}
\usepackage{amsmath,amssymb,amsthm,mathtools}
\usepackage{geometry,booktabs,array,longtable}
\usepackage[hidelinks]{hyperref}
\usepackage{url}
\geometry{margin=28mm}
\setlength{\emergencystretch}{2em}
\newtheorem{theorem}{Theorem}[section]
\newtheorem{proposition}[theorem]{Proposition}
\newtheorem{lemma}[theorem]{Lemma}
\newtheorem{corollary}[theorem]{Corollary}
\newtheorem{assumption}[theorem]{Assumption}
\theoremstyle{definition}
\newtheorem{definition}[theorem]{Definition}
\newtheorem{remark}[theorem]{Remark}
\newcommand{\Q}{\mathbb Q}
\newcommand{\Z}{\mathbb Z}
\newcommand{\A}{\mathbb A}
\newcommand{\OO}{\mathcal O}
\newcommand{\Src}{\mathcal S}
\DeclareMathOperator{\Spec}{Spec}
\DeclareMathOperator{\Frac}{Frac}
\DeclareMathOperator{\Sha}{Sha}
\title{Conditional polynomial constructions, lower bounds, and\newline
number-field obstructions for rational Diophantine integrality}
\author{Research report prepared for JUVa\thanks{Prepared with ChatGPT and
parallel mathematical and computational reviews. No claim of independent
human review or priority is made. The upstream manuscript and its arithmetic
inputs are attributed separately.}}
\date{6 November 2026\quad---\quad revised working version 3}
\begin{document}
\maketitle
\begin{abstract}
We study explicit polynomial presentations of the ordinary integers in
the rationals, starting from a pinned revision of the repository
\emph{diophantine-z-q}. Under its named arithmetic inputs and Sun's
simultaneous prime-pattern theorem, we construct an indexed polynomial
in 41 total variables, including the input, with degree at most
$2^{9566673}$. A 43-variable alternative has degree at most $2^{236}$.
The direct substitution of the Lucas polynomial into approximate
logarithmic ratios removes two retained coordinates and one finite
algebraic coordinate. It also gives quartics with 1315 variables.
An affine change reduces their coefficient height to $64\cdot3^{2^{28}}$;
an alternative with 1335 variables has coefficient height less than $2^{1020}$.
Their complete sparse coefficient tables are independently checked.
The previously implemented 558-variable construction and its
source-preserving alternatives remain available for comparison.

No construction in two through ten total variables is established.
Unconditional lower bounds exclude two total variables at every
degree. We prove that two quadratic equations, in any number of
coordinates, cannot have an infinite input image contained in a
number field's ring of integers. A quadratic system defining that
ring consequently needs at least five total coordinates and at
least three independent quadratic equations. Further results explain
precise limits of finite algebraic encoding and fixed denominator
tests without mistaking them for universal variable lower bounds.

Over every specified number field, a finite system can be combined
into one equation without adding variables. An explicit division-norm
construction gives degree at most $D\lceil\sqrt{s}\rceil$ for $s$
equations of degree at most $D$; the preferred quadratic systems yield
degree 72. These are faithful encodings, not general integrality
theorems. Exact counterexamples show that the literal newer quartics
accept all inputs over $\mathbb Q(i)$, the unshifted indexed equations
reject certain algebraic integers, and the fixed cyclic elliptic
argument fails in an explicit totally real quartic extension.
A translated square repairs the finite input contacts but not that
arithmetic obstruction. All hypotheses, proofs, computational scopes,
and unresolved dependencies are stated separately.
\end{abstract}
\tableofcontents

\section{Definitions, objectives, and the scope of the results}
\label{sec:scope}

\begin{definition}
Let $k$ be a field. A set $S\subseteq k$ is \emph{Diophantine over $k$}
if there are $r\geq0$ and $P\in k[T,Y_1,\ldots,Y_r]$ such that
\[
 t\in S\quad\Longleftrightarrow\quad
 \exists(y_1,\ldots,y_r)\in k^r:\ P(t,y_1,\ldots,y_r)=0.
\]
The variable $T$ is the input. We call $r$ the number of witnesses and
$r+1$ the total number of variables. Degrees in this report are total
degrees, including the input. For a nonzero polynomial over $\Q$, clear
numerical denominators and divide the coefficients by their greatest
common divisor. Its \emph{coefficient height} is the maximum absolute
value of the resulting primitive integer coefficients.
\end{definition}

The requested order of optimization is lexicographic: first the total
number of variables, then degree, then coefficient height. Under this
order, a smaller-variable polynomial of very high degree is preferable
to a quartic with more variables. A useful practical comparison should
nevertheless retain the alternatives with moderate coefficients and
complete coefficient tables. A circuit node used to \emph{write} a
polynomial is not a witness variable. Conversely, a newly quantified
coordinate constrained to equal a circuit value does count as a witness.

We use the upstream repository and the user's fork at revision
\begin{center}\small\ttfamily
28e87b1b62af0effd954605af7f70b2699dd5672.
\end{center}
The released package is \texttt{quartics/quartic-12504}; its manuscript
is \texttt{quartic\_12504.tex}. The repository's own documentation
distinguishes finite algebraic verification from verification of the
arithmetic theorem; its Lean material is incomplete
\cite{Repository2026}. We retain that distinction throughout.

Write $\Src$ for the rational existential system specified by the
manuscript's reciprocal comparison source, with precisions $64$ and
$16$, and its scalar order $48$. Write $S_{\Src}$ for its projected
rational solution set in the ordinary input $X$.

\begin{assumption}[The source arithmetic assertion]\label{ass:source}
The set $S_{\Src}$ is the ordinary subset $\Z\subset\Q$.
\end{assumption}

This is the assertion of the source manuscript's theorem on the
arithmetic interpretation of the revised source. It is not an additional
conjecture introduced by the optimizations. The source-preserving transformations
establish equality with $S_{\Src}$; calling their outputs definitions
of $\Z$ requires Assumption~\ref{ass:source}. Replaying a certificate
for polynomial substitution does not prove this assumption. In
particular, the present report does not claim a new independent proof
of the full arithmetic manuscript.

The new shared-point construction uses more information than the bare
equality in Assumption~\ref{ass:source}: it uses the manuscript's stated
cyclic-group, pole, formal-logarithm, height, and Lucas results, together
with Sun's simultaneous prime-pattern theorem
\cite[Theorem~4.1]{Sun2026Rational86}. Its soundness and completeness
are proved from these explicitly named inputs. This additional theorem
permits comparison points to be shared across different selected primes.
We distinguish this arithmetic modification from the unconditional
algebraic transformations of the released source. The lower-bound
theorems require neither arithmetic assumption.

The source release has 12,504 total variables, 11,939 quadratic
residuals, 135,721 monomials in their sum of squares, and maximum
coefficient
\[
 685783949153203534233601<2^{80}.
\]
These are exact algebraic statistics, independently replayed here.
The following results have different forms of explicitness.

\begin{center}
\small
\begin{tabular}{r l l p{43mm}}
\toprule
Total variables & Degree & Coefficients & Output supplied\\
\midrule
41 & $\le2^{9566673}$ & $\le2^{2^{86098382}}$ & Indexed polynomial with fixed elliptic exclusion product\\
43 & $\le2^{236}$ & $\le2^{2^{240}}$ & Indexed semilocal alternative with coefficient DAG\\
558 & $\le2^{3122}$ & $\le2^{2^{3127}}$ & Implemented source-preserving indexed polynomial\\
1315 & $4$ & $64\kappa<2^{429496736}$ & Complete table over $\mathbb Z[s]$, with $s^2=\kappa$\\
1335 & $4$ & $<2^{1020}$ & Complete integer coefficient table\\
12,504 & $4$ & $<2^{80}$ & Unchanged released package\\
\bottomrule
\end{tabular}
\end{center}

All counts include the input. The first two and the two smaller quartic
rows use the shared-point arithmetic argument and the direct Lucas
test. The 558-variable row is a source-preserving algebraic
construction. Large displayed degree and height inequalities are proved
upper bounds, not exact expansion statistics. Here
$\kappa=3^{2^{28}}$ and $s=3^{2^{27}}$ are fixed integers;
neither counts as a coordinate. The indexed polynomials specify every
matrix and finite product but do not expand their enormous tables.
Both 41- and 43-variable specifications also have separately checked
versions with $\beta=X^2+2402$, with the same displayed bounds.

The preceding 43- and 45-variable indexed formulas, the untransformed
1315-variable table, the 1338-, 1339-, 1367- and 1391-variable quartics,
source-preserving 853-variable circuit, and
12,449--12,493-variable quartics are retained as reviewable comparison
constructions. Their proofs supply several ingredients used below.
None of these counts is claimed to be globally minimal. Under the
stated arithmetic inputs, the unrestricted gap is between three and
41 total variables. For a sum of squares of quadratic residuals, the
proved gap is between five and 1315. Degree four is the least possible
degree for one defining equation if the conditional construction is
valid. A definition using at most ten total variables remains unproved.

Allowing several equations helps with degree and presentation but
does not change the minimum number of variables when degree is
unrestricted. Over $\mathbb Q$, their sum of squares combines any
finite system without a new variable. Over any specified number field,
the division-norm construction gives the same conclusion with a
controlled degree. The general and field-specific assertions are
proved separately in this report.

\section{From integer equations to rational equations}
\label{sec:reduction}

Suppose first that an integer-coefficient polynomial
$P(T,\mathbf Y)$ of degree $d$ defines $\Z$ in $\Q$, using $r$
witnesses. For $f\in\Z[X_1,\ldots,X_m]$, form
\begin{equation}\label{eq:basic-reduction}
 R_f(\mathbf X,\mathbf Y_1,\ldots,\mathbf Y_m)
 =f(\mathbf X)^2+\sum_{i=1}^m P(X_i,\mathbf Y_i)^2.
\end{equation}
\begin{proposition}\label{prop:integer-rational-reduction}
The equation $f=0$ has a solution in $\Z^m$ if and only if
$R_f=0$ has a solution over $\Q$. The construction is effective,
uses $m(r+1)$ rational variables, and has degree at most
$2\max\{\deg f,d\}$.
\end{proposition}
\begin{proof}
Given an integer zero of $f$, the defining property of $P$ supplies
a witness tuple for each coordinate, so every summand in
\eqref{eq:basic-reduction} vanishes. Conversely, a rational sum of
squares vanishes only when every summand vanishes. Thus any rational
zero of $R_f$ satisfies $f=0$ and $P(X_i,\mathbf Y_i)=0$ for every
$i$. Each $X_i$ is therefore an integer. The construction consists
only of copying, squaring, and adding finite polynomials; the bounds
follow by inspection.
\end{proof}

For example, an equation $f(a,b,c)=0$ is changed into a rational
equation with three separate integrality tests, one for each of
$a,b,c$. Its witnesses are allowed to have denominators. The tests
force the three original coordinates to be integers; they need not
force the added witnesses to be integers.

There is a better degree bound for the quartics actually supplied
here. They have the form $P=\sum_j q_j^2$, with $\deg q_j\leq2$.
In that case $P$ is nonnegative on rational inputs, so
\begin{equation}\label{eq:sos-reduction}
 f(\mathbf X)^2+\sum_{i=1}^m P(X_i,\mathbf Y_i)=0
\end{equation}
already gives the desired reduction. Its degree is at most
$\max\{2\deg f,4\}$, with the same variable count. Squaring each
copy of $P$ a second time is unnecessary.

Moreover every polynomial $f$ can be evaluated by a finite sequence
of additions and multiplications. Introduce one rational coordinate
for each intermediate value, and impose equations such as
$z-u-v=0$ or $z-uv=0$, together with the equation that the final value
is zero. These equations have degree at most two, and for fixed
$\mathbf X$ their intermediate coordinates are uniquely determined.
Combining their squares with the integrality quartics produces a
quartic rational-solvability problem for \emph{every} $f$. If $g$
intermediate coordinates are used, the variable count is
$m(r+1)+g$. These intermediate coordinates need no separate
integrality predicates: they already equal the prescribed polynomial
values. This trade preserves degree at the cost of variables, so it
is secondary under the requested priority.

\subsection{The precise contribution of MRDP}

The Matiyasevich--Robinson--Davis--Putnam theorem identifies the
computably enumerable subsets of $\Z^n$ with the integer Diophantine
sets, with effective and uniform polynomial representations
\cite{Matiyasevich1970,Jones1982}. A set is computably enumerable if
an algorithm can list its members, without necessarily deciding
nonmembership. The theorem over nonnegative integers gives the version
over all integers by the usual signed encodings and four-square
representations of nonnegative variables.

\begin{corollary}\label{cor:rational-mrdp}
If $\Z$ is Diophantine in $\Q$, then the computably enumerable
subsets of $\Q^n$ are exactly its Diophantine subsets. In particular,
rational polynomial solvability is computably enumerable complete
under many-one reductions, and hence undecidable.
\end{corollary}
\begin{proof}
Use the unique coprime numerator--positive denominator encoding of
each rational coordinate. If $S\subseteq\Q^n$ is computably
enumerable, its set of such integer encodings is computably enumerable.
MRDP supplies an integer polynomial for that set. Replace each integer
quantifier by a rational quantifier and an integrality test $P=0$;
add the equations $b_i x_i-a_i=0$. The conjunction is one polynomial
by a sum of squares, and has exactly the desired rational projection.
Conversely, enumerate all rational witness tuples and evaluate a given
polynomial exactly; this enumerates every Diophantine set.
Uniformity gives many-one hardness from any computably enumerable
set, and direct enumeration gives membership in that class.
\end{proof}

Applying this argument to a universal computably enumerable relation
also gives a universal rational Diophantine polynomial, with its index
as an input parameter. This does not make the small integer universal
polynomials into definitions over rational witnesses without modification.
Nor is an abstract Diophantine model of arithmetic automatically a
definition of the ordinary embedded subset $\Z\subset\Q$. A model,
an undecidability reduction, and a definition of that particular subset
are different assertions. All the conclusions in this section become
conditional applications to the source by
Assumption~\ref{ass:source}.

% Includable supplement. The division-algebra theorem is supplied separately
% in phase3/number_fields/division_norm_compiler_20261106_v1.tex.
\section{Several equations versus one equation}
\label{p3sys:combination}

Multiple equations can make the degree smaller and the construction
easier to inspect.  They do not lower the minimum possible number of
variables when degree is unrestricted: over the rationals, and over
every number field, a finite conjunction can be combined into one
equation without adding a variable.  The degree and field of the
combining polynomial must nevertheless be recorded.

\begin{proposition}\label{p3sys:sos-combination}
Let $K$ be a number field with a real embedding, and let
$F_1,\ldots,F_s\in K[T,Y_1,\ldots,Y_r]$ have degree at most $D$.
The system $F_1=\cdots=F_s=0$ has exactly the same $K$-solutions
as the single equation
\[
                         P=\sum_{i=1}^sF_i^2=0.
\]
No new variable is introduced, and $\deg P\le2D$.
For a single residual one may simply take $P=F_1$.
If $K=\mathbb Q$, $P\ne0$, all residuals have integer coefficients of
absolute value at most $H$, and there are $n=r+1$ total
coordinates, then
\[
               H(P)\le s\binom{n+D}{D}H^2.
\]
Dividing the coefficients by their common integer divisor can
only improve this bound.
\end{proposition}
\begin{proof}
Apply a real embedding of $K$ to a zero of the displayed sum.
Every summand is a nonnegative real number, so all of them vanish;
injectivity of the embedding gives $F_i=0$ in $K$.  The converse
and the degree bound are immediate.  There are at most
$M=\binom{n+D}{D}$ monomials in each residual.  For a fixed
monomial in $F_i^2$, choosing one exponent determines the other,
so its coefficient is a sum of at most $M$ products, each of
absolute value at most $H^2$.  Summing over $i$ proves the bound.
\end{proof}

The real-embedding condition is weaker than total reality.
It is essential to this particular combining formula: over
$\mathbb Q(i)$, the two nonzero values $1,i$ have squares summing
to zero.  A replacement over other number fields is supplied by
division algebras.

\begin{corollary}\label{p3sys:division-combination}
For any number field $K$, a system of $s\ge1$ polynomial equations
of degrees at most $D$ can be combined into one equation in the
same coordinates of degree at most
\[
                         D\lceil\sqrt{s}\rceil.
\]
For a specified $K$, a combining polynomial with rational
coefficients suffices.  This is a statement about equivalence
of solution sets over that field; it does not identify a changed
field's solution set with its ring of integers.
\end{corollary}
\begin{proof}
Put $d=\lceil\sqrt{s}\rceil$.  The division-algebra construction
of Theorem~\ref{thm:division-norm-conjunction-phase3} supplies a
homogeneous form $N$ of degree $d$ in $d^2$ variables such that
$N(\mathbf z)=0$ for $\mathbf z\in K^{d^2}$ exactly when
$\mathbf z=0$.  Substitute
\[
                         (F_1,\ldots,F_s,0,\ldots,0)
\]
into this form.  Its zero condition is precisely the original
conjunction, and its degree is at most $dD$.
\end{proof}

For example, the preferred quartic presentations in 1315 and 1335
total coordinates come from 1272 and 1292 quadratic residuals,
respectively.  In both cases $\lceil\sqrt{s}\rceil=36$.
The number-field compiler therefore combines each residual system
into an equation of degree at most 72 without changing its
coordinate count.  The supplied determinant specifications have
degree exactly 72: the input occurs only as $+X_0^2$ in the
first residual, which multiplies the identity matrix, so the
coefficient of $X_0^{72}$ is one.  Over $\mathbb Q$ the
sum-of-squares quartic is better in degree.  Over a specified
nonreal number field a field-adapted reduced-norm construction preserves the
conjunction, whereas the sum of squares need not do so.  Neither
claim says that these base-changed residuals define the new ring
of integers.

\begin{corollary}\label{p3sys:minimum-variable-equality}
For a fixed subset $S$ of a number field $K$, the least total
number of coordinates in a definition of $S$ by one existential
polynomial equation equals the least total number in a definition
by a finite system of existential polynomial equations.  The
values are understood to be infinite if there is no such
definition.
\end{corollary}
\begin{proof}
A single equation is a finite system, giving one inequality.
The preceding combination preserves every coordinate, giving
the reverse inequality.
\end{proof}

Thus a system defining $\mathbb Z$ in at most ten total variables
would immediately give a single polynomial in at most ten total
variables.  If its equations are quadratic over $\mathbb Q$,
the combined polynomial would be quartic.  The system may be
much easier to discover or display, but it does not bypass a
lower bound on the variable count.

\subsection{Finite alternatives and nonzero conditions}

Suppose $P_j(T,\mathbf Y_j)=0$ are finitely many alternative
presentations, with $r_j$ witness coordinates.  Pad the lists to
$r=\max_jr_j$ using unused coordinates, and set
\[
                          P=\prod_j P_j.
\]
Over a field its existential projection is precisely the union
of the original projections.  Its total degree is at most
$\sum_j\deg P_j$, and no selector coordinate is needed.  This
also combines a finite union of systems after combining each
conjunction.  The assertion is about finite products with fixed
index sets, not a quantified integer choice.

Nonzero conditions require separate attention.  The conjunction
$G_1\ne0,\ldots,G_t\ne0$ is equivalent over a field to
\[
                         u\prod_{j=1}^tG_j-1=0
\]
with one additional witness $u$.  The new residual has degree at
most $1+\sum_j\deg G_j$.  This generic method therefore costs
one coordinate, even though conjunction and finite disjunction
of existing equations do not.  The cost cannot simply be ignored:
the one-coordinate set $K\setminus\{0\}$ is defined by the
nonzero condition $T\ne0$, but it is not the common zero set of
any finite family of one-variable polynomials.  Such a zero set
is finite unless every polynomial is zero, in which case it is
all of $K$.  Absorbing a guard into an existing witness requires
a separate equivalence, such as the reciprocal constructions
proved below.

\subsection{The degree of a universal combining polynomial}

Here a universal combining polynomial in $s$ inputs means a
polynomial $A\in K[Z_1,\ldots,Z_s]$ with
\[
              A(\mathbf z)=0\quad\Longleftrightarrow\quad
                           \mathbf z=0\qquad(\mathbf z\in K^s).
\]
Its defining property concerns all residual-value tuples.  It
is stronger than the property needed for one particular list
of residuals, whose value tuples may satisfy additional relations.

\begin{proposition}\label{p3sys:quadratic-combiner-optimality}
For $s\ge2$, a universal combining polynomial has degree at
least two.  If its degree is at most two, it is a homogeneous
anisotropic quadratic form.  Consequently:
\begin{enumerate}
\item if $K$ has a real embedding, the minimum degree is two
for every $s\ge2$;
\item if $K$ is totally imaginary, the minimum degree is two
for $2\le s\le4$, and three for $5\le s\le9$.
\end{enumerate}
For totally imaginary $K$ and $s\ge10$, the present argument
only gives
\[
 3\le\text{minimum universal-combiner degree}
                         \le\lceil\sqrt{s}\rceil.
\]
No optimality of the last upper bound is asserted.
\end{proposition}
\begin{proof}
The constant term is zero.  A nonzero linear form in at least
two variables has a nonzero vector in its kernel, so degree one
cannot work.  Write a polynomial of degree at most two as
$A=q+\ell$, with $q$ quadratic homogeneous and $\ell$ linear.
If both parts are nonzero, choose $v\in K^s$ with
$q(v)\ell(v)\ne0$.  The field is infinite, so two proper
hypersurfaces cannot cover its affine space.  With
$t=-\ell(v)/q(v)\ne0$ one has
\[
                          A(tv)=t^2q(v)+t\ell(v)=0,
\]
contradicting universality.  If $q=0$, the linear argument
already applies.  Therefore $\ell=0$, and $q$ must be
anisotropic.

A real embedding permits the anisotropic form
$Z_1^2+\cdots+Z_s^2$ in every dimension.  For a totally
imaginary number field, the $u$-invariant is four: every
quadratic form in at least five variables is isotropic, while
a quaternion division algebra supplies an anisotropic norm
form in four variables.  The upper restriction follows from
Hasse--Minkowski and the local $u$-invariant four; quaternion
existence follows by ramifying at two finite places.  See
\cite[Theorems~14.6.1 and 14.6.9]{Voight2021} and the explicit
summary in \cite[Section~5.1]{p3sysMerkurjevParimala2025}.
Restricting an anisotropic four-variable form to a coordinate
subspace treats $s=2,3,4$.  For $5\le s\le9$, a degree-three
division norm supplies the upper bound, and the quadratic
obstruction makes it sharp.  The general upper bound is
Corollary~\ref{p3sys:division-combination}.
\end{proof}

For $s=1$ the minimum is one.  These combining-polynomial
restrictions do not impose the same degree bound on a
particular definition: dependencies among its residuals can
allow cheaper combinations.  Conversely, the general degree
lower bound for a system defining $\mathcal O_K$ is only two,
since linear systems are affine spaces and have empty, constant,
or surjective input projection.  The much stronger restrictions
for quadratic systems are proved in
Section~\ref{p3sys:lower-bounds}.


\section{The arithmetic structure retained from the source}
\label{sec:source-structure}

The polynomial is assembled from a considerably smaller list of
arithmetic requirements. It uses
\[
 K=\Q(\sqrt2),\quad d=3-\sqrt2,\quad D=d^2=11-6\sqrt2,
 \qquad E_D:y^2=x^3-Dx,
\]
and the point
\[
 P_*=\left(\frac{57}{8}-\frac72\sqrt2,
 -\frac{165}{8}+\frac{417}{32}\sqrt2\right).
\]
The source's generator assertion identifies $2E_D(K)$ with the
infinite cyclic group generated by $P_*$. Its points can therefore
encode integer indices. This assertion is a substantive arithmetic
input; the elementary chart substitutions proved below do not require
us to infer it from a computer search.

The input is converted to $\beta=X^2+2$. Four points encode indices
$n,s,t,z$, and three rational proxy coordinates satisfy a Lucas-type
equation. Formal-group comparisons at selected primes make the proxies
congruent to the indices. A common denominator and forty-one affine
contact expressions convert these congruences into height bounds.
A separate scalar batch bounds the height of $\beta$ in terms of
$|s|$. The anchor forces $|s|$ beyond the fixed cutoff. Finally Lucas
growth contradicts the height bounds unless the Lucas residual is zero
and $\beta=n\in\Z$. A rational number whose square is an integer
is an integer, so $X\in\Z$ follows.

In the other direction, the manuscript chooses sufficiently divisible
elliptic indices, uses a simultaneous ordinary-prime pattern, constructs
private comparison points, and obtains pole witnesses. The difficult
inputs include its quantitative five-contact estimate, its restricted
converse and pole-completeness argument, the prime-pattern theorem,
and the cyclic elliptic-group assertion. We preserve these inputs and
their precise witness scopes. Our changes operate on the resulting
existential presentation.

Two forms of verification should be separated. The released certificate
is a finite algebraic object: it checks lowering, elimination, invertible
row operations, variable maps, and expansion into a sum of squares.
The arithmetic manuscript must additionally prove both soundness and
completeness of its source. The former can be checked in seconds by exact
integer arithmetic. It cannot establish the latter by itself. The
proofs and evidence supplied below explicitly identify which task each
calculation performs.


\section{An economical semilocal condition}
\label{sec:lit-optimization}

Throughout this section a witness variable ranges over $\mathbb Q$ unless
another field is specified. An equivalence of polynomial presentations means
an equivalence of their projected rational solution sets in every displayed
input. Such an equivalence is independent of whether the arithmetic theorem
that identifies the original projected set with $\mathbb Z$ has been verified.
The next section matches these statements to the actual residual system.

\subsection{An explicit three-variable semilocal definition}

Write
\[
 R=\mathbb Z_{(2)}\cap\mathbb Z_{(3)}
   =\{t\in\mathbb Q:v_2(t)\ge0,\ v_3(t)\ge0\},
\]
with $v_p(0)=+\infty$. The following specialization of classical
quaternion methods is particularly economical. Daans proves three-quantifier
existential definability for every finite intersection of valuation rings
in a global field; he traces the corresponding rational cases back to
Robinson \cite[Proposition 4.2 and Remark 4.3]{Daans2024}.
Here all constants and all exceptional cases are made explicit.

\begin{proposition}\label{prop:explicit-semilocal}
For every $t\in\mathbb Q$,
\[
 t\in R\quad\Longleftrightarrow\quad
 \exists u,v,w\in\mathbb Q:\quad
 u^2+6v^2-2w^2-3t^2+2=0.
\]
\end{proposition}
\begin{proof}
Consider the quaternion algebra $B=(2,3)_{\mathbb Q}$ with generators
$i^2=2$, $j^2=3$, $ij=-ji$. Its ramification set consists exactly of
$2$ and $3$. Indeed its Hilbert symbol at $3$ equals the Legendre symbol
$(2/3)=-1$, and at $2$ the formula
$(2,u)_2=(-1)^{(u^2-1)/8}$ for odd $u$ gives $(2,3)_2=-1$.
It splits at every other finite place because both parameters are units
at those odd primes, and it splits at
infinity because both parameters are positive. In particular $B$ is a
division algebra over $\mathbb Q$.

We use the standard quadratic embedding criterion for quaternion algebras:
a quadratic field $L/\mathbb Q$ embeds in $B$ if and only if every ramified
place of $B$ remains a field in $L$. This is the quadratic case of the
Albert--Brauer--Hasse--Noether theorem and the central simple algebra
embedding criterion; equivalently, extending scalars to $L$ kills the
local invariant $1/2$ precisely when the local extension has degree two.
See \cite[Proposition 14.6.7]{Voight2021}.

Suppose first that $t\ne0$, and put
\[
 \delta=1-\frac{2}{3t^2}.
\]
This number is nonzero: $t^2=2/3$ has no rational solution.
At $2$, put $m=v_2(t)$. If $m>0$, then $v_2(\delta)=1-2m$ is odd,
so $\delta$ is not a square in $\mathbb Q_2$. If $m=0$, then
$t^2\equiv1\pmod8$ and $\delta\equiv3\pmod8$, again a nonsquare.
If $m<0$, then $\delta\in1+8\mathbb Z_2$, hence it is a square in
$\mathbb Q_2$. At $3$, if $m=v_3(t)\ge0$ then
$v_3(\delta)=-1-2m$ is odd; if $m<0$, then
$\delta\in1+3\mathbb Z_3$ and Hensel's lemma makes it a square.
Consequently
\[
 \delta\notin\mathbb Q_2^{\times2},\quad
 \delta\notin\mathbb Q_3^{\times2}
 \quad\Longleftrightarrow\quad t\in R.
\]
The embedding criterion shows that this is equivalent to the existence
of a pure quaternion $xi+yj+zij$ with square $\delta$, that is,
\[
 2x^2+3y^2-6z^2=\delta.
\]
This equivalence also excludes the case in which $\delta$ is a rational
square: a division quaternion algebra cannot contain a noncentral element
whose square is a nonzero rational square, since
$(\alpha-s)(\alpha+s)=0$ would force $\alpha=\pm s$.
A nonzero pure quaternion has reduced trace zero, so it cannot be such
a nonzero rational scalar.
Multiplying the displayed equation by $3t^2$ and setting
$u=3ty$, $v=tx$, $w=3tz$ gives the equation in the statement.
These substitutions are invertible for $t\ne0$.
Finally $t=0$ belongs to $R$ and is represented by $(u,v,w)=(0,0,1)$.
\end{proof}

\begin{lemma}\label{lem:two-semilocal-inputs}
Let $F(A,B)=A^2+AB-B^2$. For every $a,b\in\mathbb Q$ and $p=2,3$,
\[
 v_p(F(a,b))=2\min\{v_p(a),v_p(b)\}.
\]
In particular $(a,b)\in R^2$ if and only if $F(a,b)\in R$.
\end{lemma}
\begin{proof}
The assertion is immediate for $a=b=0$. Otherwise scale $a,b$ by a
common power of $p$ until both are integral and at least one is a unit.
The reduction of $A^2+AB-B^2$ has no nontrivial zero over either
$\mathbb F_2$ or $\mathbb F_3$: over $\mathbb F_2$ one checks the
three nonzero pairs, and over $\mathbb F_3$ the discriminant $5$
is a nonsquare. Thus the scaled value is a unit, proving the equality.
\end{proof}

\begin{corollary}[Five-variable replacement for the anchor]
\label{cor:five-anchor}
Fix rational inputs $f_0,f_1,C$ with $C\ne0$. The condition
$f_0/C,f_1/C\in R$ is equivalent to the existence of five rational
variables $r,D,u,v,w$ satisfying the three quadratic equations
\begin{align*}
 r&=f_0^2+f_0f_1-f_1^2,\\
 D&=C^2,\\
 u^2+6v^2-2w^2-3r^2+2D^2&=0.
\end{align*}
\end{corollary}
\begin{proof}
By Lemma~\ref{lem:two-semilocal-inputs}, the desired condition is
$r/C^2\in R$, with $r=F(f_0,f_1)$. Apply
Proposition~\ref{prop:explicit-semilocal} to $t=r/C^2$ and multiply
its equation by $C^4$. Multiplication of each of its three witnesses
by $C^2$ is a bijection of $\mathbb Q$, giving the stated equation.
\end{proof}

The released anchor uses two scaling variables and two copies of a
seven-witness trace-sum condition, hence sixteen witnesses in this module.
Corollary~\ref{cor:five-anchor} replaces that module by five witnesses.
It therefore saves eleven variables before any other change and retains
quadratic residuals, so their sum of squares remains quartic. The exact global counts and independent incidence checks are given in the next section.
The three displayed residuals have integer coefficients of absolute value
at most six. This is an exact replacement, with no change to the arithmetic
meaning or strength of the anchor.

The same method handles arbitrary finite tuples. For a finite nonempty
set $S$ of rational primes and a positive integer $m$, choose a monic
polynomial $g\in\mathbb Z[T]$ of degree $m$ whose reductions are
irreducible modulo every $p\in S$; finite-field irreducible polynomials
and the Chinese remainder theorem supply one. If $\alpha$ is a root, let
\[
 N_m(X_1,\ldots,X_m)=
 N_{\mathbb Q(\alpha)/\mathbb Q}
 (X_1+X_2\alpha+\cdots+X_m\alpha^{m-1}).
\]
For $p\in S$ the extension is unramified and inert at $p$, and its
power basis is an integral basis locally. Therefore
\[
 v_p(N_m(\mathbf x))=m\min_i v_p(x_i).
\]
This proves, with no new existential variables, that all coordinates of
$\mathbf x$ belong to $\bigcap_{p\in S}\mathbb Z_{(p)}$ if and only if
$N_m(\mathbf x)$ does. It is the norm-form version of
\cite[Proposition 4.8]{Daans2024}. In the present two-prime, two-input
case the polynomial $T^2-T-1$ gives precisely the form $F$ above.



\section{Rational relation combining and finite algebraic choices}
\subsection{Combining square conditions without assuming integrality}

Matiyasevich and Robinson's relation-combining theorem is an essential
source of small integer Diophantine representations
\cite{MatiyasevichRobinson1975,Jones1982}. Its integer-domain hypotheses
must not be silently replaced by rational-domain hypotheses. The
appropriate rational version is Zhang--Sun
\cite[Theorem 1.2]{ZhangSun2022}. Their separation parameter is
$(k+\sum A_i^2)(1+\sum A_i^{-2})$, and their theorem assumes that every
$A_i$ is nonzero. The following slightly simpler separation bound also
works and has a direct proof.

\begin{lemma}\label{lem:additive-square-separation}
Let $k\ge1$. For nonzero $A_1,\ldots,A_k\in\mathbb Q$, put
\[
 W=\sum_{i=1}^k(A_i^2+A_i^{-2}),\qquad
 J(T)=\prod_{\boldsymbol\epsilon\in\{\pm1\}^k}
 \left(T-\sum_{i=1}^k\epsilon_i W^{i-1}\sqrt{A_i}\right).
\]
Then $J\in\mathbb Q[T]$, and $J$ has a rational root if and only if
all $A_i$ are rational squares. Every complex root of $J$ is nonzero.
\end{lemma}
\begin{proof}
Independently changing the formal sign of each square root fixes the
product, so its coefficients are rational polynomials in the $A_i$ and
$W$. Let
$m=\min_i\sqrt{|A_i|}$, $M=\max_i\sqrt{|A_i|}$, and $\rho=M/m\ge1$.
The two summands corresponding to maximal and minimal absolute value
show that $W\ge2\rho^2\ge1+\rho$; the same inequality is immediate
when both extrema occur at one index. Therefore for every $j$,
\[
 \sum_{i<j}M W^{i-1}<\frac{M W^{j-1}}{W-1}
 \le mW^{j-1}.
\]
No nonempty signed subsum of the radicals can vanish, since its last
term has larger absolute value than the sum of the preceding terms.
In particular every root is nonzero. If a displayed root is rational,
every automorphism of the multiquadratic splitting field fixes it.
Subtracting its conjugate would give a nonempty signed subsum if any
radical changed sign. Thus every radical is fixed by the full Galois
group and is rational. The converse is immediate.
\end{proof}

There is also an explicit one-variable polynomial that enforces the
nonzero condition itself. Put
\[
 S=\prod_{i=1}^k A_i,\qquad D=S^2,\qquad
 N=D\sum_i A_i^2+\sum_i\prod_{j\ne i}A_j^2
\]
and define
\[
 K_k(\mathbf A,U)=\prod_{\boldsymbol\epsilon\in\{\pm1\}^k}
 \left(1-U\sum_{i=1}^k\epsilon_i
 S^{2k-2i+1}N^{i-1}\sqrt{A_i}\right).
\]
Independent sign changes show that this is an integer-coefficient
polynomial in $\mathbf A,U$. If $S=0$ it is identically $1$.
If $S\ne0$, the signed sum in each factor is
$S^{2k-1}\sum_i\epsilon_iW^{i-1}\sqrt{A_i}$, so
Lemma~\ref{lem:additive-square-separation} proves
\[
 \exists U\in\mathbb Q:K_k(\mathbf A,U)=0
 \quad\Longleftrightarrow\quad
 A_i\in\mathbb Q^{\times2}\text{ for all }i.
\]
This polynomial has degree $2^k$ in $U$; total degree grows
exponentially with $k$. To allow zero square inputs, take the finite
product, over $I\subseteq\{1,\ldots,k\}$, of
\[
 \sum_{i\notin I}A_i^2+K_{|I|}((A_i)_{i\in I},U)^2,
\]
using zero in place of $K_0$ for the empty subset. One shared rational
$U$ suffices for all branches. This treats zero inputs exactly, but its
expanded degree is correspondingly large.

Private square roots can be replaced in this way only when no other
condition needs their particular values or signs. If roots occur elsewhere,
the complete finite algebraic system must be compressed, as below.
For example, replacing $y^2=A,\ y=1$ by the bare condition that $A$
is a square would be invalid.

\subsection{One variable for finite choices and nonzero conditions}

The useful general principle is that the number of finite algebraic
choices need not equal the number of rational witnesses. Sun's recent
preprint \cite[Sections 10.3--10.6]{Sun2026Rational86} uses this principle
for a bounded-variable undecidability construction. The elementary
algebra below is independently verifiable and does not depend on that
preprint's height or prime-pattern estimates. Its undecidability result
is not itself a definition of $\mathbb Z$ inside $\mathbb Q$.

\begin{lemma}\label{lem:finite-separation}
Let $A$ be a reduced finite-dimensional $\mathbb Q$-algebra of dimension
$d\ge1$, generated by $N\ge1$ elements $a_1,\ldots,a_N$. Some integer
\[
 1\le b\le1+(N-1)\binom d2
\]
makes $z_b=a_1+ba_2+\cdots+b^{N-1}a_N$ have distinct values under all
$d$ embeddings $A\to\overline{\mathbb Q}$. For such a $b$, the
characteristic polynomial of multiplication by $z_b$ has a rational root
if and only if $A$ has a $\mathbb Q$-algebra homomorphism to $\mathbb Q$.
\end{lemma}
\begin{proof}
Two different embeddings differ on a generator, so equality of their
values on $z_b$ is a nonzero polynomial equation in $b$ of degree at most
$N-1$. At most $(N-1)\binom d2$ values are bad. For a good value the
characteristic polynomial has precisely the distinct embedding values
as its roots. If one value is rational, the absolute Galois group fixes
that value and hence its unique corresponding embedding; the values
of all generators are therefore rational. Conversely a rational
homomorphism supplies a rational root.
\end{proof}

\begin{lemma}\label{lem:reciprocal-guard}
Suppose
\[
 P_{\mathbf x}(T)=T^d+\sum_{j=1}^d
 \frac{N_j(\mathbf x)}{B(\mathbf x)}T^{d-j}
\]
is monic on $B\ne0$, and let $H(\mathbf x)$ be the product of the
numerators of all further required inequations. Then
\[
 \Psi(\mathbf x,U)=1+\sum_{j=1}^d
 N_j(\mathbf x)B(\mathbf x)^{2j-1}H(\mathbf x)^jU^j
\]
has a rational root in $U$ exactly when $BH\ne0$ and $P_{\mathbf x}$
has a nonzero rational root.
\end{lemma}
\begin{proof}
If $BH=0$, all nonconstant terms vanish and $\Psi=1$.
If $BH\ne0$, set $L=B^2H$. For $U\ne0$ the identity
\[
 \Psi(\mathbf x,U)=(LU)^dP_{\mathbf x}(1/(LU))
\]
proves the assertion, and $\Psi(\mathbf x,0)=1$ excludes $U=0$.
\end{proof}

Zero roots are covered by taking the two finite alternatives $P(T)$
and $P(T-1)$: a rational root of one of them is nonzero. Thus nonzero
conditions need no additional inverse variable once one finite-choice
variable is already present. Equations depending only on $\mathbf x$
can be imposed by adding their squares to $\Psi^2$, and finite alternatives
can be combined by multiplying the resulting polynomials, using the
same $U$ in every factor.

For completeness, this also handles nonreduced finite algebras and
additional equations involving their generators. Quotient by the ideal
of those equations, and discard factors violating an inequation by
localizing at its value; these are finite-dimensional algebra operations.
Pass to the reduced quotient. In characteristic zero its nilradical is
the kernel of the trace pairing $(a,b)\mapsto\operatorname{Tr}(ab)$:
a product decomposition into Artinian local factors reduces the assertion
to the nondegeneracy of the trace pairing of a separable field extension.
Gaussian elimination then computes this quotient. When coefficients depend
on free parameters, branch on the finitely many pivots being zero or
nonzero and record all denominators. On every branch apply
Lemma~\ref{lem:finite-separation}, enforce the nonzero discriminant of
the selected characteristic polynomial, and apply
Lemma~\ref{lem:reciprocal-guard}. This is an effective finite construction
when the original algebra has a fixed finite basis, for example for a
tower of monic equations. It preserves exceptional specializations.

Dropping the separation condition is unsound: $x^2=y^2=2$ has no
rational pair, but the characteristic polynomial for $x+y$ is
$T^2(T^2-8)$ and has the rational root zero. Its repeated roots signal
precisely the collision. Using $x+2y$ instead gives
$T^4-20T^2+36$, with no rational root.



\section{Three rational coordinates for an elliptic point over a quadratic field}
\label{sec:quadratic-point-compression}

The following elementary elimination reduces the ambient coordinate count
of a point over a quadratic field.  It trades degree for variables, which
is appropriate when the number of variables is the first optimization
criterion.  The square-root expressions below are reconstructed rational
functions, not new quantified variables.

\begin{proposition}[Eliminating one coordinate of a quadratic square root]
\label{prop:quadratic-square-root}
Let $F$ be a field of characteristic different from two, let
$\delta\in F^\times$ be a nonsquare, and write $s^2=\delta$.
Fix $j\in F$ and put $\lambda=j^2-\delta\ne0$.
For $A,B,C\in F$ with $C\ne0$, define
\[
 U=C^2-A-jB,\qquad V=\frac{U}{\lambda},\qquad
 b_1=\frac{B/2+jV}{C},\qquad b_0=C-jb_1.
\]
Then the following are equivalent:
\begin{enumerate}
 \item $(b_0+b_1s)^2=A+Bs$;
 \item $C^2V-(B/2+jV)^2=0$;
 \item $\Gamma_j(A,B,C)=0$, where
 \begin{align*}
 \Gamma_j(A,B,C)
 ={}&4\delta C^4
   -4\bigl((j^2+\delta)A+2j\delta B\bigr)C^2\\
   &+\bigl(2jA+(j^2+\delta)B\bigr)^2.
 \end{align*}
\end{enumerate}
The correspondence identifies all nonzero square roots
$b=b_0+b_1s$ satisfying $b_0+jb_1\ne0$ with the nonzero solutions
$C$ of the last equation, by $C=b_0+jb_1$.
\end{proposition}
\begin{proof}
First suppose $b^2=A+Bs$ and $C=b_0+jb_1\ne0$.
Comparing coefficients gives
\[
 A=b_0^2+\delta b_1^2,\qquad B=2b_0b_1,
 \qquad C^2=A+jB+\lambda b_1^2.
\]
Hence $V=b_1^2$ and $Cb_1=B/2+jV$, proving the reconstruction formulas
and the second assertion.

Conversely, suppose the second assertion holds and use the displayed
reconstruction formulas.  Because $C\ne0$, it follows that $b_1^2=V$.
Consequently
\[
 b_0b_1=Cb_1-jb_1^2=B/2
\]
and
\begin{align*}
 b_0^2+\delta b_1^2
 &=C^2-2jCb_1+(j^2+\delta)V\\
 &=C^2-jB-(j^2-\delta)V=A.
\end{align*}
Thus $b^2=A+Bs$, and $b\ne0$ follows from $C=b_0+jb_1\ne0$.
Finally, multiplication of the second assertion by $-4\lambda^2$
gives exactly $\Gamma_j=0$.  This also proves that the two
constructions are inverse to one another on the specified charts.
\end{proof}

For $F=\mathbb Q$ and $\delta=2$, this becomes
\[
 \Gamma_j(A,B,C)=
 8C^4-4\bigl((j^2+2)A+4jB\bigr)C^2
 +\bigl(2jA+(j^2+2)B\bigr)^2.
\]
The root coordinates can equivalently be reconstructed as
\begin{align*}
 b_0&=\frac{-4C^2+2j^2A+j(j^2+2)B}{2(j^2-2)C},\\
 b_1&=\frac{2jC^2-2jA-(j^2+2)B}{2(j^2-2)C}.
\end{align*}
Thus neither $V$, $b_0$, nor $b_1$ consumes an additional variable.

\begin{lemma}[A common chart for finitely many points]
\label{lem:common-quadratic-chart}
Let $b^{(1)},\ldots,b^{(N)}$ be nonzero elements of
$\mathbb Q(\sqrt\delta)$, where $\delta$ is a rational nonsquare.
There exists $j\in\{0,1,\ldots,N\}$ such that
\[
 b^{(r)}_0+j b^{(r)}_1\ne0\qquad(1\leq r\leq N).
\]
\end{lemma}
\begin{proof}
For a fixed nonzero pair $(b_0,b_1)$, the expression $b_0+jb_1$
vanishes for at most one rational $j$.  If $b_1=0$, it never
vanishes.  At most $N$ values are therefore excluded, and the
displayed set has $N+1$ distinct elements.
\end{proof}

\subsection{Application to elliptic doubling}

Let $K=\mathbb Q(\sqrt2)$ and let
\[
 E_D:\ b^2=a^3-Da,\qquad D\in K^\times.
\]
The assumption $D\ne0$ makes this Weierstrass cubic nonsingular.
Write $a=a_0+a_1\sqrt2$ and
$a^3-Da=A+B\sqrt2$ with $A,B\in\mathbb Q$.
An affine point with nonzero ordinate is therefore represented,
on the $j$-chart, by only three rational coordinates
$(a_0,a_1,C)$, subject to $C\ne0$ and $\Gamma_j(A,B,C)=0$.
For fixed $D$, this polynomial has total degree at most six in these
three coordinates: $A,B$ have degree at most three and the leading
terms come from a square of their linear combination.

Reconstruct $b$ as above and put
\[
 h=\frac{3a^2-D}{2b},\qquad
 \xi=h^2-2a,\qquad \eta=h(a-\xi)-b.
\]
These rational functions give the usual doubling formula
$[2](a,b)=(\xi,\eta)$.  Since $b\ne0$, the doubled point is affine.
Conversely, every nonidentity point of $2E_D(K)$ has an affine
preimage with nonzero ordinate: a preimage equal to the identity or
to a point with ordinate zero would double to the identity.
For a list of $N$ such doubled points, choose a halfpoint of each
and then apply Lemma~\ref{lem:common-quadratic-chart}.  Thus the
union of the $N+1$ common charts represents all such lists.

There is no change to the group-theoretic assertions required by the
original construction.  In particular, the identity point must be
handled separately if a surrounding formula permits it.  In the
proposed application, the 211 point occurrences are constrained to be
nonidentity by their original nonvanishing conditions; this condition
is part of the application, not a conclusion of the elimination alone.

\subsection{Guards and finite unions without selector variables}

Every rational function used in a fixed chart can be written as a
quotient of rational polynomials.  Suppose the required denominator
factors in that branch are $d_1,\ldots,d_r$, and set
$\Delta=\prod_{k=1}^r d_k$.  The single equation
\[
 Z\Delta-1=0
\]
with one rational witness $Z$ enforces all the nonvanishing conditions
at once.  After imposing this equation, clearing those denominators
from the remaining branch equations preserves both directions of
the existential assertion.  In particular, the factors $C_i$ must
be guarded; merely imposing $\Gamma_j=0$ at $C_i=0$ would introduce
points outside the displayed chart.

If a branch is given by rational polynomial residuals
$f_{j,1},\ldots,f_{j,r_j}$, including its own equation
$Z\Delta_j-1=0$, set
\[
 H_j=\sum_{k=1}^{r_j}f_{j,k}^2.
\]
Over $\mathbb Q$, the union of the branches is then represented by
\[
 \prod_{j=0}^{N}H_j=0.
\]
All branches may reuse the same list of witness-variable names: a
zero of the product belongs to at least one factor, and that factor
alone must satisfy its denominator guard.  Requiring the guards of
inactive branches would generally destroy completeness and is not
part of this construction.  Additional finite choices, such as the
256 fixed height-class choices, can be included in the same product
over pairs of branch indices without selector variables.



\section{One witness for a constructible finite-fiber projection}
\label{sec:one-finite-variable}
A constructible set is a finite union of locally closed algebraic sets.
A projection has finite geometric fibers if every fiber over every
point with coordinates in an algebraic closure is a finite set.
Finite rational fibers alone are not the hypothesis used here.

\begin{theorem}\label{thm:finite-fibre-one-variable}
Let $X\subseteq\mathbb A^{d+m}_{\mathbb Q}$ be constructible over
$\mathbb Q$, and suppose the coordinate projection
$\pi:X\to\mathbb A^d_{\mathbb Q}$ has finite geometric fibers.
There is an effectively constructible polynomial
$P\in\mathbb Z[\mathbf x,U]$ such that, for every $\mathbf a\in\mathbb Q^d$,
\[
 \mathbf a\in\pi(X(\mathbb Q))
 \quad\Longleftrightarrow\quad
 \exists U\in\mathbb Q:\ P(\mathbf a,U)=0.
\]
\end{theorem}
\begin{proof}
We give the construction including its exceptional loci. Decompose $X$
into finitely many locally closed affine pieces; represent an inequation
$h\ne0$ by the coordinate relation $hv=1$ when forming a coordinate ring.
The added coordinate is uniquely determined, so geometric fibers stay
finite. Decompose the image closure into irreducible pieces and restrict
accordingly. On one piece let $R$ be the coordinate ring of the irreducible
base closure $B$, let $L=\operatorname{Frac}(R)$, and let $A$ be the reduced
coordinate ring of the affine source. Its generic fiber algebra
$A_L$ is finite-dimensional, reduced, and separable over $L$; it is a
finite product of finite separable field extensions.

The primitive element theorem for a finite separable algebra gives
$A_L\simeq L[T]/(h)$ with $h$ monic. Clearing the denominators of a
primitive element, of expressions for the source generators in that
primitive element, and of their finitely many relations gives a nonzero
$g\in R$ for which
\[
 A_g\simeq R_g[T]/(h),\qquad h\in R_g[T]\text{ monic}.
\]
For precision, the map is surjective after clearing the generator
expressions; its kernel vanishes over $L$ and is therefore zero because
$R_g[T]/(h)$ is free over the domain $R_g$. Any intermediate relation
which vanishes only over $L$ is made zero by a further localization.
Thus this is an actual isomorphism on an open base set, not just a
bijection on generic points. Consequently a rational point of $B\cap D(g)$
has a rational lift exactly when the specialized monic $h$ has a rational
root. The residual part lies over a proper closed subset of $B$, so
induction on the dimension of the image closure gives finitely many such
charts. Each chart comes with polynomial equations $E_i(\mathbf x)=0$
for its base closure and one or more polynomial inequations; multiply the
latter together to obtain $G(\mathbf x)\ne0$. Choose polynomial
representatives in the ambient base space for all elements of $R$.

It remains to express each chart with one quantified variable. Write its
monic equation, after choosing a common coefficient denominator, as
\[
 h(T)=T^n+\sum_{j=1}^n\frac{N_j(\mathbf x)}{D(\mathbf x)}T^{n-j}.
\]
Add $D\ne0$ to the recorded chart conditions. A rational root of $h$
may be zero; to cover this exactly, use the two fixed alternatives
$h(T)$ and $h(T-1)$. If $r$ is a rational root of $h$, then the two
polynomials have roots $r$ and $r+1$, at least one of which is nonzero.
Conversely a rational root of either polynomial gives a rational root of
$h$. The shifted polynomial is still monic and has the same permissible
common denominator $D$.

For either alternative, with its coefficient numerators again denoted
$N_j$, set
\[
 \Psi(\mathbf x,U)=1+\sum_{j=1}^n
 N_j(\mathbf x)D(\mathbf x)^{2j-1}G(\mathbf x)^jU^j.
\]
At $DG=0$ this polynomial is identically $1$. At $DG\ne0$, put
$L_0=D^2G$. Then, for $U\ne0$,
\[
 \Psi(\mathbf x,U)=(L_0U)^n h(1/(L_0U)),
\]
where $h$ now denotes the selected alternative. Also
$\Psi(\mathbf x,0)=1$. Hence a rational zero of $\Psi$ exists exactly
when all inequations hold and that alternative has a nonzero rational
root. This construction cannot acquire spurious solutions at a vanished
denominator, because its value there is exactly one.

The polynomial
\[
 Q(\mathbf x,U)=\Psi(\mathbf x,U)^2+\sum_iE_i(\mathbf x)^2
\]
therefore describes that chart and alternative. Multiply $Q$ over the
finite list of charts and the two alternatives, always reusing the same
variable $U$, and clear numerical coefficient denominators. This gives
$P$. A product is zero if and only if some factor is zero, and over
$\mathbb Q$ each displayed sum of squares is zero if and only if all of
its summands vanish. The empty image is represented by $P=1$.
All the steps are effective operations in polynomial rings over
$\mathbb Q$, including ideal decomposition and localization, followed by
finite-dimensional algebra over the resulting function fields. The
induction terminates because the dimension of every exceptional base
strictly decreases. This effectivity assertion does not mean that these
computations have been executed for the source system.
\end{proof}

\begin{corollary}
A finite-fiber presentation retaining $d$ rational coordinates, one of
which is the designated input, yields a polynomial in at most $d+1$
total variables, with at most $d$ existential witnesses and the same
rational input projection. In particular, the presentation with $d=557$
proved below gives at most $558$ total variables. Its identification
with $\mathbb Z$ retains the arithmetic hypothesis of the source.
\end{corollary}

The elementary reciprocal transformation in this proof agrees with
\cite[Section 10.6, equation (10.8)]{Sun2026Rational86}; the chart argument
above extends it from the finite algebras explicitly treated there to all
constructible sets with finite geometric fibers. This citation supports
the algebraic device only. The bounded-variable undecidability theorem
in that recent preprint and its analytic arithmetic inputs are not used.
Its convention is that a monic degree-$n$ polynomial is written
$T^n+\sum_{j=1}^n(N_j/D)T^{n-j}$: the reciprocal polynomial has exponent
$U^j$ and denominator-clearing factor $D^{2j-1}$, not $D^{j-1}$.
The two shifts are essential if zero could be the only rational root.
All mathematical claims used here are proved above; no result about
the arithmetic of that preprint is needed.


\subsection{Conjunctions and the literature on existential rank}
\label{subsec:intersection-rank}

The general relation between witness number and finite fibers is part of
the theory of existential rank and essential fiber dimension developed
by Daans, Dittmann, and Fehm \cite[Corollary 4.21]{DDF2025}. Here
\emph{existential rank} means the minimum number of existential
quantifiers in a definition of a specified set. Our finite-fiber proof
is a direct constructive treatment of the special case needed here,
including exceptional strata and nonzero conditions. We make no claim
that the general quantifier-saving phenomenon is new.

For clarity, it gives the following useful known intersection bound
\cite[Theorem 1.4]{DDF2025}, also stated in
\cite[Theorem 2.5]{Daans2024}, with an explicit constructive proof over
the rationals.

\begin{proposition}\label{prop:intersection-rank}
Let $D_i\subseteq\Q^n$, for $1\le i\le r$ and $r\ge1$, be
defined by one polynomial equation with $m_i\ge1$ existential
witnesses. Their intersection has a one-polynomial definition with at
most
\[
 1+\sum_{i=1}^r(m_i-1)
\]
witnesses. Such a definition can be constructed effectively from the
given polynomials.
\end{proposition}
\begin{proof}
Write the $i$th polynomial as
\[
 F_i(\mathbf x,\mathbf u_i,Z_i)
   =\sum_{j=0}^{d_i}a_{ij}(\mathbf x,\mathbf u_i)Z_i^j,
\]
where $\mathbf u_i$ has $m_i-1$ coordinates and the witness lists
are private to their respective predicates. Let $C_i$ be the closed
condition that all coefficients $a_{ij}$ vanish. Replace the zero set
of $F_i$ by the constructible set
\[
 (\{F_i=0\}\setminus C_i)\ \cup\ (C_i\cap\{Z_i=0\}).
\]
This replacement preserves the projection to
$(\mathbf x,\mathbf u_i)$ over any field: a nonzero polynomial has
its usual roots, and an identically zero polynomial permits the choice
$Z_i=0$. Each geometric fiber in $Z_i$ is now finite. Nonzero
constant specializations give empty fibers and cause no exception.

Take the fiber product of these presentations over $\mathbf x$.
It has finite geometric fibers over
$(\mathbf x,\mathbf u_1,\ldots,\mathbf u_r)$.
Theorem~\ref{thm:finite-fibre-one-variable} replaces all the $Z_i$
by one additional variable. Quantifying the remaining
$\sum_i(m_i-1)$ private coordinates proves the formula and its
effectiveness. Genuinely shared coordinates remain in $\mathbf x$;
they are never silently duplicated or identified.
\end{proof}

For example, each of the forty-two pole predicates has a four-witness
presentation after retaining its three free pole parameters and encoding
its remaining finite choices by one variable. The proposition combines
them with $42\cdot4-41=127$ witnesses. This is exactly the $126$
retained pole coordinates and one common finite-choice variable already
used in the global finite-fiber construction. It does not provide another
forty-one-variable saving on top of that construction.


% Input fragment extending finite_fibre_strengthening.tex to other fields.
\subsection{Finite-fiber compression over other fields}
\label{subsec:finite-fibre-fields}

The one-variable compression theorem does not require an ordering. The
place where rational squares were used was the conjunction of finitely
many equations. Over any field that is not algebraically closed, a
homogenized irreducible polynomial supplies the needed replacement.

\begin{lemma}\label{lem:anisotropic-binary-conjunction}
Let $K$ be a field that is not algebraically closed. There is a homogeneous
polynomial $H(A,B)\in K[A,B]$ of degree $e\ge2$ such that
\[
 H(a,b)=0\quad\Longleftrightarrow\quad a=b=0
 \qquad(a,b\in K).
\]
Consequently any finite system of polynomial equations over $K$ can be
combined into one polynomial equation without introducing variables.
For $r$ equations of degrees at most $D$, a balanced combination has
degree at most $e^{\lceil\log_2 r\rceil}D$ when $r\ge1$.
\end{lemma}
\begin{proof}
Choose a monic irreducible polynomial $h(T)\in K[T]$ of degree $e\ge2$,
and homogenize it:
\[
 H(A,B)=B^e h(A/B).
\]
This notation denotes the polynomial obtained by homogenization, also
at $B=0$. If $b\ne0$, then $H(a,b)=0$ would give a root $a/b$ of $h$
in $K$, contradicting irreducibility. If $b=0$, monicity gives
$H(a,0)=a^e$, so $a=0$. Conversely $H(0,0)=0$.

To combine residuals $f_1,\ldots,f_r$, pad the list with zero
polynomials to a power of two, and repeatedly apply $H$ to pairs.
At each stage a combined residual vanishes exactly when both its
inputs vanish. After $\lceil\log_2 r\rceil$ stages, the one remaining
equation is equivalent to the whole system. Each stage multiplies
the degree bound by at most $e$.
\end{proof}

\begin{theorem}\label{thm:finite-fibre-arbitrary-field}
Let $K$ be a field of characteristic zero that is not algebraically
closed. Let $X\subseteq\mathbb A^{d+m}_K$ be constructible over $K$,
and suppose every geometric fiber of the coordinate projection
$\pi:X\to\mathbb A^d_K$ is finite. There is a polynomial
$P(\mathbf x,U)\in K[\mathbf x,U]$ such that
\[
 \mathbf a\in\pi(X(K))
 \quad\Longleftrightarrow\quad
 \exists U\in K:\ P(\mathbf a,U)=0
 \qquad(\mathbf a\in K^d).
\]
For a number field $K$, one may take $P\in\mathcal O_K[\mathbf x,U]$,
and the construction is effective from equations and inequations for
$X$. More generally it is effective whenever the coefficient field is
presented with the algebraic operations needed for factorization,
elimination, and primitive-element computation, and an irreducible
polynomial of degree at least two is available effectively.
\end{theorem}
\begin{proof}
The algebraic part of Theorem~\ref{thm:finite-fibre-one-variable} works
over $K$. Indeed, $K$ is infinite and perfect. After passing to reduced
affine pieces and localizing on each irreducible image closure, the
generic finite algebra is a product of finite separable field extensions
of the image's function field. The primitive element theorem gives
a monic polynomial presentation. Clearing the finitely many
denominators extends the presentation over a nonempty open subset.
The complement has smaller image dimension, so induction produces
finitely many charts, including every exceptional finite fiber.

Thus, on each chart, a point lies in the image precisely when
\[
 f_1(\mathbf x)=\cdots=f_r(\mathbf x)=0,
 \qquad g(\mathbf x)\ne0,
\]
and a certain monic polynomial
\[
 h(T)=T^n+\sum_{j=1}^n\frac{N_j(\mathbf x)}{B(\mathbf x)}T^{n-j}
\]
has a root in $K$. All coefficient denominators are nonzero on the
chart; enlarge its guard to include them. For $\epsilon\in\{0,1\}$,
let $h_\epsilon(T)=h(T-\epsilon)$ and write its coefficients with
the same common denominator $B$. Define
\[
 \Psi_\epsilon(\mathbf x,U)
 =1+\sum_{j=1}^n
      N_{\epsilon,j}(\mathbf x)
      B(\mathbf x)^{2j-1}g(\mathbf x)^jU^j.
\]
When $Bg=0$ or $U=0$, this expression is $1$. Otherwise, with
$L=B^2g$, the exact identity
\[
 \Psi_\epsilon(\mathbf x,U)
   =(LU)^n h_\epsilon(1/(LU))
\]
shows that it has a zero precisely when $h_\epsilon$ has a nonzero
root. A root of $h$ gives a nonzero root of at least one of $h(T)$
and $h(T-1)$; in particular a zero root becomes the root one after
the translation. Hence these two choices cover all roots.

For each choice of chart and $\epsilon$, apply
Lemma~\ref{lem:anisotropic-binary-conjunction} to the residuals
$f_1,\ldots,f_r,\Psi_\epsilon$. This gives one equation in
$\mathbf x,U$ equivalent over $K$ to their conjunction. Take the
product of these equations over the finitely many charts and the
two translations. A finite product vanishes if and only if one
factor vanishes, which proves the assertion without adding another
witness. For an empty constructible set take $P=1$.

If $K$ is a number field, clearing a common numerical denominator
places the finitely many coefficients in $\mathcal O_K$ without
changing the zero set. All algebraic operations in the construction
are effective over a number field. The more general effective
assertion follows under the stated algorithmic hypotheses.
\end{proof}

For number fields, Lemma~\ref{lem:four-variable-anisotropic} supplies
a quadratic form in four arguments that vanishes only at the zero
tuple. A balanced tree of that form combines $r$ residuals of degree
at most $D$ with degree at most $2^{\lceil\log_4 r\rceil}D$, again
without adding variables. Thus the same witness bound is available
even for totally imaginary number fields; the failure of sums of
squares there is not an obstruction to this compression argument.

\paragraph{Why the field hypothesis cannot simply be omitted.}
Over an algebraically closed field $K$, take $X$ to be the origin in
$\mathbb A^2_K$, with the identity projection to $\mathbb A^2_K$.
All geometric fibers are finite. Its rational image nevertheless
cannot be written as
\[
 \{(a,b)\in K^2:\exists U\in K,\ P(a,b,U)=0\}
\]
for a single polynomial $P$. If $P$ has positive degree in $U$,
then every point in the nonempty open subset where its leading
$U$-coefficient is nonzero has a lift, since every nonconstant
polynomial over $K$ has a root. The image therefore cannot be a
single point. If $P$ is independent of $U$, a nonconstant polynomial
in two variables has a zero set of dimension at least one whenever
it has a zero; a constant gives either the empty set or the whole
plane. Neither possibility gives exactly the origin. Thus
non-algebraic-closedness is essential for this general one-equation,
one-additional-variable statement.


% Proof fragment for the phase-two explicit finite-algebra construction.
\section{An explicit polynomial for a guarded finite algebra}
\label{sec:explicit-finite-algebra}

Let $R=\mathbb Z[x_1,\ldots,x_b]$. Consider a triangular monic presentation
\[
 A=R[z_1,\ldots,z_m]/(f_1,\ldots,f_m),\qquad
 f_i\in R[z_1,\ldots,z_i],
\]
where $f_i$ is monic of degree $e_i\ge1$ in $z_i$. Its standard
monomials $\prod_i z_i^{a_i}$, $0\le a_i<e_i$, form an $R$-basis.
In particular $A$ is free of rank $n=\prod_i e_i$. For $a\in A$,
write $M_a$ for its multiplication matrix in this ordered monomial basis.
The matrices $M_a$ commute, and their entries are integer polynomials in
the base coordinates. They are completely specified by repeated reduction
using the monic equations; they need not be materialized to specify a
determinant involving them.

Fix elements $G,E_1,\ldots,E_r\in A$. We assume that in every geometric
fiber, every point at which $G$ is nonzero has local length one. This
condition holds, for example, if $G$ contains the product of the diagonal
derivatives $\partial f_i/\partial z_i$ as a factor. We seek the rational
base points admitting a rational point of $A$ with
\[
 G\ne0,\qquad E_1=\cdots=E_r=0.
\]

For $0\le\ell\le m n(n+1)/2$, define
\[
 \lambda_\ell=1+\sum_{i=1}^m\ell^i z_i,\qquad
 M=M_{G\lambda_\ell},\qquad
 H(V)=\det(I-VM)=\sum_{k=0}^n a_kV^k.
\]
Here and below $a_k\in R$ depend on the fixed integer $\ell$; always
$a_0=1$. Also put
\[
 R_i(V)=\operatorname{Tr}\bigl(M_{E_i}\operatorname{adj}(I-VM)\bigr).
\]
This is a polynomial in $R[V]$ of degree at most $n-1$.

For each $1\le d\le n$, set
\begin{align*}
 H_d(V)&=\sum_{k=0}^d a_kV^k,\\
 g&=a_d\operatorname{Res}_{V}(H_d,H'_d),\qquad L=a_d^2g,\\
 \Psi(U)&=1+\sum_{j=1}^d a_{d-j}a_d^{2j-1}g^jU^j,\\
 \widetilde R_i(U)&=(LU)^{n-1}R_i\bigl(1/(LU)\bigr),\\
 Q_{\ell,d}(\mathbf x,U)&=
 \Psi(U)^2+\sum_{k=d+1}^n a_k^2+\sum_{i=1}^r\widetilde R_i(U)^2.
\end{align*}
The resultant is the Sylvester determinant with the specified degree
bounds $d$ and $d-1$. The expression for $\widetilde R_i$ denotes
polynomial coefficient reversal, hence defines a polynomial also at
$LU=0$. Every displayed expression has integer coefficients.

\begin{theorem}\label{thm:explicit-guarded-finite-algebra}
The explicitly specified polynomial
\[
 P(\mathbf x,U)=
 \prod_{\ell=0}^{m n(n+1)/2}\ \prod_{d=1}^{n} Q_{\ell,d}(\mathbf x,U)
\]
has the following property. For every $\mathbf a\in\mathbb Q^b$,
there exists $U\in\mathbb Q$ with $P(\mathbf a,U)=0$ if and only if
there is a rational point of the fiber algebra $A_{\mathbf a}$ satisfying
$G\ne0$ and all $E_i=0$. The formula uses only one coordinate in addition
to the $b$ base coordinates; its formal algebra generators are eliminated
by the determinants and are not quantified coordinates.
\end{theorem}
\begin{proof}
Fix a rational base point. First consider one factor $Q_{\ell,d}$.
At $a_dg=0$ or $U=0$, its reciprocal residual $\Psi$ equals one.
Otherwise, with $V=1/(LU)$,
\[
 \Psi(U)=(LU)^d\frac{H_d(V)}{a_d}.
\]
Thus a rational zero of $Q_{\ell,d}$ implies $a_dgU\ne0$, all
coefficients $a_k$ with $k>d$ vanish, and $H(V)=H_d(V)=0$.
The resultant condition makes this a simple root of $H$. Since
$H(0)=1$, the root $V$ is nonzero.

The rational matrix $I-VM$ has a one-dimensional kernel. Every
multiplication matrix $M_{z_i}$ preserves it and therefore acts on
it by a rational scalar $\zeta_i$. Applying the algebra relations
to this common eigenline shows that $(\zeta_1,\ldots,\zeta_m)$ is
a rational point of the fiber. Its $G\lambda_\ell$-value is $1/V$,
so its $G$-value is nonzero.

To verify the closed conditions, write
$J=\operatorname{Tr}(\operatorname{adj}(I-VM))$. At a root of $H$,
the adjugate identity gives
$M\operatorname{adj}(I-VM)=V^{-1}\operatorname{adj}(I-VM)$.
Differentiating the determinant then yields
\[
 J=-V H'(V)\ne0.
\]
The columns of the adjugate lie in the common eigenline, so
\[
 R_i(V)=E_i(\boldsymbol\zeta)J.
\]
Since $LU\ne0$, the equation $\widetilde R_i(U)=0$ is equivalent
to $R_i(V)=0$, and hence $E_i(\boldsymbol\zeta)=0$. This proves
the reverse implication without assuming that the rest of the fiber
algebra is reduced.

Conversely, suppose a suitable rational point exists. Over an
algebraic closure, multiplication by $G\lambda_\ell$ has the
value $G(p)\lambda_\ell(p)$ at every geometric point $p$, with
multiplicity equal to its local length. If $G(p)=0$, its factor
in $\det(I-VM)$ is one, even if the local algebra is nonreduced.
If $G(p)\ne0$, its local length is one by hypothesis.

At any such point $p$, the polynomial
$G(p)(1+\sum_{i=1}^m\ell^i z_i(p))$ is nonzero as a polynomial
in $\ell$, because its constant coefficient is $G(p)\ne0$.
For distinct such points $p,q$, their difference is also nonzero:
if all coefficients vanished, first $G(p)=G(q)\ne0$, and then
all their coordinates would agree. Each of these polynomials has
degree at most $m$. With at most $n$ geometric points, there are
at most $m n(n+1)/2$ bad values of $\ell$ in total, counting zero
values and collisions. Thus one of the specified integers makes
all nonzero guarded eigenvalues distinct and nonzero.

For that $\ell$, the polynomial $H$ is squarefree of some degree
$1\le d\le n$. The given rational point provides a nonzero
rational root $V$ of $H$. Its closed residuals vanish, so the trace
identities give $R_i(V)=0$. The leading coefficient and resultant
are nonzero. Taking $U=1/(LV)$ therefore gives a zero of the
corresponding factor $Q_{\ell,d}$, and hence of $P$.
\end{proof}

The degree cases are part of the literal finite formula. They include
every exceptional drop in degree of $H$. The discriminant condition
excludes bad separating choices, and the finite bound on $\ell$ proves
that this exclusion loses no valid rational point. In particular the
construction does not leave a generic primitive-element argument or an
unperformed exceptional-stratum decomposition implicit.


% Input fragment for the exact finite-algebra specification's bounds.
\subsection{Explicit degree and coefficient bounds for the specification}
\label{subsec:finite-algebra-bounds}

The following bounds concern the expanded polynomial defined by the
indexed determinant formula. They do not require expanding its matrices
or its final finite product. For a polynomial $f$, let $\|f\|_1$ denote
the sum of the absolute values of its coefficients. Then
$\|fg\|_1\le\|f\|_1\|g\|_1$ and
$\|f+g\|_1\le\|f\|_1+\|g\|_1$.

Give every base coordinate weight one, and assign a positive integer
weight $w_i$ to each algebra generator $z_i$. Choose the weights so that
every nonleading term of its monic relation has weight at most
$e_iw_i$. Reduction by these relations never increases weighted degree.
Writing $W=\sum_i(e_i-1)w_i$, the entries of a multiplication matrix
$M_f$ therefore have base degree at most $\operatorname{wdeg}(f)+W$:
the input basis monomial has weight at most $W$, and the output basis
monomial has nonnegative weight.

For a matrix over the polynomial ring, use the maximum column sum of
the coefficient norms of its entries. This norm is submultiplicative.
If
\[
 z_i^{e_i}+\sum_{k<e_i}c_{ik}z_i^k=0,
\]
the block companion description gives
\[
 \|M_{z_i}\|\le
 \max\left\{1,\sum_{k<e_i}\|M_{c_{ik}}\|\right\}.
\]
Extending to later monic generators preserves this estimate, since
multiplication by an earlier element is repeated on diagonal blocks,
up to a simultaneous permutation of rows and columns. Thus weights
and bounds $\|M_{z_i}\|\le2^{b_i}$ can be computed recursively from
the coefficient DAG. The source compiler records the entire table,
and the independent checker verifies the corresponding inequalities.

For the point quartics the weights can be taken to be two. A coefficient
bound for the companion matrix is the sum of the bounds for its two
nonzero lower coefficients. For each of the subsequent quadratic
generators, its nonzero branch has only the constant coefficient
$LM$, so one may take
\[
 w_i\ge\tfrac12\operatorname{wdeg}(LM),\qquad
 b_i\ge\log_2\|M_{LM}\|.
\]
The zero branch has relation $z_i=0$ and satisfies the same bounds.
Consequently one uniform table applies to all $2^{86}$ fixed branches.

Put $n=2^{508}$ for a uniform rank bound, $m=297$, and
$T=m n(n+1)/2$ for the largest separating parameter. Suppose the
recorded bounds give
\[
 \operatorname{wdeg}(G)\le w_G,\quad \|M_G\|\le2^{b_G},\qquad
 \operatorname{wdeg}(E_i)\le w_E,\quad \|M_{E_i}\|\le2^{b_E}.
\]
Set
\begin{align*}
 b_\lambda&=\lceil\log_2(m+1)\rceil
       +\max_i\{i\lceil\log_2 T\rceil+b_i\},\\
 \alpha&=w_G+\max_i w_i+W,\qquad \beta=b_G+b_\lambda,\\
 \alpha_E&=w_E+W.
\end{align*}
These bound the entry degrees and operator norm of
$M_{G\lambda_\ell}$, uniformly for all prescribed $\ell$ and all
source branches.

The coefficients of $H(V)=\det(I-VM)$ have base degree at most
$A_H=n\alpha$, and its full coefficient norm is at most
$2^{B_H}$ with $B_H=n(\beta+1)$. Indeed the coefficient norm of a
determinant is bounded by the product of its column sums. This bound
already includes the sum over determinant permutations; no additional
factorial is needed. By multilinearity in columns,
\[
 R_i(V)=\left.\frac{\partial}{\partial t}
       \det(I-VM+tM_{E_i})\right|_{t=0}
\]
has base coefficient degree at most
$A_R=(n-1)\alpha+\alpha_E$ and coefficient norm at most $2^{B_R}$,
where
\[
 B_R=\lceil\log_2 n\rceil+b_E+(n-1)(\beta+1).
\]
This accounts for the $n$ possible columns containing $M_{E_i}$.

For each degree case $d\le n$, the Sylvester matrix for
$\operatorname{Res}(H_d,H'_d)$ has $2d-1$ rows. Every row has
coefficient norm at most $d\|H\|_1$. Hence uniform bounds for the
reciprocal scale $L=a_d^3\operatorname{Res}(H_d,H'_d)$ are
\begin{align*}
 A_L&=(2n+2)A_H,\\
 B_L&=(2n+2)B_H+(2n-1)\lceil\log_2 n\rceil.
\end{align*}
The reciprocal residual has degree at most $n(A_L+1)$ and coefficient
norm at most $2^{nB_L+\lceil\log_2(n+1)\rceil}$. The reversed trace
residuals have degree at most $A_R+(n-1)(A_L+1)$ and coefficient norm
at most $2^{B_R+(n-1)B_L}$. Together with the high-coefficient
conditions $a_k=0$, these give a bound for the sum of squares by
doubling the largest degree, doubling the largest logarithmic norm
bound, and adding $\lceil\log_2(n+87)\rceil$ to the latter.
For the reciprocal estimate, no division of coefficient norms is used:
its $j$th coefficient is the polynomial
$a_{d-j}a_d^{3j-1}\operatorname{Res}(H_d,H'_d)^j$.
Its degree is at most $jA_L$ and its logarithmic norm bound is at most
$j(3B_H+B_{\rm Res})=jB_L$.

The number of final factors is also explicit. Write $n_0=4^{211}$.
For a mask with $k$ nonzero quadratic branches, its rank is
$n_k=n_0 2^k$. There are $\binom{86}{k}$ such masks, and its number
of separating-parameter and degree choices is
$n_k(1+m n_k(n_k+1)/2)$. Including the $212$ point charts and $256$
height choices gives exactly
\[
 B_{\mathrm{tot}}=54272\left[
 n_0 3^{86}+
 \frac{297}{2}\left(n_0^3 9^{86}+n_0^2 5^{86}\right)\right]
\]
factors. Multiplying the uniform branch degree and logarithmic
coefficient-norm bounds by $B_{\mathrm{tot}}$ bounds the complete
polynomial. All quantities here are exact integers, apart from the
displayed ceiling notation, which is evaluated by integer bit lengths.

For the emitted coefficient specification the resulting convenient
rounded bounds are
\[
 \deg P\le2^{3122},\qquad H(P)\le2^{\,2^{3127}}.
\]
The machine-readable specification records the stronger unrounded
integer degree bound and the stronger integer bound on $\log_2 H(P)$.
These numbers are rigorous upper bounds for the literal indexed
polynomial, not measurements of an expanded coefficient table.


\section{A finite-fiber presentation with 557 retained coordinates}
\label{sec:finite-fibre-application}

We now give a source-specific application of
Theorem~\ref{thm:finite-fibre-one-variable}.  The count is a theorem
about an existential polynomial obtainable by elimination.  It is not
the degree or variable count of an already expanded coefficient table.
Its identification with the ordinary integers remains conditional on
the arithmetic correctness of the source definition, exactly as for
any projection-preserving optimization of that definition.

All geometric fibers in this section are taken after extension to an
algebraic closure of $\mathbb Q$.  This qualification matters:
$\mathbb Q(\sqrt2)\otimes_{\mathbb Q}\overline{\mathbb Q}$ is the
product of two copies of $\overline{\mathbb Q}$, not a field.
Whenever a quadratic-field inverse occurs, we retain the inequation
that its rational norm is nonzero.  Its value is then uniquely determined
on both factors, including on geometric fibers.

\subsection{Point occurrences}

The root module has six nonidentity point occurrences, with labels
$A,S,T,Z$, height $J$, and common comparison $J$.  Each of the 41
comparison calls has five more, with labels $H,X_0,X_1,X_2,X_3$.
Thus the total is
\[
 6+41\cdot5=211.
\]
Replace each such point by the double of a halfpoint $(a,b)$ on
the fixed curve $E_D:b^2=a^3-Da$ over $K=\mathbb Q(\sqrt2)$.
This uses the exact nonidentity $2E_D(K)$ membership statement of the
source.  Retain $a_0,a_1$, where $a=a_0+a_1\sqrt2$, and leave
$b_0,b_1$ as coordinates to be eliminated.  Keep the original
nonzero guards on the doubled point as well as the norm guard on $b$.
All downstream point coordinates are rational functions of these
halfpoint coordinates.

This point substitution can be checked by explicit rational maps,
without a group-membership assertion as a black box.  Write $D=d^2$
for the source curve.  From a halfpoint $(a,b)$ set
\[
 u=\frac{2b}{a^2+D},\qquad
 v=\frac{a^2-2da-D}{a^2+D},\qquad
 w=\frac{a^2+2da-D}{a^2+D},\qquad h=vw.
\]
Then put $f=u^2$, $x=1/u^2$, $y=h/u^3$, and $z=-u/h$.
On the guarded domain these satisfy every old full-chart equation:
\[
 v^2+df=1,\quad w^2+v^2=2,\quad h=vw,\quad
 zh+u=0,\quad xf=1,\quad yz+x=0.
\]
In addition,
\[
 x=\frac{(a^2+D)^2}{4b^2},\qquad
 y=\frac{(a^2+D)(a^4-6Da^2+D^2)}{8b^3},
\]
so $(x,y)=[2](a,b)$ by the duplication formulas.

Conversely, the old full chart forces $u,v,w$ to be nonzero.
With $A=(1+v)(1+w)$ and $S=v+w$, define
\[
 a=\frac{A}{u^2},\qquad b=\frac{SA}{u^3}.
\]
These formulas always give a valid halfpoint.  Indeed
$v^2+w^2=2$ and $du^2=1-v^2$ imply
\[
 A(S^2-A)=-(1-v^2)^2,\qquad
 A^2+(1-v^2)^2=2SA.
\]
The first identity gives $b^2=a^3-d^2a$; the second, together with
\begin{align*}
 A^2-2(1-v^2)A-(1-v^2)^2&=2vSA,\\
 A^2+2(1-v^2)A-(1-v^2)^2&=2wSA,
\end{align*}
shows that the first set of formulas recovers the same $u,v,w$.
None of the new denominators
vanishes.  In fact $1+v=0$ or $1+w=0$ would give $du^2=0$; and
$v+w=0$ would give $v^2=w^2=1$, with the same contradiction.
Thus $A,S,b,a^2+D$ are nonzero.  This proves both directions of the
full-chart substitution directly.  The original guards on the output
exclude any remaining exceptional halfpoints in the forward map.

For every fixed geometric value of $a_0,a_1$, the equation
$b^2=a^3-Da$ has at most four geometric solutions for $(b_0,b_1)$.
Indeed, after the splitting
$K\otimes\overline{\mathbb Q}\simeq
\overline{\mathbb Q}\times\overline{\mathbb Q}$, it becomes two
independent equations $b_+^2=c_+$ and $b_-^2=c_-$, each with at most
two roots.  Restricting to the guarded open subset cannot increase
this number.  The retained point coordinates therefore number
\[
 211\cdot2=422.
\]

\subsection{The pole blocks and their zero fibers}

There is one pole block from the height call and one from each of the
41 comparison calls, hence 42 in total.  The source block has input
$b$ and private variables $t,U,e,p,q,s,h$, with $B=17b$, satisfying
\begin{align*}
 U(1+bt^2)&=b(1-bt^2),\\
 p^2&=B(1-e)(1+3e),\\
 h&=e(B-U^2),\\
 Bq^2&=(s^2-Bh)(s^2+3Bh).
\end{align*}
Retain $t,e,s$ from each block.  The factor $1+bt^2$ never vanishes
on this system, even on geometric fibers.  Otherwise $bt^2=-1$ and
the first equation would give $0=2b$, a contradiction.  Thus $U$ is
uniquely determined; $h$ is then uniquely determined, $p$ has at most
two possibilities, and $q$ has at most two possibilities when $b\ne0$.

The exceptional case $b=0$ must be treated.  The source equations then
force
\[
 U=p=h=s=0,
\]
but they do not constrain the private variable $q$.  Replace this
fiber by $q=0$.  This does not change any rational existential projection:
$q$ occurs nowhere outside this block, and replacing it by zero
preserves all its equations when $b=0$.  The same normalization makes
the geometric fiber finite.  It is a constructible condition, given
by the union of the branch $b\ne0$ and the branch $b=0,q=0$.

Consequently these blocks contribute
\[
 42\cdot3=126
\]
retained coordinates and at most four eliminated choices per block.

\subsection{The smaller anchor}

Let $C\ne0$ be the fixed source anchor constant and let $f_0,f_1$
be the two reciprocal coordinates supplied to the anchor.  Put
\[
 F=f_0^2+f_0f_1-f_1^2.
\]
By Proposition~\ref{prop:explicit-semilocal} and
Lemma~\ref{lem:two-semilocal-inputs}, the exact anchor condition is
equivalent to
\begin{equation}\label{eq:finite-anchor}
 u^2+6v^2-2w^2-3F^2+2C^4=0
\end{equation}
with three rational witnesses.  This follows by applying the
semilocal formula to $F/C^2$ and multiplying by $C^4$; multiplication
of each of its three witnesses by $C^2$ is a bijection of $\mathbb Q$.
Retain $v,w$.  Once the point coordinates and their finite choices
are fixed, $F$ is fixed and \eqref{eq:finite-anchor} is monic quadratic
in $u$.  Thus two retained coordinates suffice for this block, with
at most two eliminated geometric choices.  The old two sets of five
retained anchor coordinates have been replaced by these two.

\subsection{The remaining parameters and algebraic choices}

Retain the seven root coordinates
\[
 X,\quad s,\quad z,\quad\theta,\quad e,\quad u_c,\quad v_c.
\]
Here $X$ is the ordinary input, $s,z$ are two of the three index
proxies, $\theta$ is the common approximation to $\sqrt2$, $e$ is
the common coordinate denominator, and $u_c,v_c$ are the two common
contact parameters.  The third index proxy, denoted $\tau$ to avoid
confusing it with a retained pole variable, satisfies the monic equation
\[
 \tau^2-s(s-2)\tau z+(s-2)^2z^2-1=0.
\]
It therefore contributes at most two geometric choices and no retained
coordinate.  The eight height-class bits satisfy $\epsilon^2=\epsilon$
and contribute only finitely many choices.  Every remaining arithmetic
gate is determined by its inputs; every inverse is unique on its
retained guarded open subset.  Fixed constant constructions contribute
no free coordinates.

One nonvanishing point in this claim deserves explicit attention.
The compressed reciprocal-logarithm equations are
$\Lambda p_0=G_0$, $A_i=C_ip_i$, $A_i\Lambda=G_i$.
Their coefficients $G_i=G_J(Df_i^2)$ do not vanish for $f_i\in K$.
Here
\[
 G_J(W)=1+\sum_{j=1}^J a_jW^j,\quad
 a_1=-2/5,\quad a_j=a_{j-1}\frac{4j-5}{4j+1}\ (j\geq2),
\]
with $J\geq2$ in every application.  The extension of the two-adic
valuation to $K$ is
\[
 v(a+b\sqrt2)=\min\{v_2(a),v_2(b)+\tfrac12\}.
\]
The fixed $D$ has valuation zero and all $a_j$ have valuation one.
If $v(f)\geq0$, the constant term of $G_J(Df^2)$ is uniquely least
in valuation; if $v(f)<0$, then $v(f)\leq-1/2$ and the last
nonconstant term is uniquely least.  Thus $G_J(Df^2)\ne0$.
The displayed equations force $\Lambda,p_i\ne0$ on rational solutions.
Imposing their rational norm guards before eliminating them therefore
loses no rational solutions.  These guards are still retained on
geometric fibers: the valuation argument is not being asserted over
an algebraic closure.

For clarity, the passage from each compressed point chart back to the
full one is also explicit.  Its equations are
\[
 f_i=u_i^2,\quad v_i^2+df_i=1,\quad
 p_i^2=x_i-Df_i,\quad x_if_i=1.
\]
The nonzero $p_i$ just proved and the identity
$f_ip_i^2=v_i^2(1+df_i)$ force $v_i\ne0$.
Set $z_i=-1/p_i$, $y_i=x_ip_i$, $h_i=u_ip_i$, and
$w_i=u_ip_i/v_i$.  Substitution gives all full-chart equations.
Conversely the full chart gives $p_i=y_if_i$.
The ratio variables are then determined by
$\Lambda=G_0/p_0$, $A_i=G_i/\Lambda$, and $C_i=A_i/p_i$.
All other difference and reciprocal-coordinate guards of the original
source are retained.

\begin{proposition}\label{prop:finite-base-557}
After the exact point-membership substitution, the exact anchor
replacement, and the harmless zero-fiber normalization of the private
pole variable $q$, the source has a constructible presentation whose
projection to
\[
 422+126+2+7=557
\]
retained rational coordinates has finite geometric fibers.
Its rational existential projection onto $X$ is unchanged.
\end{proposition}
\begin{proof}
The exact equivalences of the substitutions have been established
above.  In a fiber of the retained coordinates, first choose the
ordinate coordinates of the 211 halfpoints, in finitely many ways.
The doubled points and all their downstream rational functions are
then determined wherever their guards hold.  Next choose the index
proxy $\tau$ and the eight bits.  This determines the rational input
of each pole block.  The normalized pole equations have finitely many
solutions for their remaining private variables by the case analysis
above.  Finally the anchor ordinate $u$ has at most two choices.
Additional equations or inequations only discard choices.
All remaining auxiliary coordinates are unique functions of these
choices.  Thus every geometric fiber is finite, including exceptional
and empty fibers.

For an explicit uniform upper bound on the number of choices, the
argument gives
\[
 4^{211}\cdot4^{42}\cdot2\cdot2\cdot2^8=2^{516}.
\]
This is only a bound on the size of a geometric fiber of this
presentation, not on the degree of the final polynomial after
stratification and branch multiplication.
\end{proof}

\begin{corollary}\label{cor:558-variables}
There is an effectively constructible integer-coefficient polynomial
in at most 558 total variables, one of them the input $X$, whose
projected rational zero set equals that of the source system.
Conditional on that source system defining $\mathbb Z$ in $\mathbb Q$,
the polynomial is a Diophantine definition of $\mathbb Z$ using at
most 557 witnesses.
\end{corollary}
\begin{proof}
Apply Theorem~\ref{thm:finite-fibre-one-variable} with $d=557$ to
Proposition~\ref{prop:finite-base-557}.  Its polynomial uses the 557
retained coordinates and one additional coordinate.  Existentially
quantifying all except $X$ gives the stated projection.
\end{proof}

The corollary is an existence and effective-construction bound, with
explicitly proved projection equivalence.  An actual output polynomial
of this size requires completing the finite-algebra elimination and
the exceptional-stratum computation.  No expanded coefficient table,
degree, coefficient height, or claim of minimum variable count is
asserted here.



\subsection{The emitted 558-variable source specification}
\label{subsec:emitted-558-source}

The coefficient compiler starts with the pinned $855$-variable rational
system, whose $297$ residuals consist of $211$ point equations, $84$
pole equations, one anchor equation, and one Lucas equation. Its
independent source audit is retained as a separate layer. Keep the
$557$ base coordinates of the finite-fiber presentation: $422$ point
abscissa coordinates, $126$ pole parameters, two anchor parameters,
and the seven remaining root parameters including the input.

For every point chart put $Y=2C$, write
$a^3-Da=A+B\sqrt2$, and set
\[
 K=(j^2+2)A+4jB,\qquad J=2jA+(j^2+2)B.
\]
The point equation becomes the monic integer polynomial
\[
 Y^4-2KY^2+2J^2=0.
\]
Its derivative is nonzero on every valid point chart. Indeed $Y\ne0$
there, and
\[
 K^2-2J^2=(j^2-2)^2(A^2-2B^2)\ne0.
\]
If the derivative vanished at a nonzero root, then $Y^2=K$ and
the point equation would contradict this identity. The last norm
is nonzero because the reconstructed halfpoint ordinate has
nonzero norm. Thus including all $211$ point derivatives among the
guards does not discard valid rational solutions.

Complete the Lucas equation by putting
\[
 V=2\tau-a,\qquad a=s(s-2)z,\qquad
 \Delta=a^2-4((s-2)^2z^2-1).
\]
Then $V^2-\Delta=0$ and $\tau=(V+a)/2$. In the remaining
residual numerators and guards, substitute $C_i=Y_i/2$ and this
expression for $\tau$. At each addition gate clear to the largest
power of two occurring in its children; at products add the
exponents, and at power gates multiply them. This gives exact
integer polynomial numerators, with recorded certificates
\[
 \widetilde H=2^k H\bigl(C_i=Y_i/2,\tau=(V+a)/2\bigr).
\]
No vanishing denominator is introduced by these numerical factors.

Each of the $85$ other finite witnesses occurs in just its own
quadratic residual and in no denominator guard. Its cleared equation
is $Ly^2+M=0$, where $L,M$ depend only on the base coordinates,
the point roots, and the completed Lucas root. It has the exact
finite-case replacement
\[
 \begin{array}{ll}
 \epsilon=0:& Z=0,\quad M=0,\\[2pt]
 \epsilon=1:& Z^2+LM=0,\quad LZ\ne0.
 \end{array}
\]
For the second case recover $y=Z/L$. For the first recover $y=0$.
Conversely, if $Ly^2+M=0$ and $Ly\ne0$, use the second case with
$Z=Ly$. Otherwise $M=0$, and the first case is possible. In
particular a fiber with $L=M=0$ is handled by choosing the private
witness zero, without changing any other equation. This includes
the exceptional zero pole input. Apply the same replacement to
the completed Lucas equation with $L=1$, $M=-\Delta$.

There are consequently $297$ monic generators: $211$ quartics
and $86$ linear-or-quadratic choices. For each fixed mask, their
monic triangular algebra has rank
\[
 n=4^{211}2^{\operatorname{popcount}(\epsilon)}\le2^{508}.
\]
The guard is the original guard product, multiplied by the point
derivatives and by $LZ$ in every nonzero quadratic case. It makes
all the diagonal derivatives of the monic presentation nonzero.
Every guarded geometric point is therefore simple. The only added
closed conditions are the $86$ equations
$(1-\epsilon_i)M_i=0$; in a nonzero case the corresponding equation
is identically zero. This verifies all hypotheses of
Theorem~\ref{thm:explicit-guarded-finite-algebra}, including all
zero and repeated-root cases.

Apply that theorem for each of the $212$ fixed point charts, each
of the $256$ fixed height corrections, and every one of the
$2^{86}$ masks. Their finite product is the emitted polynomial
in $558$ external variables. The $297$ formal algebra generators
are eliminated by the specified matrix determinants. They are not
extra witnesses. The input projection agrees with the pinned
rational system; its identification with $\mathbb Z$ retains the
arithmetic hypothesis of that system.

The artifact contains $101637$ integer coefficient-DAG nodes and
the complete finite determinant/coefficient/product recipe. The
independent checker imports no producer. It verifies all $211$
point identities, every private quadratic extraction, the Lucas
completion, the power-of-two substitution identities, monicity,
triangularity, coordinate inventories, and the full table of
degree and matrix-norm majorants. A separate exact small-instance
interpreter executes $177$ selector/degree factors across eight
examples, including nilpotent components, excluded double roots,
bad separators, nonrational residue fields, and zero coefficients.
It also produces a fully expanded degree-$16$ example. The enormous
source multiplication matrices and final branch product are not
materialized, and no such computation is claimed.


\section{Removing a fixed size condition without adding variables}
\label{sec:finite-anchor}

The following elementary observation is useful when the number of variables
has priority over both the degree and the size of the coefficients.  It
replaces a condition on the index of a point in a cyclic group by a finite
product of polynomials.  The length of that product can be extremely large.

\begin{proposition}[Excluding finitely many indices]\label{prop:finite-anchor}
Let $K=\mathbb Q(\sqrt2)$, let $E/K$ be an elliptic curve in a Weierstrass
equation, and let $P\in E(K)$ have infinite order.  Fix an integer $M\geq1$.
There is an explicitly computable polynomial
$A_M\in\mathbb Z[X_0,X_1]$ such that, for every nonzero integer $s$,
\[
 A_M(x_0([s]P),x_1([s]P))\ne0
 \quad\Longleftrightarrow\quad |s|>M.
\]
Here $x(Q)=x_0(Q)+\sqrt2\,x_1(Q)$.  The polynomial has degree $2M$ and
requires no additional existential variables.
\end{proposition}
\begin{proof}
For $1\leq n\leq M$, compute integers $a_n,b_n,D_n$ with $D_n>0$ and
\[
 x([n]P)=(a_n+b_n\sqrt2)/D_n.
\]
Such a representation is obtained by the usual exact group law and clearing
the two rational denominators.  Define
\[
 A_M(X_0,X_1)=\prod_{n=1}^M
 \bigl((D_nX_0-a_n)^2-2(D_nX_1-b_n)^2\bigr).
\]
Each factor is $D_n^2$ times the field norm of
$X_0+\sqrt2 X_1-x([n]P)$.  Since the norm of an element of a field vanishes
only at zero, a factor vanishes precisely when the two abscissas agree.
On a Weierstrass elliptic curve, two points with the same abscissa are
negatives of each other.  Thus $x([s]P)=x([n]P)$ holds precisely when
$[s]P=[n]P$ or $[s]P=[-n]P$.  The infinite order of $P$ gives $s=n$ or
$s=-n$.  The leading homogeneous part of the displayed product is nonzero,
so its degree is $2M$.
\end{proof}

An inequality $A_M\ne0$ is a guard, rather than an extra equation with its
own inverse variable.  In a construction which already has a common
denominator guard, one inserts $A_M$ into that product.  In the finite
algebra construction below, one inserts it into the element whose nonzero
components are selected.  In either case it costs no new coordinate.

\begin{corollary}[Conditional replacement in the repository construction]
\label{cor:556-anchor}
Assume the arithmetic soundness argument of the repository, including its
height estimate with a constant $H<2^{4783227}$, and its completeness with
arbitrarily large admissible indices $s$.  Replace the packed semilocal
anchor by the guard $A_M(x(S))\ne0$, where
\[
                    M=2^{9566463}.
\]
The resulting construction still defines $\mathbb Z$.  Applying the
finite algebra construction to it uses $556$ total variables, including the
input.
\end{corollary}
\begin{proof}
The numerical choice satisfies
\[
 (20H)^2<400\,2^{9566454}<2^{9566463}=M.
\]
Proposition~\ref{prop:finite-anchor} therefore supplies the strict size
condition $|s|>(20H)^2$ used in the soundness argument.  Completeness is
preserved by choosing the admissible index larger than this fixed $M$.
The former packed anchor used three coordinates: two retained coordinates
and one coordinate finite over them.  Removing them lowers the number of
retained coordinates from $557$ to $555$.  The finite algebra construction
adds one coordinate to the retained base, giving $556$ total coordinates.
\end{proof}

This corollary is an arithmetic modification of the defining system.
Unlike elimination of finite fibres, it is not asserted to preserve the
projection of an arbitrary system independently of the source arithmetic
theorem.  It also does not make the enormous finite product practical to
expand.  Its integer coefficients are specified exactly by a terminating
indexed construction.  If the contacts or height theorem are changed, the
constant $M$ must be chosen again from the new height estimate; the
numerical value above must not be reused without that check.



% Includable lemma and proof. Independently checked; no expansion of the
% fixed exclusion product is claimed.
\section{Coefficient bounds for a fixed elliptic exclusion product}
\label{sec:fixed-anchor-height}

The finite size condition in the shared-point construction uses fixed
multiples of a fixed elliptic point.  The number of factors is extremely
large, but their coefficients admit elementary uniform bounds.  This
section gives an exact algorithm defining those coefficients and proves
bounds without computing the full product.

Put $k=\mathbb Q(\sqrt2)$.  Every $x\in k$ has a unique expression
\[
 x=\frac{a+b\sqrt2}{c},\qquad
 a,b,c\in\mathbb Z,\quad c>0,\quad \gcd(a,b,c)=1.
\]
We call $(a,b,c)$ its primitive common-denominator triple and set
\[
 h_0(x)=\left\lceil\log_2\max(1,|a|,|b|,c)\right\rceil.
\]
Thus the triple of zero is $(0,0,1)$ and $h_0(0)=0$.  This elementary
coefficient height is used only in this section; it is not the canonical
elliptic height or the absolute logarithmic Weil height.

\begin{lemma}\label{lem:fixed-triple-height}
For $x,y\in k$ one has
\[
 \begin{aligned}
 h_0(x+y)&\le h_0(x)+h_0(y)+1,\\
 h_0(xy)&\le h_0(x)+h_0(y)+2,
 \end{aligned}
\]
and, for $x\ne0$,
\[
                         h_0(x^{-1})\le2h_0(x)+2.
\]
\end{lemma}
\begin{proof}
Write the primitive triples of $x,y$ as $(a,b,c)$ and
$(a',b',c')$.  Before cancellation their sum has triple
\[
 (ac'+a'c,\,bc'+b'c,\,cc')
\]
and their product has triple
\[
 (aa'+2bb',\,ab'+a'b,\,cc').
\]
The coordinates are bounded respectively by
$2^{h_0(x)+h_0(y)+1}$ and $3\cdot2^{h_0(x)+h_0(y)}$.
Passing to a primitive triple can only decrease their maximum.  For
the inverse use
\[
                  x^{-1}=\frac{ca-cb\sqrt2}{a^2-2b^2}.
\]
The denominator is nonzero because $x\ne0$.  Its absolute value is
at most $3\cdot2^{2h_0(x)}$, and the numerator coefficients are at
most $2^{2h_0(x)}$.  Changing all signs if necessary and cancelling
the common divisor proves the last bound.
\end{proof}

Consider the fixed elliptic curve and point
\[
 \begin{gathered}
 E:\quad y^2=x^3-D_E x,\qquad D_E=11-6\sqrt2,\\
 P=\left(\frac{57-28\sqrt2}{8},
          \frac{-660+417\sqrt2}{32}\right).
 \end{gathered}
\]
Direct substitution verifies $P\in E(k)$.  For a finite point
$R=(x,y)$ put $H_0(R)=\max(h_0(x),h_0(y))$, and set
$H_0(\mathcal O)=0$ for the identity.  Here $h_0(D_E)=4$,
$h_0(x(P))=6$, $h_0(y(P))=10$, and $H_0(P)=10$.

\begin{lemma}\label{lem:fixed-point-binary-height}
For every integer $n\ge1$, the point $nP$ is defined by exact
arithmetic in $k$ and
\[
                         H_0(nP)\le17n^8.
\]
Consequently, if $nP\ne\mathcal O$ and
$x(nP)=(a_n+b_n\sqrt2)/c_n$ is written in its primitive
common-denominator triple, then
\[
               \max(|a_n|,|b_n|,c_n)\le2^{17n^8}.
\]
No nontorsion assumption is needed for these inequalities.
\end{lemma}
\begin{proof}
We first establish bounds for doubling and addition of the fixed point.
Let $R=(x,y)$ be finite and $H=H_0(R)$.  If $y\ne0$, the usual
doubling formulas are
\[
 m=\frac{3x^2-D_E}{2y},\qquad
 x_2=m^2-2x,\qquad y_2=m(x-x_2)-y.
\]
Lemma~\ref{lem:fixed-triple-height}, together with
$h_0(2)=1$ and $h_0(3)=2$, gives successively
\[
 \begin{array}{c|c}
 \text{quantity}&\text{bound for its coefficient height}\\ \hline
 3x^2-D_E&2H+11\\
 (2y)^{-1}&2H+8\\
 m&4H+21\\
 x_2&9H+48\\
 x-x_2&10H+49\\
 y_2&15H+73.
 \end{array}
\]
Thus $H_0(2R)\le15H+80$.  When $R=\mathcal O$ or $y=0$,
the result is the identity, so the same bound holds using its
specified height zero.

For $R\ne\mathcal O,P,-P$, write $P=(x_P,y_P)$ and use
\[
 m=\frac{y-y_P}{x-x_P},\qquad
 x_+=m^2-x-x_P,\qquad y_+=m(x-x_+)-y.
\]
The denominator is nonzero: points of a short Weierstrass curve with
the same abscissa are opposite.  Using the uniform bounds
$h_0(x_P),h_0(y_P)\le10$ gives
\[
 \begin{array}{c|c}
 \text{quantity}&\text{bound for its coefficient height}\\ \hline
 y-y_P,\ x-x_P&H+11\\
 (x-x_P)^{-1}&2H+24\\
 m&3H+37\\
 x_+&7H+88\\
 x-x_+&8H+89\\
 y_+&12H+129.
 \end{array}
\]
It follows that $H_0(R+P)\le12H+129$.  The exceptional cases
obey this bound as well.  If $R=\mathcal O$ the output has height
10, and if $R=-P$ it has height zero.  If $R=P$, the doubling
estimate above with $H_0(P)=10$ bounds $H_0(2P)$ by 223,
which is smaller than $12\cdot10+129$.

To compute $nP$, begin with the leading binary digit of $n$, so the
initial point is $P$.  Read each subsequent digit from left to right,
doubling and then, for a digit one, adding $P$.  Every step is a
well-defined group operation; use the identity and the exceptional
cases just specified instead of a formula with zero denominator.
One binary step replaces a height bound $H$ by at most
\[
                     12(15H+80)+129=180H+1089.
\]
The same upper bound covers a digit zero.  Therefore the next height
$H'$ satisfies $H'+7\le180(H+7)$.  There are
$r=\lfloor\log_2 n\rfloor$ digits after the leading digit, and hence
\[
 H_0(nP)+7\le17\cdot180^r
       \le17\cdot(2^r)^8\le17n^8,
\]
since $180<2^8$.  This proves the asserted weaker bound
$H_0(nP)\le17n^8$ and the bound for the primitive triple.
\end{proof}

\subsection{A bound in a finite algebra}

Let $R=\mathbb Z[X_1,\ldots,X_d]$ and let $A$ be a commutative
$R$-algebra free of positive finite rank, with a specified basis.  For
$a\in A$ let $M_a$ be its multiplication matrix.  If
$\|f\|_1$ denotes the sum of the absolute values of the coefficients
of a polynomial $f\in R$, define
\[
              \|M\|=\max_j\sum_i\|M_{ij}\|_1.
\]
This norm satisfies the triangle inequality, is submultiplicative,
and gives $\|M_c\|=|c|$ for an integer scalar $c$.
The same arguments apply over rational coefficients with the
corresponding coefficient norm.

Suppose $N_0,N_1,D_0\in A$ and, for a real number $b\ge0$,
\[
        \|M_{N_0}\|,\ \|M_{N_1}\|,\ \|M_{D_0}\|\le2^b.
\]
For each $n\ge1$ with $nP\ne\mathcal O$, define
\[
 \mathcal E_n=(c_nN_0-a_nD_0)^2
                       -2(c_nN_1-b_nD_0)^2\in A;
\]
when $nP=\mathcal O$, define $\mathcal E_n=1$.  For an integer
$M\ge1$ put $\mathcal E_{\le M}=\prod_{n=1}^M\mathcal E_n$.

\begin{proposition}\label{prop:fixed-exclusion-norm}
The multiplication matrices of these factors satisfy
\[
 \begin{aligned}
  \|M_{\mathcal E_n}\|&\le2^{34n^8+2b+4},\\
  \|M_{\mathcal E_{\le M}}\|&\le2^{(2b+38)M^9}.
 \end{aligned}
\]
The fixed coefficients in the product are specified by the binary
algorithm of Lemma~\ref{lem:fixed-point-binary-height}; they are not
additional quantified variables.
\end{proposition}
\begin{proof}
For a finite $nP$, let $h_n=h_0(x(nP))$.  The triangle inequality
gives
\[
 \|M_{c_nN_0-a_nD_0}\|,\ \|M_{c_nN_1-b_nD_0}\|
       \le2^{h_n+b+1}.
\]
Submultiplicativity and the factor's two terms now give
\[
 \|M_{\mathcal E_n}\|
     \le3\cdot2^{2h_n+2b+2}
     \le2^{34n^8+2b+4}.
\]
The identity factor when $nP=\mathcal O$ obeys the same bound.
Multiplying and using $\sum_{n=1}^M n^8\le M^9$ and $M\le M^9$
gives
\[
 \|M_{\mathcal E_{\le M}}\|
 \le2^{34\sum_{n=1}^M n^8+(2b+4)M}
 \le2^{(2b+38)M^9},
\]
as claimed.  This reasoning does not require expansion of the product.
\end{proof}

At a rational specialization for which $D_0\ne0$, the value
$\mathcal E_n$ for a finite $nP$ is zero precisely when
\[
               \frac{N_0+N_1\sqrt2}{D_0}=x(nP).
\]
Indeed $u^2-2v^2=0$ with $u,v\in\mathbb Q$ forces $u=v=0$.
If the left side is the abscissa of a finite point $Q\in E(k)$,
the equality is equivalent to $Q=\pm nP$.  Thus the product,
together with the existing denominator guard, excludes exactly these
finitely many abscissas.  The convention for $nP=\mathcal O$ makes
the indexed algorithm total and introduces no unwanted exclusion.


% Working input fragment. This is a new arithmetic construction, not an exact
% equivalence of the complete old witness space. No compilation claim is made.
\section{A shared-point construction}
\label{sec:new-shared-source}

\paragraph{Scope.}
The construction below assumes the source manuscript's cyclic-group,
five-contact height, scalar-height, and pole-completeness results, and also
Theorem~4.1 of Sun, \emph{86 unknowns for Hilbert's tenth problem over the
rationals}, arXiv:2610.20691v2. The last theorem is an additional arithmetic
input. The number of variables is not inferred from Sun's undecidability
theorem. All algebraic operations below have rational coefficients; field
operations in $K=\mathbb Q(\sqrt2)$ mean their two rational components.

\subsection{Fixed data and the complete rational system}
Use
\[
 d=3-\sqrt2,\quad D=d^2=11-6\sqrt2,\quad
 E:y^2=x^3-Dx,
\]
\[
 P=(57/8-7\sqrt2/2,-165/8+417\sqrt2/32),\qquad
 2E(K)=\mathbb ZP.
\]
For a nonzero point put $f(Q)=1/x(Q)$ and $z(Q)=-x(Q)/y(Q)$.
Write $\sigma$ for conjugation. Every displayed quotient is accompanied
by the inequation that its denominator, or its field norm, is nonzero.
The old point charts are replaced by doubles of affine halfpoints.

The ten point occurrences are
\[
 A,S,T,Z,J_0,H,U_A,U_S,U_T,U_Z.
\]
Every occurrence is nonidentity and lies in $2E(K)$. The points $A,S,T,Z$
have integer labels $(n,a,t,b)$, respectively; these labels are used only
in the proof and are not additional rational coordinates.

Set $\beta=X^2+2$. Use rational proxies $\xi,\tau,\zeta$ with
\[
 \tau^2-\xi(\xi-2)\tau\zeta+(\xi-2)^2\zeta^2-1=0.
\]
Only $\xi,\zeta$ are retained; $\tau$ is a finite algebraic coordinate.
Let $Q_0=P$, $(Q_1,Q_2,Q_3,Q_4)=(A,S,T,Z)$ and
$(U_0,U_1,U_2,U_3,U_4)=(H,U_A,U_S,U_T,U_Z)$.
Form $V_i=U_i-[4]Q_i$ for $0\le i\le4$. These five differences are
required to be nonidentity. Affine subtraction also retains the guard
$x(U_i)\ne x([4]Q_i)$; completeness will choose a common index different
from $0,4,-4$.

For $e\in\{16,64\}$ put
\[
 L_e(Q)=z(Q)G_{e/4-1}(Df(Q)^2),\quad
 G_m(W)=1+\sum_{j=1}^m c_jW^j,
\]
where $c_1=-2/5$ and $c_j=c_{j-1}(4j-5)/(4j+1)$.
The source proves these logarithms nonzero at every required nonidentity
point. For $k=(16,64,64,64)$ and $b=(\beta,\xi,\tau,\zeta)$ define
\[
 L_i=\frac{L_{k_i}(U_i)}{L_{64}(H)},\qquad
 L'_i=\sigma\!\left(\frac{L_{k_i}(V_i)}{L_{64}(V_0)}\right).
\]
The eighteen field tests $D_j$, with orders $m_j$, are
\[
 \begin{array}{c|c}
 f(U_0),\ldots,f(U_4),\sigma f(V_0),\ldots,\sigma f(V_4)&1\ (10\text{ entries})\\
 L_i-L'_i\quad(1\le i\le4)& k_i\\
 L_i-b_i\quad(1\le i\le4)&k_i.
 \end{array}
\]
There is no approximate square root of two and no common comparison
point $J$ in this construction.

Enumerate the eight rational components of $x(A),x(S),x(T),x(Z)$ as
$x_1,\ldots,x_8$. On finite branches set $A_0=1$ and, for $1\le r\le8$,
set $A_r=x_r^2$ if $x_r\ne0$, and $A_r=1$ if $x_r=0$. These are
expressions, not new variables. The common slope is one rational
coordinate $\lambda$. Put
\[
 \Theta_0=\{0,-1/2,-1,1/2399,-32/407\},\qquad
 F_5(T)=\prod_{\theta\in\Theta_0}(T-\theta),\quad t_r=A_r\lambda.
\]
Require $t_r\notin\Theta_0\cup\{1\}$, and define, for $0\le r\le8$,
\[
 R_r=A_r^{-5}\frac{F_5(t_r)}{(t_r-1)^6},\qquad
 \widetilde D_{r,j}=\frac{D_j}{(t_r-1)^{m_j+1}},
\]
\[
 W_r=\prod_{j=1}^{18}\left(1+R_r
       \left(\frac{\widetilde D_{r,j}}{R_r^{m_j}}\right)^2\right),
 \quad T_r=\operatorname{Tr}_{K/\mathbb Q}(W_r^{-1}),
 \quad B_r=R_r/T_r^2.
\]
Require $W_r,T_r\ne0$ and the following pole condition for each $B_r$.

Retain two rational coordinates $e_r,s_r$, set
\[
 C_r=17B_r,\qquad h_r=e_r(C_r-B_r^2),
\]
and impose the two monic equations
\[
 p_r^2=C_r(1-e_r)(1+3e_r),\qquad
 q_r^2=\frac{(s_r^2-C_rh_r)(s_r^2+3C_rh_r)}{C_r}.
\]
The roots $p_r,q_r$ are finite algebraic coordinates. This is precisely
the old pole predicate with preparation variable zero. No new finite
multiplier choice is needed. The prime-pattern theorem will be invoked
with its odd squarefree avoidance modulus divisible by $3\cdot17$.
It then makes each selected prime square in $\mathbb Q_2,\mathbb Q_3$
and $\mathbb Q_{17}$. Thus the fixed multiplier $17$ works throughout.

For use in the scalar-height block alone, retain the finite set
\[
 \mathcal C=\left\{\pm\prod_{p\in I}p:
 I\subseteq\{2,3,5,7,11,13,17\}\right\}.
\]

\paragraph{The unchanged scalar-height block.}
Put
\[
 r_\beta=\beta(\beta-1)(\beta-2401)(32\beta+343)(2\beta-1).
\]
Take the six rational coefficients of
$f(J_0),f(S),z(S)/z(J_0)$, with orders $(1,1,1,1,48,48)$.
Use the source's fixed binary norm construction $N_r(a,b)=a^2+rb^2$
on the constant leaf $1$ and the six normalized tests, and denote its
root by $W$. On a finite branch choose $c\in\mathcal C$ and set
$b_*=cr_\beta W^2$, $B_*=17b_*$. Retain $t_*,e_*,s_*$ and define
\[
 U_* =b_*\frac{1-b_*t_*^2}{1+b_*t_*^2},\qquad
 h_*=e_*(B_*-U_*^2).
\]
Require $b_*\ne0$, $1+b_*t_*^2\ne0$, and the two monic equations
\[
 p_*^2=B_*(1-e_*)(1+3e_*),\qquad
 q_*^2=\frac{(s_*^2-B_*h_*)(s_*^2+3B_*h_*)}{B_*}.
\]
The extra guard $b_*\ne0$ removes no required integral-input witnesses:
the source constructs $W$ as a unit at a prime dividing the nonzero
integer $r_\beta$, which has absolute value greater than one.

\paragraph{A fixed finite size condition.}
Let $H$ be the original five-contact constant, with
$1\le H<2^{4783227}$, and put $M=2^{9566463}$. Require
\[
 \mathcal A_M(S):=\prod_{j=1}^M
 N_{K/\mathbb Q}\bigl(x(S)-x([j]P)\bigr)\ne0.
\]
This is an indexed finite expression with fixed rational coefficients.
It forces $|a|>M>(20H)^2$. It has not been expanded. An alternative
source replaces this guard by the proved three-witness semilocal anchor;
that alternative adds two retained coordinates and one quadratic root.

\subsection{The finite algebra and the variable count}
For a point occurrence retain $a_0,a_1$ and set $a=a_0+a_1\sqrt2$,
$a^3-Da=A+B\sqrt2$. Introduce a root $C$ of the monic quartic
\[
 C^4-AC^2+B^2/2=0,
\]
require $C\ne0$ and $A^2-2B^2\ne0$, and put
\[
 b=C+\frac{B}{2C}\sqrt2.
\]
The quartic gives $b^2=a^3-Da$. Use $[2](a,b)$ for the point,
retaining all necessary norm guards. This single chart suffices for
the specific curve used here; there is no finite chart parameter.

\begin{lemma}[A single ordinate chart over a quadratic field]
Let $F$ have characteristic different from two, let
$K=F(\sqrt\delta)$ be a quadratic field, and let $d\in K^\times$
satisfy $d^3\notin F\cup\sqrt\delta F$. On
$E_d:y^2=x^3-d^2x$, every nonidentity $Q\in2E_d(K)$ has a
halfpoint whose ordinate has nonzero $F$-component. In particular
this applies to our curve. More generally it applies over a real
quadratic field whenever $d=u+v\sqrt\delta$ with $uv\ne0$.
\end{lemma}
\begin{proof}
Start with any affine halfpoint, so $b\ne0$ and $a\ne0$.
Translation by $T_0=(0,0)$ gives
\[
 (a,b)+T_0=(-d^2/a,d^2b/a^2)
\]
and does not change its double. If both ordinates had $F$-component
zero, then
$b$ and $d^2b/a^2$ would be nonzero $F$-multiples of $\sqrt\delta$,
so $d^2/a^2\in F$. Thus $w=a/d$ has square in $F$ and belongs to
$F$ or $\sqrt\delta F$. But $b^2=d^3w(w^2-1)\in F^\times$.
In the first case this forces $d^3\in F$, and in the second it
forces $d^3\in\sqrt\delta F$, both contradictions. The nonzero
factor $w(w^2-1)$ follows from $b\ne0$.

Over a real quadratic field, the two components of
$d^3$ are $u(u^2+3\delta v^2)$ and $v(3u^2+\delta v^2)$,
both nonzero when $uv\ne0$. For the present curve, explicitly,
\[
 d^3=45-29\sqrt2,\qquad
 d^3\sqrt2=-58+45\sqrt2.
\]
The addition formula has the displayed positive sign: its slope is
$b/a$, its new abscissa is $-d^2/a$, and its new ordinate is
$(b/a)(a+d^2/a)-b=d^2b/a^2$.
\end{proof}

The quartic derivative is nonzero at the permitted roots. Indeed,
with $Y=C^2$, the equation is $Y^2-AY+B^2/2=0$, of discriminant
$A^2-2B^2\ne0$, and $C\ne0$. Thus its derivative may be recorded
explicitly as a guard without excluding any permitted root.

The retained coordinates number
\[
 1+20+2+1+18+3=45.
\]
The remaining coordinates are ten quartic roots, eighteen pole roots,
two scalar-height roots and the Lucas root. Their monic triangular
degree product is $4^{10}2^{21}=2^{41}$. Vanishing square radicands
are treated by the linear relation that their root is zero; on each
nonzero branch its root is guarded nonzero. The finite-fibre theorem
therefore supplies one additional rational coordinate, for a total of
46. Finite choices and their strata are combined by products, not by
new existential selectors.

\subsection{Local tests and the fixed preparation-zero predicate}

\begin{lemma}[A restricted completeness domain with preparation zero]
The preparation-zero pole just specified excludes negative odd
valuations outside $\{2,3,17\}$. It has rational witnesses whenever
$b$ has square reduced denominator, is positive and square in
$\mathbb Q_2,\mathbb Q_3,\mathbb Q_{17}$, and satisfies
\[
 b<17,\qquad v_2(b)\ge3,\quad v_3(b)\ge1,\quad v_{17}(b)\ge2.
\]
This assertion uses only the pole-completeness arithmetic input already
used by the source.
\end{lemma}
\begin{proof}
The equations are the old equations with preparation variable zero and
$U=b$, so the old one-sided valuation obstruction applies. For
completeness put $B=17b$ and $V=B(B-b^2)=17^2b^2(1-b/17)$.
The stated inequalities make $V$ positive and square at the three
specified local fields. They are exactly the sufficient preparation
conditions in the source's Lemma
\texttt{original-seven-parity-completeness}. Its remaining argument
constructs rational points in square class $B$ on the two twists and
then rescales them to the displayed equations. No further preparation
variable is used. Equivalently one may repeat that argument with
$U=b$: at a support prime of the square class of $B$ outside the three
exceptions, $v_p(b)$ is positive odd, so $1-b/17$ is a square unit;
at a support prime of $V$ the valuation of $B$ is even and cancellation
forces $B/b^2$ to be a square unit. Thus the same disjoint-support and
mutual-squareness hypotheses for the simultaneous two-twist construction
hold. Its conclusions supply the required witnesses.
\end{proof}

\begin{lemma}[One place for all local tests]
Let $R\ne0$ be rational, let $D_j\in K$, and let $m_j\ge1$.
Set
\[
 W=\prod_j(1+R(D_j/R^{m_j})^2),\quad
 T=\operatorname{Tr}_{K/\mathbb Q}(W^{-1}),\quad b=R/T^2,
\]
with the displayed nonzero guards and a valid preparation-zero pole
certificate for $b$. If $p>76176$ and $v_p(R)$ is odd, then it is
positive and some place $\mathfrak p\mid p$ satisfies
$v_{\mathfrak p}(D_j)\ge m_j$ for every $j$.
\end{lemma}
\begin{proof}
Put $\alpha=v_p(R)$. All valuations in $K_{\mathfrak p}$ are integral,
because $p$ is unramified. Each nonzero nonconstant term of a factor
has odd valuation
$b_{j,\mathfrak p}=\alpha+2(v_{\mathfrak p}(D_j)-m_j\alpha)$,
so its factor has valuation $\min(0,b_{j,\mathfrak p})$. Put
$\lambda_{\mathfrak p}=v_{\mathfrak p}(W)\le0$ and
$\lambda=\max_{\mathfrak p\mid p}\lambda_{\mathfrak p}\le0$.
Integrality of trace gives $v_p(T)\ge-\lambda$. Therefore
\[
 v_p(b)=\alpha-2v_p(T)\le\alpha+2\lambda.
\]
The left side is odd and cannot be negative by the pole obstruction.
It is consequently positive, so $\alpha>0$ and
$\lambda>-\alpha/2$. At a place attaining $\lambda$, each nonpositive
summand of $v(W)$ is at least $\lambda$. For each nonzero $D_j$ this gives
\[
 \alpha+2(v(D_j)-m_j\alpha)>-\alpha/2,
 \qquad v(D_j)>(m_j-3/4)\alpha\ge m_j-3/4.
\]
Integral valuations yield $v(D_j)\ge m_j$. Zero terms meet the same
conclusion automatically. This is the reciprocal-trace argument of
Sun's Lemma~3.1, written here for one shared group.
\end{proof}

\begin{lemma}[Two anchors compare the four integer labels]
At a prime $p>76176$, suppose the eighteen field tests have their stated
orders at a place $\mathfrak p\mid p$. Then every proxy $b_i$ is
$p$-integral and
\[
 b_i\equiv n_i\pmod{p^{k_i}},\qquad
 (n_1,n_2,n_3,n_4)=(n,a,t,b).
\]
This statement does not assume that $p$ splits or that the label of
$H$ is a $p$-unit.
\end{lemma}
\begin{proof}
Write $H=[h]P$, $U_i=[u_i]P$. The first ten tests place $H,U_i$ in
the formal group at $\mathfrak p$, and $V_0,V_i$ in the formal group
at its conjugate place. The guarded denominators ensure $h\ne0,4$.
The source's formal logarithm identities give the exact ratios
\[
 r_i=u_i/h,\qquad r'_i=(u_i-4n_i)/(h-4).
\]
The reciprocal truncations have the same valuations as these ratios
and agree with them to relative precision $k_i$. At least one of $h$
and $h-4$ is a $p$-unit; its corresponding exact ratio is integral.
The test $v(L_i-L'_i)\ge k_i$ then makes the other ratio integral too.
Consequently the relative precision becomes absolute precision for
both ratios. The remaining test gives
$b_i\equiv r_i\equiv r'_i\pmod{p^{k_i}}$. Multiplying by the integer
denominators and subtracting yields
$4(b_i-n_i)\equiv0\pmod{p^{k_i}}$. Since $p\ne2$, the assertion follows.
If $p$ is inert, conjugation acts within the same local field and the
same proof applies. No splitting assertion is needed for soundness.
\end{proof}

\subsection{Height transfer with the original numerical constant}
Let $M_0(t)$ denote the odd-contact radical for the five new contacts,
with primes at most $76176$ omitted. The map $\phi(t)=2+1/t$ carries
these contacts to the original five contacts. Its matrix and inverse
are integral with determinant $-1$, and the coordinate forms differ
by fixed rational factors. The prime divisors of the new nonzero
coefficients and contact differences, including differences from $1$,
are
\[
 2,3,5,7,11,37,109,439,2399.
\]
Thus no new exceptional prime is introduced. For $t$ avoiding the
five contacts,
\[
 M_0(t)=M_5(\phi(t)),\qquad
 h(t)\le h(\phi(t))+\log3
       \le(H+\log3)M_0(t)^{24}.
\]

At a prime in $M_0(t_r)$, $t_r-1$ is a unit and exactly one factor
of $F_5(t_r)$ has positive odd valuation. Because $A_r^{-5}$ is a
rational square, $v_p(R_r)$ is odd even when $A_r$ is not a unit.
The preceding trace lemma applied to $\widetilde D_{r,j}$ therefore
gives the original, unscaled tests to their stated precision. Hence
the two-anchor lemma applies at every prime in every $M_0(t_r)$.

Suppose the $K$th powers of all nine radicals divide a nonzero integer
$N$. Then
\[
 h(\lambda),h(A_r\lambda)
 \le(H+\log3)|N|^{24/K}.
\]
For $x_r\ne0$, the identity $A_r=x_r^2$ gives
\[
 2h(x_r)=h(A_r)\le h(\lambda)+h(A_r\lambda).
\]
The same inequality for $h(x_r)$ is immediate at $x_r=0$. Finally,
\[
 h(x(Q_i))\le h(x_{2i-1})+h(x_{2i})+\tfrac32\log2
 <7H|N|^{24/K},
\]
because $H\ge1$, $|N|\ge1$, and $2\log3+\tfrac32\log2<5$.
The original constant $7H$ has therefore been preserved, rather than
replaced by an uncontrolled constant.

\subsection{Soundness of the definition}
Take a rational solution. Write $a_0=|a|$ and
$B=\max(|n|,|a|,|t|,|b|)$. The finite exclusion gives
$a_0>M>(20H)^2$. The unchanged scalar-height argument gives
\[
 h(\beta)\le H\sqrt{a_0}.
\]
For completeness about its use, its constant leaf is $1$ on the
all-left path of the binary norm tree, and the other six leaves are
the stated tests divided by $r_\beta$ to their assigned powers.
At every odd five-contact of $\beta$, the old norm and pole argument
gives order 48 for the parameter ratio, hence divisibility of the
integer label $a$ by the 48th power of that radical. The old
five-contact height theorem gives exactly the displayed estimate.

Set
\[
 \Delta=t^2-a(a-2)tb+(a-2)^2b^2-1\in\mathbb Z.
\]
If $\Delta\ne0$, the proxy equation and the precision-64 congruences
show that each of the nine radicals has 64th power dividing $\Delta$.
The source index-height inequality $h(x([m]P))\ge m^2-6$ and the
previous subsection give
\[
 B^2-6<7H|\Delta|^{3/8}<14HB^{3/2},
\]
where $|\Delta|<5B^4$ since $B\ge a_0>7$. Hence
$B<(14H+6)^2\le(20H)^2$, contradicting $B\ge a_0$.
Thus $\Delta=0$. The source's elementary Lucas-growth lemma, applied
to the nonzero label $b$, now gives
\[
 \log B>a_0-3.
\]

If $\beta\ne n$, let $N$ be the absolute numerator of $\beta-n$ in
lowest terms. The precision-16 comparisons show that every radical
has 16th power dividing $N$. As $N\le2e^{h(\beta)}B$, one obtains
\[
 B^2<7HN^{3/2}+6,
 \qquad \log B<2\log(32H)+3h(\beta).
\]
The bounds $a_0>(20H)^2$ and $H\ge1$ imply both
$h(\beta)<a_0/20$ and $\log(32H)<a_0/20$. Consequently
$\log B<a_0/4<a_0-3$, a contradiction. Therefore $\beta=n\in\mathbb Z$.
Since $\beta=X^2+2$ and $X\in\mathbb Q$, reduction of the denominator
of $X$ gives $X\in\mathbb Z$.

\subsection{Completeness and compatible simultaneous orders}

Fix an integer $X$ and put $n=X^2+2$. Use the source's scalar-height
completeness construction: choose $J_0$ in sufficiently deep formal
subgroups at the finitely many relevant primes, and then choose a
positive label $a$ sufficiently divisible relative to $J_0$. There
are arbitrarily large such labels, so require $a>M$. The source's
Lucas lemma supplies positive integers $t,b$ satisfying $\Delta=0$.
Set the proxies to $(n,a,t,b)$. This fixes the four root points, their
eight abscissa components and all positive square scales $A_r$.
The scalar norm root is a unit at every prime dividing $r_\beta$;
as $|r_\beta|>1$, it is nonzero. Thus the additional scalar guard
is satisfied and the original scalar pole witnesses exist.

Temporarily prescribe $H=[h]P$ and $U_i=[hn_i]P$. Each comparison
test is now a rational function of the point $[h]P$, obtained using
fixed multiplications and translations. At $h=64$ all its denominators
are nonzero: the relevant integer multiples are nonzero, $64\ne\pm4$,
and the source reciprocal logarithms never vanish there. Continuity
in the finitely many groups $E(K_{\mathfrak p})$, for
$\mathfrak p\mid2\cdot3\cdot17$, gives a positive integer $L$,
divisible by 64, and integers $c_p\ge0$ such that
\[
 h\equiv64\pmod L
 \quad\Longrightarrow\quad
 v_{\mathfrak p}(D_j)\ge-c_p
 \quad(\mathfrak p\mid p,\ p=2,3,17).
\]
One obtains $L$ by taking a common multiple of the indices of small
open formal subgroups contained in the required neighborhoods.

There is also a real bound $C$, independent of all future progression
moduli, attainable by arbitrarily large positive $h$ in every
nonconstant progression. Here is the compact-group argument. The
closure $G$ of $\{([h]P,[h]P^\sigma):h\in\mathbb Z\}$ in the
product of the real elliptic groups is a compact group with finitely
many components. Each component projects onto the identity circle,
up to translation, in each real elliptic factor: an infinite cyclic
projection is infinite, and its connected closure is the circle.
The finitely many poles of the required rational functions therefore
remove a set of Haar measure zero from each component. Choose a
neighborhood avoiding them in each component and bound all tests
there by one $C$. Multiplication by any positive integer maps $G^0$
onto itself, so the closure of any nonconstant progression contains
one component. Its positive tail is dense in that closure, proving
the assertion. This is the argument in Sun's Lemma~8.1; in particular
the bound is chosen before the prime pattern and its order moduli.

Choose a positive integer $D_0$ for which $D_0^2\theta/A_r$ is integral
for every contact $\theta$ and every $r$, and put
\[
 v=D_0^2w^2,\qquad \lambda=u/v.
\]
The contact positions are $u-(\theta/A_r)v$, a finite collection of
forms $u+c_jw^2$ with integer $c_j$. Delete repeated positions;
within each fixed $r$ its five positions are distinct. The zero
contact gives the common position $u$. Enlarge the prescribed lower
bound for the primes beyond $76176$ and beyond all prime factors of
the fixed coefficients and scales.

Invoke Sun's Theorem~4.1 with arbitrary prescribed divisibility of $w$,
arbitrarily large $u/w^2$, and the odd squarefree avoidance modulus
containing every odd prime factor of $3\cdot17L$ other than seven.
It supplies positive prime values at all positions, Gaussian
generators $\pi\equiv9\pmod{16}$, $\pi\equiv1+3i\pmod7$, and
$\pi\equiv i$ modulo that avoidance modulus. Its two reduction orders
are $d_+=N(\pi+1)$ and $d_-=N(\pi-1)$, with
\[
 v_2(d_+)=2,\qquad v_2(d_-)=6,\qquad
 \gcd(d_-(p),d_+(q))=4
\]
for all selected primes $p,q$. Both orders are prime to the odd part
of $L$: the prime seven follows separately from
$d_+\equiv6$, $d_-\equiv2\pmod7$.

Put $U=\operatorname{lcm}d_-$ and $V=\operatorname{lcm}d_+$.
The three congruences
\[
 h\equiv0\pmod U,\qquad h\equiv4\pmod V,\qquad
 h\equiv64\pmod L
\]
are compatible. Their pairwise common odd factors have been removed,
and the two-parts are compatible because
$v_2(U)=6$, $v_2(V)=2$, $64\mid L$, and $4\mid60$.
The generalized Chinese remainder theorem supplies a nonconstant
progression. Choose a sufficiently large positive $h$ in it attaining
the real bound $C$. This avoids $0,4,-4$ and every affine-chart
exception. At a selected prime choose the place with reduction order
$d_-$. Then $H,U_i$ are formal there and $V_0,V_i$ are formal at the
conjugate place of order $d_+$. Their exact two logarithm ratios equal
$n_i$, so all eighteen tests have their required local orders.

\subsection{Completeness of the nine pole occurrences}
The determinant of the contact coordinate change is a unit outside
the fixed exceptions. More explicitly, for each $r$,
\[
 R_r=\frac{\prod_{\theta\in\Theta_0}(u-(\theta/A_r)v)}
             {v^5(t_r-1)^6}.
\]
The denominator is a rational square, since $v$ is a square.
At its five distinct contact primes $R_r$ has valuation one, and
at every other prime it has even valuation. The contact primes do
not divide $v$; otherwise two distinct prime positions would share
that prime. They do not divide the denominator of a fixed scale or
a fixed difference, by the chosen lower bound. Thus $t_r-1$ is a
unit at these primes.

At either place above a contact prime, each factor of $W_r$ has
valuation zero or negative. Its inverse has residue one or zero,
respectively. At the successful place all factors have residue one.
Therefore $T_r$ has residue one or two and is a unit there. It follows
that $B_r=R_r/T_r^2$ has positive odd valuation at its contact primes
and even valuation at all other primes. In particular its reduced
denominator is a square.

The imposed Gaussian congruences give each contact prime
$p\equiv1\pmod8$ and $p\equiv1\pmod{3\cdot17}$. Such a prime is
square in each of $\mathbb Q_2,\mathbb Q_3,\mathbb Q_{17}$.
The displayed product formula for $R_r$, and then division by $T_r^2$,
shows that $B_r$ has the same local square properties. This is the
reason a single fixed multiplier $17$ suffices for all nine blocks.

It remains to arrange their quantitative preparation-zero conditions.
The real bound $C$ was uniform in the later congruence moduli, so for
$\lambda\to+\infty$ the fixed scales give, uniformly in the chosen $h$,
\[
 R_r\sim A_r^{-6}\lambda^{-1},\quad
 W_r=1+O(\lambda^{-3}),\quad T_r=2+O(\lambda^{-3}),\quad
 B_r\sim(4A_r^6\lambda)^{-1}.
\]
For sufficiently large $\lambda$ all factors at both real embeddings
are positive and $0<B_r<17$. At a small prime put
$\nu_p=v_p(v)=-v_p(\lambda)$ and $\eta_{p,r}=v_p(A_r)$.
For sufficiently large $\nu_p$, direct valuation calculation gives
\[
 v_p(R_r)=\nu_p-6\eta_{p,r},\qquad
 v_{\mathfrak p}(\widetilde D_{r,j})
 \ge-c_p+(m_j+1)(\nu_p-\eta_{p,r}).
\]
The nonconstant term of its factor consequently has valuation at least
\[
 3\nu_p+(10m_j-8)\eta_{p,r}-2c_p.
\]
These bounds tend to infinity with $\nu_p$. Thus $W_r$ and its inverse
are arbitrarily close to one at each place above $p$. Integrality of
trace, also at the ramified prime two, gives
$v_p(T_r)=v_p(2)$ once the error is sufficiently small. Hence
\[
 v_p(B_r)=\nu_p-6\eta_{p,r}-2v_p(2).
\]
Arbitrarily high prescribed divisibility of $w$ makes these valuations
meet $v_2(B_r)\ge3$, $v_3(B_r)\ge1$, $v_{17}(B_r)\ge2$ for every $r$.
All such denominator requirements and the real lower bound were chosen
before the application of the prime-pattern theorem. The restricted
pole-completeness lemma now supplies all nine structural pairs and
their square roots. This completes the rational solution for every
integral input.

\begin{theorem}[Conditional forty-six-variable construction]
Assume the arithmetic inputs explicitly stated at the beginning of
this section, including Sun's simultaneous prime-pattern theorem.
The rational system specified above has input projection precisely
$\mathbb Z$. It gives, by the finite-algebra construction, an
effectively specified polynomial over $\mathbb Z$ with at most
46 total variables, including its ordinary rational input. With the
three-witness semilocal anchor in place of the indexed finite exclusion,
the same conclusion holds with at most 48 total variables.
\end{theorem}
\begin{proof}
Soundness and completeness were proved in the preceding subsections.
The point substitution is exact by the single-chart lemma. On each
finite branch all denominators are guarded, all non-base choices lie
in the stated finite triangular algebra, and its zero-radicand and
repeated-root specializations have the prescribed finite treatment.
There are 45 retained coordinates and one final coordinate, or two
additional retained anchor coordinates in the semilocal version.
Clearing fixed rational coefficients gives integer coefficients.
The finite anchor is a specified indexed polynomial, not a claim that
its astronomically many factors have been expanded.
\end{proof}

\subsection{Preserving the scalar certificate when its point is shared}
\label{subsec:independent-height-sharing}

We verify that the scalar-height point $J_0$ can be identified with the
common comparison point $H$.  This is an arithmetic modification of the
system, with the same hypotheses as the shared-point construction.  The
scalar pole witnesses must be chosen after the final common index has
been selected.

\begin{proposition}
Under the stated source arithmetic inputs and Sun's simultaneous
prime-pattern theorem, imposing $J_0=H$ preserves the input projection
of the shared-point construction.  Consequently the indexed-exclusion
version uses at most $44$ total coordinates after finite-algebra
compression, and the semilocal-anchor version uses at most $46$.
\end{proposition}
\begin{proof}
Soundness follows from the construction with separate points by assigning
the same point to both occurrences.  For completeness fix an integer input,
write $r=r_\beta$, and choose the signed squarefree correction $c$ so that
$cr>0$ and $cr$ is square in $\mathbb Q_2,\mathbb Q_3,\mathbb Q_{17}$.
The number $cr$ is a nonzero integer.  Let $\mathcal S$ be its finite prime
support.

Choose a positive integer $h_0$ divisible by $64$ such that
$J_*= [h_0]P$ is sufficiently deep in the formal subgroups at all places
above $\mathcal S$.  Next choose the positive label $a$ of $S$ sufficiently
divisible relative to $h_0$, and sufficiently large for the size anchor.
The original scalar-completeness argument permits these successive choices.
It gives the following explicit conditions on the six scalar tests,
normalized by their assigned powers of $r$: every normalized test is
$p$-integral for $p\mid r$, and every nonconstant normalized test is divisible
by $p$ for $p\mid c$ with $p\nmid r$.  The denominators used in those tests
are nonzero.  The source Lucas construction now fixes the two other root
labels and all root points, and all four proxies are their integer labels.

These scalar conditions persist when $J_*$ is replaced by $[h]P$ in a
sufficiently small neighborhood at every place above $\mathcal S$.
Indeed each test is a rational function defined at $J_*$; the coefficient
maps from $K$ to its two rational components are continuous in the product
of the completions above $p$.  An integral lower valuation bound defines
an open condition at a point satisfying it.  Simultaneously impose the
uniform lower bounds for the eighteen comparison tests at the places
above $2,3,17$, using continuity at the nonzero indices $h_0$ and $h_0-4$.
Their denominators are nonzero by the reciprocal-logarithm lemma and
$64\mid h_0$, $h_0>0$.

The resulting finitely many local neighborhoods are preserved by one
integer progression
\[
                       h\equiv h_0\pmod L,
                       \qquad 64\mid L.
\]
To see this, put a sufficiently small open formal subgroup inside each
translated neighborhood, take a common multiple of its finite group
indices, and further multiply by $64$.

Use the compact real-group argument to obtain bounds for the comparison
tests uniform in every future progression modulus.  Choose the fixed
denominator divisibilities and real lower bound for the common slope as
in the separate-point construction.  In the simultaneous prime theorem,
include all odd prime factors of $3\cdot17L$, except seven, in its avoidance
modulus.  Its Gaussian class modulo seven already excludes seven from
both reduction orders.  With $U=\operatorname{lcm}d_-$ and
$V=\operatorname{lcm}d_+$, the hypotheses therefore give
\[
 \gcd(U,V)=4,\qquad \gcd(U,L)=64,\qquad \gcd(V,L)=4.
\]
The three congruences
\[
 h\equiv0\pmod U,\qquad h\equiv4\pmod V,
 \qquad h\equiv h_0\pmod L
\]
are compatible: $64\mid h_0$ and $4\mid h_0-4$ supply the two remaining
conditions.  Choose a sufficiently large positive $h$ in this progression
attaining the uniform real bounds, and put $J_0=H=[h]P$.
The comparison witnesses and the nine preparation-zero pole blocks are
then constructed exactly as before.

It remains to choose the scalar pole witnesses for this final $h$.
Let $W$ be its scalar norm root.  At a prime dividing $r$, the normalized
leaves are integral and the constant leaf $1$ on the all-left path makes
$W$ a unit.  At a prime dividing $c$ but not $r$, the nonconstant leaves
are divisible by that prime, so the same tree again has unit root.
These conditions hold because the progression preserved them.  Since
$|r|>1$, at least one such prime exists and $W\ne0$.

Thus $b_*=crW^2$ is positive.  At primes in $\mathcal S$ its valuation is
the nonnegative valuation of the integer $cr$; at all other primes its
valuation is even.  Its reduced denominator is therefore a square.
Its square classes at $2,3,17$ equal those of $cr$, so it is locally square
at all three places.  The original scalar pole-completeness theorem now
supplies its preparation and curve witnesses.  No real bound for $W$,
and no new rational coordinate, was required.

Sharing removes one point occurrence, hence two retained abscissa
coordinates and one quartic algebra generator.  The retained count becomes
$43$ (or $45$ with the semilocal anchor); the finite-algebra construction
adds one coordinate.  This proves the stated counts.
\end{proof}


% Additional arithmetic refinement. Counts require the finite-algebra compiler.
\subsection{Removing the scalar preparation coordinate}
\label{subsec:scalar-preparation-zero}

The scalar-height certificate can also use preparation zero, after a
different normalization of its contact parameter and a larger fixed set
of squareclass corrections.  This is a change of arithmetic construction;
it is not an identity on the old full witness space.  The following
argument uses the same explicitly stated arithmetic inputs as the
shared-point definition.

Put
\[
 \delta(T)=T^2+T+23,\qquad \epsilon=\frac{204}{23^{12}},
 \qquad R_\beta=\epsilon^2\frac{r_\beta}{\delta(\beta)^4},
\]
where $\beta=X^2+2$ and $r_\beta$ is the previous scalar contact
polynomial.  Its nonvanishing and $\delta(\beta)>0$ hold for every
rational input.  Define the two fixed prime sets
\[
 \mathcal B=\{5,19,23,349,661,5209\},\qquad
 \mathcal L=\{2,3,11,17,37,67,71\},
\]
and the finite correction set
\[
 \mathcal C'=\left\{\pm\prod_{p\in I}p:
                   I\subseteq\mathcal B\cup\mathcal L\right\}.
\]
It has $2^{14}$ elements.  All these primes belong to the original
exceptional set $S_*$; the change introduces no new exceptional prime
in the height theorem.

Use the same six rational scalar tests and their same orders
$(1,1,1,1,48,48)$, with $J_0=H$.  In their fixed binary norm tree
replace $r_\beta$ by $R_\beta$, including in every leaf denominator.
Thus the six leaves are divided by the assigned powers of $R_\beta$;
the constant leaf remains $1$ on the all-left path.  Denote the root
by $W$, and require $W\ne0$.  Choose $c\in\mathcal C'$ on a finite
branch, and put
\[
 b=cR_\beta W^2,\qquad B=17b,\qquad h=e(B-b^2).
\]
Keep just the two structural coordinates $e,s$, together with the
two finite square roots
\[
 p^2=B(1-e)(1+3e),\qquad
 q^2=\frac{(s^2-Bh)(s^2+3Bh)}B.
\]
The old preparation coordinate $t_*$ and its rational expression
$U_*$ are absent.  The guard $b\ne0$ is retained.

\begin{lemma}[The fixed quadratic and the local correction basis]
\label{lem:scalar-fixed-quadratic}
The polynomial $\delta$ has no root modulo any prime of $\mathcal L$.
The squareclasses of the seven primes in $\mathcal L$ form a basis
of
\[
 \mathbb Q_2^*/\mathbb Q_2^{*2}\ \times\
 \mathbb Q_3^*/\mathbb Q_3^{*2}\ \times\
 \mathbb Q_{17}^*/\mathbb Q_{17}^{*2}.
\]
If an integer $\beta$ makes a prime divide both $r_\beta$ and
$\delta(\beta)$, that prime belongs to $\mathcal B$.  At any finite
scalar contact of a rational $\beta$ outside $S_*$, $\delta(\beta)$
is a unit.
\end{lemma}
\begin{proof}
Modulo two, $\delta=T^2+T+1$.  At the odd primes of $\mathcal L$
its discriminant $-91$ is a quadratic nonresidue.  These are finite
residue computations.  For clarity the local squareclass vectors
are displayed below.  The columns are the valuation at two, the
two unit coordinates for the basis $(-1,5)$ in $\mathbb Z_2^*$,
the valuation and nonsquare-unit coordinate at three, and the
valuation and nonsquare-unit coordinate at seventeen:
\[
\begin{array}{c|rrrrrrr}
2 &1&0&0&0&1&0&0\\
3 &0&1&1&1&0&0&1\\
11&0&1&1&0&1&0&1\\
17&0&0&0&0&1&1&0\\
37&0&0&1&0&0&0&1\\
67&0&1&1&0&0&0&0\\
71&0&1&0&0&1&0&1
\end{array}
\]
Their matrix has rank seven over $\mathbb F_2$, proving the basis
claim by the usual local squareclass descriptions.

The relevant evaluations of the quadratic are
\[
 \begin{split}
 \delta(0)&=23,\qquad \delta(1)=5^2,\\
 \delta(2401)&=5^2\cdot349\cdot661,\\
 32^2\delta(-343/32)&=5^2\cdot5209,\\
 4\delta(1/2)&=5\cdot19.
 \end{split}
\]
For an odd prime dividing the scalar contact polynomial, reduction
therefore gives exactly the asserted list of possible common primes.
The prime two cannot divide $\delta(\beta)$ for integer $\beta$.
The same evaluation argument shows that at a rational finite
contact outside $S_*$ the quadratic is a unit: the contact value
and its fixed denominators are integral there, and its displayed
evaluation is a unit.  This proves both remaining statements.
\end{proof}

\begin{lemma}[Preservation of scalar soundness]
\label{lem:normalized-scalar-soundness}
The modified scalar certificate implies the same estimate
\[
                      h(\beta)\le H\sqrt{|a|}
\]
as the previous scalar-height block.
\end{lemma}
\begin{proof}
First, $\beta=X^2+2$ has no odd contact with infinity at any finite
prime.  Indeed, if $v_p(\beta)<0$, then $v_p(X)<0$ and
$v_p(\beta)=2v_p(X)$ is even.  It therefore suffices to consider
the finite contacts counted by the source's radical $M_5(\beta)$.

At each such odd contact outside $S_*$, the preceding lemma makes
$\delta(\beta)$ a unit.  The fixed factor $\epsilon$ and every
$c\in\mathcal C'$ are also units there.  Consequently
$v_p(R_\beta)=v_p(r_\beta)$ is positive and odd.  The source's
no-cancellation argument for $N_R(u,v)=u^2+Rv^2$ now applies
unchanged, with $R=R_\beta$.  The pole condition on $cRW^2$ forces
the six original rational tests to their assigned orders.  In
particular it gives the order-$48$ divisibility of the scalar index
$a$ at every prime in $M_5(\beta)$.  Thus
$M_5(\beta)^{48}\mid a$, and the original five-contact height
theorem gives the stated estimate, with its original constant $H$.
No test has been divided by an extra, possibly nonunit, normalizer:
only the contact parameter of the norm construction has changed.
\end{proof}

\begin{lemma}[Completeness with preparation zero]
\label{lem:normalized-scalar-completeness}
For every integral input the shared-point construction can satisfy
the modified scalar block together with all of its comparison,
height, and size conditions.
\end{lemma}
\begin{proof}
Fix an integer input, so $\beta\ge2$ is an integer.  First choose
the subset of $\mathcal B$ in $c$ so that
\[
               v_p(cr_\beta)\equiv0\pmod2\quad(p\in\mathcal B).
\]
The factors $\epsilon^2$ and $\delta(\beta)^4$ do not change these
parities.  Choose the sign of $c$ so that $cr_\beta>0$.  Finally
the independent basis $\mathcal L$ supplies a subset making
$cr_\beta$, and hence $cR_\beta$, square in
$\mathbb Q_2,\mathbb Q_3,\mathbb Q_{17}$.  The two prime sets
are disjoint, so this last choice preserves all prescribed parities
at $\mathcal B$.

The reduced denominator of $cR_\beta$ is a square.  At the primes
of $\mathcal B$ its valuation is even.  Away from $\mathcal B$,
if $p\mid\delta(\beta)$, then $p$ divides neither $r_\beta$ nor
$c$: this follows from the resultant calculation and the absence
of roots modulo primes of $\mathcal L$.  Its valuation is then a
multiple of four.  At the remaining primes the valuation is
nonnegative, since the only fixed denominator of $\epsilon$ is
supported at $23\in\mathcal B$.

At two, three, and seventeen the quadratic is a unit and
the factor $204^2=(2^2\cdot3\cdot17)^2$ gives
\[
 v_2(cR_\beta)\ge4,\qquad
 v_3(cR_\beta)\ge2,\qquad
 v_{17}(cR_\beta)\ge2.
\]
The largest absolute value of a correction is
\[
 C_{\max}=\prod_{p\in\mathcal B\cup\mathcal L}p
          =518513917206770054130.
\]
For real $\beta\ge2$ one has
\[
 \frac{|r_\beta|}{\delta(\beta)^4}
 <\frac{1202\cdot204\cdot2}{\beta^3}
 \le61302<2^{16}.
\]
The exact integer inequality
\[
              2^{16}C_{\max}204^2<23^{24}
\]
therefore gives $0<cR_\beta<1$.  All numerical and finite-field
computations in this paragraph are reproduced by the companion
exact checker.

It remains to choose the common point and the anchor point so
that their norm root $W$ is a unit at the finite support of
$cR_\beta$, and $|W-1|<1/2$ in the real absolute value.  The local
condition is the same formal-group construction as before: take
$H_*=[h_0]P$ sufficiently deep at those finitely many places,
then take $S$ sufficiently deeper relative to $H_*$.  Every
nonconstant normalized leaf can consequently have an arbitrarily
large positive valuation.  Even when $v_p(R_\beta)<0$, this makes
each nonconstant summand of every norm node divisible by $p$;
the constant leaf on the all-left path then gives $W\equiv1\pmod p$.

We give the simultaneous real choice explicitly, since it is needed
to remove the preparation coordinate.  Let $G$ be the closure of
the cyclic subgroup generated by the pair of real embeddings of
$P$, let $G^0$ be its identity component, and let $m$ be the order
of $G/G^0$.  The closure of any subgroup $k\mathbb ZP$ contains
$G^0$; if $m\mid k$, it is exactly $G^0$.  The projection of
$G^0$ to either real elliptic curve is nontrivial, since that
projection of $P$ has infinite order.  Thus there are small
nonempty open sets in $G^0$ near its identity which avoid the
identity in both projections.

Choose such a relatively compact open set $\Omega$, so close to
the identity that the two rational coefficients of $f(H)$ are
as small as needed for all $H\in\Omega$.  Shrink it to have
$z(H)$ bounded away from zero at both real embeddings.  Choose
$h_0$ positive, divisible by $64m$, satisfying the local formal
conditions and with $[h_0]P\in\Omega$.  Such a choice follows
from the density of the required finite-index subgroup in $G^0$.
Choose the positive index of $S$ in the required finite-index
subgroup, satisfying its anchor and local conditions, and with
$S$ sufficiently close to the real identity.  Then all six scalar
tests, including the two coefficients of $z(S)/z(H)$, are as small
as prescribed, uniformly for $H\in\Omega$.  At zero values of
all six tests the fixed norm root is exactly one.  Its continuity
therefore gives $|W-1|<1/2$, and in particular $0<W<2$, throughout
this real choice.  The real
conditions on $S$ are compatible with any required size lower
bound because a positive orbit in the finite-index subgroup
returns infinitely often to the chosen open set.

After fixing $S$ and the Lucas root indices, prescribe the common
index modulo $L$, where $64m\mid L$, so that all scalar local
conditions persist near $h_0$ and all small-prime comparison
conditions hold.  Inside $\Omega$ remove the pole loci of the
finitely many comparison functions and then choose a nonempty
open set with compact closure avoiding those loci.  Each locus
has empty interior: its projection is a finite preimage of a
point under a nonzero multiplication map on a real elliptic
curve; all the root labels are nonzero by the completeness choice.
The comparison functions are uniformly bounded on the
chosen compact set.

The simultaneous prime-pattern theorem is now applied as in the
previous sharing proof, including every odd factor of $L$ except
seven in its avoidance modulus.  Its two order moduli $U,V$
satisfy
\[
 \gcd(U,V)=4,\quad \gcd(U,L)=64,\quad \gcd(V,L)=4.
\]
Thus $h\equiv0\pmod U$, $h\equiv4\pmod V$, and
$h\equiv h_0\pmod L$ are compatible.  The final progression
consists of multiples of $m$ and has closure $G^0$ in the real
group.  It therefore contains arbitrarily large positive indices
whose real points lie in the chosen compact-open neighborhood.
The complete comparison and detector construction goes through,
and the scalar inequalities and local units are preserved.

For this final common point,
\[
 b=cR_\beta W^2\quad\hbox{satisfies}\quad
 0<b<4<17.
\]
It has square reduced denominator: at every prime in the support
of $cR_\beta$ the norm root is a unit, and at all other primes
its additional valuation is even.  It is square at two, three,
and seventeen and has valuations there at least $4,2,2$,
respectively.  The preparation-zero pole-completeness lemma
therefore supplies $e,s,p,q$.  All added guards hold, and no
preparation coordinate has been used.
\end{proof}

\begin{theorem}[Conditional forty-three-variable construction]
\label{thm:shared43}
Under the same arithmetic inputs as the shared-point construction,
the scalar modification above preserves the integer input projection.
After finite-algebra compression it uses at most $43$ total rational
coordinates with the indexed finite exclusion, or $45$ with the
semilocal anchor.  The corresponding witness counts are $42$ and
$44$.  These are exact construction bounds, not claims that a
coefficient table of the final polynomial has been expanded.
\end{theorem}
\begin{proof}
Soundness changes only the scalar-height input, which was proved
unchanged.  The preceding lemma proves simultaneous completeness.
The scalar preparation coordinate is deleted, while its two square
roots are retained as finite algebra generators.  Thus the retained
base dimension decreases from $43$ to $42$, or from $45$ to $44$
with the semilocal anchor.  The triangular algebra still has rank
at most $2^{39}$, or $2^{40}$, respectively.  The enlarged fixed
correction set is combined by a product over its branches and
adds no rational coordinate.  The finite-algebra theorem supplies
one final coordinate, giving the asserted totals.
\end{proof}


\subsection{Fewer finite choices and smaller algebra rank}
\label{subsec:shared-degree-refinements}

The following two refinements preserve the retained-coordinate counts.
They improve the subsequent degree bound by removing unnecessary
finite choices and finite algebra generators.  They apply in
particular to the forty-three-coordinate construction.

\begin{lemma}[Fixed nonzero height scales]
\label{lem:fixed-height-scales}
Write $x(Q)=x_0+x_1\sqrt2$ for any nonidentity root point of the
shared construction.  One may replace the two piecewise height
scales attached to these components by
\[
                         A_0=x_0^2,
              \qquad A_1=(1+x_1^2)^2.
\]
Both are positive nonzero rational squares.  This removes all
eight zero-component case parameters and their closed equations,
without changing the arithmetic soundness and completeness proofs.
The original point-height bound with constant $7H$ still holds.
\end{lemma}
\begin{proof}
The point $Q$ is the double of a halfpoint $(a,b)$ on
$y^2=x^3-Dx$, with $D=(3-\sqrt2)^2$ and $b\ne0$.
The duplication formula gives
\[
                   x(Q)=\left(\frac{a^2+D}{2b}\right)^2.
\]
At each real embedding the numerator $a^2+D$ is strictly positive
and the denominator is nonzero.  Thus both real embeddings of
$x(Q)$ are positive, and their average $x_0$ is positive.
The assertions about the two scales follow.

Put $T=(H+\log3)|N|^\alpha$, where $\alpha$ is the radical
exponent used in the applicable height comparison.  The common
slope detector and each scaled detector give
\[
               h(\lambda),h(A_i\lambda)\le T,
               \qquad h(A_i)\le2T.
\]
For the first component $h(A_0)=2h(x_0)$, so $h(x_0)\le T$.
For a rational number $x_1=a/b$ in lowest terms with $b>0$,
\[
 h(1+x_1^2)=\log(a^2+b^2)\ge2h(x_1),
 \qquad h(A_1)\ge4h(x_1).
\]
The displayed equality uses
$\gcd(a^2+b^2,b^2)=1$.  Therefore $h(x_1)\le T/2$.  The
standard height inequality for the fixed quadratic basis gives
\[
 h(x(Q))\le h(x_0)+h(x_1)+\tfrac32\log2
 \le\tfrac32(H+\log3)|N|^\alpha+\tfrac32\log2
 <7H|N|^\alpha,
\]
because $H\ge1$ and $|N|\ge1$.  Thus the soundness estimates are
preserved, with room in the same constant.  Every completeness
argument uses only that the scales are fixed positive rational
squares once the root points are chosen; the new choices have
exactly those properties.  All nine detector constructions and
their simultaneous prime patterns remain available.
\end{proof}

\begin{lemma}[Parametrizing the first pole conic]
\label{lem:pole-conic-parametrization}
Let $B\in\mathbb Q^*$ and consider
\[
                        p^2=B(1-e)(1+3e).
\]
Its rational points are covered by two rational formulas with
one parameter.  For a finite choice $\sigma\in\{0,1\}$ retain
one rational coordinate $t$, put $D_t=t^2+3B$, require $D_t\ne0$,
and set
\[
 e=\frac{t^2+(4\sigma-1)B}{D_t},
 \qquad p=-\frac{4Bt(1-\sigma)}{D_t}.
\]
Every permitted assignment satisfies the conic equation, and every
old rational point occurs in one of the two choices.  In particular
the square root coordinate $p$ can be eliminated exactly.
\end{lemma}
\begin{proof}
For $\sigma=0$,
\[
                 1-e=\frac{4B}{D_t},\qquad
                 1+3e=\frac{4t^2}{D_t},
\]
which gives the equation.  For $\sigma=1$ the formulas give
$e=1,p=0$ directly.  Conversely, if an old point has $e\ne1$,
put $t=p/(e-1)$.  Substitution of its conic equation gives
\[
                       t^2+3B=-\frac{4B}{e-1}\ne0.
\]
The formulas with $\sigma=0$ recover both $e$ and $p$.  If
$e=1$, then $p=0$; the choice $\sigma=1,t=0$ gives this point
with denominator $3B\ne0$.  The other point with $p=0$,
namely $e=-1/3$, is already covered by $\sigma=0,t=0$.
Accordingly no nonzero guard on $p$ or on $t$ is imposed.
\end{proof}

Apply this lemma separately to the first conic in each of the
nine coordinate detectors and in the scalar block.  Their
$B$ values are nonzero by the existing guards.  Each block now
retains its conic parameter $t$ and its coordinate $s$, and
keeps only the second square root $q$ as a finite algebra
generator.  The parameter $t$ here parametrizes a conic; it
is distinct from the scalar preparation coordinate that was
removed in the preceding subsection.  The exceptional conic
choices are fixed finite parameters, not quantified selectors.

The preferred finite-anchor source therefore has $42$ retained
coordinates and twenty finite generators: nine quartic point
generators, ten quadratic pole generators, and one quadratic
Lucas generator.  Its algebra rank is at most
\[
                         4^9 2^{11}=2^{29}.
\]
The semilocal source has $44$ retained coordinates and one
additional quadratic anchor generator, hence rank at most
$2^{30}$.  Finite-algebra compression still gives totals
$43$ and $45$, respectively.  The ten conic choices replace
the ten now-absent first-root zero/nonzero strata, whereas
all eight zero-component choices have disappeared.  The
generic determinant compiler must use these updated ranks
and finite branch lists for its final degree bounds.


% Source-to-polynomial conversion for the preferred shared-point systems.
\section{The compiled shared-point polynomials with Lucas proxies}
\label{sec:compiled-shared-polynomials}

We apply the explicit finite-algebra construction to the preferred
sources \texttt{source-43-conic.json.gz} and
\texttt{source-45-conic-semilocal.json.gz}.  This section proves the
algebraic conversion performed by \texttt{compile\_shared\_source.py}.
The arithmetic argument identifying the source projection with
$\mathbb Z$ is a separate input; it uses the pinned manuscript,
Sun's stated theorem, the scalar normalization of
Section~\ref{subsec:scalar-preparation-zero}, and the degree
refinements of Section~\ref{subsec:shared-degree-refinements}.
The compiler preserves the rational projection of each supplied
source, including every exceptional quadratic case.

\subsection{Retained coordinates and rational parametrization}

The finite-anchor source retains 42 rational coordinates: the
input, eighteen abscissa coordinates for nine points over
$\mathbb Q(\sqrt2)$, two Lucas parameters, one common slope,
eighteen parameters for the nine coordinate poles, and two
parameters for the scalar pole.  Its remaining finite choices
consist of nine point roots, one Lucas root, and ten second
pole roots.  There are twenty formal algebra generators.
The semilocal source retains two additional anchor coordinates
and has one additional finite anchor root.

The first pole roots have already been eliminated by
Lemma~\ref{lem:pole-conic-parametrization}.  More explicitly,
for each nonzero $B$ the conic
\[
                     p^2=B(1-e)(1+3e)
\]
is parametrized using a retained coordinate $t$ and a fixed
choice $\sigma\in\{0,1\}$:
\[
 e=\frac{t^2+(4\sigma-1)B}{t^2+3B},\qquad
 p=-\frac{4Bt(1-\sigma)}{t^2+3B},\qquad t^2+3B\ne0.
\]
The choice zero covers all points with $e\ne1$: its inverse is
$t=p/(e-1)$, and the conic equation gives
$t^2+3B=-4B/(e-1)\ne0$.  The choice one covers $e=1,p=0$;
$t=0$ is permitted because $B\ne0$.  The point $e=-1/3,p=0$
is covered by the choice zero with $t=0$.  No nonzero guard on
$p$ or $t$ is imposed.  Thus this operation preserves the rational
input projection while removing one finite quadratic root in
each of the ten pole blocks.

The fixed nonzero square height scales of
Lemma~\ref{lem:fixed-height-scales} are
\[
                       x_0^2,\qquad (1+x_1^2)^2.
\]
They eliminate the eight component-zero finite choices and their
closed conditions.  Consequently each preferred source has
$16384\cdot2^{10}=2^{24}$ original finite parameter choices:
the scalar correction and the ten conic charts.  These are
fixed product indices, not quantified coordinates.

\subsection{Integral triangular monic equations}

Write $a=a_0+a_1\sqrt2$ and
$a^3-(11-6\sqrt2)a=A+B\sqrt2$.  In the fixed point chart, a
halfpoint ordinate is reconstructed as
\[
 b=C+\frac{B}{2C}\sqrt2,\qquad C\ne0,
 \qquad C^4-AC^2+\frac{B^2}{2}=0.
\]
The source variable $Y=2C$ satisfies the integral monic equation
\[
                         Y^4-4AY^2+8B^2=0.
\]
Its derivative $4Y^3-8AY$ is an explicit source guard.  Hence
all retained point roots are simple, including after base change
to an algebraic closure.  Equivalently, the discriminant of the
quadratic in $Y^2$ is $16(A^2-2B^2)$, and the source guards
imply both $Y\ne0$ and $A^2-2B^2\ne0$.

The Lucas equation has the form
\[
                         \tau^2-a\tau+b=0,
\]
where $a,b$ are integer polynomials in retained parameters.
Set $Z=2\tau-a$ and $\Delta=a^2-4b$.  The equation becomes
$Z^2-\Delta=0$, with inverse $\tau=(Z+a)/2$.  This coordinate
change is invertible over $\mathbb Q$, including when
$\Delta=0$.  A repeated Lucas root is therefore not discarded.

Every other finite equation is already monic and quadratic with
no linear term.  Earlier substitutions can introduce powers of
two in its coefficient.  Suppose its current form is
\[
                          y^2+2^{-k}N=0,
\]
where $N$ is an integer polynomial in retained coordinates and
earlier formal generators.  With $s=\lceil k/2\rceil$ and
$Z=2^s y$, it becomes
\[
                          Z^2+2^{2s-k}N=0.
\]
The exponent $2s-k$ is zero or one.  These changes preserve
triangularity and are invertible over $\mathbb Q$.  The producer
records each exact numerator and power of two.  The independent
checker verifies the identities by induction through the
coefficient DAG, without expanding large polynomials.

\subsection{Quadratic exceptions and geometric simplicity}

For every resulting equation $Z_i^2+c_i=0$, take a fixed binary
choice between
\[
 \begin{array}{c|c|c}
 \text{choice}&\text{monic algebra relation}&\text{additional condition}\\ \hline
 0&Z_i=0&c_i=0,\\
 1&Z_i^2+c_i=0&Z_i\ne0.
 \end{array}
\]
The condition $c_i=0$ is imposed as a closed residual, not as a
quotient relation, so it does not destroy freeness of the
presented algebra.  The second case contributes $Z_i$ to the
guard.  These cases cover precisely all rational roots.  In the
zero case the diagonal derivative is one; in the nonzero case it
is $2Z_i$, nonzero on the guarded locus.

The compiler forms the original guard from its distinct
structural factors, discarding repeated factors and nonzero
integer constants.  The source rational-square clearings $Z=Dq$,
$Z^2=ND$ for $q^2=N/D$ are checked exactly, including the
guard on $D$.  The closed-condition and indexed-anchor argument
denominators are also checked against the recorded guards.
After the substitutions, the
coefficient numerators are integer polynomials and the discarded
numerical denominators are powers of two.

The full guard is the product of the original factors, the
nonzero-case quadratic generators, and the indexed anchor when
present.  Each point derivative is among the original factors.
Thus every diagonal derivative of the triangular monic
presentation is nonzero on the guarded open subset.  At every
geometric point of that subset the fiber has local length one.
This is the exact hypothesis of
Theorem~\ref{thm:explicit-guarded-finite-algebra}; generic
separability alone is not being used.

There are eleven quadratic choices for the finite-anchor source
and twelve for the semilocal source.  A mask with $k$ nonzero
choices has free rank $4^9 2^k$, with respective maxima $2^{29}$
and $2^{30}$.  All source parameters, quadratic masks, separating
integers, and degree cases in the theorem are literal finite
product indices.  Its formal algebra generators are eliminated
by the determinant and coefficient operations.

\subsection{The fixed elliptic exclusion product}

For the finite-anchor source, the guard also contains a fixed
product excluding small elliptic indices.  The integer
\[
                             M=2^{9566463}
\]
and the fixed point $P$ are part of the specification.  Write the
variable abscissa as $(N_0+N_1\sqrt2)/D$ after exact denominator
conversion.  For each finite $nP$, compute its primitive triple
\[
                        x(nP)=\frac{a_n+b_n\sqrt2}{c_n}.
\]
The exact integer factor is
\[
                  (c_nN_0-a_nD)^2-2(c_nN_1-b_nD)^2.
\]
For an identity multiple the factor is one.  The product runs
from $n=1$ to $M$, and $D$ is already guarded.  Since
$\sqrt2\notin\mathbb Q$, a factor vanishes at a rational
specialization exactly when the abscissa equals $x(nP)$.
The triples are defined by the exact binary multiplication
algorithm of Section~\ref{sec:fixed-anchor-height}.  Thus all
integer coefficients and index bounds are specified, with no
additional quantified variables.  The enormous collection of
individual factors is not materialized.

The semilocal source uses its quadratic anchor instead and has
no indexed elliptic exclusion product.  Its complete coefficient
DAG is materialized.  Both variants use the outer indexed
determinants and finite products of the finite-algebra theorem.

\begin{proposition}\label{prop:shared-source-compilation}
The preferred emitted specifications define integer polynomials
$P_{43}$ and $P_{45}$ in respectively 43 and 45 coordinates,
including the input.  Their rational input projections equal
those of their corresponding supplied sources.  They have 42
and 44 existential witnesses, respectively.  Subject to the
arithmetic characterization of those sources, each is a
Diophantine definition of $\mathbb Z$ in $\mathbb Q$.

With $H(P)$ denoting the largest absolute coefficient, the
verified bounds are
\[
 \begin{array}{c|c|c|c}
 &\text{maximum algebra rank}&\deg P&H(P)\\ \hline
 P_{43}&2^{29}&2^{9566680}&2^{\,2^{86098389}}\\
 P_{45}&2^{30}&2^{243}&2^{\,2^{246}}.
 \end{array}
\]
The degree and height entries are upper bounds, not asserted
exact values.
\end{proposition}
\begin{proof}
The conic parametrization covers every rational source solution,
and all subsequent rational coordinate changes are invertible.
The two quadratic cases cover all roots and impose exactly the
original equations.  Denominator coverage makes the cleared
source equations and guards equivalent to the original ones.
The fixed anchor clearing is equivalent by the norm identity
just proved.  Hence every finite branch has its intended
rational projection.  The diagonal derivative check permits
application of Theorem~\ref{thm:explicit-guarded-finite-algebra}.
Its one added coordinate gives the stated counts, and a finite
product realizes the union of all source choices and masks.

For the bounds, apply the weighted reduction and multiplication
matrix norm estimates of Section~\ref{subsec:finite-algebra-bounds}
with nine quartics and respectively eleven or twelve quadratic
choices.  Proposition~\ref{prop:fixed-exclusion-norm} bounds the
indexed fixed exclusion product.  Every generator weight,
operator norm estimate, determinant and resultant bound, and
finite factor count is recorded in the specification and checked
independently.  They give the displayed rounded estimates.  For
the indexed exclusion product the estimates are stored as exact
powers of two, so the bound computation itself does not require
materializing integers with millions of digits.
\end{proof}

The specifications contain 6224 and 6245 coefficient-DAG nodes.
They are exact compact descriptions using integer polynomial
operations, coefficient extraction, determinants, and explicitly
bounded finite products.  They are not expanded coefficient
tables or circuits listing every gate of rank-sized matrices.
An independent exact interpreter executes the complete theorem
on small algebras, including nilpotents, root collisions, degree
drops, and zero quadratic cases.  The theorem and source conversion
proof justify the unexpanded large instances.

The other emitted source variants are preserved in separate
packages for independent review of the arithmetic replacements.
All use the same compiler and the same proof of algebraic
conversion.  The package index records their individual counts,
bounds, hashes, and passing verification receipts.


\section{Testing the integer relation without proxy coordinates}
\label{sec:direct-integer-relation}

The first new improvement removes the three rational Lucas proxies
from the shared-point construction.  Two were retained coordinates;
the third was a finite quadratic root.  Their role was to impose
one integer polynomial relation on three elliptic-group labels.
The relation can instead be tested directly on the approximate
logarithmic ratios.  This gives a conditional definition with
$41$ total rational coordinates, including the input, or $43$
with the semilocal anchor.

\subsection{An elementary substitution principle}

\begin{lemma}\label{lem:direct-polynomial-congruence}
Let $L$ be a number field, let $\mathfrak p$ lie over a rational
prime $p$, and normalize its valuation by $v_{\mathfrak p}(p)=1$.
Suppose $z_i\in L$ and $n_i\in\mathbb Z$ satisfy
\[
 v_{\mathfrak p}(z_i)\ge0,
 \qquad v_{\mathfrak p}(z_i-n_i)\ge k
 \quad(1\le i\le r).
\]
For every $F\in\mathbb Z[T_1,\ldots,T_r]$,
\[
 v_{\mathfrak p}\bigl(F(\mathbf z)-F(\mathbf n)\bigr)\ge k.
\]
Consequently a requirement $v_{\mathfrak p}(F(\mathbf z))\ge k$
is equivalent to $p^k\mid F(\mathbf n)$ when $k$ is a positive
integer.  The same congruence assertion holds for rational-coefficient
polynomials whose coefficient denominators are units at $p$; in that
case the final condition is $v_p(F(\mathbf n))\ge k$.
\end{lemma}
\begin{proof}
For each monomial, replace its factors one at a time.  Every term
of the resulting telescoping difference contains some $z_i-n_i$;
all its other factors, including the coefficient, are integral at
$\mathfrak p$.  Thus every term has valuation at least $k$.
The nonarchimedean triangle inequality proves the assertion.
The final equivalence follows because $F(\mathbf n)$ is a rational
integer and the normalized valuation restricts to $v_p$ on it.
\end{proof}

This principle does not make the ratios equal to integer labels.
It transfers a polynomial relation modulo the selected prime powers.
The global height comparison remains the step that forces the
integer residual to vanish exactly.

For a quantitative general version, write
$\|F\|_1$ for the sum of the absolute values of its coefficients.
If $\deg F=d$, $|n_i|\le B$, and $B\ge1$, then
\[
                    |F(\mathbf n)|\le\|F\|_1B^d.
\]
Thus a height argument of the form
\[
 B^2-6<7H|F(\mathbf n)|^{24/k}
\]
forces $F(\mathbf n)=0$ whenever
\[
 k>12d,\qquad
 B>\left((7\|F\|_1^{24/k}+6)H\right)^{1/(2-24d/k)}.
\]
Indeed, if the integer residual were nonzero, the first inequality
and $B,H\ge1$ would imply the reverse strict inequality on $B$.
For several integer relations one can either test them separately,
or use the single relation $\sum_jF_j^2=0$.  The latter doubles
the largest polynomial degree and its precision requirement must
be adjusted accordingly.

\subsection{The replacement in the shared construction}

Retain the nine nonidentity points
\[
              A,S,T,Z,H,U_A,U_S,U_T,U_Z\in2E(K),
\]
where $K=\mathbb Q(\sqrt2)$, and all the point, subtraction,
logarithm and denominator guards of the preferred shared source.
Their root labels are $(n,a,t,b)$ for $(A,S,T,Z)$, respectively.
These are integers used in the proof; they are not rational
coordinates of the polynomial.  Put $\beta=X^2+2$.

Let $L_i,L'_i$, $1\le i\le4$, be the primary and conjugate
alternate logarithmic ratios from the previous construction,
with orders $(16,64,64,64)$.  Delete the rational coordinates
$\xi,\zeta$ and the monic quadratic root $\tau$, together with
the equation involving them.  Define
\[
 F(U,V,W)=V^2-U(U-2)VW+(U-2)^2W^2-1.
\]
Use the following sixteen field tests in every coordinate detector:
\[
\begin{array}{c|c}
 f(U_0),\ldots,f(U_4),\sigma f(V_0),\ldots,\sigma f(V_4)
       &1\quad(10\text{ tests})\\
 L_i-L'_i\quad(1\le i\le4)&16,64,64,64\\
 L_1-\beta&16\\
 F(L_2,L_3,L_4)&64.
\end{array}
\]
Here $U_0=H$, $Q_0=P$, $(Q_1,Q_2,Q_3,Q_4)=(A,S,T,Z)$,
and $V_i=U_i-[4]Q_i$, exactly as before.  There are ten tests
of order one, two of order sixteen, and four of order sixty-four.
Every other part of the preferred source is retained: its fixed
positive square scales, nine reciprocal-trace detectors, normalized
scalar-height block sharing $H$, conic parametrizations, and the
chosen size anchor.

\begin{proposition}[Preservation of the integer input projection]
\label{prop:direct-Lucas-preservation}
Under the explicitly stated arithmetic inputs of the shared
construction, the modified system has input projection exactly
$\mathbb Z$.
\end{proposition}
\begin{proof}
Consider first a rational solution and a prime counted by one
of the nine height radicals.  The detector selects a place
above that prime at which all sixteen tests have their assigned
orders.  The two-anchor argument uses only the first ten tests
and the four differences $L_i-L'_i$.  Hence it still proves
\[
 L_i\text{ is integral},\qquad
 L_i\equiv n_i\pmod{p^{k_i}},
 \qquad(n_1,n_2,n_3,n_4)=(n,a,t,b).
\]
In particular this conclusion does not require any proxy equation.

Lemma~\ref{lem:direct-polynomial-congruence} now gives
\[
 F(L_2,L_3,L_4)\equiv
 \Delta:=t^2-a(a-2)tb+(a-2)^2b^2-1\pmod{p^{64}}.
\]
The new order-$64$ test therefore says exactly that the integer
$\Delta$ is divisible by $p^{64}$.  The radical divisibilities
used in the original height argument are unchanged.  If
$B=\max(|n|,|a|,|t|,|b|)$ and $\Delta\ne0$, that argument
gives
\[
 B^2-6<7H|\Delta|^{3/8},\qquad |\Delta|<5B^4.
\]
The unchanged anchor contradicts these estimates.  Therefore
$\Delta=0$, and the nonzero root label $b$ yields the same
Lucas growth inequality $\log B>|a|-3$.

The normalized scalar block is unchanged and still gives
$h(\beta)\le H\sqrt{|a|}$.  The order-$16$ test $L_1-\beta$
and the corresponding cross comparison give the same divisibility
of the numerator of $\beta-n$.  If $\beta\ne n$, the resulting
bound
\[
                    \log B<2\log(32H)+3h(\beta)
\]
contradicts the anchor and the Lucas growth inequality as before.
Thus $\beta=n\in\mathbb Z$, and $X\in\mathbb Z$.

For completeness, start with an integer input.  Choose the same
integer root labels as in the previous completeness proof, so
$F(a,t,b)=0$.  At every selected place, choose the common index
and points so that $L_i\equiv n_i\pmod{p^{k_i}}$ and the same
holds for the alternate ratios.  The lemma then gives
$v_{\mathfrak p}(F(L_2,L_3,L_4))\ge64$, which is the only
requirement on this new test.  Exact vanishing of $F(L_2,L_3,L_4)$
is not required.

The new test is a fixed polynomial in the old ratios.  Its real
values are bounded on the same compact sets avoiding the old
denominators, and its valuations have finite lower bounds in the
same small-prime neighborhoods.  Consequently the normalization,
prime-pattern, local-square and preparation-zero completeness
arguments apply without alteration.  No new denominator or
exceptional prime is introduced by $F$ itself.  This constructs
all the remaining witnesses for every integral input.
\end{proof}

\subsection{Exact coordinate and finite-algebra counts}

With the finite size exclusion, the retained coordinates are
\[
 \underbrace{1}_{X}+
 \underbrace{18}_{\text{nine halfpoint abscissa pairs}}+
 \underbrace{1}_{\lambda}+
 \underbrace{18}_{\text{nine pole pairs}}+
 \underbrace{2}_{\text{scalar pole pair}}=40.
\]
The finite algebra has nine quartic point generators and ten
quadratic pole generators.  Its rank is at most
\[
                            4^9 2^{10}=2^{28}.
\]
The finite-algebra construction adds one final rational coordinate,
giving $41$ total coordinates and $40$ existential witnesses.
The semilocal anchor adds two retained coordinates and one
quadratic generator, giving rank at most $2^{29}$ and totals
$43$ and $42$, respectively.  The finite correction and conic
choices are unchanged.  The removed Lucas root has no remaining
equation, repeated-root stratum, or denominator guard.

The files \texttt{source-41-direct.json.gz} and
\texttt{source-43-direct-semilocal.json.gz} record the complete
coefficient DAGs and finite parameters.  The copied producer is
\texttt{build\_direct\_source.py}.  The frozen earlier source
files are unchanged.  Final degree and coefficient-height bounds
come from the separately checked finite-algebra compilation of
these new inputs.


\section{Changing the input shift without adding coordinates}
\label{sec:p3-beta-shift}

The substitution $\beta=X^2+2$ is convenient over $\mathbb Q$, but its
interaction with the five scalar contacts is not stable under extension
of the ground field. We first describe a uniform rational modification
and then a field-dependent way to avoid every member of a prescribed
finite exceptional set. Neither statement changes the elliptic
cyclic-group hypothesis.

\begin{proposition}\label{p3:integer-shift}
Let $C$ be an integer with $C\ge2$ such that $2401-C$ is not a
nonnegative integer square. Replace $\beta=X^2+2$ by
$\beta=X^2+C$ throughout the shared-point construction and the direct
Lucas-polynomial refinement. Under their stated arithmetic inputs,
the resulting rational input projection is still $\mathbb Z$.
The retained-coordinate counts and finite-algebra ranks are unchanged.
In particular every integer $C\ge2402$ is permitted.
\end{proposition}
\begin{proof}
Write $X=a/b$ in lowest terms, with $b>0$. Then
\[
                \beta=\frac{a^2+Cb^2}{b^2},
          \qquad \gcd(a^2+Cb^2,b^2)=1.
\]
Thus $\beta$ is an integer if and only if $X$ is an integer.
Moreover, at any finite prime where $\beta$ has negative valuation,
$v_p(X)<0$ and $v_p(\beta)=2v_p(X)$ is even. These are precisely
the input-specific facts used in scalar soundness: the finite
five-contact radical accounts for every odd contact, without a new
contact at infinity. The scalar height argument and the comparisons
which force $\beta$ to equal an integer label therefore remain valid.

For integral $X$, the new value of $\beta$ is an integer at least
two. It cannot equal $0,1,-343/32,$ or $1/2$. It cannot equal
$2401$, by the hypothesis on $C$. Hence the scalar contact product
\[
 r_\beta=\beta(\beta-1)(\beta-2401)(32\beta+343)(2\beta-1)
\]
is nonzero. The scalar normalization, local squareclass choices,
prime-pattern construction, and simultaneous point choices apply
to this integer value of $\beta$ just as before. Their estimates
were uniform for all integral $\beta\ge2$ away from these contacts;
they did not require the special numerical value of the constant
two. In particular the same height constant and precision choices
are valid. This proves completeness.

Changing a fixed coefficient introduces no coordinate and no new
algebra generator. It also does not change the formal degrees of the
source expressions. The independently computed upper bounds for the
emitted coefficient formulas must nevertheless be recorded separately,
since their numerical coefficient sizes can change.
\end{proof}

If a number field $K$ has a real embedding and $C\ge2402$, then
$X^2+C$ avoids every scalar contact for every $X\in K$. Indeed at
that embedding its value is at least $2402$, whereas every contact
is real and at most $2401$. The normalizing denominator
$\beta^2+\beta+23$ is also nonzero: its real discriminant is
$-91$. Thus $C=2402$ removes these two specific domain obstructions
over every number field with a real embedding. This observation
does not imply that the resulting formula defines $\mathcal O_K$:
the cyclic point group and the arithmetic comparison arguments
require their own field-dependent replacements.

\begin{lemma}[Avoiding a finite set by a translated square]
\label{p3:finite-square-avoidance}
Let $K$ be a number field and let $\Theta\subset K$ be finite.
There is an effectively searchable arithmetic progression of positive
integers $C$ such that
\[
                  \theta-C\notin K^2
                       \qquad(\theta\in\Theta).
\]
In particular $X^2+C$ avoids $\Theta$ for every $X\in K$. The
progression contains integers greater than any prescribed bound.
\end{lemma}
\begin{proof}
For each $\theta\in\Theta$, choose a different odd rational prime
$\ell_\theta$ which splits completely in the normal closure of
$K/\mathbb Q$ and at which $\theta$ is integral. Infinitely many
such primes exist by Chebotarev's theorem
\cite[Section~2.1, Theorem~1]{p2lbSerre1981}. Excluding the finitely
many ramified primes, denominators, and previously used primes
does not affect existence. Choose a prime $\mathfrak p_\theta$
of $K$ above $\ell_\theta$. Its residue field is
$\mathbb F_{\ell_\theta}$.

Choose a nonzero nonsquare
$u_\theta\in\mathbb F_{\ell_\theta}$ and impose the integer
congruence
\[
                C\equiv\overline\theta-u_\theta
                       \pmod{\ell_\theta}.
\]
Every such $C$ makes $\theta-C$ a unit at
$\mathfrak p_\theta$ with nonsquare residue. It cannot be a
square in $K$: a square which is a unit has a unit square root
and square residue. The Chinese remainder theorem imposes all
these congruences simultaneously and supplies the desired
arithmetic progression.

The construction is effective given exact number-field data:
enumerate rational primes, test complete splitting by finite
factorization outside the discriminant, compute residues, and
test squares in the finite residue fields. Chebotarev guarantees
that each required search terminates. No conjectural effective
prime bound is needed for this assertion.
\end{proof}

For the scalar formulas one may take $\Theta$ to consist of the
five rational contacts together with those roots of
$T^2+T+23$ which lie in $K$. Choose $C\ge2402$ in the
resulting progression. Then both $r_\beta$ and the normalizing
denominator are nonzero for every $X\in K$, and the conditional
definition over $\mathbb Q$ remains valid by
Proposition~\ref{p3:integer-shift}. The shift is allowed to depend
on $K$, but remains an ordinary fixed integer coefficient, not
an extra variable. The field-independent choice $2402$ suffices
when a real embedding exists.

No fixed integer shift can remove these contacts over every number
field. For any fixed $C$, the algebraic integer $X=\sqrt{-C}$ in
$\mathbb Q(\sqrt{-C})$ satisfies $X^2+C=0$. The field-dependent
choice in Lemma~\ref{p3:finite-square-avoidance} is therefore
essential if all number fields are to be considered.


\section{The compiled polynomials with forty-one and forty-three coordinates}
\label{sec:p3-compiled-polynomials}

We apply the guarded finite-algebra theorem to the direct
Lucas-polynomial source of Section~\ref{sec:direct-integer-relation}.
All formal generators, monic relations, coefficient conversions,
denominator guards, and finite index ranges are supplied in the
corresponding polynomial specification.

\begin{theorem}\label{thm:p3-polynomial41}
Under the stated source arithmetic inputs and Sun's simultaneous
prime-pattern theorem, the supplied polynomial $P_{41}$ defines
$\mathbb Z$ in $\mathbb Q$ using forty-one total coordinates,
including the input. It has forty existential witnesses and satisfies
\[
        \deg P_{41}\le 2^{9566673},
        \qquad H(P_{41})\le 2^{\,2^{86098382}}.
\]
The semilocal-anchor alternative $P_{43}^{\mathrm{loc}}$ uses
forty-three total coordinates, forty-two existential witnesses, and
satisfies
\[
        \deg P_{43}^{\mathrm{loc}}\le 2^{236},
        \qquad H(P_{43}^{\mathrm{loc}})\le 2^{\,2^{240}}.
\]
Both are explicit indexed integer polynomials. These are upper bounds
for degree and height, not exact expansion statistics.
\end{theorem}
\begin{proof}
The source has forty retained rational coordinates and nineteen
formal algebra generators, nine of degree four and ten of degree
two. Its maximum finite rank is $2^{28}$. The semilocal alternative
adds two retained coordinates and one quadratic generator, giving
maximum rank $2^{29}$. Their rational input projections are
$\mathbb Z$ by Proposition~\ref{prop:direct-Lucas-preservation}.

The source-to-integer conversion is the same exact procedure proved
in Section~\ref{sec:compiled-shared-polynomials}. For every rational
square equation $q^2=N/D$, its denominator is guarded and the
substitution $Z=Dq$ gives $Z^2=ND$. Subsequent powers-of-two
clearings preserve the triangular monic presentation. Every
quadratic generator is split into its zero case and nonzero case,
with the appropriate closed equation or derivative guard. The nine
point quartics retain their derivative guards. Thus the required
geometric simplicity holds at every guarded point, including
exceptional parameter specializations.

Theorem~\ref{thm:explicit-guarded-finite-algebra} eliminates these
formal generators with one additional coordinate. It gives the
stated totals and preserves the rational projection of each source.
The emitted specifications have respectively 6018 and 6039 integer
coefficient-DAG nodes. The independent checker verifies the source
hash, all affine and rational-square changes, guard coverage, monicity,
triangularity, every finite index range, and the recursive degree and
coefficient majorants. The bounds are exactly those displayed above.
The fixed elliptic exclusion product is treated by the canonical
clearing and height estimate already proved for that product.
\end{proof}

The packages are
\nolinkurl{research/phase3/finite_fibre/polynomial-41/} and
\nolinkurl{research/phase3/finite_fibre/polynomial-43-semilocal/}.
Each contains its complete indexed specification, construction
report, and passing independent verification receipt. The large
multiplication matrices and outer products are not expanded.
The semilocal variant has a materialized coefficient DAG; the
forty-one-coordinate variant additionally contains the explicitly
indexed elliptic exclusion product.

The same producer also emits the two versions with input shift
$C=2402$ from Proposition~\ref{p3:integer-shift}. Their retained
counts, maximum ranks, and displayed degree and height bounds
are the same. They have 6019 and 6040 coefficient-DAG nodes,
respectively, and separate passing receipts. The additional
coefficient node is a fixed integer, not a quantified coordinate.
These variants remove the scalar-contact obstruction over fields
with a real embedding; they do not resolve the other arithmetic
obstructions over extensions.



\section{The quadratic compilation with Lucas proxy coordinates}
\label{sec:new-quartic}

The finite-algebra construction minimizes the number of coordinates by
allowing a very high degree.  For a degree-four presentation, auxiliary
coordinates serve a different purpose: they keep each arithmetic relation
quadratic.  We now apply this method to the shared-point arithmetic system,
using the common height point, the fixed square scales, comparison order
$52$, scalar order $32$, and the semilocal size condition.  The scalar pole
in this quartic retains its preparation coordinate.  The separate
normalization which removes that coordinate in the smallest-variable
construction is not used here; its extra arithmetic would cost more
quadratic intermediate coordinates.

% Input fragment: refinement of the shared-point construction.
\subsection{Comparison order fifty-two with an explicit larger anchor}
\label{subsec:shared-precision52}

The high comparison order can be lowered from $64$ to $52$ while
keeping beta order $16$ and scalar-height order $48$.  This does not
change the number of retained coordinates of the shared construction.
It shortens eight reciprocal logarithms, and can therefore reduce the
number of variables in a quartic circuit presentation.  Such a circuit
count is a separate, implementation-dependent assertion.

\begin{proposition}\label{prop:shared-precision52}
Assume the arithmetic inputs of the shared-point construction and
$1\le H<2^{4783227}$.  Replace every order-$64$ comparison by order
$52$, and replace the corresponding logarithms $zG_{15}(Df^2)$ by
$zG_{12}(Df^2)$.  Retain the beta order $16$, the scalar order $48$,
and all other equations and guards.  Either of the following size
conditions is sufficient:
\[
 |a|>(21H)^{13/2},
\]
or the semilocal anchor with fixed constant
\[
 C=3^{2^{26}},\qquad
 u^2+6v^2-2w^2-3(f_0^2+f_0f_1-f_1^2)^2+2C^4=0.
\]
Here $f_0+f_1\sqrt2=f(S)$.  Alternatively the finite exclusion
condition can use $M_{52}=2^{31091008}$.  Each version has the same
integer input projection as the order-$64$ construction under the
stated arithmetic inputs.
\end{proposition}

\begin{proof}
We first check that the changed logarithm has the required properties.
The formal identity
\[
 \log_E(z)=z\sum_{j\ge0}c_j(Df^2)^j,
 \qquad c_1=-2/5,\quad
 c_j=c_{j-1}(4j-5)/(4j+1),
\]
is independent of the chosen truncation.  At a place above
$p>76176$, normalized by $v(p)=1$, a nonzero point in the formal
group has $r=v(z)\ge1$ and $v(Df^2)=4r$.  The denominator of $c_j$
divides $(4j+1)!$.  For $j\ge13$ with $4j+1<p$, the $j$th term
therefore has valuation at least $52r$.  Otherwise $j\ge19044$ and
\[
 v\bigl(c_j(Df^2)^j\bigr)
 \ge4jr-\frac{4j+1}{p-1}
 \ge4jr-j\ge52r.
\]
These lower bounds tend to infinity, so they control the whole tail.
Consequently
\[
 v\bigl(\log_E(z)-zG_{12}(Df^2)\bigr)\ge53r,
 \qquad
 \frac{zG_{12}(Df^2)}{\log_E(z)}\in1+p^{52}\mathcal O_{K,\mathfrak p}.
\]
This proof applies to either a split or an inert good prime, since
both local extensions are unramified.  The source's global
nonvanishing argument applies to every $G_J$ with $J\ge2$ and hence
to $G_{12}$.  Thus all logarithm denominators remain nonzero.

The two-anchor comparison and the detector proof use only the
specified precision; they now give order $52$ for the three Lucas
labels and order $16$ for beta.  The transformed contact-height
estimate is unchanged.  Put $a_0=|a|$ and
$B=\max(|n|,|a|,|t|,|b|)$.  If the integer Lucas residual
$\Delta$ is nonzero, the resulting estimates are
\[
 B^2-6<7H|\Delta|^{24/52},\qquad |\Delta|<5B^4.
\]
Since $B,H\ge1$, they imply
\[
 B^{2/13}<\bigl(7\,5^{6/13}+6\bigr)H<21H,
 \qquad B<(21H)^{13/2}.
\]
The strict constant inequality follows by raising
$5^{6/13}<15/7$ to the thirteenth power; equivalently,
$5^6 7^{13}<15^{13}$.  In particular replacing the factor by
$14H$ would not be justified: $7\,5^{6/13}>14$.
The size hypothesis contradicts $B\ge a_0$, so $\Delta=0$ and
the unchanged Lucas lemma gives $\log B>a_0-3$.

The scalar estimate remains $h(\beta)\le H\sqrt{a_0}$.
If $\beta\ne n$, the order-$16$ argument remains
\[
 \log B<2\log(32H)+3h(\beta).
\]
The new size hypothesis implies $a_0>(20H)^2$, because
$21^{13}>20^4$ and $H\ge1$.  Hence
$h(\beta)<a_0/20$.  Also
$\log(32H)<20H^2<a_0/20$ for $H\ge1$.
It follows that $\log B<a_0/4<a_0-3$, a contradiction.  Thus
$\beta=n$ and $X\in\mathbb Z$ as before.

For the explicit constants,
\[
 \log_2\bigl((21H)^{13/2}\bigr)
 <\frac{13}{2}(5+4783227)=31091008.
\]
The exact source subgroup calculation for an anchor
$C=3^{2^q}$ gives
\[
 0\ne a\in2\,3^{2^{q-1}-1}\mathbb Z.
\]
For $q=26$, using $3^2>2^3$,
\[
 \log_2|a|\ge1+(2^{25}-1)\log_2 3
 >1+\frac32(2^{25}-1)=50331647.5>31091008.
\]
The semilocal anchor is therefore sufficient.  The alternative
indexed exclusion with $M_{52}=2^{31091008}$ is sufficient by the
same inequality.  Its enormous number of factors is part of the
explicit indexed specification, not an expanded coefficient table.

Completeness still chooses the anchor index in an arbitrarily
divisible nonzero subgroup, or larger than a fixed finite exclusion.
All the prime-pattern, congruence, local-square, and real-smallness
arguments are unchanged: lowering a required comparison precision
only weakens the necessary formal divisibilities.  This proves the
claimed equality of input projections.
\end{proof}

For the detector arithmetic, with $F=F_5(A\lambda)$ and
$\delta=A\lambda-1$, the three possible weights become
\[
 q_1=\frac{A^5\delta^2}{F},\qquad
 q_{16}=\frac{A^{155}\delta^{152}}{F^{31}},\qquad
 q_{52}=\frac{A^{515}\delta^{512}}{F^{103}}.
\]
These follow from $q_m=A^{10m-5}\delta^{10m-8}/F^{2m-1}$.
There are eight order-$52$ logarithms and two order-$16$
logarithms.  Each changed logarithm loses exactly three terms.
The fixed coefficient in the anchor is $2C^4=2\cdot3^{2^{28}}$;
using it as a literal coefficient adds no variable.

\begin{corollary}[Scalar order thirty-two]
\label{cor:shared-scalar32}
Under the hypotheses and the enlarged size condition of
Proposition~\ref{prop:shared-precision52}, the two order-$48$ scalar-height
tests can be replaced by order-$32$ tests.  This preserves the integer input
projection and does not change the point inventory.
\end{corollary}
\begin{proof}
The scalar-height argument with order $m$ gives
$h(\beta)\le H|a|^{24/m}$.  At $m=32$ it gives
$h(\beta)\le H|a|^{3/4}$.  The size condition implies
\[
 |a|>(21H)^{13/2}>(20H)^4,
\]
because $H\ge1$ and $21^{13}>20^8$.  Thus
$h(\beta)<|a|/20$, exactly the estimate used in the last contradiction
of Proposition~\ref{prop:shared-precision52}.  The Lucas part uses only the
order-$52$ comparisons and is unchanged.  Completeness chooses $S$ deeper
than a fixed scalar-height point, and decreasing the requested order
preserves that construction, also when the height point is shared as
proved above.  This proves both implications.
\end{proof}

\subsection{The domains of the division equations}
\label{subsec:direct-field-quotients}
Put $k=\mathbb Q(\sqrt2)$, $d=3-\sqrt2$, and $D=d^2$.
\begin{lemma}\label{lem:doubling-quotient-domain}
On $E:y^2=x^3-Dx$, the equation
\[
                           (2y)m=3x^2-D
\]
forces $y\ne0$ at every affine point, and hence defines the usual
doubling slope.  Every finite point of $2E(k)$ has nonzero abscissa
and nonzero ordinate.
\end{lemma}
\begin{proof}
If $y=0$, the curve equation gives either $x=0$ or $x^2=D$.
The numerator $3x^2-D$ is then respectively $-D$ or $2D$,
both nonzero.  Thus the slope equation cannot be satisfied when
$y=0$.  This argument uses the curve equation, and does not need
a separate assertion that $3x^2-D$ never vanishes elsewhere.

Let $Q=2R$ be finite.  The point $R$ is affine with nonzero
ordinate, and its doubling formula gives
\[
                  x(Q)=\left(\frac{x(R)^2+D}{2y(R)}\right)^2.
\]
Since $D=d^2$ and $k$ embeds in $\mathbb R$, the numerator is
nonzero.  Thus $x(Q)$ is a nonzero square in $k$.  A point with
zero ordinate has abscissa $0,d$, or $-d$.  The last two cannot
be squares in $k$: their norm to $\mathbb Q$ is $7$, whereas
the norm of a square in $k$ is a rational square.  Consequently
$y(Q)\ne0$ as well.
\end{proof}

The lemma shows that a halfpoint on the curve followed by the direct
doubling-slope equations defines a finite point of $2E(k)$ with the
required nonzero coordinates.  Both usual coordinates of the doubled
point are then given by their standard polynomial expressions in the
slope, so they lie on $E$.  Repeating this operation is valid by the
same argument.  If the usual subtraction slope is defined by its
existing abscissa-inequality guard, the resulting finite difference
of two points in $2E(k)$ also belongs to $2E(k)$ and has the same
nonvanishing properties.

Therefore the parameter $z(Q)=-x(Q)/y(Q)$ may be specified by a
direct two-coordinate quotient for every such point $Q$.  It is
itself nonzero.  For the reciprocal logarithms used in the source,
\[
              \ell_J(Q)=z(Q)G_J\bigl(D/x(Q)^2\bigr),
\]
the separately proved global nonvanishing of $G_J$ for $J\ge2$
makes $\ell_J(Q)\ne0$.  All truncations used here have
$J\in\{3,12,15\}$.  Ratios with these logarithms as denominators
may likewise use direct quotient equations.  This last step depends
on the source's stated nonvanishing lemma, not on numerical evaluation.

The abscissa inequality for the subtraction slope is retained: if two
points have the same abscissa, an unguarded quotient equation could have
both sides zero.  None of the preceding arguments removes that guard.


\subsection{The exact arithmetic conventions}

Represent a rational intermediate value by an affine linear form in the
coordinates already introduced.  Addition, subtraction, and multiplication
by a fixed rational coefficient then require no new coordinate.  A product
of two nonconstant affine forms $a,b$ is represented by one coordinate $v$
and the quadratic equation $v-ab=0$.  A rational inverse uses $av-1=0$.
More generally, an arbitrary quadratic expression $q$ can be assigned to
one new coordinate with $v-q=0$.  Common expressions are reused.

An element of $K=\mathbb Q(\sqrt2)$ is a pair of rational affine forms.
Its product with another pair is assigned to two coordinates, using
\[
 v_0-a_0b_0-2a_1b_1=0,\qquad
 v_1-a_0b_1-a_1b_0=0.
\]
This costs two coordinates, although an arithmetic algorithm minimizing
scalar multiplications might instead use three products and free additions.
The objective here is the number of existential coordinates.  An inverse
uses two coordinates and
\[
 a_0v_0+2a_1v_1=1,\qquad a_0v_1+a_1v_0=0.
\]
These equations hold exactly when the represented element is nonzero and
$(v_0,v_1)$ is its inverse.  They enforce the necessary domain without
introducing a separate norm and a separate inverse of that norm.

When the domain lemmas above already prove that $b\ne0$, a quotient
$z=a/b$ uses just the two equations
\[
 b_0z_0+2b_1z_1=a_0,\qquad b_0z_1+b_1z_0=a_1.
\]
Their determinant is $b_0^2-2b_1^2\ne0$, so their solution is unique.
This rule is used for doubling slopes, local parameters, logarithm ratios,
and the scalar local-parameter ratio, with the domain proof supplied in
the preceding subsection.  The subtraction slope retains an explicit
inverse of the difference of abscissas.  These domain arguments are part
of the proof of the quartic presentation; a coefficient-table check alone
would not justify dropping a denominator condition.

\subsection{Reusing polynomial evaluations}

Two elementary reorganizations save many auxiliary coordinates.  First,
write a polynomial with rational coefficients in blocks:
\[
 G(W)=\sum_{r=0}^s W^{br}
               \left(\sum_{j=0}^{b-1}c_{br+j}W^j\right).
\]
After computing $1,W,\ldots,W^b$, each parenthesized expression is an
affine linear combination of their two rational components.  Horner
evaluation in $W^b$ finishes the calculation.  This is the
Paterson--Stockmeyer principle \cite{PatersonStockmeyer1973}.  Here $b=3$
for $G_{12}$ and $b=4$ for $G_3$.  The identity is obtained simply by
grouping the coefficients according to their exponent modulo $b$; no
approximation of the logarithm is introduced by this evaluation rule.

Second, the eighteen field tests $D_j$ are the same at every coordinate
detector.  If $\delta=A\lambda-1$, the corresponding factor is
\[
 1+q_mD_j^2,\qquad q_m=R^{1-2m}\delta^{-2m-2},
 \qquad m\in\{1,16,52\}.
\]
For each order $m$, form once the polynomial
\[
 E_m(U)=\prod_{j:m_j=m}(1+UD_j^2)\in K[U].
\]
Its coefficients are the elementary symmetric polynomials in the listed
$D_j^2$.  There are respectively ten, two, and six factors.  For each of
the nine detectors, evaluate $E_1(q_1)E_{16}(q_{16})E_{52}(q_{52})$.
This is exactly the product defining $W$, with the field coefficients
reused across all nine detectors.

Finally, for $W=w_0+w_1\sqrt2$,
\[
 \operatorname{Tr}(W^{-1})=\frac{2w_0}{w_0^2-2w_1^2},\qquad
 B=\frac{R(w_0^2-2w_1^2)^2}{4w_0^2}.
\]
The inverse of $w_0$ enforces the trace condition.  Since $w_0\ne0$,
the norm of $W$ is nonzero in the quadratic field.  Together with
$R\ne0$, this also gives $B\ne0$, so an additional inverse of $B$
would be redundant.

\subsection{Finite choices at degree four}

The point occurrences use four rational halfpoint coordinates and two
curve equations.  The three-coordinate quartic chart is more economical
for finite-algebra elimination, but is not needed in this degree-four
construction.  The nine point occurrences are the four root points and
the five shared comparison points; the first comparison point is also
the scalar-height point.

The fixed square scales proved above remove all zero-component case
choices.  Only the scalar correction remains.  For each
$p\in\{2,3,5,7,11,13,17\}$ introduce a Boolean coordinate $\epsilon_p$,
and introduce one sign coordinate $\epsilon_-$.  Impose
$\epsilon^2-\epsilon=0$ for all eight coordinates and compute
\[
 c=(1-2\epsilon_-)\prod_p\bigl(1+(p-1)\epsilon_p\bigr).
\]
Over $\mathbb Q$, this realizes exactly the required 256 signed squarefree
corrections.  Unlike a finite product of alternative polynomials, this
keeps the total degree four after quadratic compilation, at the cost of
eight coordinates.

\subsection{Removing redundant quadratic coordinates}

The compiler next performs two reversible operations.  Replacing a
residual $q_i$ by $q_i+cq_j$, while retaining $q_j$, does not change the
common zero set.  If a residual has the form
\[
                  av+q=0,\qquad a\in\mathbb Q^\times,
                  \quad v\notin q,
\]
we may substitute $v=-q/a$ in all the other residuals and remove this
coordinate and equation.  Every zero of the new system reconstructs one
and only one value of $v$; the converse follows by substitution.  The
compiler accepts such a substitution only if every resulting residual
still has degree at most two.  In particular a quadratic replacement is
permitted only when that coordinate occurs linearly elsewhere.  Its exact
row operations and substitutions are retained as a replayable certificate.

After these operations, clear the rational denominator of each residual
and divide its integer coefficients by their common divisor.  Neither
operation changes its rational zero set.  If the resulting residuals are
$q_1,\ldots,q_s$, define
\[
                         P=\sum_{i=1}^s q_i^2.
\]
The rational zero set of $P$ is exactly their common zero set, because
each summand is nonnegative.  All these equivalences concern rational
solutions and do not require numerical sampling.

\begin{theorem}[The supplied quartics]\label{thm:quartic-1367}
Assume the source arithmetic inputs and Sun's simultaneous prime-pattern
theorem used in the shared-point construction.  The supplied polynomial
$P_{1367}$ has 1367 total variables, 1366 existential witnesses, and total
degree four, and defines $\mathbb Z$ in $\mathbb Q$.  It is specified by
a complete sparse coefficient table over $\mathbb Z[\kappa]$, where
\[
                         \kappa=3^{2^{28}}
\]
is a fixed integer, not a coordinate.  There is also a fully expanded
ordinary integer quartic $P_{1391}$ with 1391 total variables and primitive
coefficient height less than $2^{1020}$, with the same defining property.
\end{theorem}
\begin{proof}
Apply the just-described rules to the shared-point source with the
precision and size choices specified at the start of this section.
Each arithmetic coordinate has its prescribed value, except where the
explicit finite Boolean choices are intended; the paired quotient domain
lemmas establish the same assertion for the divisions.  The reversible
quadratic operations preserve the rational input projection.  The source
arithmetic theorem therefore identifies the final projection with
$\mathbb Z$.

The exact implementation introduces 1367 coordinates after sixty
quadratic-preserving eliminations.  Its 1322 primitive residuals expand
to 29,419 nonzero encoded entries over $\mathbb Z[\kappa]$, giving
29,415 distinct coordinate monomials after collecting the fixed
constant's powers. The independent checker
replays every arithmetic relation, row operation, and elimination and
independently convolves every pair of residual terms.  This checks the
entire sparse table, including mixed coefficients.  Its maximum variable
degree is exactly four.

For the ordinary integer alternative, start with the fixed integer
$3^{16}$ and use twenty-four successive squarings to compute
$3^{16\cdot2^{24}}=3^{2^{28}}$.  The associated coordinates are retained
to avoid expanding this integer into the coefficients.  Their defining
equations have degree at most two and uniquely determine their values.
The resulting 1391-coordinate system has 1346 residuals and 29,489
nonzero expanded integer terms.  The independent expansion check gives
coefficient bit length at most 1020.  The coefficient content is one.
These facts prove the degree, height, and count assertions.
\end{proof}

The supplied data are in
\nolinkurl{research/phase2/quartic/quartic-best/polynomial.json.gz} and
\nolinkurl{research/phase2/quartic/quartic-small-coefficients/polynomial.json.gz}.
An entry $[\mathbf m,k,c]$ means $c\kappa^k\prod_{i\in\mathbf m}X_i$,
with repetitions in $\mathbf m$ recording powers.  The index zero is the
input.  In the ordinary integer table every $k$ is zero.  This convention,
the complete tables, and the exact construction give ordinary finite
integer polynomials.  No gate or coefficient symbol is silently omitted
from the variable count.

The program \nolinkurl{build_shared_quartic.py} is a producer; it is not
trusted by the verification program
\nolinkurl{verify_shared_quartic.py}, which does not import it.
The separate mathematical review checks the source formulas and the
domain lemmas.  Neither computational check is claimed to verify the
external arithmetic theorems.  No global minimality claim is made for
1367 variables: the theorem supplies a concrete upper bound for the
quartic problem, while the unconditional lower bound for this
sum-of-squares architecture is five total variables.


\section{Quartics obtained from the direct Lucas test}
\label{sec:p3-quartics}

The direct test also improves the degree-four construction.
For this purpose we retain the scalar block with its preparation
coordinate, because its quadratic evaluation is less expensive
than the normalized scalar source used for finite-algebra
compression. We retain the proved comparison order $52$,
scalar order $32$, and fixed size constant
$\kappa=3^{2^{28}}$.

Replace the three proxy comparisons and the exact Lucas equation
by the one field-valued test
\[
 F(L_2,L_3,L_4),\qquad
 F(U,V,W)=V^2-U(U-2)VW+(U-2)^2W^2-1,
\]
of order $52$. Keep $L_1-\beta$ of order $16$ and all four
two-anchor cross comparisons. The direct-polynomial congruence
lemma applies with $k=52$ just as with $k=64$.
The precision-$52$ height proof uses the same integer residual
$\Delta$ and its bound $|\Delta|<5B^4$, so its threshold and
the chosen anchor are unchanged. The scalar-order-$32$ proof
is unaffected. Completeness uses the same fixed-point
neighborhoods: a polynomial in the original logarithm ratios
remains bounded there, introduces no new denominator, and
has the required prime-power precision.

The quadratic compiler evaluates the new field test by computing
\[
 A=U(U-2),\qquad B=VW,\qquad C=(U-2)W,
 \qquad F=V^2-AB+C^2-1.
\]
Each field multiplication uses its two rational component equations
as before. The coordinate detectors now have ten tests of order
one, two of order sixteen, and four of order fifty-two.
Their symmetric-product coefficient lists are shared across all
nine coordinate detectors. This removes more arithmetic
coordinates than merely deleting the three proxies would do.
All reciprocal-domain lemmas, Boolean scalar choices, and
quadratic-preserving substitutions from the preceding quartic
construction are retained.

\begin{theorem}\label{thm:p3-quartic1315}
Under the arithmetic inputs stated for the shared construction,
the supplied polynomial $P_{1315}$ defines $\mathbb Z$ in
$\mathbb Q$, has 1315 total variables and degree exactly four.
It is the sum of squares of 1272 primitive quadratic residuals
and has 28,991 nonzero encoded entries in its complete table over
$\mathbb Z[\kappa]$, with the specified fixed integer
$\kappa=3^{2^{28}}$. After collecting powers of this constant,
there are 28,987 distinct monomials in the quantified coordinates.

The alternative $P_{1339}$ has 1339 total variables, degree
exactly four, 1296 quadratic residuals, and 29,061 expanded
ordinary integer terms. Its primitive coefficient height is
less than $2^{1020}$.
\end{theorem}
\begin{proof}
The preceding argument proves the arithmetic source replacement
at precision $52$. The exact quadratic compilation and the
domain equivalences were proved in Section~\ref{sec:new-quartic}.
For the present source the independent checker verifies 1285
arithmetic gates, eleven row additions, and sixty coordinate
eliminations. It then independently squares and convolves every
residual and compares the complete resulting coefficient table.
The active coordinate count is 1315, the number of primitive
residuals is 1272, and the number of nonzero table entries is
28,991. The encoding separates different powers of $\kappa$
in a coefficient. Collecting these powers merges four pairs of
entries: the constant coefficient is positive, of the form
$c_0+1024\kappa^2$, and three further coefficients are
$1-192\kappa\ne0$. Thus there are 28,987 monomials in the
actual coordinates, with no cancellation after substituting the
specified integer. A nonzero degree-four coefficient occurs, so the degree
is exactly four. Over $\mathbb Q$ a sum of squares vanishes
exactly when all its residuals vanish, giving the asserted
input projection.

For the ordinary integer alternative, retain the same twenty-four
successive squaring coordinates computing $\kappa$ from $3^{16}$.
The count becomes 1339, and the checker verifies 1309 arithmetic
gates with the same eleven row additions and sixty eliminations.
Its 1296 primitive residuals give exactly 29,061 nonzero integer
terms. Their largest coefficient has bit length at most 1020,
and their content is one. This proves the degree and height
claims without leaving any coefficient expansion implicit.
\end{proof}

The two tables are in
\nolinkurl{research/phase3/quartic/quartic-direct/} and
\nolinkurl{research/phase3/quartic/quartic-direct-small-coefficients/}.
Their producer is \nolinkurl{build_direct_quartic.py}.
The independently written checker
\nolinkurl{research/phase2/quartic/verify_shared_quartic.py}
is reused without modification and does not import this producer.
These algebraic checks do not certify the external arithmetic
lemmas. The number-field section separately examines both
literal quartics and faithful combinations of their residual systems.


\section{Reducing the anchor coefficients and shortening its constant chain}
\label{sec:p3-affine-anchor-refinement}

The following refinements are exact transformations of the supplied
quadratic residual systems.  They require no additional arithmetic
theorem.  The original tables remain available; the transformed
tables are separate artifacts with certificates referring to their
frozen parents.

\subsection{An affine change in the three private anchor coordinates}

Put
\[
 \kappa=3^{2^{28}},\qquad s=3^{2^{27}},\qquad s^2=\kappa.
\]
In the $1315$-coordinate table the anchor residual is
\[
 R=-3A^2+16u^2+96v^2-32w^2+32\kappa,
 \qquad A=X_{52}+X_{53}-X_{54},
\]
where $(u,v,w)=(X_{55},X_{56},X_{57})$.  These three coordinates
occur in no other residual.  The other three displayed coordinates
are constrained arithmetic values, so they cannot be treated as
independent anchor witnesses.

Make the rational affine substitution
\begin{equation}\label{eq:p3-anchor-affine-map}
 u=s+\frac{U-3V+3W}{4},\qquad
 v=\frac{s}{2}+\frac V4,\qquad
 w=\frac{3s}{2}+\frac W4.
\end{equation}
Its linear determinant is $1/64$, and its inverse is
\[
 U=4u+12v-12w+8s,\qquad V=4v-2s,\qquad W=4w-6s.
\]
It is therefore a bijection over every characteristic-zero field,
fixing the input and every other coordinate.  Direct expansion gives
\begin{equation}\label{eq:p3-translated-anchor}
 \widetilde R=-3A^2+U^2-6UV+6UW+15V^2-18VW+7W^2+8sU.
\end{equation}
Indeed the constant terms sum to
$16s^2+24s^2-72s^2+32s^2=0$, and the linear terms sum to $8sU$.
All coefficients of the new residual are integers, with $s$ regarded
as the specified fixed integer, not a coordinate.

\begin{proposition}\label{prop:p3-translated-quartic}
Replacing $R$ by $\widetilde R$ in the frozen $1315$-coordinate
residual list gives a polynomial $\widetilde P_{1315}$ with exactly
the same input projection as $P_{1315}$ over every characteristic-zero
field.  Its degree is exactly four, its primitive integer coefficient
height is
\[
                  H(\widetilde P_{1315})=64\kappa
                                      <2^{429496736},
\]
and its complete table has $29021$ nonzero ordinary coordinate
monomials.  It has $1272$ quadratic residuals and $1314$ existential
witnesses.  In particular, under the already stated arithmetic inputs,
it defines $\mathbb Z$ in $\mathbb Q$.
\end{proposition}
\begin{proof}
The polynomial identity
\[
 \widetilde P_{1315}
       =P_{1315}\circ\text{the substitution in
                           \eqref{eq:p3-anchor-affine-map}}
\]
and the explicit inverse prove equality of input projections.
This conclusion concerns the literal single polynomial as well as
the common residual zero set; it does not use positivity of squares.
No new coordinate or equation is introduced, and every residual
remains quadratic.

For the height assertion write
\[
 B=-3A^2+U^2-6UV+6UW+15V^2-18VW+7W^2.
\]
Then $\widetilde R^2=B^2+16sUB+64s^2U^2$.  Since $U,V,W$
occur nowhere else, the coefficient of the monomial $U^2$ in
the sum of squares is exactly $64s^2$.  It is the only coefficient
involving $s^2$.  The coefficients involving $s$ have absolute value
at most $288s$.  The independently expanded table records that
every remaining integer coefficient has absolute value less than
$2^{1020}$, and no monomial receives two different powers of $s$.
The value $s=3^{2^{27}}$ makes $64s^2$ larger than both other bounds.
The input coefficient of $X^4$ is exactly one.  Thus the integer
content is one, the height is exactly $64s^2=64\kappa$, and the
degree is exactly four.  Finally $3^5<2^8$ gives
\[
 \log_2(64\kappa)<6+\frac85\,2^{28}<429496736.
\]
The displayed residual list and its exact convolution give the
stated monomial count.
\end{proof}

\subsection{Eliminating the last coordinate in the constant chain}

In the ordinary-coefficient $1339$-coordinate table, the coordinate
$X_{81}$ occurs in exactly two residuals:
\[
                   X_{80}^2-X_{81},\qquad R_0+32X_{81},
\]
where
\[
 R_0=-3(X_{52}+X_{53}-X_{54})^2
                       +16X_{55}^2+96X_{56}^2-32X_{57}^2.
\]
Substitute $X_{81}=X_{80}^2$, delete its defining residual, and
renumber all subsequent coordinates by subtracting one.  The second
residual becomes $R_0+32X_{80}^2$ and is still quadratic.

\begin{proposition}\label{prop:p3-small-coefficient-1338}
The resulting system has $1338$ total coordinates and $1295$
quadratic residuals.  Its sum of squares $P_{1338}$ has degree exactly
four, $29059$ nonzero ordinary integer monomials, primitive integer
content one, and coefficient height less than $2^{1020}$.
Its rational input projection equals that of $P_{1339}$.
\end{proposition}
\begin{proof}
On a common zero of the parent residuals, the first displayed equation
uniquely determines $X_{81}=X_{80}^2$.  Deleting this coordinate
therefore gives a solution of the new system.  Conversely, inserting
that square in a new solution restores all parent equations.  This
is a bijection of common residual zero sets over every
characteristic-zero field and fixes the input.

Over a formally real field, in particular over $\mathbb Q$, a sum
of squares vanishes precisely when all its residuals vanish.  Hence
the same equivalence holds for the literal sum-of-squares input
projections over those fields.  No such converse for the literal
polynomials over nonreal fields is inferred from this elimination.

The substitution keeps every residual quadratic.  Exact convolution
of the remaining residuals gives $29059$ nonzero coefficients, all
of absolute value below $2^{1020}$.  The unchanged coefficient of
the input monomial $X^4$ is one, proving content one and degree four.
The complete maximum coefficient is recorded in the construction
report.  The coordinate and residual counts each decrease by one.
\end{proof}

\subsection{Inlining the three initial constants}

The first three coordinates of the remaining constant chain are
uniquely determined as
\[
                  X_{58}=3^{32},\qquad X_{59}=3^{64},
                  \qquad X_{60}=3^{128}.
\]
Their three defining residuals are
\[
       3^{32}-X_{58},\qquad X_{58}^2-X_{59},
       \qquad X_{59}^2-X_{60}.
\]
They have no other occurrences except in the next equation of
the chain.  Substitute the three constants and delete these
coordinates and residuals.  The first remaining chain equation
becomes $3^{256}-X_{61}=0$ before renumbering, or
$3^{256}-X_{58}=0$ afterwards.

\begin{theorem}[The preferred ordinary-coefficient quartic]
\label{thm:p3-small-coefficient-1335}
The supplied polynomial $P_{1335}$ has $1335$ total coordinates,
$1334$ existential witnesses, $1292$ quadratic residuals, degree
exactly four, and $29050$ nonzero ordinary integer monomials.
Its primitive coefficient height is unchanged from $P_{1338}$ and
is less than $2^{1020}$.  Under the same arithmetic hypotheses,
it defines $\mathbb Z$ in $\mathbb Q$.
\end{theorem}
\begin{proof}
Each eliminated coordinate equals the displayed fixed integer
at every common zero of the parent residuals.  Substitution and
reinsertion are inverse maps on the corresponding common zero
sets, over every characteristic-zero field, and fix the input.
All remaining residuals have degree at most two.  Over $\mathbb Q$
their sum of squares therefore has the same input projection.

The first remaining chain residual contributes the constant
$3^{512}$ on squaring, and
\[
                             3^{512}<2^{812}<2^{1020}.
\]
For an elementary check of the first inequality, $3^{41}<2^{65}$
and $512\cdot65<812\cdot41$.  The only changed constant
coefficient is the old one minus $3^{64}$ plus $3^{512}$; it
has bit length at most $826$.  The new linear coefficient has
absolute value $2\cdot3^{256}<2^{407}$.  Every other coefficient
is inherited, up to coordinate renumbering, or removed.  The
previous coefficient of maximum height occurs in an unaffected
monomial, so the exact height is unchanged.  These assertions
are also checked against the complete coefficient table.

Twenty constant-chain coordinates remain, representing
$3^{2^8},3^{2^9},\ldots,3^{2^{27}}$.  The final square is
used directly in the anchor.  Thus the ordinary-coefficient
presentation adds twenty coordinates to the $1315$-coordinate
literal-constant presentation.  Exact convolution gives $29050$
nonzero terms.  The unchanged coefficient of the input monomial
$X^4$ is one, proving primitive content one and degree four.
\end{proof}

The preferred fresh packages are
\nolinkurl{research/phase3/quartic/quartic-direct-translated/} and
\nolinkurl{research/phase3/quartic/quartic-direct-small-coefficients-1335/}.
The intermediate $1338$-coordinate package is retained separately
under \nolinkurl{quartic-direct-small-coefficients-1338/}.
Each contains the full residual list, complete coefficient table,
parent hash, and exact transformation certificate.  Their producer
is \nolinkurl{transform_direct_quartics.py}.  It verifies both the
residual substitutions and the full polynomial identities against
the frozen parent tables.  The separate program
\nolinkurl{verify_final_transforms_independent.py} imports no producer
code and independently checks every transformed residual, every
expanded coefficient, all coordinate maps, and the stated heights
and supports.  All three checks pass.  The translated table uses
the coefficient symbol $S$ for the fixed integer $s$; an entry
$[\mathbf m,k,c]$ means $cs^k\prod_{i\in\mathbf m}X_i$.
This symbol contributes no variable and no degree in the quantified
coordinates.


\section{Exact quartics and their computational verification}
\label{sec:quartic-artifacts}

The preceding anchor formula is applied to the actual retained residual
system, not only to an informal count. The old private anchor uses two
scaling coordinates $a_i$ with $Ca_i=f_i$, and, for each $i=0,1$, the
seven-witness equations
\begin{align*}
 b_{i0}^2-2b_{i1}^2-3b_{i2}^2+6b_{i3}^2&=1,\\
 (a_i-2b_{i0})^2-8c_{i1}^2-12c_{i2}^2+24c_{i3}^2&=4.
\end{align*}
The norm-one trace-sum theorem for the quaternion algebra $(2,3)$
identifies this condition with $a_i\in R$; its ramification set is
$\{2,3\}$ and it splits at infinity
\cite[Proposition 6]{Koenigsmann2016}. The second equation has scalar
part $(a_i-2b_{i0})/2$, so the trace convention agrees exactly.
The retained constant is $C=3^{2^{24}}\ne0$.

\begin{proposition}\label{prop:computed-quartics}
There are exact integer quartics $P_{12493}$ and $P_{12465}$ with
the same rational input projection as the released source. Their
statistics are
\[
\begin{array}{c|r|r|r|r}
 &\text{total variables}&\text{residuals}&\text{monomials}&\log_2 H\ \text{bound}\\\hline
 P_{12493}&12493&11936&135667&<80\\
 P_{12465}&12465&11908&130427&<3247.
\end{array}
\]
The height of $P_{12493}$ equals the height of the release.
Both complete expanded integer coefficient tables are supplied.
\end{proposition}
\begin{proof}
For the first polynomial replace the six old anchor equations by
the three equations of Corollary~\ref{cor:five-anchor}. Their sixteen
private coordinates occur nowhere else. The five new coordinates
retain quadratic residuals and have exactly the same projection at the
interface $(C,f_0,f_1)$. Eleven variables and three residuals are removed.
The sum of the squares of all retained residuals is a quartic.

For the second polynomial first perform twenty-eight additional exact
eliminations and thirteen invertible row operations. An elimination
uses a row $cY+g=0$, with $c\in\Q^\times$ and $g$ independent of
$Y$, substitutes $Y=-g/c$ into every occurrence, and removes that row
and coordinate. Both rational solution maps are immediate. A row
operation replaces $f$ by $af+bg$ while retaining $g$, with
$a\in\Q^\times$; again both solution sets agree. The certificate
records every affected row, all denominator clearings, and the variable
remapping. Every resulting residual is still quadratic. Apply the same
eleven-variable anchor replacement afterward.

The independent checker replays the pinned baseline certificate, checks
every extra operation using exact fractions without importing the
producer, checks the literal six anchor rows and their full incidence,
verifies the nonzero defining chain of $C$, and expands every square
by ordered products. It compares every coefficient with the supplied
table and verifies that each reported variable occurs. This gives the
displayed finite statistics. The semantic anchor equivalence is proved
above; the checker verifies its exact placement. Neither part uses an
unproved arithmetic optimization of the source.
\end{proof}

The small-coefficient package provides a useful alternative when file
size matters. Under the requested variable-first ordering,
$P_{12465}$ is preferable despite its coefficient increase. Its largest
coefficient has 978 decimal digits; the exact integer is recorded in the
verification receipt, rather than reproduced here.

\subsection{Removing constants without expanding their decimal digits}

There are still fifteen fixed coordinates in $P_{12465}$, beginning
with $c_0=3^{1024}$ and continuing by $c_{k+1}=c_k^2$ for
$0\leq k<14$. Thus $c_{14}=C=3^{2^{24}}$. The new anchor coordinate
$D_0=C^2$ is fixed as well. Substitute these sixteen constants into
every remaining residual and delete their defining rows. This cannot
increase degree. The only remaining large constant is
\[
 M=C^4=3^{2^{26}},
\]
in the anchor residual
\[
 u^2+6v^2-2w^2-3r^2+2M.
\]
Here $M$ is a fixed integer, not a variable.

\begin{proposition}\label{prop:constant-quartic}
The constant substitution gives a quartic in exactly 12,449 total
variables, with 11,892 residuals and 130,378 monomials, having the
same rational input projection. Its primitive integer coefficient
height is exactly
\begin{equation}\label{eq:constant-height}
 H=4\cdot3^{134217728}+H_0,
\end{equation}
where
\begin{align*}
H_0={}&373684652\cdot10^{144}\\
&+ 355249371111055415460965317040747589616777785974\cdot10^{96}\\
&+ 770831188591770167322175457505857070439628239726\cdot10^{48}\\
&+ 443171972269170621942059432904883611120549682273.
\end{align*}
Its coefficients are completely specified by a table over $\Z[M]$
of degree at most two in the fixed constant $M$.
\end{proposition}
\begin{proof}
Each deleted coordinate is uniquely fixed by its defining equation.
Direct incidence checking shows that, outside the deleted chain, these
coordinates occur only in the displayed anchor places. Thus substitution
and reconstruction give inverse maps on the retained coordinates.
The exact expansion over $\Z[M]$ gives the stated counts.
No nonconstant monomial disappears on specializing the fixed constant:
its coefficient is $a+bM$ with $|a|<2^{3247}<M$, so it is nonzero
when $b\ne0$; a nonzero constant coefficient remains nonzero as well.
Only the coefficient of the constant monomial has an $M^2$ term.
The constant coefficient is $4M^2+H_0$. Every other coefficient has
absolute value at most $2^{3247}+24M$, which is strictly smaller.
The coefficient of $u^4$ is one, so the integer content is one and the
degree is four. Each remaining variable occurs in a residual, and a
sum of squares of real polynomials cannot lose its dependence on such
a variable: its leading coefficient in that variable is a nonzero
sum of squares. This verifies the exact variable count.
\end{proof}

The integer in \eqref{eq:constant-height} is intentionally not converted
to a decimal string. This is an exact polynomial specification by an
ordinary integer power, not a numerical approximation. An elementary
upper bound is $H<2^{268435459}$. The smaller variable count is purchased
with very large coefficients, in accordance with the specified priority.

\subsection{Verification boundaries}

The files separate producer output from independent verification. The
constant table is checked as a polynomial over $\Z[M]$, so there is no
floating-point arithmetic or gigantic decimal conversion in that check.
For the preceding elimination suffix, seven deliberate mutations of
pivots, solved numerators, affected-row lists, substitution patches,
and row operations were rejected. These checks test the advertised
finite transformations.
They neither test all rational inputs nor certify the manuscript's
restricted converse, height estimate, or prime-pattern argument.


\section{A fully specified polynomial with 854 variables}
\label{sec:explicit-circuit}

This section gives a finite formula for an actual polynomial, not just
an existence argument. Its degree is much higher than four. Arithmetic
expressions in this section introduce no quantified coordinates unless
explicitly listed. The accompanying JSON circuit makes every addition,
multiplication, power, residual numerator, and nonzero condition literal.

\subsection{The finite list of rational equations}

Use the curve and constants of Section~\ref{sec:source-structure}, and
write $\sigma$ for the automorphism of $K=\Q(\sqrt2)$ sending
$\sqrt2$ to $-\sqrt2$. For a nonidentity point $Q=(x,y)$ define
\[
 f(Q)=1/x,\qquad z(Q)=-x/y.
\]
All the denominators occurring here and below are required to be
nonzero. A field inverse is written using its rational norm, so both
rational coefficient formulas are specified.

There are six root points $A,S,T,Z,J_0,J$, and five private points
$Q_{\ell,0},\ldots,Q_{\ell,4}$ for each $1\leq\ell\leq41$.
Their total number is $211$. For each point use the three rational
coordinates $(a_0,a_1,C_Q)$ and the common chart parameter
$j\in\{0,\ldots,211\}$ of
Proposition~\ref{prop:quadratic-square-root}. Namely put
$a=a_0+a_1\sqrt2$, write $a^3-Da=A_Q+B_Q\sqrt2$, reconstruct
$b=b_0+b_1\sqrt2$, impose
\[
 C_Q^2 V_Q-(B_Q/2+jV_Q)^2=0,
 \qquad V_Q=\frac{C_Q^2-A_Q-jB_Q}{j^2-2},
\]
and use the double $Q=[2](a,b)$ in every subsequent expression.
Require $C_Q\ne0$, $N_{K/\Q}(b)\ne0$, and nonzero norms for
the abscissa and ordinate of the output point. The source-chart
equivalence and its inverse maps are proved in
Section~\ref{sec:finite-fibre-application}.

Set $\beta=X^2+2$. Introduce rational proxies $s,\tau,z$ and impose
\begin{equation}\label{eq:circuit-lucas}
 \tau^2-s(s-2)\tau z+(s-2)^2z^2-1=0.
\end{equation}
Use one rational $\theta$, with $\varepsilon=\theta^2-2$.
The letter $z$ without an argument in \eqref{eq:circuit-lucas}
is a proxy; $z(Q)$ is the specified elliptic rational function.

The finite reciprocal logarithms are
\[
 \mathcal L_L(Q)=z(Q)G_{L/4-1}(Df(Q)^2),
 \qquad L\in\{16,64\},
\]
where $G_J(W)=1+\sum_{r=1}^J a_rW^r$,
$a_1=-2/5$, and $a_r=a_{r-1}(4r-5)/(4r+1)$.
Their nonvanishing on the required rational point charts is proved
in the finite-fiber application.

We next specify the norm operation used in every contact. Given a
nonzero rational $r$, an ordered list of tests $t_i$ with positive
orders $L_i$, and $L_1=1$, form the normalized leaves
$1,t_1/r,t_2/r^{L_2},\ldots$. First pair $1,t_1/r$ by
$N_r(a,b)=a^2+rb^2$. Group all remaining tests by increasing order,
preserving their original order within each group; pair successive
normalized leaves and carry an unpaired last leaf. At each higher
level pair successive nodes by $N_r$ and carry any last node until
one root $W_r$ remains. The input of the pole equations is $b=rW_r^2$.
This prescription is finite and fixes the order of every operation.

For a rational input $b$, a pole occurrence has five coordinates
$t,e_p,p,q,h_p$. Define
\[
 U=\frac{b(1-bt^2)}{1+bt^2},\qquad B=17b,
 \qquad H=e_p(B-U^2),
\]
and impose
\begin{equation}\label{eq:circuit-pole}
 p^2=B(1-e_p)(1+3e_p),\qquad
 Bq^2=(h_p^2-BH)(h_p^2+3BH).
\end{equation}
The denominator $1+bt^2$ cannot vanish on the original pole equations;
its explicit nonzero condition is retained here.

For the height occurrence take
\[
 r_0=\beta(\beta-1)(\beta-2401)(32\beta+343)(2\beta-1)
\]
and the six rational coefficients of
$f(J_0),f(S),z(S)/z(J_0)$, in that order. Their orders are
$(1,1,1,1,48,48)$. Replace the resulting pole input by $c r_0W_{r_0}^2$,
where the fixed finite choice is
\[
 c\in\mathcal C=\left\{\pm\prod_{p\in I}p:
 I\subseteq\{2,3,5,7,11,13,17\}\right\}.
\]
There are $256$ such integers, covering the source's eight Boolean
choices exactly.

Introduce a nonzero coordinate $e$, and let $a_1,\ldots,a_9$ be
$e$ followed by the eight rational coefficients of
$ex(A),ex(S),ex(T),ex(Z)$. Introduce two contact coordinates $u_c,v_c$.
With $W_0=25878772920$, the forty-one contacts are
\[
 \frac{v_c}{32W_0},\quad u_c,\quad u_c-v_c,\quad
 u_c-2401v_c,\quad u_c+\frac{343}{32}v_c,
\]
and, for each $i=1,\ldots,9$, on putting $A_i=47027a_i+3$,
\[
 u_c-A_iv_c,\quad u_c-A_iv_c+v_c,\quad
 u_c-A_iv_c+2401v_c,\quad
 u_c-A_iv_c-\frac{343}{32}v_c.
\]
Write these contacts as $r_1,\ldots,r_{41}$ in the displayed order.

For contact $\ell$, the five target points are
$(4P_*,4A,4S,4T,4Z)$. Subtract them respectively from the five
private points to obtain $Y_{\ell,i}$, using the affine subtraction
formula and its nonzero abscissa-difference denominator. The ordered
twenty field tests are the following:
\begin{enumerate}
\item $f(Q_{\ell,i})$ for $0\leq i\leq4$;
\item $f(J)^\sigma$;
\item $f(Y_{\ell,i})^\sigma$ for $0\leq i\leq4$;
\item $\mathcal L_{16}(Q_{\ell,1})/\mathcal L_{64}(Q_{\ell,0})-\beta$,
and the analogous three ratios with numerator precision $64$ and
proxies $s,\tau,z$;
\item $z(Y_{\ell,i})^\sigma/z(J)^\sigma$ for $0\leq i\leq4$.
\end{enumerate}
Replace each field test $A+B\sqrt2$ by the two rational tests
$\varepsilon B,A+\theta B$, and prepend $\varepsilon$. The
resulting forty-one orders are
\[
 (\underbrace{1,\ldots,1}_{23},16,16,
 \underbrace{64,\ldots,64}_{8},16,16,
 \underbrace{64,\ldots,64}_{6}).
\]
Apply the norm prescription at $r_\ell$ and then one copy of
\eqref{eq:circuit-pole}. This makes $42$ pole occurrences in total.

Finally, write $f(S)=f_0+f_1\sqrt2$, set
$F=f_0^2+f_0f_1-f_1^2$, and let $C=3^{2^{24}}$ be the fixed
anchor constant. The three-witness anchor is
\[
 u_a^2+6v_a^2-2w_a^2-3F^2+2C^4=0.
\]
We next absorb its private coordinate $u_a$ and all denominator
conditions into one coordinate, without increasing their combined count.

\subsection{Eliminating the denominator witness}

For fixed $(j,c)$, clear the denominators of all rational equations
other than the anchor, and call their polynomial numerators
$E_1,\ldots,E_{296}$. There are $211$ point equations, $84$ pole
equations, and one Lucas equation. During the construction record every
nonzero numerator required for a division, including rational norms
for divisions in $K$, and the explicit point and coordinate guards.
Let $G$ be their product, with repeated factors omitted if desired.
All the rational functions used are defined whenever $G\ne0$.

Write
\[
 6v_a^2-2w_a^2-3F^2+2C^4=N/B
\]
with integer polynomial numerator $N$ and denominator $B$. Choose
$F=R_0/S_0$ in the rational-expression construction. Then
\[
 B=S_0^2,\qquad
 N=S_0^2(6v_a^2-2w_a^2+2C^4)-3R_0^2.
\]
The actual circuit's cleared anchor has exactly this form.
All denominator factors are included among the recorded conditions.
The three anchor coordinates occur nowhere in $G$ or the other
residuals, and $B$ is independent of them.

The anchor always has a solution with $u_a\ne0$ whenever it has any
solution. Indeed its three witnesses cannot all vanish, since
$3F^2=2C^4$ is impossible for rational $F/C^2$. If $u_a=0$ and
$w_a\ne0$, replace $(u_a,w_a)$ by $(4w_a,3w_a)$; this preserves
$u_a^2-2w_a^2$. If $u_a=w_a=0$ and $v_a\ne0$, replace
$(u_a,v_a)$ by $(-12v_a/7,-5v_a/7)$; this preserves
$u_a^2+6v_a^2$. In either case $u_a$ becomes nonzero without
changing any outside equation.

Rescale $t_a=S_0u_a$, so that the cleared anchor is $t_a^2+N=0$.
Replace its private coordinate by $U$ and put
\[
 \Psi=1+NG^2U^2.
\]
Define the branch polynomial
\begin{equation}\label{eq:854-branch}
 F_{j,c}=\sum_{i=1}^{296}E_i^2+\Psi^2,
\end{equation}
and the required integer polynomial
\begin{equation}\label{eq:854-polynomial}
 P_{854}=\prod_{j=0}^{211}\ \prod_{c\in\mathcal C}F_{j,c}.
\end{equation}

\begin{proposition}\label{prop:854-explicit}
The polynomial \eqref{eq:854-polynomial} uses 854 total variables
and has the same rational input projection as $\Src$. A supplied
integer arithmetic circuit specifies it completely. Its degree is
at most $63570781474816$, and its coefficient height is at most
$2^{757866260454400}$.
\end{proposition}
\begin{proof}
If $GU=0$, then $\Psi=1$, so \eqref{eq:854-branch} has no rational
zero. Otherwise $G\ne0$ also implies $S_0\ne0$, and
\[
 \Psi=(GU)^2\bigl((1/(GU))^2+N\bigr).
\]
A zero therefore gives an original anchor witness
$u_a=1/(GS_0U)$. Conversely, use the preceding transformations to
choose $u_a\ne0$, then put $U=1/(GS_0u_a)$. This preserves the
existential input projection. It need not preserve a fiber in which
both of the other private anchor witnesses have been fixed, and no
such claim is needed. The sum of squares imposes all the outside
equations and all the divisions used before taking numerators.

The point substitutions and the pole substitutions have the inverse
maps proved above and below. Among the $212$ fixed chart choices
there is a valid common choice for any list of $211$ source halfpoints.
The $256$ fixed corrections cover the eight Boolean source choices.
A product of branch polynomials permits reuse of the same witness
coordinates, since only one branch must vanish. In particular, no
denominator condition from an inactive branch is imposed.

The coordinate count is
\[
 633+210+3+3+1+1+2+1=854:
\]
point coordinates, pole coordinates, anchor coordinates $(U,v_a,w_a)$,
proxies, $\theta$, $e$, two contact coordinates, and the input.
The final $U$ also enforces the denominators, so no inverse coordinate
is added. All 854 coordinates occur in the supplied specification.

For the bounds, assign each integer constant its absolute coefficient
sum, each variable degree one, and each fixed parameter its maximum
absolute value over its stated finite list. At a sum, degrees take a
maximum and coefficient sums add; at a product, degrees add and
coefficient sums multiply; at a power both bounds scale accordingly.
These exact integer recurrences bound every node. Finally multiply the
branch bounds over $212\cdot256=54272$ branches. Replaying these
recurrences on the supplied circuit gives the displayed bounds.
They are upper bounds, not exact expansion statistics.
\end{proof}

This explicit circuit and the 558-variable indexed construction
serve different purposes. The former can be regenerated directly from
the finite formulas and inspected as a circuit; the latter uses further
finite-algebra elimination whose indexed output is specified and checked,
although its rank-sized matrices and full expansion are not materialized.
Neither is presented as a practical solver for Diophantine equations.


\section{An explicit reduction to 853 variables}
\label{sec:explicit-853}

Two further elementary steps combine the private anchor root with the
Lucas root. Their proof concerns rational solutions. It is not a claim
about geometric fibers over an algebraic closure.

\subsection{The pole denominators are automatic over the rationals}

\begin{lemma}\label{lem:pole-denominator-automatic}
After clearing the denominator $\delta=1+bt^2$ from
\eqref{eq:circuit-pole}, its two equations already imply
$\delta\ne0$ over $\Q$.
\end{lemma}
\begin{proof}
Put $B_p=17b$ and
$A=B_p\delta^2-b^2(1-bt^2)^2$. The cleared second equation is
\[
 B_pq^2\delta^4
 =(h_p^2\delta^2-B_pe_pA)(h_p^2\delta^2+3B_pe_pA).
\]
If $\delta=0$, then $bt^2=-1$, so $b,t,B_p$ are nonzero and
$A=-4b^2$. The equation becomes $0=-48B_p^2b^4e_p^2$, whence
$e_p=0$. The first equation now gives $p^2=17b=-17/t^2$,
or $(pt)^2=-17$, impossible in $\Q$.

This also applies to the reduced rational numerators used by the
circuit. In each actual second conic, the denominator contains the
numerator of $\delta$ to exactly the fourth power, and every other
factor is among the retained nonzero conditions. The first-conic
denominator uses only retained factors. Multiplying back yields the
displayed polynomial identity even when $\delta=0$. These exact
denominator-factor statements are checked in the independent verification.
\end{proof}

More generally the same proof works with $B_p=cb$, over a field of
characteristic different from two and three, whenever $-c$ is a
nonsquare. It deliberately does not give the corresponding assertion
over every number field.

Remove precisely the $42$ pole-denominator conditions from the product
$G$ of the preceding section, and call the remaining product $G_0$.
Its factors depend only on the input, point coordinates, coordinate
denominator, and contact coordinates. In particular $G_0$ is independent
of all three proxies, all pole coordinates, $\theta$, and the private
anchor coordinates. This follows from the displayed construction and
was checked on the entire circuit. All other rational denominators
are covered by $G_0$; the omitted ones are enforced by
Lemma~\ref{lem:pole-denominator-automatic}.

\subsection{Combining two square roots and the nonzero conditions}

Write the Lucas equation as
\[
 \tau^2-a\tau+b_0=0,\quad
 a=s(s-2)z,\quad b_0=(s-2)^2z^2-1,\quad
 \Delta=a^2-4b_0.
\]
Thus $v=2\tau-a$ satisfies $v^2=\Delta$. The scaled anchor witness
satisfies $t_a^2=-N$, where $N$ is independent of $\tau$.
Choose its nonzero representative as in the preceding section.
Its sign can be selected so that $W=t_a+v\ne0$: the two choices
$t_a+v$ and $-t_a+v$ cannot both vanish when $t_a\ne0$.
Consequently no new finite selector is necessary.

The two square-root equations have the exact single relation
\begin{equation}\label{eq:joint-quartic}
 h(W)=W^4+2(N-\Delta)W^2+(N+\Delta)^2=0,
\end{equation}
with reconstruction
\[
 v=\frac{W^2+N+\Delta}{2W},\qquad
 t_a=\frac{W^2-N-\Delta}{2W},\qquad
 \tau=\frac{a+v}{2}.
\]
Indeed $4W^2(v^2-\Delta)=h(W)=4W^2(t_a^2+N)$.
Conversely, $W=t_a+v$ and the two square equations imply
\eqref{eq:joint-quartic} and these reconstruction formulas.

Set $W=1/(G_0U)$ and define
\[
 \Psi_*=1+2(N-\Delta)G_0^2U^2+(N+\Delta)^2G_0^4U^4,
 \qquad
 \tau=\frac{1+2aG_0U+(N+\Delta)G_0^2U^2}{4G_0U}.
\]
The equation $\Psi_*=0$ forces $G_0U\ne0$ and is the reciprocal
form of \eqref{eq:joint-quartic}. Substitute the displayed rational
formula for $\tau$ into all the other residuals. Choose a syntactic
upper bound $k$ for each residual's degree in $\tau$, and multiply
its substituted expression by $(4G_0U)^k$ to obtain a polynomial
numerator. At addition gates take the maximum of the children's
bounds; at multiplication gates add them; at power gates multiply
by the exponent. This specifies the emitted circuit even if
cancellation lowers the actual degree. The operation is reversible
on $G_0U\ne0$.
Delete the old Lucas residual and the anchor residual; denote the
remaining $295$ polynomial numerators by $E'_i$. The branch polynomial
and the full polynomial are
\[
 F^*_{j,c}=\sum_{i=1}^{295}(E'_i)^2+\Psi_*^2,
 \qquad
 P_{853}=\prod_{j=0}^{211}\prod_{c\in\mathcal C}F^*_{j,c}.
\]

\begin{proposition}\label{prop:853-explicit}
The integer polynomial $P_{853}$ has 853 total variables and the
same rational input projection as $\Src$. The supplied finite
arithmetic circuit has degree at most $894090701148160$ and
coefficient height at most $2^{29154205340326912}$.
\end{proposition}
\begin{proof}
The forward construction may change only the three private anchor
witnesses, making $t_a\ne0$ and then choosing its sign to make
$W\ne0$. Set $U=1/(G_0W)$; every transformed residual vanishes.
Conversely, $\Psi_*=0$ gives $G_0U\ne0$. Reconstruct $W,v,t_a$
and $\tau$ by the formulas above, then $u_a=t_a/S_0$. They satisfy
the old anchor and Lucas equations. Undoing each homogenization gives
the other original numerator equations. The retained guards and
Lemma~\ref{lem:pole-denominator-automatic} make every original rational
division valid. The source point-chart equivalences then apply as
before. These operations prove equality of rational input projections.

Two coordinates, $u_a$ and $\tau$, have been replaced by the single
coordinate $U$. The count is
\[
 633+210+2+2+1+1+2+1+1=853,
\]
respectively point coordinates, pole coordinates, two remaining anchor
coordinates, two remaining proxies, $\theta$, the coordinate denominator,
two contact coordinates, the input, and $U$. There are still $54272$
finite branches. Exact circuit bounds are propagated as in
Proposition~\ref{prop:854-explicit}; the degree and height statements
are upper bounds. The maximum homogenizing degree in $\tau$ is $1280$.
\end{proof}

The omission of the pole guards is justified here over $\Q$.
The 558-variable geometric-fiber argument retains those guards and
does not use this omission. This distinction prevents exceptional
geometric fibers from being silently reintroduced into that argument.


% Includable section. Requires amsmath, amssymb, amsthm, hyperref.
% Bibliography entries are supplied in lower_bounds.bib.
\section{Unconditional lower bounds}
\label{sec:lower-bounds}

The lower bounds in this section do not assume the correctness of any proposed
Diophantine definition.  In fact, they hold over every number field.
Let $K$ be a number field and let $\mathcal O_K$ be its ring of integers.
For a polynomial
\[
 F(T,Y_1,\ldots,Y_r)\in K[T,Y_1,\ldots,Y_r]
\]
write
\[
 S_F=\{t\in K:\text{there exist }y_1,\ldots,y_r\in K
                    \text{ with }F(t,y_1,\ldots,y_r)=0\}.
\]
The variable $T$ is the input; the $r$ variables $Y_i$ are the witnesses.
The degree is the total degree, including $T$.

\begin{theorem}\label{thm:lower-bounds}
If $S_F=\mathcal O_K$, then
\[
                    r\geq 2\qquad\text{and}\qquad\deg F\geq 4.
\]
In particular, a Diophantine definition of $\mathbb Z$ in $\mathbb Q$
by one polynomial needs at least three total variables and has degree at least
four.  Thus a valid quartic construction has the smallest possible degree,
regardless of its number of variables.
\end{theorem}

We prove the two assertions separately.  The variable bound uses the integral
form of Hilbert irreducibility.  The cubic obstruction uses Koll\'ar's theorem
on unirational cubic hypersurfaces; the treatment of cones and reducible
polynomials is included explicitly.

\subsection{One witness and Hilbert irreducibility}

We use the following standard form of Hilbert irreducibility.  Given finitely
many irreducible polynomials $f_i(T,Y)\in K(T)[Y]$, each of positive degree in
$Y$, and finitely many forbidden values of $T$, there are infinitely many
$a\in\mathcal O_K$ outside those values for which all the specializations are
defined, preserve their $Y$-degrees, and remain irreducible in $K[Y]$.
See \cite[Chapter~4, Section~4]{lbSchinzel2000}; the integral number-field
form is also recalled in \cite[Introduction]{lbMuller2001}.
The restriction to integral specializations is essential here: the
Hilbertianity of $K$ alone, without that strengthening, is not the assertion
being applied.

\begin{proposition}\label{prop:one-witness}
Let $G\in K[T,Y]$.  If $S_G$ contains all but finitely many elements of
$\mathcal O_K$, then $S_G$ contains all but finitely many elements of $K$.
\end{proposition}
\begin{proof}
The conclusion is immediate if $G=0$.  If $G\neq0$ and has degree zero in
$Y$, then $S_G$ is finite, contrary to the hypothesis.  Otherwise factor in
$K(T)[Y]$:
\[
                G=c(T)\prod_{i=1}^{s}f_i(T,Y)^{e_i},
\]
where $c(T)\neq0$, the $e_i$ are positive integers, and the $f_i$ are monic
irreducible polynomials of positive $Y$-degree.
If some $f_i$ is linear, it has a root $u(T)\in K(T)$.  The identity
$G(T,u(T))=0$ then specializes for all but finitely many $t\in K$,
proving the conclusion.

If every $f_i$ has $Y$-degree at least two, integral Hilbert irreducibility
provides infinitely many $a\in\mathcal O_K$ at which $c(a)\neq0$ and all
$f_i(a,Y)$ remain irreducible of degree at least two.  None has a root in
$K$, so $G(a,Y)$ has no root in $K$.  This contradicts the hypothesis.
\end{proof}

Since $K\setminus\mathcal O_K$ is infinite, the proposition rules out one
witness.  For example, $1/n\notin\mathcal O_K$ for every integer $n\ge2$,
because a rational algebraic integer is an integer.  With no witness a
nonzero polynomial in $T$ has only finitely many roots, whereas the zero
polynomial defines all of $K$.  This proves the variable assertion of
Theorem~\ref{thm:lower-bounds}.

\subsection{A topological obstruction to rational parametrizations}

Let $r_1$ and $r_2$ denote the numbers of real embeddings and conjugate pairs
of complex embeddings of $K$.  Put
\[
 K_{\infty}=\mathbb R^{r_1}\times\mathbb C^{r_2},
\]
viewed as a real vector space.  The diagonal embedding
$\iota:K\longrightarrow K_{\infty}$ has dense image, while
$\iota(\mathcal O_K)$ is a discrete lattice.  Indeed, an integral basis of
$K$ identifies these images with the images of $\mathbb Q^{[K:\mathbb Q]}$
and $\mathbb Z^{[K:\mathbb Q]}$ under an invertible real linear map.

\begin{lemma}\label{lem:rational-function-integral}
Let $R\in K(U_1,\ldots,U_m)$ be nonconstant, and let $U$ be a nonempty
Zariski-open subset of $\mathbb A^m_K$ on which $R$ is defined.
Then $R(U(K))$ is not contained in $\mathcal O_K$.
\end{lemma}
\begin{proof}
Assume to the contrary that all these values lie in $\mathcal O_K$.
After shrinking $U$ we may take $U$ to be a principal open subset, defined by
one nonzero polynomial.  Apply every archimedean embedding of $K$ to its
coefficients and to those of $R$.  This gives a continuous map, on a nonempty
Euclidean open subset of $K_{\infty}^{m}$,
\[
 R_{\infty}:\prod_{v\mid\infty}K_v^m
       \longrightarrow\prod_{v\mid\infty}K_v=K_{\infty}.
\]
Choose a small connected open ball $B$ in its domain.  The diagonal image of
$K^m$ is dense in $B$.  By hypothesis its image under $R_{\infty}$ lies in
the closed discrete set $\iota(\mathcal O_K)$; by continuity the same is true
on all of $B$.  A continuous map from a connected set to a discrete set is
constant.  Each archimedean component of $R$ is therefore constant on an
open set, so its numerator minus that constant times its denominator is the
zero polynomial.  This forces $R$ to be constant, a contradiction.
\end{proof}

For $K=\mathbb Q$, this argument is simply the density of $\mathbb Q$ in
$\mathbb R$ and the discreteness of $\mathbb Z$.  For a general number field
one must use all archimedean embeddings together; $\mathcal O_K$ need not be
discrete in one selected embedding.

An integral variety $X/K$ is called \emph{unirational over $K$} if some
projective space admits a dominant rational map to $X$, defined over $K$.

\begin{corollary}\label{cor:unirational-integral}
If $X$ is unirational over $K$ and $f\in K(X)$ is nonconstant, then the values
of $f$ at the $K$-points where it is defined cannot all belong to
$\mathcal O_K$.
\end{corollary}
\begin{proof}
Compose $f$ with a dominant rational parametrization of $X$, and restrict
the source to a suitable nonempty open subset of affine space.  Dominance
makes the resulting rational function nonconstant.  Apply
Lemma~\ref{lem:rational-function-integral}.
\end{proof}

\subsection{Quadratic polynomials}

\begin{lemma}\label{lem:quadratic-integral}
If $\deg F\le2$ and $S_F\subseteq\mathcal O_K$, then $S_F$ has at most two
elements.
\end{lemma}
\begin{proof}
Suppose that $a,b,c\in K^{r+1}$ are zeros of $F$ with three distinct input
coordinates.  If these points are collinear, the restriction of $F$ to their
line has degree at most two and three distinct roots.  Thus the entire
line is contained in $F=0$.  The input is nonconstant on that line and takes
all values in $K$, a contradiction.

Otherwise the three points span an affine plane $L$ over $K$.  The input
coordinate is a nonconstant affine linear function on $L$.  If
$F|_L=0$, a nonvertical line in $L$ again gives a contradiction.  If
$F|_L$ factors into affine linear factors over $K$, one of its at most two
lines contains two of the selected points, hence has nonconstant input;
this gives the same contradiction.  A nonzero linear polynomial cannot
vanish at three noncollinear points.

The remaining case is an irreducible quadratic polynomial on $L$.
It defines an irreducible conic with a smooth $K$-point: a singular conic
has at most one singular point, so the three distinct points cannot all be
singular.  Projection from a smooth $K$-point parametrizes the conic by
rational functions over $K$.  The input is nonconstant on this conic,
for otherwise an irreducible quadratic curve would be contained in an
affine line.  Its composition with the parametrization is a nonconstant
rational function taking only values in $\mathcal O_K$, contrary to
Lemma~\ref{lem:rational-function-integral}.
\end{proof}

\subsection{Cubic cones and the sharp degree obstruction}

We use the following published theorem, in its characteristic-zero case.
\begin{theorem}[Koll\'ar]\label{thm:kollar-input}
Let $K$ be a perfect field and let $X\subseteq\mathbb P^{n+1}_K$ be a
geometrically integral cubic hypersurface of dimension $n\ge2$ which is not a cone.
If $X(K)\neq\varnothing$, then $X$ is unirational over $K$.
\end{theorem}
This is the geometrically integral case of Theorem~2 of
\cite{lbKollar2002}; the theorem there is stated for irreducible cubics.
We handle geometric reducibility separately below, so no convention about
\emph{irreducible} is needed.  Its cone exclusion must not be omitted: a cone over a smooth plane cubic need not be unirational.

For a homogeneous cubic $H$ on a finite-dimensional $K$-vector space $E$,
define its vertex subspace by
\[
 W=\{v\in E:D_vH=0\text{ as a polynomial on }E\},
 \qquad D_vH(x)=\left.\frac{d}{ds}H(x+sv)\right|_{s=0}.
\]
The equations defining $W$ are linear over $K$.  In characteristic zero,
$D_vH=0$ is equivalent to $H(x+sv)=H(x)$ for all $x,s$.
Choosing a basis of $W$ and extending it to a basis of $E$ therefore makes
$H$ independent of the coordinates along $W$.
Consequently $H$ descends to a homogeneous cubic $\overline H$ on $E/W$.
By maximality of $W$, this quotient cubic has no nonzero vertex direction
and is not a cone.  These statements remain true after extending scalars:
the kernel of a matrix over $K$ commutes with field extension.

\begin{proof}[Proof of the degree assertion in Theorem~\ref{thm:lower-bounds}]
Suppose $S_F=\mathcal O_K$ and $\deg F\le3$.
The polynomial $F$ is neither zero nor a nonzero constant.
If $F$ is reducible, each of its irreducible factors has degree at most two.
In that case, the projection is the union of the
projections of finitely many linear or quadratic factors.  Every such
projection is contained in $\mathcal O_K$, hence finite by
Lemma~\ref{lem:quadratic-integral}.  This cannot give $\mathcal O_K$.
The same lemma excludes $\deg F\le2$.
We may therefore suppose that $F$ is irreducible of degree three.

Homogenize $F$ to an irreducible cubic
\[
 H(X_0,X_1,\ldots,X_{r+1})
      =X_0^3 F(X_1/X_0,\ldots,X_{r+1}/X_0).
\]
Here $X_1/X_0$ is the input function.
We first exclude geometric reducibility.  Since the characteristic is zero,
$H$ has no repeated geometric components.  Irreducibility over $K$ makes
the Galois group act transitively on them.  Their degrees are equal and
sum to three; thus if $H$ is geometrically reducible, it is a product of
three conjugate linear forms.  A $K$-rational zero belongs to every one of
these hyperplanes, because it is fixed by Galois.  Consequently the entire
line joining any two $K$-rational zeros lies in all three hyperplanes.
Taking affine zeros with inputs $0$ and $1$ would produce every input in
$K$, a contradiction.  We may therefore assume $H$ is geometrically
integral.

Let $W$ be its vertex subspace.
For any $v\in W$, any affine rational zero $x$ normalized by $x_0=1$,
and any $s\in K$ with $1+s v_0\neq0$, the point $x+sv$ is a zero of $H$.
Its input is
\[
                       \frac{x_1+s v_1}{1+s v_0}.
\]
If $v_1-x_1v_0\neq0$, this is a nonconstant M\"obius function of $s$;
its image omits at most one element of $K$.  It cannot be contained in
$\mathcal O_K$, because $K\setminus\mathcal O_K$ is infinite.
Hence $v_1=x_1v_0$ at every such point.  Apply this to zeros with input
$0$ and input $1$, whose existence follows from $S_F=\mathcal O_K$.
It follows that
\[
                         v_0=v_1=0\qquad(v\in W).
\]

Thus both $X_0$ and $X_1$ descend to independent linear forms on $E/W$.
Choose coordinates on the quotient beginning with these forms.  The
irreducible quotient cubic $\overline H$ is not a cone, and its affine
chart $X_0=1$ has exactly the same input projection as $H$.
Indeed, every rational point on the quotient lifts through a $K$-linear
splitting of $E\to E/W$, preserving $X_0$ and $X_1$, and $H$ is constant
along the fibers.

If the projective quotient hypersurface has dimension at least two,
Theorem~\ref{thm:kollar-input} makes it unirational, since it has a
$K$-point with $X_0=1$.  Its input function $X_1/X_0$ is nonconstant,
as its image contains all of the infinite set $\mathcal O_K$.
Corollary~\ref{cor:unirational-integral} now gives a nonintegral input,
a contradiction.

If the quotient has dimension one, its affine chart is defined by a
polynomial in the input and just one further coordinate.  Its projection
cannot be $\mathcal O_K$ by Proposition~\ref{prop:one-witness}.
If the quotient has dimension zero, the affine chart has only finitely
many points and likewise cannot project onto $\mathcal O_K$.
These cases exhaust the possibilities, proving $\deg F\ge4$.
\end{proof}

\begin{remark}[Scope of the optimality statement]
The conclusion concerns one polynomial equation with all witnesses in the
field, and total degree in all variables.  It does not prohibit a system
of quadratic equations: a sum of their squares over $\mathbb Q$ is quartic.
It also does not determine the smallest possible number of witnesses at
degree four, or the smallest degree among definitions attaining the
unknown minimum number of witnesses.  The unconditional numerical
bounds established here are two witnesses and degree four; no example
meeting both bounds is asserted.
\end{remark}

\subsection{A stronger one-witness obstruction}

For completeness there is a useful strengthening which uses Siegel's
theorem rather than just Hilbert irreducibility.
In its rational-function form, Siegel's theorem says that if $C/K$ is a
smooth projective geometrically integral curve of positive genus and
$f\in K(C)$ is nonconstant, then
\[
              \{P\in C(K):f(P)\text{ is defined and belongs to }
                       \mathcal O_K\}
\]
is finite.  One can derive this from the usual theorem on affine integral
points by taking $C$ minus the poles of $f$: the map to $\mathbb A^1$
defined by $f$ is finite, and after inverting finitely many primes its
coordinate functions are integral over $\mathcal O_{K,S}[f]$.
This reduction is also used in
\cite[Section~2.1]{lbMuller2001}; see \cite[Chapter~8]{lbLang1983}.

\begin{proposition}\label{prop:one-witness-finite}
If $G\in K[T,Y]$ satisfies $S_G\subseteq\mathcal O_K$, then $S_G$ is finite.
Thus no infinite subset of $\mathcal O_K$ is the projection of the
$K$-zeros of a single polynomial with just one witness.
\end{proposition}
\begin{proof}
Factor $G$ and treat its finitely many irreducible factors separately.
Factors depending only on $T$ give finite projections.  For another factor,
let $C$ be its plane curve.  If $C$ is not geometrically irreducible, its
$K$-points lie in the intersections of distinct geometric components,
a finite set.  Otherwise pass to its smooth projective normalization
$\widetilde C$.  Away from the finitely many singular points of $C$,
rational points lift to $\widetilde C(K)$, and the input defines a
nonconstant rational function $t$ on $\widetilde C$.

If $\widetilde C$ has positive genus, Siegel's theorem gives only finitely
many points for which $t$ is integral.  If its genus is zero and it has
any rational point, then $\widetilde C\simeq\mathbb P^1_K$.
The input then becomes a nonconstant rational function on
$\mathbb P^1_K$, and on a nonempty open subset all its values would be
integral, contradicting Lemma~\ref{lem:rational-function-integral}.
If it has no rational point, only the finitely many previously excluded
singular points remain.  Each case gives a finite projection.
\end{proof}

\begin{corollary}\label{cor:cubic-finite}
Let $F\in K[T,Y_1,\ldots,Y_r]$ have total degree at most three.
If $S_F\subseteq\mathcal O_K$, then $S_F$ is finite.
\end{corollary}
\begin{proof}
For an infinite projection the proof of
Theorem~\ref{thm:lower-bounds} applies unchanged: choose any two distinct
input values in the geometric-reducibility argument and in the vertex
elimination.  In the curve case use
Proposition~\ref{prop:one-witness-finite} instead of
Proposition~\ref{prop:one-witness}.  The other cases already give a
contradiction to infinitude.
\end{proof}


% Includable research supplement. Theorem environments and bibliography are
% provided by paper.tex. All variable counts include the input unless stated.

\section{Stronger obstructions for small polynomial presentations}
\label{p2lb:scope}

The unrestricted bounds established in the preceding lower-bound section are two existential
witnesses and total degree four.  The arguments below give stronger bounds
for important classes of presentations.  They do not prove that every
definition needs more than two witnesses.  In particular, the general case
of a sign-changing quartic in three total variables remains unresolved here.

For a number field $K$, write $\mathcal O_K$ for its ring of integers.
If $V\subseteq\mathbb A^n_K$ is an algebraic set with first coordinate $T$,
put
\[
 \pi(V(K))=\{t\in K:\text{some }(t,y_1,\ldots,y_{n-1})\in V(K)\}.
\]
Thus $n$ counts all coordinates, and $n-1$ counts existential witnesses.
All assertions in this supplement are independent of the proposed
arithmetic proof in the repository.

\subsection{Curves and rational parametrizations}

We use the following consequences of Siegel's theorem and the topology
of a number field.  They apply to open subsets, which is important when
using birational maps.

\begin{lemma}\label{p2lb:curves}
Let $V$ be an algebraic set over $K$ of dimension at most one, and let
$f:V\to\mathbb A^1_K$ be a regular function.  If
$f(V(K))\subseteq\mathcal O_K$, then $f(V(K))$ is finite.
The same conclusion holds for a rational function restricted to any
Zariski-open subset on which it is defined.
\end{lemma}
\begin{proof}
There are finitely many irreducible components.  A zero-dimensional
component contributes finitely many points, and a component on which
$f$ is constant contributes at most one value.  On an irreducible curve
which is not geometrically irreducible, every $K$-point lies on the
intersection of two distinct conjugate geometric components.  There
are only finitely many such points.

For a geometrically integral curve, normalization and projective
completion discard or add only finitely many points.  Let $C$ be the
resulting smooth projective curve and let $f\in K(C)$ be nonconstant.
If $C$ has positive genus, Siegel's theorem in its rational-function
form gives only finitely many $P\in C(K)$ for which $f(P)$ is defined
and integral.  One obtains this form from the usual theorem by removing
the poles of $f$ and enlarging a finite set of primes so that the
coordinate functions of that affine curve are integral over
$\mathcal O_{K,S}[f]$; see \cite[Chapter~8]{lbLang1983}.

If $C$ has genus zero and has a $K$-point, it is $K$-isomorphic to
$\mathbb P^1$.  A nonconstant rational function on a nonempty open
subset of $\mathbb P^1(K)$ cannot take only integral values.  Indeed,
embed $K$ diagonally in
$K_\infty=\mathbb R^{r_1}\times\mathbb C^{r_2}$.  The image of $K$
is dense, while that of $\mathcal O_K$ is a discrete lattice.  Continuity
would force the rational function to be constant on a small connected
open ball, and hence identically constant.  If $C(K)$ is empty, only
the previously discarded singular points remain.  These cases prove
the assertion.
\end{proof}

\begin{lemma}\label{p2lb:unirational}
Let $X/K$ be an integral unirational variety, let $U\subseteq X$ be
nonempty and open, and let $f\in K(X)$ be nonconstant and defined on $U$.
Then $f(U(K))$ is not contained in $\mathcal O_K$.
\end{lemma}
\begin{proof}
Compose $f$ with a dominant rational map from projective space to $X$.
After restricting to an affine open set, the composition is a
nonconstant rational function in finitely many variables.  The density
and discreteness argument in the preceding proof, now on
$K_\infty^m$, gives a contradiction if all its values are integral.
\end{proof}

\begin{lemma}\label{p2lb:ruled}
Let $C/K$ be a geometrically integral curve, let $m\geq1$, and let
$X$ be birational over $K$ to $C\times\mathbb P^m$.
Suppose $U\subseteq X$ is nonempty and open and $U(K)$ is Zariski
dense.  A nonconstant $f\in K(X)$, defined on $U$, cannot satisfy
$f(U(K))\subseteq\mathcal O_K$.
\end{lemma}
\begin{proof}
Restrict the birational correspondence to isomorphic nonempty opens
and replace $C$ by its smooth projective model.  We may work on an
open subset of $C\times\mathbb A^m$.  If $f$ is nonconstant in the
affine coordinates over $K(C)$, then outside finitely many points of
$C$ its restriction to a fiber is nonconstant.  Outside another finite
set that fiber meets the chosen open subset in a nonempty open subset
of $\mathbb A^m$.  Density of $U(K)$ supplies a rational base point
outside these exceptional sets.  The restriction to this fiber violates
Lemma~\ref{p2lb:unirational}.

Otherwise $f$ belongs to $K(C)$ and is nonconstant there.  For every
$c\in C(K)$ outside a finite set, the open fiber has a $K$-point,
so $f(c)$ is integral.  Lemma~\ref{p2lb:curves} gives only finitely
many such values, hence only finitely many such $c$, because a
nonconstant map of projective curves has finite fibers.  This
contradicts the density of $U(K)$ in $C\times\mathbb A^m$.
\end{proof}

\begin{remark}
The preceding lemma concerns a variety birational over $K$ to a
product with projective space.  Geometric ruledness after extending
scalars is not sufficient for its proof.  In particular, it must not
be applied without further argument to an arbitrary conic bundle
over $\mathbb P^1_K$ which has no rational section.
\end{remark}

\begin{lemma}\label{p2lb:dense-component}
Let $V/K$ be an algebraic set and $f:V\to\mathbb A^1_K$ a regular
function.  If $f(V(K))$ is infinite and contained in $\mathcal O_K$,
the Zariski closure of $V(K)$ has a geometrically integral component
$X$ of dimension at least two such that $X(K)$ is Zariski dense
and $f|_X$ is nonconstant.
\end{lemma}
\begin{proof}
The closure has finitely many irreducible components, each with dense
$K$-points on the open set obtained by deleting the other components.
If such a component were geometrically reducible, its rational points
would lie in the proper closed intersections of its conjugate
geometric components, a contradiction.  Components on which $f$ is
constant contribute finitely many values, and components of dimension
at most one contribute finitely many by Lemma~\ref{p2lb:curves}.
Thus at least one component has the required properties.
\end{proof}

\subsection{Two quadrics in projective four-space}

The next argument uses only the elementary degree bound
$\deg X\geq1+\operatorname{codim}(X)$ for a nondegenerate projective
integral variety, B\'ezout's theorem, and the following theorem of
Koll\'ar: a geometrically integral cubic hypersurface of dimension at
least two over a perfect field is unirational if it has a rational
point and is not a cone \cite[Theorem~2]{lbKollar2002}.
The elementary degree bound is recalled in
\cite[Introduction]{p2lbEGHP2006}.  We give the projection argument
explicitly in order to include reducible and nonreduced intersections.

\begin{proposition}\label{p2lb:two-quadrics-surface}
Let $Q_1,Q_2\in K[X_0,\ldots,X_4]$ be coprime homogeneous
quadratics, and let $X$ be a geometrically integral surface component
of $V(Q_1,Q_2)$.  Suppose $X(K)$ is Zariski dense.
Then $X$ is either unirational over $K$ or birational over $K$ to
$C\times\mathbb P^1$ for a geometrically integral curve $C/K$.
\end{proposition}
\begin{proof}
The complete intersection $W=V(Q_1,Q_2)\subseteq\mathbb P^4$
has degree four.  More precisely, if $m_X$ is its multiplicity along
$X$, then $m_X\deg X\leq4$.  This follows either from B\'ezout's
theorem with multiplicities or from the complete-intersection Hilbert
series $(1-t^2)^2/(1-t)^5$ and additivity of degree.

If $\deg X\leq2$, the elementary degree bound shows that its linear
span has dimension at most three.  It is therefore a plane or an
integral quadric surface in its span.  Its linear span is defined over
$K$, and density supplies a smooth $K$-point in the quadric case.
Projection from that point makes the quadric rational; a plane is
rational as well.

Suppose $\deg X\geq3$.  The degree inequality forces $m_X=1$.
Since the characteristic is zero, a nonempty open subset of $X$ is
smooth in the complete intersection $W$.  Choose a $K$-point there,
using density, and change coordinates so that it is
$p=(1:0:0:0:0)$.  The quadratics have the form
\[
 Q_1=X_0L_1+q_1,\qquad Q_2=X_0L_2+q_2,
\]
where $L_i$ are linear and $q_i$ are quadratic in $X_1,\ldots,X_4$.
Smoothness of $W$ at $p$ means that $L_1,L_2$ are independent.
The cubic
\[
 H=L_1q_2-L_2q_1
\]
is not identically zero.  Otherwise coprimality of $L_1,L_2$ gives
$q_1=L_1L_3$ and $q_2=L_2L_3$ for a linear form $L_3$, contradicting
coprimality of $Q_1,Q_2$.

The surface $X$ cannot be contained in the plane $L_1=L_2=0$,
because its degree is at least three.  Interchange the two quadratics
if necessary so that $L_1$ does not vanish identically on $X$.
Projection from $p$ then has the rational inverse
\[
 X_0=-q_1/L_1
\]
on the corresponding open set.  Thus $X$ is birational to a
geometrically integral surface component $Y$ of the cubic
$V(H)\subseteq\mathbb P^3$.  The rational points of $Y$ are dense.

If $Y$ is a plane or a quadric, it is rational as above.  If it is
a cubic which is not a cone, Koll\'ar's theorem makes it unirational.
Finally, a geometrically integral cubic surface which is a cone has
a single vertex point defined over $K$.  Indeed the vertex directions
form the kernel of a linear map over $K$ given by the directional
derivatives of its cubic equation; more than one independent vertex
direction would leave a binary cubic, which is geometrically reducible.
The complement of the vertex is birational over $K$ to the product
of its plane cubic base with $\mathbb P^1$.  This proves the result.
\end{proof}

\begin{theorem}\label{p2lb:quadratic-system-four}
Let $q_1,\ldots,q_s\in K[T,Y_1,\ldots,Y_{n-1}]$ have total
degree at most two, where $n\leq4$.  If the rational projection
\[
 S=\{t\in K:\exists\mathbf y\in K^{n-1},\quad
                   q_1(t,\mathbf y)=\cdots=q_s(t,\mathbf y)=0\}
\]
is contained in $\mathcal O_K$, then $S$ is finite.
\end{theorem}
\begin{proof}
Adding unused coordinates reduces the assertion to $n=4$.
Zero residuals can be removed; if all residuals vanish identically,
then $S=K$, contrary to the hypothesis.  A nonzero constant residual
gives the empty set.  Let $V$ be their common zero set.

Suppose $S$ is infinite, and choose $X\subseteq V$ as in
Lemma~\ref{p2lb:dense-component}.  Any common polynomial factor
$g$ of the residuals has degree at most two.  The entire hypersurface
$g=0$ lies in $V$, and its input projection is therefore contained
in $\mathcal O_K$.  The quadratic obstruction proved in the first
report says this projection has at most two elements.  Consequently
$X$ is not contained in $g=0$.  Dividing the residuals by their common
greatest divisor does not remove $X$ and leaves quadratics with
greatest common divisor one.

If their $K$-linear span contains a nonzero affine linear polynomial,
restrict to its affine hyperplane, isomorphic to $\mathbb A^3_K$.
Not all restricted residuals vanish identically, since that would
give the original quadratics a common linear factor.  The resulting
surface $X$ is a hypersurface component of a nonzero quadratic in
three variables, hence is a plane or a quadric.  Density gives a
rational parametrization, contradicting Lemma~\ref{p2lb:unirational}.

Otherwise choose two coprime quadratics in the linear span of the
residuals.  To see that this is possible, first choose a nonzero
quadratic $q$ and then a linear combination not divisible by any
irreducible factor of $q$; finitely many proper linear subspaces
cannot cover a vector space over the infinite field $K$.
Their affine common zero set has dimension two.  Thus $X$ is a
surface component.  Their degree-two homogenizations remain coprime,
and the projective closure of $X$ is a component of their intersection
in $\mathbb P^4$.

Proposition~\ref{p2lb:two-quadrics-surface} makes this component
unirational or birational to a product of a curve and $\mathbb P^1$.
The input is a nonconstant rational function and is integral on
the rational points of its affine open subset.  This contradicts
Lemma~\ref{p2lb:unirational} or Lemma~\ref{p2lb:ruled}.
\end{proof}

\begin{corollary}\label{p2lb:sos-five}
Suppose
\[
 P(T,\mathbf Y)=\sum_{j=1}^s c_jq_j(T,\mathbf Y)^2,
 \qquad c_j\in\mathbb Q_{>0},\quad\deg q_j\leq2.
\]
If the rational existential projection of $P=0$ is $\mathbb Z$,
then $P$ uses at least five total variables, hence at least four
existential witnesses.
The same lower bound holds over a number field for any quartic
obtained by applying an anisotropic quadratic form to quadratic
residuals and defining its ring of integers.
\end{corollary}
\begin{proof}
In the rational case, the sum is zero exactly when all $q_j$ vanish.
For an anisotropic combining form the same equivalence holds by
definition of anisotropy.  Apply Theorem~\ref{p2lb:quadratic-system-four}.
\end{proof}

This improves the unrestricted three-total-variable bound for the
sum-of-squares quartic architecture used by the released polynomial
and all the expanded quartic variants constructed above.  It does
not rule out a sign-changing quartic in three or four total variables.

\subsection{Nonnegative polynomials in three variables}

\begin{proposition}\label{p2lb:nonnegative-squarefree}
Let $P\in\mathbb Q[T,Y,Z]$ be squarefree and nonnegative on
$\mathbb R^3$.  If its rational input projection is contained in
$\mathbb Z$, that projection is finite.
\end{proposition}
\begin{proof}
Every real zero of a differentiable nonnegative function is a local
minimum, so all three partial derivatives of $P$ vanish there.
The singular locus of a squarefree hypersurface in characteristic
zero has codimension at least one in that hypersurface.  In this
case it has dimension at most one.  All rational zeros lie in this
singular locus.  Apply Lemma~\ref{p2lb:curves} to its input function.
\end{proof}

\begin{theorem}\label{p2lb:nonnegative-degree-eight}
Let $P\in\mathbb Q[T,Y,Z]$ be nonnegative on $\mathbb R^3$.
If its rational input projection is an infinite subset of $\mathbb Z$,
then $P$ has an irreducible factor of degree at least four and
multiplicity at least two.  In particular $\deg P\geq8$.
\end{theorem}
\begin{proof}
Choose a component $X$ of the closure of its rational zeros as in
Lemma~\ref{p2lb:dense-component}.  It has dimension two, so its prime
ideal is generated by an irreducible factor $g$ of $P$.  Write
$P=g^m h$ with $g\nmid h$.  Density supplies a rational point on
$X$ which is smooth on $g=0$ and does not lie on $h=0$.
If $m=1$, then $\nabla P=h\nabla g$ is nonzero there, contrary
to nonnegativity of $P$.  Thus $m\geq2$.

The projection of the whole hypersurface $g=0$ is contained in
the projection of $P=0$, and is infinite because the input is
nonconstant on $X$ and $X(\mathbb Q)$ is dense.  The unconditional
degree-three obstruction of the preceding lower-bound section gives $\deg g\geq4$.
Therefore $\deg P\geq m\deg g\geq8$.
\end{proof}

Nonnegativity only on rational inputs gives the same result, because
$\mathbb Q^3$ is dense in $\mathbb R^3$ and $P$ is continuous.
The degree-eight conclusion is restricted to nonnegative polynomials;
it is not an unconditional degree-eight lower bound for two witnesses.

\subsection{Further necessary conditions}

\begin{proposition}\label{p2lb:linear-coordinate}
If a polynomial $P(T,Y_1,\ldots,Y_r)\in K[T,\mathbf Y]$ defines
$\mathcal O_K$ in $K$, its degree in each witness which occurs is
at least two.  Its degree in $T$ is also at least two.
\end{proposition}
\begin{proof}
Suppose $P=AY_r+B$, where $A\ne0$ is independent of $Y_r$.
For all but finitely many $t\in K$, the specialization
$A(t,Y_1,\ldots,Y_{r-1})$ is a nonzero polynomial.  Since $K$ is
infinite, choose the other witnesses so that it does not vanish and
then solve uniquely for $Y_r$.  The projection is cofinite in $K$,
which is incompatible with $\mathcal O_K$.

If $P=AT+B$ with $A\ne0$ independent of $T$, its projection includes
the values of $-B/A$ on the nonempty open set $A\ne0$.  Integrality
of all these values forces this rational function to be constant,
by the density and discreteness argument.  Thus $B=-cA$ for some
$c\in K$, and $P=A(T-c)$.  If $A$ has a $K$-zero, every input is
accepted; otherwise only $c$ is accepted.  Neither is $\mathcal O_K$.
If $P$ is independent of $T$, its projection is empty or all of $K$.
\end{proof}

\begin{proposition}\label{p2lb:quartic-surface-necessary}
Suppose a quartic $P\in\mathbb Q[T,Y,Z]$ defines $\mathbb Z$ in
$\mathbb Q$.  Then $P$ is irreducible and geometrically irreducible,
its rational points are Zariski dense on its affine surface, and
$P$ takes both positive and negative values on $\mathbb R^3$.
The input function is nonconstant on that surface.  Every curve
birational over $\mathbb Q$ to $\mathbb P^1$ on which the input is defined is vertical
(the input is constant), and the same holds for every geometrically
integral genus-one curve on the surface with infinitely many rational
points.
\end{proposition}
\begin{proof}
If $P$ factors nontrivially, every factor has degree at most three.
Each factor's rational projection is contained in $\mathbb Z$ and
is finite by the earlier degree-three obstruction.  Their finite
union cannot be $\mathbb Z$.  If the irreducible $P$ is geometrically
reducible, its rational points lie in intersections of distinct
conjugate components, of dimension at most one; Lemma~\ref{p2lb:curves}
again gives a finite projection.  The same lemma shows that a proper
Zariski closure of the rational points would give a finite projection.
Thus the surface is geometrically integral and its rational points
are dense.  The input is nonconstant because its image is infinite.

If $P$ had only one sign, apply
Theorem~\ref{p2lb:nonnegative-degree-eight} to $P$ or $-P$.
Finally, a nonvertical rational curve with a rational point would
violate Lemma~\ref{p2lb:unirational}.  On a genus-one curve with
infinitely many rational points a nonconstant input function would
take only finitely many integral values at rational points by
Siegel's theorem, whereas every rational point where the input is
defined would have to give such a value.  This is impossible.
\end{proof}

\subsection{Why the unrestricted two-witness case remains difficult}

A smooth quartic surface in $\mathbb P^3$ has trivial canonical
bundle and is a K3 surface.  It is not covered by the unirational
cubic theorem.  Zariski density of rational points, even when known,
does not itself imply density in the real topology or the existence
of a nonintegral value of a specified rational function.  Likewise,
potential density after a finite field extension cannot be substituted
for a statement about rational points over $\mathbb Q$.
For examples and the distinction between density over the ground field
and potential density, see \cite{p2lbHarrisTschinkel2000}.

The same limitation appears naturally in the improved bound for
quadratic systems.  A smooth complete intersection of three quadrics
in $\mathbb P^5$ is a K3 surface: adjunction gives
$K_X=\mathcal O_X(2+2+2-6)=\mathcal O_X$, and the
complete-intersection cohomology gives $H^1(X,\mathcal O_X)=0$.
Thus five affine coordinates already permit a class of surfaces not
handled by the two-quadric argument.  This observation is a limit of
the proof, not an existence claim for a five-variable definition.

\begin{proposition}[Conditional obstruction for general-type surfaces]
\label{p2lb:lang-surface}
Assume the weak Bombieri--Lang conjecture for surfaces over $\mathbb Q$:
the rational points of a smooth projective surface of general type are
not Zariski dense.  Then no polynomial in three total variables whose
geometrically integral projective surface has a smooth projective
model of general type can have an infinite rational input projection
contained in $\mathbb Z$.
In particular, under this conjecture a polynomial in three variables
with smooth projective closure of degree at least five cannot define
$\mathbb Z$.
\end{proposition}
\begin{proof}
Lemma~\ref{p2lb:curves} forces Zariski density of rational points
on a surface with such a projection.  Passing to isomorphic dense
opens of a smooth projective model preserves density, contradicting
the conjecture.  A smooth degree-$d$ surface in $\mathbb P^3$ has
canonical bundle $\mathcal O_X(d-4)$, which is ample when $d\geq5$.
\end{proof}

Mazur's conjecture about finitely many real connected components of
the closure of rational points would rule out every Diophantine
definition of $\mathbb Z$ in $\mathbb Q$, in any number of variables.
It therefore cannot supply a conditional numerical improvement
compatible with simultaneously assuming the proposed arithmetic
theorem.  Proposition~\ref{p2lb:lang-surface}, by contrast, is a
restriction on one class of possible two-witness surfaces.


% Includable supplement. Independent of the proposed arithmetic construction.
\section{A restriction on the degree in two witnesses}
\label{sec:p2lb-conics}

The total degree can conceal an important distinction: a polynomial may
have arbitrarily high degree in the input, while remaining quadratic in
its witnesses.  With only two witnesses, such a polynomial cannot define
the ring of integers.  We prove a simultaneous version which also rules
out a finite disjunction of such equations.

\begin{theorem}\label{p2lb:quadratic-witnesses}
Let $K$ be a number field and let
$F_1,\ldots,F_s\in K[T,Y,Z]$ have total degree at most two in the
pair $(Y,Z)$, with no restriction on their degree in $T$.  Put
\[
 S=\{t\in K:\text{for some }i\text{ and }y,z\in K,
                         \ F_i(t,y,z)=0\}.
\]
Exactly one of the following alternatives holds:
\begin{enumerate}
\item at least one equation has a solution
$(Y,Z)\in K(T)^2$, and $S$ is cofinite in $K$;
\item no equation has such a solution, and $K\setminus S$ contains
infinitely many elements of every coset $a+\mathfrak a$, where
$a\in\mathcal O_K$ and $\mathfrak a$ is a nonzero ideal of
$\mathcal O_K$.
\end{enumerate}
In particular, no single polynomial quadratic in two witnesses, or
finite disjunction of such polynomials, defines $\mathcal O_K$ in $K$.
\end{theorem}

The cofinite assertion in the first alternative follows by specializing
the rational functions in a generic solution away from their finitely
many poles.  The substantive issue is to find integral specializations
which simultaneously preserve the obstructions in the second alternative.
We use Faddeev's residue sequence, Chebotarev's theorem and integral
Hilbert irreducibility.  The argument is a consequence of these classical
tools; no claim of historical priority is made.

\subsection{Quaternion obstructions at integral specializations}

For a field $F$ of characteristic different from two, $(a,b)$ denotes the
Brauer class of the quaternion algebra generated by $i,j$ with
$i^2=a$, $j^2=b$ and $ij=-ji$.  A smooth conic is split, or equivalently
has an $F$-point, precisely when its associated quaternion class is zero.
For this correspondence and the residue sequence below see
\cite[Theorem~5.1 and Section~6.1]{p2lbBergEtAl2024}.

We recall explicitly the residue convention.  At the closed point of
$\mathbb A^1_K$ given by a monic irreducible $f(T)$, write
\[
 a=f^r u,\qquad b=f^s v,
\]
where $u,v$ are units in $K[T]_{(f)}$.  If $\theta$ is a root of $f$
and $L=K(\theta)$, the residue of $(a,b)$ is the square class of
\[
             (-1)^{rs}u(\theta)^s/v(\theta)^r
             \quad\text{in }L^\times/L^{\times2}.
\]
Faddeev's theorem implies that a class in $\operatorname{Br}K(T)$
with all finite residues zero comes from $\operatorname{Br}K$.
Indeed the sum-of-corestrictions relation first makes the residue at
infinity zero, since its residue field is $K$; exactness then gives a
constant class.  In characteristic zero the sequence applies to the
whole Brauer group because $\operatorname{Br}\overline K(T)=0$.

We shall also use the following elementary local formula.  If $F$ is a
nonarchimedean local field with odd residue field $\kappa$, if $\varpi$
is a uniformizer, and if $u,v$ are units, the quaternion
$(\varpi^r u,\varpi^s v)$ is split exactly when
\[
                    (-1)^{rs}\overline u^{\,s}/\overline v^{\,r}
                    \quad\text{is a square in }\kappa^	imes.
\]
One obtains this by bilinearity of the Hilbert symbol: two unit
parameters give a split algebra, $(\varpi,u)$ is detected by the square
class of $\overline u$, and $(\varpi,\varpi)=(\varpi,-1)$.
These facts also follow from the norm criterion for the unramified
quadratic extension; see \cite[Proposition~5.5]{p2lbVoight2019}.

\begin{lemma}\label{p2lb:brauer-progression}
Let $\alpha\in\operatorname{Br}K(T)$ be a nonzero quaternion class.
If $\alpha$ is constant, every specialization outside a finite set
is nonzero.  If $\alpha$ is not constant, then outside any prescribed
finite set of primes of $K$ there is a prime $\mathfrak p$ and a
residue class $c\bmod\mathfrak p^2$ such that, outside a finite set
of inputs, every $t\in\mathcal O_K$ in that residue class has
\[
                   \alpha(t)\ne0\quad\text{in }
                   \operatorname{Br}K_{\mathfrak p}.
\]
\end{lemma}
\begin{proof}
The assertion for a constant class follows from specialization on an
open set where its equality with the chosen quaternion presentation is
defined.  Otherwise choose a nonzero finite residue using Faddeev's
theorem.  With the notation just introduced, let
\[
 d=(-1)^{rs}u(\theta)^s/v(\theta)^r\in L^\times
\]
be its nonsquare representative.  Choose a finite Galois extension
$M/K$ containing $L(\sqrt d)$, and an element
$g\in\operatorname{Gal}(M/L)$ which sends $\sqrt d$ to $-\sqrt d$.
Such an element exists because $d$ is not a square in $L$.

Chebotarev's theorem gives infinitely many primes of $K$ with Frobenius
conjugate to $g$; see \cite[Section~2.1, Theorem~1]{p2lbSerre1981}.
Choose one outside the prescribed finite set and all further primes of
bad reduction for the displayed algebraic data, including primes above
two.  Choose a prime of $M$ whose Frobenius is exactly $g$, and let
$\mathfrak q$ be its restriction to $L$.  Since $g$ fixes $L$, this
prime of $L$ has residue degree one over $\mathfrak p$.  Since $g$
acts nontrivially on $\sqrt d$, the reduction of $d$ at
$\mathfrak q$ is a nonsquare.  The excluded bad primes ensure that
it is a nonzero unit.

The corresponding completion of $L$ is $K_{\mathfrak p}$, so
$f$ has a simple root $\theta_{\mathfrak p}$ in
$\mathcal O_{K_{\mathfrak p}}$.  Choose a residue class
$c\bmod\mathfrak p^2$ for which
\[
 v_{\mathfrak p}(t-\theta_{\mathfrak p})=1
 \qquad(t\equiv c\bmod\mathfrak p^2).
\]
There is such a class because the reduction map from $\mathcal O_K$
onto $\mathcal O_{K_{\mathfrak p}}/\mathfrak p^2$ is surjective.
For all these inputs, $v_{\mathfrak p}(f(t))=1$, and $u(t),v(t)$
are units whose reductions equal those at $\theta$ and
$\mathfrak q$.  The local formula therefore makes the specialized
quaternion nonsplit, since its residue square class is that of $d$.
The finitely many inputs where a rational presentation is undefined may
be discarded.  This proves the assertion.
\end{proof}

\begin{corollary}\label{p2lb:simultaneous-conics}
Let $C_1,\ldots,C_m$ be smooth conics over $K(T)$ with no
$K(T)$-point.  Given any coset $a+\mathfrak a$ as in
Theorem~\ref{p2lb:quadratic-witnesses}, it contains a smaller nonempty
coset of a nonzero ideal such that, outside finitely many inputs, none
of the specialized conics has a $K$-point.
\end{corollary}
\begin{proof}
Constant nonzero quaternion classes already exclude every good
specialization.  For each remaining conic, use the preceding lemma
with a fresh prime, distinct from all previously chosen primes and
coprime to $\mathfrak a$.  The Chinese remainder theorem simultaneously
imposes its bad residue class modulo the square of that prime and the
original coset.  Every resulting good specialization is locally insoluble
at the selected prime for each corresponding conic.
\end{proof}

\subsection{Degenerate conics and the simultaneous theorem}

\begin{proof}[Proof of Theorem~\ref{p2lb:quadratic-witnesses}]
Write $F=K(T)$.  If one equation has a point in $F^2$, its rational
functions specialize to a point at every input outside a finite set.
This gives the first alternative.  Suppose henceforth that none does.
We classify each equation over $F$, and discard finitely many inputs
so that all rational coordinate changes, ranks and nonzero coefficients
in this classification specialize correctly.

An equation of witness degree zero is a nonzero element of $F$, and
has no point at any good input.  Witness degree one always gives an
affine $F$-point, so cannot occur in the case under consideration.
For witness degree two, homogenize in $Y,Z$ with a third coordinate
$W$, giving a ternary quadratic form and its projective conic $C$.
If the form has rank three and its quaternion class is zero, then
$C\simeq\mathbb P^1_F$.  The affine open $W\ne0$ is nonempty
and hence has an $F$-point, a contradiction.  Thus every rank-three
case has a nonzero quaternion class, and
Corollary~\ref{p2lb:simultaneous-conics} simultaneously excludes all
of them on a smaller coset of the given integral progression.

A rank-one ternary form is a scalar times the square of an
$F$-linear form.  That line cannot be the line at infinity, since the
original polynomial has witness degree two.  It therefore has an
affine $F$-point, again a contradiction.  A split rank-two form is the
product of two distinct $F$-linear forms.  At least one line is not
the line at infinity, and supplies the same contradiction.

The remaining possibility is a nonsplit rank-two form.  After a linear
change over $F$ it is a nonzero scalar times
\[
                         U^2-d(T)V^2,
                         \qquad d(T)\notin F^{\times2}.
\]
Its only projective $F$-point is its radical point, defined by
$U=V=0$.  If that point were affine, it would give an $F$-point of
the original equation.  Therefore its original $W$ coordinate is
zero.  At every good specialization for which $d(t)$ remains a
nonsquare, the same reasoning leaves only this radical point, still
at infinity, and the specialized affine equation has no solution.

There are only finitely many nonsquare functions $d(T)$ to treat.
The smaller coset already obtained contains a progression
$t=c+bz$ with $c\in\mathcal O_K$, $0\ne b\in\mathcal O_K$
and $z\in\mathcal O_K$: choose $b$ in its defining nonzero ideal.
The polynomials
\[
                         Y^2-d(c+bZ)
\]
remain irreducible over $K(Z)$, because $T\mapsto c+bZ$ is an
isomorphism of rational function fields.  The integral Hilbert
irreducibility theorem, simultaneously for this finite list, gives
infinitely many integral $z$ for which all these specializations
remain irreducible and all previously excluded finite conditions are
avoided.  See the integral form recalled in the original lower-bound
section and \cite[Introduction]{lbMuller2001}.  If the list is empty,
any infinitely many good inputs in the coset suffice.

The resulting inputs lie in the original coset $a+\mathfrak a$,
and none admits a solution of any equation.  This proves the second
alternative and the theorem.
\end{proof}

\begin{corollary}\label{p2lb:two-witness-degree-three}
If $P(T,Y,Z)$ defines $\mathbb Z$ in $\mathbb Q$, then its total
degree in the two witnesses $Y,Z$ is at least three.  More generally,
after factoring over $\mathbb Q(T)$, at least one factor has total
witness degree at least three.  This remains true after every
invertible affine change of the witness coordinates over
$\mathbb Q(T)$, provided the transformed equation and inverse change
are defined at all but finitely many input values.
\end{corollary}
\begin{proof}
The first assertion is the theorem with one equation.  If all factors
had witness degree at most two, clear their rational-function
denominators.  Outside finitely many inputs, the original zero condition
is their finite disjunction, contradicting the theorem.  An invertible
affine change over $\mathbb Q(T)$ preserves the rational solution sets
at every good specialization, so the same argument applies after it.
\end{proof}

\begin{corollary}\label{p2lb:conic-integral-superset}
Let $S\subset K$ contain $\mathcal O_K$, and suppose that $K\setminus S$
is infinite.  Then $S$ is not a finite union of projections of equations
in two witnesses which are jointly quadratic in those witnesses, even if
the degree in the input is unrestricted.
\end{corollary}
\begin{proof}
The second alternative of Theorem~\ref{p2lb:quadratic-witnesses} would
omit infinitely many elements of $\mathcal O_K$, while the first would
make $K\setminus S$ finite.  Both contradict the stated assumptions.
\end{proof}

For example, take
\[
 S=\mathbb Z_{(2)}\cap\mathbb Z_{(3)}
   =\{a/b\in\mathbb Q:\gcd(a,b)=1,\ b>0,\ 2\nmid b,\ 3\nmid b\}.
\]
This set contains $\mathbb Z$ and has infinite complement, for instance
the distinct elements $2^{-n}$ for $n\ge1$.  Thus a quadratic
three-witness definition of this semilocal ring cannot be replaced by a
quadratic two-witness definition, or a finite union of such definitions.
The obstruction concerns joint witness degree; it allows arbitrary input
degree and does not assert a lower bound for equations of higher witness
degree.

This is an unconditional restriction on a possible two-witness
definition at every input degree.  It does not exclude general quartic
surfaces, whose generic fibers over the input can have positive genus.
It also does not extend by this proof to quadratic equations in three
witnesses: their fibers are quadric surfaces, and are not controlled by
a single quaternion obstruction of the kind used above.


% Includable phase-three supplement. All results here are unconditional.
\section{Lower bounds for systems of polynomial equations}
\label{p3sys:lower-bounds}

Let $K$ be a number field, and let $\mathcal O_K$ be its ring of
integers.  A polynomial system in $n$ total coordinates has the form
\[
 F_1(T,Y_1,\ldots,Y_{n-1})=\cdots
   =F_s(T,Y_1,\ldots,Y_{n-1})=0.
\]
Its degree is $\max_i\deg F_i$, including the input $T$.
Allowing several equations does not change the number of quantified
coordinates.  We also allow finitely many nonzero conditions when
stating obstructions: these specify a Zariski open subset of the common
zero set, and strengthen rather than weaken the conclusions for closed
systems.  A finite union of such presentations is treated componentwise.

The main new restriction is that \emph{two quadratic equations cannot
suffice in any number of coordinates}.  The proof includes singular,
reducible, and nonreduced intersections, rather than assuming a smooth
complete intersection.

\subsection{Integral-valued functions on simple varieties}

\begin{lemma}\label{p3sys:unirational-function}
Let $X$ be an integral $K$-variety which is $K$-unirational, let
$U\subset X$ be nonempty and open, and let $f:U\to\mathbb A^1_K$
be a regular function.  If $f(U(K))\subseteq\mathcal O_K$, then
$f$ is constant.  In particular, its image is finite.
\end{lemma}
\begin{proof}
A dominant rational parametrization reduces the assertion to a rational
function on a nonempty open subset of affine space.  If this function
were nonconstant, one of its partial derivatives would be nonzero.
Choose a $K$-point at which that derivative and every relevant
denominator are nonzero.  Such a point exists because $K$ is infinite.

Put $K_\infty=\prod_{v\mid\infty}K_v$, with one factor for each real
place and one complex factor for each pair of complex embeddings.
The diagonal copy of $K$ is dense in $K_\infty$: a $\mathbb Q$-basis
of $K$ identifies this assertion with the density of
$\mathbb Q^{[K:\mathbb Q]}$ in its real vector space.  In contrast,
$\mathcal O_K$ is a lattice, hence is closed and discrete there.
The inverse-function theorem at the chosen point, at every
archimedean place, gives a nonempty open subset of $K_\infty$ in
the local image of the rational function.  Density of $K$ makes its
$K$-values dense in this open subset.  They cannot all belong to the
discrete set $\mathcal O_K$.  This is the required contradiction.
\end{proof}

\begin{lemma}\label{p3sys:curve-function}
Let $X$ be an algebraic set over $K$ of dimension at most one,
let $U\subseteq X$ be open, and let $f:U\to\mathbb A^1_K$ be
regular.  If $f(U(K))\subseteq\mathcal O_K$, then its image is
finite.
\end{lemma}
\begin{proof}
Reduce to the finitely many irreducible components.  Zero-dimensional
components give finite sets.  If an irreducible curve is not geometrically
integral, all its $K$-points lie in the intersections of distinct
geometric components, a finite set.  Otherwise take its smooth
projective normalization $C$.  Except for finitely many points, rational
points lift to $C(K)$, and $f$ becomes a rational function on $C$.
There is nothing to prove when this function is constant.

If $C$ has genus zero and a $K$-point, it is isomorphic to
$\mathbb P^1_K$, and Lemma~\ref{p3sys:unirational-function} excludes
a nonconstant integral-valued function on its permitted open subset.
If it has no $K$-point, only the finitely many excluded points remain.
If the genus is positive, Siegel's theorem makes the set of $K$-points
where a fixed nonconstant rational function is integral finite.
For this last formulation, remove the poles of the function: the
resulting affine curve is finite over the affine line.  After inverting
a fixed finite set of primes, its coordinate ring is integral over
the ring generated by that function.  Integral function values
therefore give points integral outside this fixed set, to which
Siegel's theorem applies; see \cite[Chapter~8]{lbLang1983}.
\end{proof}

\begin{lemma}\label{p3sys:cubic-function}
Let $X\subseteq\mathbb P^N_K$ be a hypersurface defined by a
nonzero homogeneous polynomial of degree at most three.  For an open
subset $U\subseteq X$ and a regular function
$f:U\to\mathbb A^1_K$, the containment
$f(U(K))\subseteq\mathcal O_K$ implies that $f(U(K))$ is finite.
The same conclusion holds for $X=\mathbb P^N_K$.
\end{lemma}
\begin{proof}
The assertion for projective space follows from
Lemma~\ref{p3sys:unirational-function}.  Factor the defining
polynomial over $K$ and treat each irreducible factor separately;
multiplicities do not change the rational point set.  If such a
factor of degree two or three is not geometrically integral, its
geometric factors are conjugate linear forms, since these degrees
are prime.  Every $K$-point lies in their common intersection.
That intersection is a $K$-linear subspace, so the preceding lemma
for unirational varieties applies.  Linear factors are handled in
the same way.

A geometrically integral quadric with a smooth $K$-point is rational
by projection from that point.  If it has no smooth $K$-point, all
its $K$-points lie in its linear vertex.  Both cases give the result.
For a geometrically integral cubic of dimension at most one, use
Lemma~\ref{p3sys:curve-function}.  In dimension at least two,
Koll\'ar's theorem says that a cubic which is not a cone and has
a $K$-point is $K$-unirational
\cite[Theorem~2]{lbKollar2002}.  A number field is perfect,
as required in that theorem.

It remains to treat a cubic cone.  Its maximal linear vertex is
defined over $K$: in characteristic zero it is the projectivization
of the vector subspace of directions along which its defining form
is constant.  Quotient by these directions and choose a complementary
$K$-linear space.  This gives a geometrically integral cubic base
$B$ which is not a cone.  Its dimension is at least one; a
geometrically integral zero-dimensional cubic would be impossible.
Outside the vertex, projection to $B$ has affine-space fibers, whose
projective closures are the linear spaces joining a point of $B$
to the vertex.  The cone is birational to
$B\times\mathbb P^{r+1}$ if the vertex has dimension $r$.

If $B(K)$ is empty, all rational points of the cone lie in the
vertex.  If $B(K)$ is finite, each permitted fiber is an open
subset of a linear space; the function is constant on it by
Lemma~\ref{p3sys:unirational-function}.  Together with the vertex
this gives a finite image.  If $\dim B\ge2$ and $B(K)$ is
nonempty, Koll\'ar's theorem makes $B$, and hence the cone,
unirational.

The only remaining case is a cubic curve $B$ with infinitely many
$K$-points.  Pass to its normalization $C$, discarding finitely
many exceptional base points.  On a dense open of the cone we may
write $f$ as a rational function on
$C\times\mathbb P^{r+1}$.  If it depends nontrivially on the
projective-space coordinates, its specialization is nonconstant
for all base points outside a finite subset of $C$.  A rational
base point outside that subset then contradicts
Lemma~\ref{p3sys:unirational-function} on the fiber.  Thus
$f\in K(C)$.  For all but finitely many points of $C(K)$, the
permitted fiber is a nonempty open subset of projective space and
has a $K$-point.  Consequently $f$ is integral at all these base
points.  Lemma~\ref{p3sys:curve-function} forces $f$ to be
constant, because a nonconstant function on a curve has finite
fibers and $C(K)$ is infinite.  Being constant in the function
field, it is constant wherever it is regular on the irreducible
cone.  This also treats all points omitted from the birational
chart and completes the proof.
\end{proof}

\subsection{Two quadrics in any dimension}

\begin{theorem}\label{p3sys:two-quadrics-any-dimension}
Let $X\subseteq\mathbb P^N_K$ be the common zero set of at most
two homogeneous quadratic forms, let $U\subseteq X$ be open,
and let $f:U\to\mathbb A^1_K$ be regular.  If
$f(U(K))\subseteq\mathcal O_K$, then $f(U(K))$ is finite.
No smoothness, irreducibility, or reducedness assumption is made.
\end{theorem}
\begin{proof}
We induct on $N$.  The cases $N\le1$ follow from the curve lemma.
A zero equation may be discarded.  With at most one nonzero equation,
use Lemma~\ref{p3sys:cubic-function}.  If the two quadratic forms
have a common factor of degree two, they are proportional.  If their
greatest common divisor is a linear form $L$, write them as
$LL_1,LL_2$; their common zero set is the union of the linear
spaces $L=0$ and $L_1=L_2=0$.  The assertion again follows from
the unirational lemma.  We may therefore suppose that the two
forms $Q_1,Q_2$ are coprime.

If $U(K)$ is empty there is nothing to prove.  Choose $p\in U(K)$
and coordinates with $p=(1:0:\cdots:0)$.  Since $Q_i(p)=0$,
we can write
\[
 Q_1=X_0L_1(X_1,\ldots,X_N)+q_1(X_1,\ldots,X_N),\qquad
 Q_2=X_0L_2(X_1,\ldots,X_N)+q_2(X_1,\ldots,X_N),
\]
where the $L_i$ are linear and the $q_i$ are quadratic.  We consider
the rank of the span of $L_1,L_2$.

If the rank is two, projection from $p$ is an isomorphism between
\[
 X\setminus V(L_1,L_2)
 \quad\text{and}\quad
 V(L_1q_2-L_2q_1)\setminus V(L_1,L_2)
          \ \subseteq\mathbb P^{N-1}_K.
\]
On the chart $L_i\ne0$ the inverse is given by
$X_0=-q_i/L_i$; these formulas agree on the overlap.  The cubic
$L_1q_2-L_2q_1$ is nonzero.  Otherwise coprimality of the two
linear forms would give $q_1=L_1M$, $q_2=L_2M$ for a linear
form $M$, making $X_0+M$ a common factor of $Q_1,Q_2$.
The image of $f$ on this open part is finite by
Lemma~\ref{p3sys:cubic-function}.  The exceptional part
$X\cap V(L_1,L_2)$ is cut out by two quadrics in the smaller
linear space $V(L_1,L_2)\simeq\mathbb P^{N-2}_K$.
The induction hypothesis applies to the restriction of $f$ and
the open subset inherited from $U$.

If the rank is one, replace the two forms by an invertible linear
combination so that $L_1\ne0$ and $L_2=0$.  Away from $L_1=0$,
projection from $p$ is an isomorphism to the open subset
\[
                         V(q_2)\setminus V(L_1)
                          \subseteq\mathbb P^{N-1}_K,
\]
with inverse $X_0=-q_1/L_1$.  The polynomial $q_2$ is nonzero,
since the second equation was nonzero.  The quadric lemma gives
a finite image here.  On the exceptional linear space $L_1=0$
we again have two quadrics in a smaller projective space, so
induction gives a finite image there as well.

If the rank is zero, both equations are independent of $X_0$,
and $p$ is a vertex of the intersection.  For every $q\in U(K)$,
the line joining $p$ to $q$ lies in $X$.  The restriction of $f$
to the open part of this line has only integral values at its
$K$-points, hence is constant by the unirational lemma.  It is
defined at both $p$ and $q$, so $f(q)=f(p)$.  The image is a
singleton.  The three cases prove the induction.
\end{proof}

\begin{corollary}\label{p3sys:two-quadratics-affine}
A system of at most two equations of total degree at most two,
in any number of rational coordinates over $K$, cannot define
$\mathcal O_K$ in its input.  More generally, whenever its
projection, possibly restricted by finitely many polynomial
nonzero conditions, lies in $\mathcal O_K$, that projection is
finite.  The same holds for a finite union of such presentations.
\end{corollary}
\begin{proof}
Homogenize each equation to degree two and take the open chart
where the homogenizing coordinate is nonzero.  The input becomes
a regular rational function on that chart.  Apply the theorem;
a finite union of finite images is finite.
\end{proof}

Linear equations can first be solved by Gaussian elimination.  On the
resulting affine subspace, the input is an affine linear function and
the remaining residuals are still quadratic.  Thus the same obstruction
applies whenever those residuals span a vector space of dimension at
most two.  A quadratic system defining $\mathcal O_K$ must have at
least three independent quadratic residuals after the linear equations
have been removed.  This assertion counts equations, not variables.

\subsection{Restrictions on the number of coordinates}

\begin{lemma}\label{p3sys:horizontal-closure}
Let $U$ be a locally closed subset of affine space over $K$ and
let $f:U\to\mathbb A^1_K$ be regular.  If $f(U(K))$ is an
infinite subset of $\mathcal O_K$, the Zariski closure of $U(K)$
has an irreducible component $Z$ of dimension at least two,
with Zariski-dense $K$-points and nonconstant input function.
The component is geometrically integral.
\end{lemma}
\begin{proof}
Take the finitely many irreducible components of the closure.  The
rational points are dense on each component after removing its
intersections with the others.  A component that is not geometrically
integral has all its $K$-points in the intersections of distinct
geometric components; this contradicts density in characteristic
zero.  If every component had dimension at most one or constant
input, the curve lemma would make the union of their input images
finite.  At least one component has the stated properties.
\end{proof}

\begin{theorem}\label{p3sys:small-coordinate-bounds}
Suppose a finite polynomial system, possibly with nonzero guards,
has an infinite rational input projection contained in
$\mathcal O_K$.
\begin{enumerate}
\item It uses at least three total coordinates.  This holds at every
degree and for a finite union of presentations.
\item If it uses three total coordinates, every nonzero residual
has degree at least four.  In fact the residuals have a common
irreducible factor of degree at least four with nonconstant input.
\item If it uses at most four total coordinates, its residuals
cannot all be quadratic.  More strongly, in four coordinates it
cannot contain two coprime residuals of degree at most two,
regardless of the degrees of its other residuals.
\end{enumerate}
\end{theorem}
\begin{proof}
For the first assertion, a proper algebraic subset of affine
two-space has dimension at most one, so the curve lemma applies.
If the closure is all of affine two-space, a locally closed subset
with that closure contains a nonempty open subset; its input cannot
be integral-valued there by the unirational lemma.  The cases with
fewer coordinates are included.  A finite union does not help,
since an infinite image would require an infinite image from one
member.

In three coordinates, choose the component $Z$ in
Lemma~\ref{p3sys:horizontal-closure}.  If there is a nonzero
residual, its common zero set is proper, so $\dim Z=2$.  Its
height-one prime ideal is generated by an irreducible polynomial
$g$.  Every residual vanishes on $Z$, hence is divisible by $g$.
The permitted open subset of $V(g)$ is contained in the source
system and has dense rational points and nonconstant input.
Lemma~\ref{p3sys:cubic-function} excludes $\deg g\le3$.
This proves the second assertion.  If all residuals are zero,
the source is open in affine space and is excluded by the
unirational lemma.

For the stronger statement in four coordinates, suppose two
residuals $Q_1,Q_2$ of degree at most two are coprime.  Their
common zero set has dimension at most two.  Thus the horizontal
component $Z$ again has dimension two and is an irreducible
component of this common zero set.  Homogenize to quadrics,
restrict to the affine chart, and remove all other components.
Intersect further with the source's nonzero guards.  This gives
an open subset of the intersection of two projective quadrics
whose rational points lie in the source system.  The image is
finite by Theorem~\ref{p3sys:two-quadrics-any-dimension}, whereas
density of the points on $Z$ and nonconstancy of the input make
that image infinite.  This is a contradiction.

Finally suppose all residuals have degree at most two.  A nonzero
constant makes the source empty.  A nonzero linear equation reduces
the ambient affine dimension to at most three and is excluded by
the second assertion.  If the quadratic residuals have a common
factor, either they are proportional quadrics or they are a common
linear factor times linear forms.  Their common zero set is then
a quadric or a union of linear spaces, and the cubic-function
lemma applies on every permitted open part.  Otherwise choose
one residual $Q_1$ and a linear combination $Q_2$ of the others
which avoids every irreducible factor of $Q_1$.  Such a choice
exists over the infinite field $K$: each forbidden divisibility
condition is a proper linear subspace of the coefficient choices.
The resulting pair is coprime, so the preceding paragraph applies.
\end{proof}

In particular, a quadratic system defining the integers over the
rationals needs at least five total coordinates and at least three
independent quadratic equations after removal of linear constraints.
A finite system of cubic equations in three total coordinates cannot
define the integers either.  These restrictions leave open systems in
the requested range of five to ten coordinates.  They do not provide
a construction in that range and do not exclude it.


% Restrictions on a compression method, not absolute optimality claims.
\section{What generic projection can and cannot compress}
\label{p3sys:projection-limits}

The finite-fiber method eliminates algebraic choices over a retained
base.  It cannot automatically replace freely varying base coordinates
by a single primitive element.  The distinction is geometric, and
it must not be confused with a lower bound for every possible
Diophantine definition.

\begin{proposition}\label{p3sys:finite-projection-dimension}
Let $X$ be an integral $K$-variety of dimension $d$.  If a rational
map from a nonempty open subset of $X$ to $\mathbb A^b_K$ has
finite geometric fibers, then $b\ge d$.
In particular, rational changes of coordinates followed by finite-fiber
projection cannot reduce the dimension of the retained base below $d$.
\end{proposition}
\begin{proof}
Let $Y$ be the closure of the image.  A generically finite map makes
the function-field extension $K(X)/K(Y)$ finite.  Their transcendence
degrees over $K$ are therefore equal.  Thus
$d=\dim Y\le b$.  Restricting to a nonempty open or making a
birational change does not change the function field or its
transcendence degree.
\end{proof}

\begin{corollary}\label{p3sys:monic-base-barrier}
Suppose a source is described by a positive-rank finite free algebra
over $K[X_1,\ldots,X_b]$, followed only by nonzero conditions.
Every nonempty irreducible component of its geometric open locus
has dimension $b$.  A compression which preserves such a component
up to birational equivalence and finite-fiber projection must retain
at least $b$ base coordinates.
\end{corollary}
\begin{proof}
A finite algebra is integral over its base, and every minimal
prime of a finite flat algebra over a domain lies over the zero
prime.  To see the latter assertion, use going down for a flat
extension: a nonzero contracted prime would admit a strictly
smaller prime below the proposed minimal prime.  The induced
finite dominant map from each irreducible component to the
base preserves dimension.  A nonempty open in that component
still has dimension $b$.  Apply the preceding proposition.
\end{proof}

The unstratified 43-coordinate source before the direct Lucas-residual
refinement has 42 retained base coordinates and a triangular monic
finite algebra, followed only by guards.  The new direct-residual
source has 40 retained coordinates and the same property.  Thus a
generic projection does not, by itself, turn either source into a
finite-fiber presentation over nine retained coordinates.  Such an
operation would have positive-dimensional geometric fibers.  A proof
that those fibers retain the required rational points would be new
arithmetic information, not a consequence of the primitive-element
theorem.  The subsequent decomposition into zero and nonzero quadratic
root cases can introduce closed conditions and smaller components;
the dimension assertion just made concerns the original unstratified
source, not each of those later cases separately.

\begin{proposition}\label{p3sys:no-generic-rational-section}
Let a source variety with regular input $T$ project to
$\mathbb A^b_K$ with coordinates $(T,Z_2,\ldots,Z_b)$.
If that projection has a rational section on a nonempty open
subset of $\mathbb A^b_K$, then the source's input projection
contains all but finitely many elements of $K$.
Consequently a source whose rational input projection is
$\mathcal O_K$ has no such section.
\end{proposition}
\begin{proof}
Shrink the domain of the section to a nonempty principal open
$G(T,\mathbf Z)\ne0$.  For all but finitely many $t\in K$,
the specialization $G(t,\mathbf Z)$ is a nonzero polynomial:
at least one of its nonzero coefficient polynomials in $T$
has only finitely many roots.  Since $K$ is infinite, one can
choose $\mathbf Z\in K^{b-1}$ with $G(t,\mathbf Z)\ne0$.
The section then supplies a rational source point at input $t$.
The ring $\mathcal O_K$ has infinite complement in $K$, so it
cannot contain such a cofinite projection.
\end{proof}

The role of the base in this statement matters.  A rational section
over an already constrained subvariety is compatible with an
arithmetic image.  For example, in the presentation
\[
                         Y^2=T
\]
of rational squares one may add arbitrarily many completely free
coordinates.  They can later be deleted with no change of rational
projection, although the deleting map has positive-dimensional fibers.
Its section is over the square-cover variety, not over the entire
input affine line.  Thus dimension alone is not a lower bound on the
number of coordinates in every equivalent presentation.

There is a second reason for caution.  The geometric dimension of
one equation need not be the dimension of the closure of its rational
solutions.  The equation
\[
                         U^2+V^2=0
\]
defines a one-dimensional hypersurface over an algebraic closure,
but its rational solution set is just $(0,0)$.  Therefore a large
geometric source can in principle have its rational points confined
to a much smaller algebraic set.  The dimension restriction must
be applied to a component actually preserved by the proposed
compression, not silently to all equivalent rational presentations.

\begin{corollary}\label{p3sys:arithmetic-dimension-floor}
For every polynomial-system definition of $\mathcal O_K$, the
Zariski closure of its rational witness points has a horizontal
component of dimension at least two.  On that component the input
map has positive-dimensional generic geometric fiber.
\end{corollary}
\begin{proof}
Use Lemma~\ref{p3sys:horizontal-closure}.  Its nonconstant input
function is dominant onto the affine line, so the generic fiber
has dimension $\dim X-1\ge1$.
\end{proof}

This last restriction is intrinsic to the rational-point locus,
unlike the preceding bound for a particular finite-flat presentation.
It rules out compressing all witnesses to finitely many algebraic
choices over the input alone.  It does not rule out a two-dimensional
arithmetic witness variety, nor a definition with three to ten total
coordinates.  No such new construction is furnished by these
projection arguments.


\section{Geometric limitations of compressing the existing witnesses}
\label{sec:geometric-compression-limits}

The bound of $41$ total coordinates is not a lower bound for
definitions of $\mathbb Z$.  This section proves a more limited
statement: certain proposed ways to compress the existing elliptic
point modules cannot work.  A new arithmetic construction can avoid
these limitations by restricting or replacing those modules while
preserving only the input projection.

\subsection{Finite algebra requires geometric finiteness}

\begin{proposition}\label{prop:finite-algebra-dimension-obstruction}
Let $V$ be a geometrically irreducible variety of dimension $d$ over
a field $F$, and let $S\subseteq V(F)$ be Zariski dense.  Suppose
there are finitely many varieties $Y_j$, morphisms
$\pi_j:Y_j\longrightarrow\mathbb A_F^m$ with finite geometric
fibres, and rational maps $g_j:Y_j\dashrightarrow V$, defined
where they are used, such that
\[
                         S\subseteq\bigcup_jg_j(Y_j(F)).
\]
Then $m\ge d$.
\end{proposition}
\begin{proof}
The dimension theorem gives $\dim Y_j\le m$ because each
geometric fibre of $\pi_j$ is finite.  Every irreducible component
of the closure of $g_j(Y_j)$ therefore has dimension at most $m$.
If $m<d$, the union of these finitely many closures is a proper
closed subset of the irreducible variety $V$.  It cannot contain
the Zariski dense set $S$.
\end{proof}

In particular, a finite triangular monic algebra over $m$ rational
parameters, followed by rational decoding formulas, cannot cover
a Zariski dense $d$-dimensional family when $m<d$.  A finite list
of exceptional charts does not change this conclusion.

This observation also distinguishes finite rational fibres from
finite geometric fibres.  Pasten proves that, for a positive-rank
elliptic curve over a number field, there is an affine open set
$U$ and a morphism $U\times U\to\mathbb A^1$ which is injective
on rational points, although those rational points are Zariski
dense in the surface \cite[Corollary 1.2]{P3Pasten2020}.
Its rational fibres have at most one point.  Its generic geometric
fibre has dimension one.  Proposition~\ref{prop:finite-algebra-dimension-obstruction}
shows that this injection does not supply a finite-algebra decoder
for the pair from one rational coordinate.

\subsection{The rank-one point group is dense in a surface}

We now establish the obstruction for the exact elliptic curve
used in the construction.  The proof needs no height estimate,
prime-pattern theorem, or restricted converse for elliptic curves.

\begin{lemma}\label{lem:conjugate-twist-density}
Let $F$ be a number field, let $K/F$ be quadratic with involution
$\sigma$, and let $d\in K^*$.  Put $N=d\sigma(d)\in F^*$ and
suppose $\sqrt N\notin K(i)$.  On
\[
                         E_d:y^2=x^3-d^2x
\]
let $P\in E_d(K)$ be nontorsion.  The cyclic group generated by
$P$, viewed in the Weil restriction
$A=\operatorname{Res}_{K/F}E_d$, is geometrically Zariski dense
in the two-dimensional variety $A$.
\end{lemma}
\begin{proof}
After extension to an algebraic closure,
$A\simeq E_d\times E_{\sigma(d)}$.  Put $u=d/\sqrt N$.
The change of variables
\[
 E_{\sigma(d)}\longrightarrow E_d,
 \qquad (x,y)\longmapsto(u^2x,u^3y)
\]
is an isomorphism because
$u^4\sigma(d)^2=d^2$.  Write $R$ for the image of
$\sigma(P)$.  There is an automorphism $\tau$ fixing $K(i)$
and changing the sign of $\sqrt N$.  Consequently
\[
                            \tau(P)=P,\qquad\tau(R)=-R.
\]

The geometric endomorphism ring of $E_d$ is $\mathbb Z[i]$,
and all its endomorphisms are defined over $K(i)$.  One can see
this from complex uniformization: the order-four automorphism
$(x,y)\mapsto(-x,iy)$ makes the period lattice a rank-one
$\mathbb Z[i]$-module; since $\mathbb Z[i]$ is a principal ideal
domain, its multiplier ring is exactly $\mathbb Z[i]$.
In particular $\tau$ fixes every such endomorphism.

Suppose $\alpha(P)+\beta(R)=0$ for endomorphisms
$\alpha,\beta\in\mathbb Z[i]$.  Applying $\tau$ gives
$\alpha(P)-\beta(R)=0$.  Thus $2\alpha(P)=2\beta(R)=0$.
A nonzero endomorphism of an elliptic curve is an isogeny and
has finite kernel.  Since $P$ and $R$ are nontorsion, it follows
that $\alpha=\beta=0$.

For completeness, this absence of relations implies the claimed
density.  The Zariski closure of a subgroup of an algebraic group
is an algebraic subgroup.  If the closure of $\mathbb Z(P,R)$
were proper in $E_d^2$, its identity component would have
dimension zero or one.  Dimension zero would make $(P,R)$
torsion.  In dimension one, let $C$ be that identity component.
Both projections $C\to E_d$ are nonzero isogenies: otherwise
a nonzero multiple of $P$ or $R$ would vanish.  If these maps
are $\phi,\psi$, a dual isogeny to $\phi$ gives a nonzero
relation
\[
                       [\deg\phi]y-\alpha x=0
                       \quad\hbox{on }C
\]
for an endomorphism $\alpha$ of $E_d$.  A nonzero multiple of
$(P,R)$ lies in $C$, so this is a nontrivial endomorphism
relation between $P$ and $R$, contradicting the previous
paragraph.  The standard isogeny facts used here may be found
in \cite[Sections 8 and 12]{P3MilneAbelian}.
\end{proof}

\begin{corollary}[The point modules used here]
\label{cor:point-modules-dimension}
For $K=\mathbb Q(\sqrt2)$, $d=3-\sqrt2$, and the explicit
point
\[
 P=(57/8-7\sqrt2/2,-165/8+417\sqrt2/32),
\]
the set $\mathbb ZP$ is Zariski dense in
$\operatorname{Res}_{K/\mathbb Q}E_d$.  Every nonzero
finite-index subgroup has the same property.  Nine independent
copies of the point module consequently form a Zariski dense
subset of a variety of dimension eighteen.

A construction which covers all those independent point choices
using finitely many rational decoding charts over a finite
algebra needs at least eighteen retained rational parameters.
In particular, replacing each point module separately by one
rational parameter and finitely many algebraic choices is
impossible.
\end{corollary}
\begin{proof}
Here $N(d)=7$ and the quadratic subfields of
$\mathbb Q(\sqrt2,i)$ are
$\mathbb Q(\sqrt2)$, $\mathbb Q(i)$, and
$\mathbb Q(\sqrt{-2})$.  Thus $\sqrt7\notin K(i)$.

The stated point is nontorsion by an elementary local calculation.
The curve has good reduction at three, the reduction of $P$ is
the nonzero two-torsion point $(\sqrt2,0)$, and
\[
 z(2P)=
 \frac{3396356159088-4135204317756\sqrt2}{3287811625151},
 \qquad z=-x/y,
\]
has valuation one at the unramified quadratic place above three.
For a point of positive integral depth there, multiplication by
a three-adic unit preserves depth, while multiplication by three
increases depth by exactly one.  Indeed the integral formal
group has odd multiplication series, whose first term is $nz$;
the higher terms in $[3]z$ have valuation at least $3v(z)$,
strictly greater than $1+v(z)$.  No nonzero integer can therefore
annihilate $2P$.  The displayed exact calculation is reproduced
by \texttt{check\_elliptic\_density\_data.py}.

Apply Lemma~\ref{lem:conjugate-twist-density}.  Replacing $P$
by a nonzero multiple preserves its nontorsion property and
the same argument.  A Cartesian product of Zariski dense sets
is dense in the product variety, which has dimension eighteen.
The final assertion follows from
Proposition~\ref{prop:finite-algebra-dimension-obstruction}.
\end{proof}

\paragraph{Scope of this obstruction.}
The nine point choices are not independent after every arithmetic
condition of the final definition has been imposed.  The corollary
therefore gives a limitation on separate replacement of the point
modules and on witness-preserving finite-algebra parametrization.
It is not a lower bound of eighteen witnesses for $\mathbb Z$,
and it does not exclude a new definition with at most ten total
coordinates.  Such a reduction must exploit additional arithmetic
relations among the points, replace the point construction, or use
a different kind of representation.


\section{Limits of finite tests on denominator powers}
\label{sec:denominator-obstructions}

The local obstruction behind the pole construction excludes odd negative
valuations.  It is natural to try to replace its many occurrences by
one product, or to impose finitely many such conditions directly on the
input.  The following results rule out two precise versions of those
strategies.  They do not identify the image of every pole equation with
the sets defined below; that identification would require an additional
converse theorem.

\begin{definition}
For an integer $m\ge2$, let
\[
 \mathcal D_m=\{q\in\mathbb Q_{>0}:
       \text{the positive reduced denominator of $q$ is an $m$th power}\}.
\]
Equivalently, $q\in\mathcal D_m$ if and only if $q>0$ and
$m\mid v_p(q)$ whenever $v_p(q)<0$.  In particular,
$\mathcal D_2$ consists of the positive rationals with square
denominator.
\end{definition}

\subsection{Multiplying local obstructions loses information}

\begin{proposition}\label{prop:denominator-monomial-rank}
Let $c_1,\ldots,c_r\in\mathcal D_2$, let
$A=(a_{ji})\in\operatorname{Mat}_{r\times n}(\mathbb Z)$, and put
\[
 M_j(b_1,\ldots,b_n)=c_j\prod_{i=1}^n b_i^{a_{ji}}.
\]
If $A$ has rank less than $n$ over $\mathbb F_2$, there are positive
rationals $b_1,\ldots,b_n$, some outside $\mathcal D_2$, such that
every $M_j(\mathbf b)$ belongs to $\mathcal D_2$.
Consequently an equivalence
\[
 \bigl(b_i\in\mathcal D_2\text{ for every }i\bigr)
 \quad\Longleftrightarrow\quad
 \bigl(M_j(\mathbf b)\in\mathcal D_2\text{ for every }j\bigr)
\]
on all positive rational inputs requires $r\ge n$ and full column
rank of $A$ modulo two.  These necessary conditions are not asserted
to be sufficient.
\end{proposition}
\begin{proof}
Choose a nonzero vector $\epsilon\in\ker(A:\mathbb F_2^n\to
\mathbb F_2^r)$ and a prime $p$ outside the numerator and denominator
supports of all $c_j$.  Set $b_i=p^{-\epsilon_i}$.  At least one
$b_i$ has reduced denominator $p$ and is not in $\mathcal D_2$.
But
\[
 v_p(M_j)=-\sum_i a_{ji}\epsilon_i
\]
is even.  Every other valuation of $M_j$ is the corresponding
valuation of $c_j$, so every negative valuation is even.  Positivity
is automatic.
\end{proof}

For example, the product of two inputs with denominator $p$ has square
denominator.  The proposition shows that choosing several fixed
monomial weights does not repair this cancellation unless the parity
conditions retain at least as much independent information as the
original conditions.

\subsection{Fixed rational pullbacks cannot characterize integrality}

A positive Boolean combination means a finite expression formed using
only conjunction and disjunction.  Rational functions are evaluated
only where their denominators are nonzero.  The result below permits
arbitrary fixed rational functions and different denominator exponents.

\begin{theorem}\label{thm:finite-denominator-tests}
Let $S\subseteq\mathbb Q$ be defined by a finite positive Boolean
combination of predicates of the following forms:
\[
 f(X)\in\mathcal D_m\quad(f\in\mathbb Q(T),\ m\ge2),
 \qquad g(X)=0,\qquad g(X)\ne0
 \quad(g\in\mathbb Q[T]).
\]
If $S$ contains all but finitely many rational integers, then there is
an integer $n\in S$ and infinitely many nonintegral elements of $S$
converging to $n$ in the real topology.  In particular no such finite
combination defines $\mathbb Z$ in $\mathbb Q$.
\end{theorem}

We first record the elementary prime-divisor fact needed in the proof.

\begin{lemma}\label{lem:simple-prime-values}
Let $h\in\mathbb Z[T]$ be nonconstant and squarefree over $\mathbb Q$.
Outside any prescribed finite set of primes, there is a prime $p$ and
an integer residue class $a\pmod{p^2}$ such that
\[
                        v_p(h(n))=1
              \quad\text{for every }n\equiv a\pmod{p^2}.
\]
One may simultaneously require finitely many polynomials coprime to
$h$ over $\mathbb Q$ to be nonzero modulo $p$ on this class.
\end{lemma}
\begin{proof}
A nonconstant integer polynomial has infinitely many prime divisors
of its nonzero integer values.  Here is a proof.  Translate its
argument so its constant term $c$ is nonzero.  If all such prime
divisors belonged to a finite set $\Pi$, choose $M$ divisible by
every prime in $\Pi$.  The nonconstant integer polynomial
\[
                         H(T)=h(cMT)/c
\]
has constant term one and satisfies $H(t)\equiv1\pmod M$.
Every prime divisor of a nonzero value of $H$ would belong to $\Pi$,
but none could divide that value.  Thus every nonzero value would
be $1$ or $-1$, a contradiction.

Now exclude the prescribed primes, the prime divisors of the leading
coefficient and discriminant of $h$, and those of the resultants with
the other polynomials.  A remaining prime divisor of a value supplies
a simple root $a_0$ modulo $p$.  Among its $p$ lifts modulo $p^2$,
exactly one is a root modulo $p^2$.  Any other lift has valuation one.
The resultant conditions give the last assertion.
\end{proof}

\begin{proof}[Proof of Theorem~\ref{thm:finite-denominator-tests}]
Write the defining expression in disjunctive normal form, as a finite
union of conjunctions, called branches in this proof.  A branch
containing a nonzero polynomial equality can contain only finitely
many integers; branches containing an identically false predicate
are empty.  Discard these branches and enlarge the finite exceptional
set of integers accordingly.  All remaining polynomial conditions
are identities or nonvanishing conditions.

Write each nonzero rational function in reduced form $P/Q$ over
$\mathbb Q[T]$.  Call a remaining branch bad if one of its predicates
$P(X)/Q(X)\in\mathcal D_m$ has an irreducible factor of $Q$ whose
multiplicity $e$ is not divisible by $m$.  For each bad branch choose
one such predicate and factor.  Exclude the prime supports of the
rational scalar factors and of all coefficient denominators in these
factorizations.  By
Lemma~\ref{lem:simple-prime-values}, choose a distinct prime $p$ and
a residue class modulo $p^2$ on which this factor has valuation one,
while $P$ and every other denominator factor are units.  Then
\[
                         v_p(P(n)/Q(n))=-e
\]
on that class, so the chosen predicate, and hence that branch, fails.
The Chinese remainder theorem combines these finitely many classes
into an infinite arithmetic progression on which all bad branches
fail.

Choose an integer $n$ in this progression outside the finite
exceptional set.  Since $n\in S$, it belongs to a remaining branch
which is not bad.  Every denominator factor in each predicate of
this branch has multiplicity divisible by its exponent $m_j$.
After absorbing a nonzero rational constant into the numerator,
write all its denominator predicates as
\[
              f_j(X)=P_j(X)/B_j(X)^{m_j}\in\mathcal D_{m_j},
              \qquad\gcd(P_j,B_j)=1.
\]
These rational functions are defined and positive at $n$.

Let $\Pi$ be a finite set of primes containing all coefficient
denominators, leading-coefficient supports, nonzero resultant
supports, and the numerator and denominator supports of the
nonzero rational numbers $f_j(n)$.  Enlarge it if needed so that
outside $\Pi$ all the displayed polynomials are integral with
unit leading coefficient and each pair $(P_j,B_j)$ has unit
resultant.  For each $q\in\Pi$, continuity in $\mathbb Q_q$
gives an integer $e_q\ge1$ such that
\[
 x-n\in q^{e_q}\mathbb Z_q
       \quad\Longrightarrow\quad
 v_q(f_j(x))=v_q(f_j(n))\quad\text{for every }j.
\]
Choose a positive integer $M$ divisible by every $q^{e_q}$,
and put $L=\operatorname{lcm}(m_1,\ldots,m_s)$.  If there are
no denominator predicates, set $L=2$.  For any sufficiently large
prime $p\notin\Pi$ not dividing $M$, define
\[
                              x_p=n+M/p^L.
\]
These are distinct nonintegers and converge to $n$ in $\mathbb R$.

We check every denominator condition of the chosen branch.  At
$q\in\Pi$ the valuation agrees with its value at $n$, which
satisfies that condition.  At $p$ we have $v_p(x_p)=-L$; the unit
leading coefficients give
\[
 v_p(f_j(x_p))=-L\deg P_j+m_jL\deg B_j,
\]
a multiple of $m_j$.  At a prime $q\notin\Pi\cup\{p\}$,
the number $x_p$ is integral.  If $B_j(x_p)$ is a unit, then
$v_q(f_j(x_p))\ge0$.  Otherwise the unit resultant makes
$P_j(x_p)$ a unit and
\[
                      v_q(f_j(x_p))=-m_jv_q(B_j(x_p)),
\]
again a permitted valuation.  Real continuity makes all $f_j(x_p)$
positive for sufficiently large $p$, and preserves the finitely
many polynomial nonvanishing conditions.  Thus all sufficiently
large choices of $p$ give $x_p\in S\setminus\mathbb Z$.
\end{proof}

\begin{corollary}\label{cor:boolean-denominator-tests}
The conclusion of Theorem~\ref{thm:finite-denominator-tests} also
holds for arbitrary finite Boolean combinations of the listed
predicates, allowing negation.
\end{corollary}
\begin{proof}
Use disjunctive normal form with both positive and negative denominator
literals.  Negated polynomial equalities and inequalities are already
among the allowed predicates.  Discard finite branches as before, and
use the same Chinese remainder construction to eliminate every branch
with a bad positive denominator literal.  Choose the resulting
admitted integer $n$ also away from the finitely many zeros and poles
of all nonzero rational functions in the formula.

Consider a negative literal $f(n)\notin\mathcal D_m$ in the chosen
branch.  An identically zero $f$ makes it automatically true.  Otherwise
either $f(n)<0$, which persists in a real neighborhood, or $f(n)>0$
and there is a prime $q$ with $v_q(f(n))<0$ not divisible by $m$.
The latter failure persists in a sufficiently small $q$-adic
neighborhood of $n$.  Add one such prime for each relevant negative
literal to $\Pi$, and choose $M$ sufficiently divisible to preserve
all those valuations as well.  The same numbers $x_p=n+M/p^L$
then satisfy every positive and negative literal for all sufficiently
large permitted primes $p$.  They are nonintegral and converge to $n$.
\end{proof}

The theorem is deliberately about predicates on fixed rational
functions of the input, with no extra existential coordinates.
It does not rule out auxiliary elliptic points, variable rational
functions, or other arithmetic predicates.  It also does not replace
the independent proof needed to identify the exact image of a given
pole system.  Its conclusion is that even arbitrarily complicated
fixed rational pullbacks of perfect-denominator conditions cannot
alone remove every nonintegral rational number.


\section{What rational pairing and existential rank do, and do not, supply}
\label{sec:packing-literature}

The finite-algebra method removes algebraic choices over retained
parameters.  A proposed further compression must justify the decoding
of the remaining rational parameters.  Set-theoretic countability and
injectivity of a coding function do not give this justification.

\subsection{Integer pairing is not rational pairing}

The familiar polynomial
\[
 C(u,v)=\frac{(u+v)(u+v+1)}2+v
\]
pairs nonnegative integers, but has infinitely many rational points
in every rational fibre.  Indeed, for any $c,s\in\mathbb Q$, set
\[
 v=c-\frac{s(s+1)}2,\qquad u=s-v.
\]
Then $C(u,v)=c$.  Thus using this pairing on rational witnesses
introduces spurious decodings.  Restricting its arguments to integers
would require precisely the integrality predicates whose cost is being
optimized.

Poonen proves that a suitable non-density consequence of the
Bombieri--Lang conjecture implies a polynomial injection
$k^2\to k$ over a number field
\cite[Theorem 1.1 and Remark 1.2]{P3Poonen2010}.  Even such an
injection would not supply rational functions recovering the two
inputs, nor a finite geometric fibre for primitive-element
compression.  Every nonconstant morphism
$\mathbb A^2\to\mathbb A^1$ has a one-dimensional generic
geometric fibre.  The elliptic-surface example of Pasten in
Section~\ref{sec:geometric-compression-limits} makes the distinction
between rational injectivity and geometric finiteness unconditional.

\subsection{A recent explicit result has a restricted domain}

Hirvensalo's Theorem 2 gives an explicit polynomial injection on
\[
 \Lambda_\Pi=\mathbb Z[1/p:p\in\Pi]\cap[0,1],
\]
where $\Pi$ is a fixed finite nonempty set of odd primes
\cite{P3Hirvensalo2026}.  With
$M=\operatorname{lcm}_{p\in\Pi}(p-1)$ and $N\ge2$ coprime to
every $p(p-1)$, its construction is a positive constant multiple of
\[
                       (x^M+y^M)^N+y^M.
\]
The degree is $MN$; for $\Pi=\{5\}$ one can take $M=4,N=3$,
giving degree twelve.  The theorem does not assert an injection
on all rational pairs.  The variable denominators in the present
elliptic construction are not restricted to such a fixed set of primes,
and decoding would still need an independent existential definition.
It therefore gives no unconditional reduction of the present count.

\subsection{Existential rank gives the appropriate distinction}

For a Diophantine set $D\subseteq K^r$, its existential rank is the
least number of quantified field variables in an existential definition
of $D$.  The rank can be smaller than the dimension of a particular
witness variety, since a different definition need only preserve the
input projection.  Daans--Dittmann--Fehm's essential fibre dimension
formalizes this distinction, and their characteristic-zero bound is
\cite[Corollary 4.21]{DDF2025}
\[
             \operatorname{rk}^{\exists,+}_K(D)
                    \le\operatorname{efd}_K(D)+1.
\]
The explicit finite-algebra construction used here is an effective
special case of that general principle, with all exceptional algebraic
fibres handled in the specification.  It does not show that positive
dimensional arithmetic fibres have finite algebraic decoders.

The same paper's Remark 8.8 excludes a one-witness existential definition
of the ring of integers of any number field.  It does not provide an
upper bound of nine witnesses for $\mathbb Z$ in $\mathbb Q$.
Its MRDP-based bounds are conditional on an existing integrality
definition and multiply its witness cost; they do not remove that
cost \cite[Remark 8.23]{DDF2025}.

\subsection{What remains necessary for a count of at most ten}

The $41$-coordinate definition has one input, eighteen retained
coordinates for nine elliptic points, one common slope, twenty
coordinates for ten pole pairs, and one final algebraic coordinate.
Replacing the finite algebra by another finite primitive-element
construction does not change these forty retained coordinates.
Proposition~\ref{prop:finite-algebra-dimension-obstruction} rules out
separately encoding the full nine-point family over fewer than eighteen
parameters.  Proposition~\ref{prop:denominator-monomial-rank} rules out
one proposed direct aggregation of the local denominator conditions.
Corollary~\ref{cor:boolean-denominator-tests} further rules out every
finite Boolean combination of fixed rational pullbacks of
perfect-denominator predicates as a direct definition of integrality.

Neither statement is a lower bound for every definition of
$\mathbb Z$.  A substantially smaller construction would need new
arithmetic relations that reduce the independent point choices, a
different method of enforcing integrality, or a proved simultaneous
construction of pole witnesses.  Merely requiring the ten independent
pole conics to use a common parameter is not justified by the separate
existence theorems used in completeness.  No argument in this work
establishes a definition with two through ten total variables.


\section{Quantitative consequences for rational Diophantine problems}
\label{subsec:quantitative-mrdp}

There is a useful refinement of the reduction in
Section~\ref{sec:reduction}. The literature gives a bound of $12N$
witnesses for computably enumerable subsets of $\Q$ when $\Z$ has
a definition with $N$ witnesses \cite[Remark 8.23]{DDF2025}.
Using a small positivity condition together with the constructive
intersection bound gives the following alternative.

\begin{lemma}\label{lem:integer-positivity-rational}
For an integer $t$,
\[
 t\ge0\quad\Longleftrightarrow\quad
 \exists u,v,w\in\Q:\ t=u^2+v^2+w^2+w.
\]
\end{lemma}
\begin{proof}
If $t\ge0$, the three-square theorem applies to $4t+1$.
Its three squares have one odd and two even entries, since
$4t+1\equiv1\pmod4$. Write them as $(2w+1)^2,(2u)^2,(2v)^2$
with integers $u,v,w$ and divide by four. Conversely, a rational
solution gives
$t+1/4=u^2+v^2+(w+1/2)^2\ge0$. The assumption $t\in\Z$
then gives $t\ge0$. This is the integer positivity identity used
in \cite[Equation (1.1)]{Sun2021}, with its witnesses now allowed
to be rational. The integer hypothesis must be retained.
\end{proof}

\begin{proposition}\label{prop:quantitative-mrdp}
Suppose $\Z$ is defined in $\Q$ with $N\ge1$ existential witnesses.
Every computably enumerable subset of $\Q$ then has a one-polynomial
definition with at most $11N+1$ existential witnesses. Furthermore,
polynomial solvability over $\Q$ restricted to at most $9N+2$
variables is undecidable.
\end{proposition}
\begin{proof}
We use the published integer representation theorem
\cite[Theorem 1.1(i)]{Sun2021}: for every computably enumerable
$B\subseteq\mathbb Z_{\ge0}$ there is an integer polynomial $g_B$
such that, for $c\ge0$,
\[
 c\in B\quad\Longleftrightarrow\quad
 \exists y_1,\ldots,y_9\in\Z:\
 y_9\ge0,\quad g_B(c,y_1,\ldots,y_9)=0.
\]
For a computably enumerable $A\subseteq\Q$, encode integer pairs
by the injective polynomial
\[
 q(a)=2a^2-a,\qquad
 f(a,b)=(q(a)+q(b))(q(a)+q(b)+1)+2q(b).
\]
The map $q:\Z\to\mathbb Z_{\ge0}$ is injective, and the second
formula is twice the usual pairing polynomial on nonnegative integers.
For a fixed sum of the two arguments its values are distinct, and
successive sums have disjoint ranges, proving injectivity.
The set
\[
 B=\{f(a,b):a,b\in\Z,\ b\ne0,\ a/b\in A\}
\]
is computably enumerable and contained in $\mathbb Z_{\ge0}$.
Thus $x\in A$ is equivalent to the existence of
$y_0,\ldots,y_9\in\Q$ satisfying
\[
 xy_0,y_0,y_1,\ldots,y_9\in\Z,\quad
 y_0\ne0,\quad y_9\ge0,\quad
 g_B(f(xy_0,y_0),y_1,\ldots,y_9)=0.
\]
There are eleven integrality predicates, each using $N$ private
witnesses. Use Lemma~\ref{lem:integer-positivity-rational} for
$y_9\ge0$, writing $y_9=u^2+v^2+w^2+w$, and use
$y_0\eta-1=0$ for $y_0\ne0$.

Regard $g_B(f(xy_0,y_0),y_1,\ldots,y_8,Y)$ as a polynomial in $Y$.
On the closed stratum where all its coefficients vanish, impose
$y_9=0$; on the complement retain its root equation. This preserves
the final projection to $x$. Indeed, on the former stratum the only
other restrictions on $y_9$ are integrality and nonnegativity, both
satisfied by zero. Its private integrality witnesses can be chosen
afresh and we may set $u=v=w=0$. We do not claim preservation of
the projection to every larger tuple of private coordinates.

Normalize each integrality polynomial in its last private witness as
in the proof of Proposition~\ref{prop:intersection-rank}. Retain
$y_0,\ldots,y_8$, the first $N-1$ witnesses of each integrality
predicate, and $u,v$. Besides $x$, this retains
\[
 9+11(N-1)+2=11N
\]
coordinates. Over every geometric retained tuple, $y_9$ has finitely
many choices; then each last integrality witness has finitely many
choices, $w$ satisfies the monic quadratic
$w^2+w+u^2+v^2-y_9=0$, and $\eta$ is unique whenever it exists.
The resulting constructible system has finite geometric fibers.
Theorem~\ref{thm:finite-fibre-one-variable} replaces all these
remaining choices by one witness, giving $11N+1$ in total.

For the last assertion fix a nonrecursive computably enumerable
$B\subseteq\mathbb Z_{\ge0}$ and use $g_B$ directly, with its
input $c$ supplied as a nonnegative integer parameter. Only the nine
integrality predicates on $y_1,\ldots,y_9$ and the three-witness
positivity predicate are needed. Apply the same normalization to
$g_B(c,y_1,\ldots,y_8,Y)$. Retaining $y_1,\ldots,y_8$, the first
$N-1$ private coordinates in each predicate, and $u,v$ leaves
$8+9(N-1)+2=9N+1$ coordinates. One final finite-choice witness
gives $9N+2$. An algorithm deciding rational solvability in this
many variables would
therefore decide $c\in B$, a contradiction. This step concerns an
integer input parameter; it makes no assertion that the same formula
defines $B$ on arbitrary rational inputs without an integrality test
on that parameter.
\end{proof}

With $N=40$, the conditional 41-variable construction therefore gives
$441$ existential witnesses for each computably enumerable subset of
$\Q$, and undecidability of rational polynomial solvability in at most
$362$ variables. These conclusions have the same arithmetic assumptions
as that construction. Using only the source-preserving 558-variable
construction instead, $N=557$ gives the bounds $6128$ and $5015$.
These are quantitative consequences of the defining-polynomial method,
not claims of optimality among all undecidability reductions. In
particular a specialized undecidability theorem and a definition of the
ordinary subset $\Z\subset\Q$ are different objectives. The first
bound improves $12N$ when $N>1$. No polynomial achieving these
MRDP-derived counts has been expanded in this report.


% Research contribution: number fields.  This is an input fragment for the
% principal report; theorem/proposition/lemma/remark environments are assumed.

\section{Number fields: what transfers, and what does not}
\label{sec:number-fields}

Let $K$ be a number field, that is, a finite extension of $\mathbb Q$, and let
$\mathcal O_K$ be its ring of algebraic integers.  Coefficients in $K$ are
allowed in a Diophantine definition over $K$.  This section distinguishes
the existence of a definition for each $K$ from the much stronger assertion
that one fixed equation works after every extension of scalars.

\subsection{An exact counterexample to extending the released polynomial}

Write $P(X_0,\ldots,X_{12503})$ for the released quartic, with the variable
indices used in its monomial table, and put
\[
 C=5937913971887799417367.
\]
The following calculation is independent of the arithmetic argument that
$P$ defines $\mathbb Z$ over $\mathbb Q$.

\begin{proposition}\label{prop:quartic-gaussian}
Over $\mathbb Q(i)$, the existential projection of the released polynomial
$P$ onto its input coordinate is all of $\mathbb Q(i)$.  In particular,
this polynomial does not define $\mathbb Z[i]$ in $\mathbb Q(i)$.
\end{proposition}
\begin{proof}
Set every witness variable other than $X_{2242}=u$ and $X_{2243}=v$ equal
to zero.  The exact resulting polynomial is
\begin{equation}\label{eq:gaussian-restriction}
 P(t,0,\ldots,u,\ldots,v,\ldots,0)=C+4t^2+t^4+u^2+v^2.
\end{equation}
For any $t$ in a field containing $\mathbb Q(i)$, put
\[
 R=C+4t^2+t^4,\qquad u=\frac{1-R}{2},\qquad
 v=i\frac{1+R}{2}.
\]
Then
\[
 R+u^2+v^2
 =R+\frac{(1-R)^2-(1+R)^2}{4}=0.
\]
Thus every $t\in\mathbb Q(i)$ is accepted.  For example, $t=1/2$ is not
in $\mathbb Z[i]$ and has the witnesses
\[
 X_{2242}=-\frac{95006623550204790677873}{32},\qquad
 X_{2243}=i\frac{95006623550204790677905}{32},
\]
with all other witness coordinates zero.

For reproducibility, \texttt{check\_base\_change.py} checks
\eqref{eq:gaussian-restriction} in two independent ways: by retaining the
relevant monomials of the expanded table, and by specializing and squaring
the quadratic residuals.  It also substitutes the displayed generic
witnesses directly in both representations, using exact arithmetic in
$\mathbb Q(i)[t]$.  The source table has SHA-256
\begin{center}\small\ttfamily
7f5aa3117af82a0b323f0a4e74a18599204f5540408106d3ce99612b5ed830ca.
\end{center}
\end{proof}

\begin{remark}
Proposition~\ref{prop:quartic-gaussian} is a counterexample to direct
extension of the particular equation.  It is not a counterexample to
the existence of a different Diophantine definition of
$\mathbb Z[i]$ in $\mathbb Q(i)$.
The obstruction is already present in the final sum-of-squares operation:
$u^2+v^2=0$ no longer forces $u=v=0$ over $\mathbb Q(i)$.
Even replacing this operation by a suitable norm form would not, by
itself, establish that the original arithmetic system characterizes
integrality over a new field.
\end{remark}

\subsection{A general obstruction to uniform scalar extension}

\begin{proposition}\label{prop:bounded-extension-false-positive}
Suppose that $P\in\mathbb Q[T,W_1,\ldots,W_m]$ has total degree $D$ and
defines $\mathbb Z$ in $\mathbb Q$.  There is a finite set
$E\subset\mathbb Q$ with $|E|\leq D-1$ such that, for every
$q\in\mathbb Q\setminus(\mathbb Z\cup E)$, there exist a number field
$K_q$ and $\mathbf w\in K_q^m$ satisfying
\[
 2\leq[K_q:\mathbb Q]\leq D,\qquad P(q,\mathbf w)=0,
 \qquad q\notin\mathcal O_{K_q}.
\]
Consequently no fixed polynomial defining $\mathbb Z$ in $\mathbb Q$
defines $\mathcal O_K$ in every number field $K$ by scalar extension.
\end{proposition}
\begin{proof}
The polynomial must depend on at least one $W_j$: otherwise its rational
zero set is finite or all of $\mathbb Q$.  Expand it as
\[
 P(T,\mathbf W)=\sum_{\alpha}c_{\alpha}(T)\mathbf W^{\alpha}
\]
and select $\alpha\ne0$ with $c_{\alpha}\ne0$.  Let $E$ be the rational
zero set of $c_{\alpha}$.  Since $\deg c_{\alpha}\leq D-|\alpha|$,
we have $|E|\leq D-1$.  If $q\notin E$, the polynomial
$P(q,\mathbf W)$ is nonconstant.

Choose a variable on which this polynomial depends.  Its leading
coefficient in that variable is a nonzero polynomial in the remaining
variables.  A nonzero polynomial over an infinite field cannot vanish
at every rational point; by induction on the number of variables,
we can therefore specialize all the remaining variables to rational
numbers so that the resulting polynomial $h(U)\in\mathbb Q[U]$ is
nonconstant.  Its degree is at most $D$.
Choose an irreducible factor of $h$, let $\theta$ be a root, and set
$K_q=\mathbb Q(\theta)$.  The specialized coordinates, with $U=\theta$,
give the required zero of $P$.  The degree of $K_q$ is at most $D$.
It cannot be one because $q\notin\mathbb Z$ and $P$ defines
$\mathbb Z$ over $\mathbb Q$.

Finally, $\mathcal O_{K_q}\cap\mathbb Q=\mathbb Z$: a rational
algebraic integer has denominator one, by the rational-root theorem
applied to a monic integral equation.  Hence $q$ is not integral.
\end{proof}

\begin{remark}
If one chosen witness has degree at most $d$ in $P$ and genuinely occurs
after specialization, the same proof gives $[K_q:\mathbb Q]\leq d$.
For a sum of squares whose residuals are affine linear in that witness,
a nonconstant quadratic specialization with no rational zero has
strictly negative discriminant.  Indeed it is a sum of squares of real
affine linear polynomials, so its discriminant is nonpositive, and
discriminant zero would give a rational root.  In this case an imaginary
quadratic extension already supplies a false positive.
\end{remark}

\subsection{Combining equations over an arbitrary number field}

The use of one polynomial rather than a finite system causes no
definability obstruction over a number field, although it can affect
the degree severely.

\begin{lemma}\label{lem:norm-combine}
Let $K$ be a number field, let $s\geq1$, and let
$f_1,\ldots,f_s\in K[\mathbf X]$ have degree at most $d$.
There is a polynomial $F\in K[\mathbf X]$, with no additional variables
and degree at most $sd$, such that
\[
 F(\mathbf x)=0\quad\Longleftrightarrow\quad
 f_1(\mathbf x)=\cdots=f_s(\mathbf x)=0
 \qquad(\mathbf x\in K^n).
\]
\end{lemma}
\begin{proof}
Choose a field extension $L/K$ of degree $s$ and a $K$-basis
$\beta_1,\ldots,\beta_s$ of $L$.  Such an extension exists: choose a
discrete valuation of $K$ with a uniformizer $\pi\in K$; the polynomial
$U^s-\pi$ is Eisenstein.  The norm form
\[
 N(Z_1,\ldots,Z_s)=
 N_{L/K}\!\left(\sum_{j=1}^s\beta_j Z_j\right)
\]
is homogeneous of degree $s$ over $K$, and it vanishes on a vector in
$K^s$ precisely when that vector is zero.  Take
$F=N(f_1,\ldots,f_s)$.
\end{proof}

For the degree objective, the preceding elementary norm bound can be
substantially improved.  The following construction uses only a single
nonarchimedean place.

\begin{lemma}\label{lem:four-variable-anisotropic}
Every number field $K$ admits an anisotropic quadratic form in four
variables with coefficients in $\mathcal O_K$.  Consequently a system
of $s\geq2$ polynomial equations of degree at most $d$ over $K$ can be
combined, without additional variables, into one equation of degree at
most
\[
 2^{\lceil\log_4 s\rceil}d.
\]
If $K$ admits a real embedding, degree at most $2d$ suffices for every
$s$.
\end{lemma}
\begin{proof}
Choose a prime ideal $\mathfrak p$ of $\mathcal O_K$ above an odd
rational prime, and let $v$ be its normalized discrete valuation.
Choose $a\in\mathcal O_K$ whose image in
$\mathcal O_K/\mathfrak p$ is a nonsquare, and choose
$\pi\in\mathfrak p\setminus\mathfrak p^2$.  For any
$(x,y)\ne(0,0)$ in $K^2$, reduction modulo $\mathfrak p$ after
division by $\pi^{\min\{v(x),v(y)\}}$ gives
\[
 v(x^2-ay^2)=2\min\{v(x),v(y)\}.
\]
Indeed the two leading residues cannot cancel because $a$ is a
nonsquare in the residue field.  It follows that
\[
 q(X_1,X_2,X_3,X_4)
 =X_1^2-aX_2^2-\pi X_3^2+a\pi X_4^2
\]
is anisotropic: equality of its two binary parts would equate an even
valuation with an odd one, unless both parts vanish, in which case all
four coordinates are zero.

Place the $s$ residuals at the leaves of a complete tree of branching
degree four and height $h=\lceil\log_4 s\rceil$, padding the unused
leaves by zero.  At every internal vertex apply $q$ to its four child
polynomials.  Anisotropy shows inductively that the resulting polynomial
vanishes exactly when every original residual vanishes.  Each level
at most doubles the degree, giving the stated bound.  If $K$ has a
real embedding, the sum of the squares of all the residuals works
directly, since a sum of real squares is zero only when every summand
is zero.
\end{proof}

\subsection{Precise transfer principles}

\begin{lemma}[Transitivity]\label{lem:diophantine-transitivity-nf}
Let $A\subset B\subset C$ be rings of characteristic zero whose
fraction fields are not algebraically closed.  If $A$ is Diophantine in
$B$ and $B$ is Diophantine in $C$, then $A$ is Diophantine in $C$.
Here the coefficients used to define $A$ are regarded as constants in $C$.
\end{lemma}
\begin{proof}
Take a defining existential system for $A$ over $B$.  Regard its
variables as ranging over $C$, and append a defining system for
membership in $B$ to every variable, including the free variable.
The resulting finite existential system defines exactly $A$ in $C$.
It can be combined into one equation by an anisotropic norm form,
as in Lemma~\ref{lem:norm-combine} for the number-field cases used here.
For the stated more general rings, the same construction uses the
homogenization of an irreducible polynomial over their fraction fields
and clears denominators.
\end{proof}

\begin{proposition}\label{prop:integral-basis-transfer}
If $\mathbb Z$ is Diophantine in a number field $K$ of degree $n$, then
$\mathcal O_K$ is Diophantine in $K$.
\end{proposition}
\begin{proof}
Choose an integral basis $\omega_1,\ldots,\omega_n$ of $\mathcal O_K$.
Then
\[
 x\in\mathcal O_K\quad\Longleftrightarrow\quad
 \exists z_1,\ldots,z_n\in\mathbb Z:\quad
 x=\sum_{j=1}^n\omega_jz_j.
\]
Replace each condition $z_j\in\mathbb Z$ by its Diophantine
definition over $K$, and combine the finitely many equations by
Lemma~\ref{lem:norm-combine}.
\end{proof}

\begin{proposition}[Quantitative integral-basis construction]
\label{prop:quantitative-integral-basis}
Suppose a polynomial $G(T,\mathbf W)$ of degree $d$ defines
$\mathbb Z$ over $K$, using $m$ witnesses.  If $[K:\mathbb Q]=n$,
then $\mathcal O_K$ has a definition using at most
\[
 n(m+1)-1
\]
witnesses.  For $n\geq2$ the total degree can be at most
$2^{\lceil\log_4 n\rceil}d$, and at most $2d$ when $K$ admits a
real embedding.  For $n=1$, the original polynomial suffices.
\end{proposition}
\begin{proof}
In the integral-basis construction, eliminate the last coefficient by
\[
 z_n=\frac{x-\sum_{j<n}\omega_j z_j}{\omega_n}.
\]
Apply $G$ to each of $z_1,\ldots,z_n$, using separate witness tuples.
The remaining coefficient variables number $n-1$ and the witnesses
inside the $n$ copies of $G$ number $nm$.  Substitution of the displayed
linear expression for $z_n$ does not increase degree, and multiplication
by a nonzero element of $K$ clears any coefficient denominator.  There
are now $n$ equations of degree at most $d$; apply
Lemma~\ref{lem:four-variable-anisotropic}.
\end{proof}

\begin{proposition}\label{prop:descent-integrality}
Let $L/K$ be an extension of number fields.  If $\mathcal O_L$ is
Diophantine in $L$, then $\mathcal O_K$ is Diophantine in $K$.
\end{proposition}
\begin{proof}
Choose a $K$-basis $\beta_1=1,\beta_2,\ldots,\beta_e$ of $L$.
In a defining polynomial for $\mathcal O_L$, replace each quantified
$L$-variable by $\sum_{j=1}^e\beta_j y_j$, with the $y_j$ quantified
over $K$.  Expand every coefficient in the same basis and equate the
$e$ coordinates of the resulting expression to zero.  For a free
variable $x\in K$, this system is equivalent to $x\in\mathcal O_L$,
which is equivalent to $x\in\mathcal O_K$.  Combine the equations
using Lemma~\ref{lem:norm-combine}.
\end{proof}

\subsection{The remaining arithmetic question}

The theorem of Koymans--Pagano and, independently,
Alp\"oge--Bhargava--Ho--Shnidman proves that
\begin{equation}\label{eq:dl-theorem}
 \mathbb Z\text{ is Diophantine in }\mathcal O_K
 \quad\text{for every number field }K.
\end{equation}
This is Theorem~1.2 of \cite{KoymansPagano2026} and Corollary~1.2 of
\cite{AlpogeBhargavaHoShnidman2026}.  These statements concern witnesses
in $\mathcal O_K$; they do not supply unrestricted witnesses in $K$.

\begin{theorem}[Conditional equivalence of the extension questions]
\label{thm:nf-equivalence}
Assume that $\mathbb Z$ is Diophantine in $\mathbb Q$.
For a number field $K$, the following are equivalent:
\begin{enumerate}
 \item $\mathcal O_K$ is Diophantine in $K$;
 \item $\mathbb Z$ is Diophantine in $K$;
 \item $\mathbb Q$ is Diophantine in $K$.
\end{enumerate}
The equivalence of the first two assertions uses the established
theorem~\eqref{eq:dl-theorem} but does not require the assumption about
$\mathbb Q$.
\end{theorem}
\begin{proof}
The first assertion implies the second by
\eqref{eq:dl-theorem} and transitivity.  The second implies the first
by Proposition~\ref{prop:integral-basis-transfer}.
For the implication from the second to the third, use
\[
 x\in\mathbb Q\quad\Longleftrightarrow\quad
 \exists a,b\in\mathbb Z\ \exists c\in K:
 \quad bx=a,\quad bc=1.
\]
Finally, if $\mathbb Q$ is Diophantine in $K$, transitivity transports
the assumed definition of $\mathbb Z$ over $\mathbb Q$ to $K$.
\end{proof}

These transfer statements do not construct a Diophantine definition of
$\mathbb Q$ in an arbitrary number field.  The rank-stability theorems
behind~\eqref{eq:dl-theorem} establish definitions over rings of integers,
not that missing field-level assertion.  The results proved here give
no counterexample to the existence of a field-dependent definition of
$\mathcal O_K$, and do not settle that general question.  They do show
that an extension argument must contain new arithmetic input: neither
the reuse of the released equation nor an unrestricted change of the
witness field suffices.


% Fresh phase-three supplement.  All variable counts include the input.
\section{What the newer constructions do over number fields}
\label{sec:new-constructions-number-fields-phase3}

Three different questions must be separated.  A fixed polynomial
over $\mathbb Q$ has an existential projection over every extension
field.  An equation-combining procedure may preserve the projection
of a specified source system over that extension.  Neither fact
identifies the resulting set with the ring of integers of the
extension.  The following results locate explicit failures of the
first and third assertions, and give an effective repair of the
second in Section~\ref{sec:division-norm-compiler-phase3}.
The counterexamples and algebraic encoding results here do not
depend on accepting the source manuscript's arithmetic theorem.

\subsection{The four literal quartics accept every input over $\mathbb Q(i)$}

Let $P_N$ denote the recorded sum-of-squares quartic with $N$
total coordinates, indexed from zero, where $X_0$ is the input.
Set all witnesses to zero except the following two coordinates:
\[
\begin{array}{r|r|r|r}
N&\text{quadratic residuals}&u\text{ coordinate}&v\text{ coordinate}\\\hline
1315&1272&72&73\\
1339&1296&96&97\\
1367&1322&62&78\\
1391&1346&86&102
\end{array}
\]
The exact restriction in each case is
\begin{equation}\label{eq:quartic-two-squares-restriction-phase3}
                     P_N(T,0,\ldots,u,\ldots,v,\ldots,0)
                        =C_N+T^4+u^2+v^2.
\end{equation}
Here $C_N$ is the constant obtained by the same restriction at
$T=u=v=0$.  For $N=1315,1367$ it includes the term
$1024\kappa^2$, with the fixed integer $\kappa=3^{2^{28}}$.
The full integer constants, coordinate names, and hashes of the
four complete coefficient tables are recorded in
\nolinkurl{quartic-base-change-verification_20261106_v1.json}.
Defining $C_N$ by this specialization avoids any ambiguity from
rounding or decimal transcription.

\begin{proposition}\label{prop:quartic-collapse-Qi-phase3}
Each of these four literal quartics has existential input
projection equal to $K$ over every characteristic-zero field
$K$ containing $i$.  In particular none defines
$\mathcal O_{\mathbb Q(i)}$ in $\mathbb Q(i)$.
\end{proposition}
\begin{proof}
For any $T\in K$ put $R=C_N+T^4$ and take
\[
                         u=\frac{1-R}{2},\qquad
                         v=i\frac{1+R}{2}.
\]
Then $u^2+v^2=-R$, so
\eqref{eq:quartic-two-squares-restriction-phase3} vanishes.
Taking $T=1/2$ gives a nonintegral input with witnesses already
in $\mathbb Q(i)$.  This identity is exact for the fixed integer
$\kappa$ as well as before its substitution.
\end{proof}

The independent script
\nolinkurl{check_quartic_base_change_20261106_v1.py}
obtains the restriction twice: from the expanded sparse polynomial
and by squaring the independently restricted primitive residuals.
It compares the two results and substitutes the displayed generic
witnesses using exact arithmetic in
$\mathbb Q(i)[T,\kappa]$.  The conclusion is therefore stronger
than the general observation that sums of squares can cancel in a
nonreal field.

\subsection{An exact integral input rejected by the unshifted sources}

The unshifted source uses $\beta=X^2+2$ and guards the scalar
contact expression
\begin{equation}\label{eq:scalar-contact-polynomial-phase3}
       r_\beta=\beta(\beta-1)(\beta-2401)
                         (32\beta+343)(2\beta-1)
\end{equation}
away from zero.  Its later nonzero rational-square rescaling does
not remove these zeros.  The individual factors are present as
nonvanishing guards in the indexed coefficient specifications.
Consequently $X=\sqrt{2399}$, $i$, and $\sqrt{-2}$ are rejected
by the unshifted guarded source, even though each is an algebraic
integer in the corresponding quadratic field.

For the indexed polynomials the failure is particularly exact.
Recall the notation of the finite-algebra specification:
$G$ is its product of guards, $n$ is its algebra rank,
\[
 H(V)=\det(I-VM_{G\lambda})=\sum_{k=0}^n a_kV^k,
 \qquad
 R_E(V)=\operatorname{Tr}\bigl(M_E\operatorname{adj}
                       (I-VM_{G\lambda})\bigr).
\]
Each branch has an index $1\le d\le n$ and uses
$g=a_d\operatorname{Res}(H_d,H_d')$, $L=a_d^2g$,
\[
 \Psi=1+\sum_{j=1}^d a_{d-j}a_d^{2j-1}g^jU^j,
 \qquad
 \widetilde R_E=\sum_{k=0}^{n-1}[V^k]R_E(V)(LU)^{n-1-k}.
\]
Its final residuals are $\Psi$, the high coefficients
$a_{d+1},\ldots,a_n$, and the $\widetilde R_E$.

\begin{proposition}\label{prop:indexed-contact-value-one-phase3}
For each unshifted indexed polynomial with respectively
41, 43, 43, and 45 total coordinates from the direct and proxy constructions,
the polynomial equals $1$ identically in all remaining
coordinates after imposing any of
\[
                         X^2=2399,\qquad X^2=-1,
                              \qquad X^2=-2.
\]
The assertion holds over every characteristic-zero field.
\end{proposition}
\begin{proof}
The indicated contact factor is zero as a polynomial in the
retained base coordinates, so $G=0$ in every specialized
finite algebra, for every value of every finite parameter.
Thus $M_{G\lambda}=0$, $H=1$, and all $a_k$ with $k\ge1$
vanish.  Hence $g=L=0$ and $\Psi=1$.
Moreover $R_E(V)=\operatorname{Tr}(M_E)$ is constant.
Every branch retains nine quartic point generators, so
$n\ge4^9>1$.  Therefore its reversed trace is zero:
the sole possibly nonzero coefficient has the factor
$(LU)^{n-1}=0$.  Every branch sum of squares is consequently
$1$, and so is their finite product.  This calculation does
not use anisotropy of a sum of squares.
\end{proof}

The contact checker verifies a literal multiplicative path from
the zero factor to the materialized guard, evaluates that factor
in $\mathbb Q[X]/(X^2-a)$, and checks the branch-rank assertion.
Its separate checks of the four quadratic systems follow exact
input-dependent gate equations until an inverse equation becomes
$-1=0$.  The recorded reversible row operations transport this
contradiction to each final quadratic residual system.  Thus a
faithful norm combination of the unshifted residuals still rejects
these inputs.  Over $\mathbb Q(\sqrt{2399})$ the literal sum of
squares is faithful as well, and therefore rejects that integral
input.  Over $\mathbb Q(i)$ the literal quartic instead accepts
every input by Proposition~\ref{prop:quartic-collapse-Qi-phase3};
this is precisely the distinction between a system and its
sum-of-squares encoding.

These finite contact obstructions admit the field-dependent shift
repair of Section~\ref{sec:p3-beta-shift}.  The assertion above
is about the unshifted artifacts and does not apply to the
separately emitted $\beta=X^2+2402$ artifacts.  Removing a finite
list of excluded inputs does not restore the elliptic arithmetic
argument considered next.

\subsection{A totally real quartic obstruction to the finite elliptic anchor}

Put $F=\mathbb Q(\sqrt2)$, $d=3-\sqrt2$ and
$D=d^2=11-6\sqrt2$.  Let
\[
                         E:\ y^2=x^3-Dx,
\]
and let $P\in E(F)$ be the fixed point used in the finite
anchor.  Its particular coordinates are immaterial below;
only $P\in E(F)$ is used.  The anchor excludes equality with
the finitely many abscissas $x([n]P)$, $1\le n\le M$, by
taking the product of their nonzero rational norms.

\begin{proposition}\label{prop:totally-real-anchor-counterexample-phase3}
There is a totally real quartic number field $K$, with
$\sqrt2\notin K$, and a point $Q\in2E(K(\sqrt2))$ of
fixed height such that
\[
              x(Q)\in F,\qquad x(Q)\ne x([n]P)
                       \quad\text{for every }n\in\mathbb Z
\]
whenever the abscissa on the right is defined.  The same fixed
point therefore satisfies every finite anchor exclusion, for
every positive bound $M$.
\end{proposition}
\begin{proof}
Let $s$ be a root of
\[
                         s^4-1696s^2+18432=0
\]
and put $K=\mathbb Q(s)$.  Set $F'=\mathbb Q(\sqrt{2737})$,
where $2737=7\cdot17\cdot23$.  The two values of $s^2$ are
$848\pm16\sqrt{2737}$, both strictly positive.  They have
norm $18432=2\cdot96^2$, which is not a rational square.
Thus neither is a square in $F'$, and $[K:\mathbb Q]=4$.
All four conjugates of $s$ are real, so $K$ is totally real.

If $K$ contained $\sqrt2$, it would equal $F'(\sqrt2)$.
Writing $s=a+b\sqrt2$ with $a,b\in F'$ and using
$s^2\in F'$ gives $ab=0$.  In either case the norm of
$s^2$ from $F'$ to $\mathbb Q$ would be a rational square:
it would be $N(a)^2$ or $4N(b)^2$.  This contradicts the
computed norm.  Hence $\sqrt2\notin K$.

In $K(\sqrt2)$ set
\[
                  b=\frac{s}{2}+\frac{48}{s}\sqrt2,
                          \qquad R=(8,b).
\]
The defining equation for $s$ gives
$b^2=424+48\sqrt2=8^3-8D$, so $R\in E(K(\sqrt2))$.
It has nonzero ordinate.  The doubling formula gives
\begin{align*}
 x(Q)=x([2]R)
   &=\frac{(8^2+D)^2}{4b^2}
     =\frac{5697-900\sqrt2}{1696+192\sqrt2}\\
   &=\frac{312741}{87584}
                          -\frac{40941}{43792}\sqrt2\in F.
\end{align*}
This value is not $0,d$, or $-d$, as is already seen from
its coefficient of $\sqrt2$.  Thus $y(Q)\ne0$.

On the other hand $b\notin F$.  If it belonged to $F$, the
norm of its square would be a square in $\mathbb Q$; but
\[
           N_{F/\mathbb Q}(b^2)=424^2-2\cdot48^2
                           =175168=2^6\cdot7\cdot17\cdot23
\]
is not a square.  The doubling slope is
$(3\cdot8^2-D)/(2b)$.  Since $b^2\in F$, the resulting
ordinate $y(Q)$ is a scalar in $F$ times $b$.  That scalar
is nonzero because $y(Q)\ne0$, so $y(Q)\notin F$.

Every multiple $[n]P$ has coordinates in $F$.  Equality of
its abscissa with $x(Q)$ would imply
$Q=[n]P$ or $Q=-[n]P$, and would put $y(Q)$ in $F$, a
contradiction.  Hence each defined difference
$x(Q)-x([n]P)$ is a nonzero element of $F$ and has nonzero
rational norm.  Every finite product is nonzero.
\end{proof}

The halfpoint is represented in the literal two-coordinate
model by $a_0=8$, $a_1=0$, $b_0=s/2$, $b_1=48/s$ in $K$.
Thus the example applies to the extended point chart itself;
it is not an artifact of a split quadratic algebra.
The script \nolinkurl{check_real_anchor_20261106_v1.py}
checks all curve and doubling identities by exact arithmetic
in the degree-eight compositum.

The conclusion disproves the extension of the inference
``finite exclusion forces a large integer elliptic label.''
Over the original field that inference depends on a cyclic
group description.  Enlarging the field introduces points
outside that fixed cyclic group.  We have not used this
single point to construct witnesses for every other source
equation at a nonintegral input; the proposition is a
counterexample to an essential arithmetic inference, not a
claim of a complete false positive for the entire source.

\subsection{A domain assertion that does extend to fields with a real place}

One useful component survives more generally.  Let
\[
 G_J(W)=1+\sum_{j=1}^J c_jW^j,\qquad
 c_1=-\frac25,\qquad
 \frac{c_j}{c_{j-1}}=\frac{4j-5}{4j+1}\quad(j\ge2).
\]
All $c_j$ are negative, and induction gives the exact identity
\[
                        G_J(1)=-\frac{4J-1}{2}c_J>0.
\]
Consequently $G_J(w)\ge G_J(1)>0$ for $0\le w\le1$.
At a real embedding of the curve above, $D>0$.  For an
affine double $Q=[2](a,b)$ with the required denominators
and $y(Q)$ nonzero,
\[
                  x(Q)=\left(\frac{a^2+D}{2b}\right)^2>0.
\]
The curve equation then implies $x(Q)>\sqrt D$, and hence
$0<D/x(Q)^2<1$.  Thus
\[
             -\frac{x(Q)}{y(Q)}G_J\!\left(\frac D{x(Q)^2}\right)
                          \ne0.
\]
An embedding is injective, so this proves the nonvanishing
in any number field with a real embedding.  If the formal
quadratic algebra splits, apply the argument to each component;
the source denominator guards ensure both are defined.

The analogous assertion is false over unrestricted number
fields.  Choose an algebraic root $w$ of $G_J$.  It is neither
zero nor one.  Choose $x=d/\sqrt w$ and a nonzero square root
$y$ of $x^3-Dx$.  This gives an algebraic point $Q$ where the
displayed logarithm truncation vanishes.  Adjoin a halfpoint
of $Q$, which exists over a finite extension because the
multiplication-by-two morphism is surjective, and, if desired,
adjoin $i$.  The resulting number field supplies the asserted
failure even on $2E$.  Therefore a final norm combination
repairs conjunctions, but does not automatically justify all
quotient-domain simplifications originally proved over
$\mathbb Q$.

\subsection{No one fixed existential formula works over every number field}

\begin{theorem}\label{thm:no-uniform-existential-integer-rings-phase3}
Let $k_0$ be a number field.  Fix an existential formula over
$k_0$ whose quantifier-free part is a finite union of systems
of polynomial equations and nonvanishing conditions.
Suppose its input projection over $k_0$ contains $\mathbb Z$.
There is a finite extension $K/k_0$ over which its input
projection contains a rational noninteger.  In particular no
one such fixed formula defines $\mathcal O_K$ in every finite
extension $K/k_0$.
\end{theorem}
\begin{proof}
Over an algebraic closure $\overline{k_0}$ the formula's
witness locus is constructible.  Its image in the affine
input line is constructible by Chevalley's theorem
\cite[Tag 00FE]{StacksChevalley}.  A constructible subset of
an affine line over an algebraically closed field is finite
or cofinite.  This image contains $\mathbb Z$, so it is
cofinite.  Choose $q\in\mathbb Q\setminus\mathbb Z$ in it.
The finitely many algebraic witness coordinates belong to
some finite extension $K/k_0$.  The formula therefore holds
at $q$ over $K$.  But a rational algebraic integer is an
integer: if $q=a/b$ in lowest terms satisfies a monic
integer polynomial, clearing denominators shows $b\mid a^m$
and hence $b=1$.  Thus $q\notin\mathcal O_K$.
\end{proof}

This includes any finite number of simultaneous equations,
any finite alternatives, and any finite list of denominator
guards.  Increasing these counts cannot overcome the
uniformity obstruction.  It does not preclude choosing a
different formula for each number field.

There is a useful real refinement when the coefficient field
is $\mathbb Q$, or more generally when $k_0$ has a chosen
real embedding.  The real input projection is semialgebraic
by quantifier elimination and contains all positive and
negative integers.  A semialgebraic subset of the line is
a finite union of points and intervals; it therefore contains
both sufficiently distant tails.  Choose a nonintegral
rational input in one tail.  The transfer principle for real
closed fields supplies real algebraic witnesses.  They,
together with $k_0$, generate a number field $K$ with a real
embedding where that input is accepted.  Quantifier
elimination and this transfer principle are stated, for
example, in \cite[Theorem 2.1]{Basu2014} and
\cite[Corollary 4.5 and Example 4.6]{Plaumann2023}.
The resulting field need not be totally real; only a real
embedding is asserted.

\subsection{What remains open after these checks}

None of these counterexamples proves that some
$\mathcal O_K$ is not Diophantine in $K$.  They prove that
the literal formulas, fixed cyclic arithmetic argument,
or fixed field-independent coefficient choice cannot supply
the desired universal conclusion.  A field-adapted norm form
faithfully combines an already correct finite system; it
does not produce the missing arithmetic definition.
Any extension of the rational construction must additionally
control the new field-valued point coordinates and arithmetic
witnesses, or supply an independent existential definition
of the embedded rational field.  Under the claimed rational
definition, such rational isolation would allow an integral
basis construction of $\mathcal O_K$, but that additional
existential assertion has not been established here for
every number field.


% Fresh phase-three supplement; phase-two artifacts are not modified.
\section{Combining equations over number fields by division norms}
\label{sec:division-norm-compiler-phase3}

This section concerns the faithful encoding of polynomial systems.
It does not assert that extending the arithmetic source to a number
field defines that field's ring of integers.  The distinction is
essential: a correct equation-combining operation preserves whatever
solutions the extended source has, including any arithmetic defects.

\subsection{A rational norm form adapted to a specified number field}

The reduced norm of a central division algebra of degree $d$ is a
homogeneous form of degree $d$ in $d^2$ coordinates, vanishing only
at the origin.  We give an explicit cyclic construction and a local
proof of division, so that the effectivity and the coefficients are
specified.  Standard background on cyclic algebras is available in
\cite[Chapter IV]{MilneCFT2020}; the prime-selection step uses
Chebotarev's theorem, for example \cite[Section 2.4]{KedlayaCFT}.

\begin{theorem}\label{thm:division-norm-conjunction-phase3}
Let $K$ be any specified number field and let $r\ge1$.
Put $d=\lceil\sqrt r\rceil$.  There is an effectively specified
homogeneous polynomial
\[
                 N\in\mathbb Z[Z_1,\ldots,Z_{d^2}]
\]
of degree $d$ such that, for $z_j\in K$,
\[
                  N(z_1,\ldots,z_{d^2})=0
                  \quad\Longleftrightarrow\quad z_1=\cdots=z_{d^2}=0.
\]
Consequently a system of $r$ equations over $K$, of degrees at
most $D$, can be combined into one equation of degree at most
$dD$, with no additional variables.  The combining form can have
rational integer coefficients even though its choice depends on $K$.
\end{theorem}
\begin{proof}
For $r=1$ use the equation itself.  Assume $d\ge2$.  Let $K'$ be
the normal closure of $K/\mathbb Q$.  Choose a rational prime
$q\equiv1\pmod d$ unramified in $K'$, and let $L$ be the unique
degree-$d$ subfield of $\mathbb Q(\zeta_q)$.  Such a prime exists
by Dirichlet's theorem.  Every nontrivial subfield of $L$ is
ramified at $q$, whereas $K'$ is unramified there.  Thus
$L\cap K'=\mathbb Q$.  Both extensions are Galois, so their
compositum has the product Galois group.

Chebotarev's theorem supplies a rational prime $p\ne q$ which
splits completely in $K'$ and whose Frobenius generates
$\operatorname{Gal}(L/\mathbb Q)$.  Hence $p$ is inert in $L$.
Fix a generator $\sigma$ of this cyclic Galois group and form
\[
 A=(L/\mathbb Q,\sigma,p)
   =\bigoplus_{j=0}^{d-1}L u^j,
 \qquad u^d=p,\qquad ua=\sigma(a)u.
\]
These relations give the usual associative cyclic algebra.

At a place of $K$ above $p$, the completion is $\mathbb Q_p$,
and the corresponding completion of $L$ is the unramified
degree-$d$ extension $L_p/\mathbb Q_p$.  On the local cyclic
algebra define
\[
 w\left(\sum_{j=0}^{d-1}a_j u^j\right)
       =\min_{a_j\ne0}\left(v_{L_p}(a_j)+\frac jd\right).
\]
The minimum is attained at exactly one index, since the
fractional parts $j/d$ are distinct.  Multiplying the unique
lowest terms shows that $w(ab)=w(a)+w(b)$: the automorphism
preserves the valuation, and reducing $u^d$ to $p$ preserves
the weight.  No other product term has this lowest weight.
The local algebra therefore has no zero divisors and is a
division algebra, because it is finite dimensional over a field.
It follows that $A\otimes_{\mathbb Q}K$ is also a division
algebra: a zero divisor would remain one after extension to
that completion.  The same argument proves that $A$ is division.

Here is a concrete basis of $L$ consisting of algebraic integers
for writing its norm form.
Let $g$ generate $(\mathbb Z/q\mathbb Z)^*$, let
$H=\langle g^d\rangle$, and put
\[
 \eta_i=\sum_{h\in H}\zeta_q^{g^ih},\qquad 0\le i<d.
\]
The periods are linearly independent over $\mathbb Q$.
Indeed a relation among them gives a polynomial in $\zeta_q$
of degree at most $q-1$ and constant coefficient zero.
It must be a multiple of $1+T+\cdots+T^{q-1}$, whose constant
coefficient is one, and therefore must be zero.  Since the
periods are fixed by $H$, they form a basis of $L$.  Their sum
is $-1$, so $1,\eta_1,\ldots,\eta_{d-1}$ is also a basis.
Choose
$\sigma$ so that $\sigma(\eta_i)=\eta_{i+1}$, with indices
modulo $d$.

Write $a_j=Z_{j,0}+\sum_{i=1}^{d-1}\eta_i Z_{j,i}$.  Let
$D(a)=\operatorname{diag}(a,\sigma(a),\ldots,\sigma^{d-1}(a))$
and define the matrix $U$ by
\[
 Ue_j=e_{j-1}\ (1\le j<d),\qquad Ue_0=p e_{d-1}.
\]
Then $U^d=pI$ and $UD(a)=D(\sigma(a))U$.  The norm form is
\begin{equation}\label{eq:explicit-cyclic-norm-phase3}
 N(Z)=\det\left(\sum_{j=0}^{d-1}D(a_j)U^j\right).
\end{equation}
Applying $\sigma$ to this matrix conjugates it by $U$, so
its determinant has rational coefficients.  The periods are
algebraic integers and $p$ is an integer, so those rational
coefficients are integers.  The form is homogeneous of degree
$d$.  A nonzero coordinate vector over $K$ represents a
nonzero, hence invertible, element of $A\otimes K$; its matrix
is invertible and its norm is nonzero.  This proves anisotropy.

The construction is effective.  One can search for $q$, compute
its cyclotomic periods, and search primes $p$ testing complete
splitting in $K'$ and the indicated Frobenius class in $L$.
The cited existence theorems ensure termination.  Finally
substitute the $r$ residual polynomials for the first $r$
coordinates of $N$, and zeros for all remaining coordinates.
Anisotropy proves equivalence, and homogeneity gives the degree
bound without introducing any quantified coordinate.
\end{proof}

\begin{remark}[Coefficient bound]
In every complex embedding, $|\eta_i|\le(q-1)/d$.
Every entry of the matrix in
\eqref{eq:explicit-cyclic-norm-phase3} is a linear form whose
coefficient absolute sum is at most $(q-1)p$.  Expanding the
determinant and applying the triangle inequality therefore gives
\[
                       \|N\|_1\le d!\bigl((q-1)p\bigr)^d.
\]
This includes all determinant permutations and all monomials.
The form is primitive: assigning $Z_{0,0}=1$ and all other
coordinates zero represents the identity and gives $N=1$.
If all input residuals have
coefficient norm at most $M$, their combined polynomial has
coefficient norm at most $\|N\|_1M^d$.
\end{remark}

\subsection{The degree-36 form and the untransformed direct systems}

The untransformed direct-Lucas systems have respectively 1315 and 1339
total coordinates, with 1272 and 1296 quadratic residuals.
Both residual counts are at most $36^2=1296$.  For these systems
take the completely explicit choices
\[
                    d=36,\qquad q=37,\qquad p=2.
\]
The element 2 generates $(\mathbb Z/37\mathbb Z)^*$: indeed
$2^{18}\equiv-1$ and $2^{12}\equiv26\not\equiv1\pmod{37}$,
which exclude the maximal proper divisors of 36.
Thus $L=\mathbb Q(\zeta_{37})$ has degree 36 and two is inert.
The preceding local valuation proof shows that
\[
                    A=(L/\mathbb Q,\sigma,2),
             \qquad \sigma(\zeta_{37})=\zeta_{37}^2,
\]
is division.  Formula~\eqref{eq:explicit-cyclic-norm-phase3}
therefore gives an integer norm form of degree 36 in 1296
coordinates.  Here the period basis is simply
$1,\zeta_{37}^{2},\zeta_{37}^{2^2},\ldots,
\zeta_{37}^{2^{35}}$, with exponents reduced modulo 37.

\begin{lemma}\label{lem:division-coprime-base-change-phase3}
Let $A$ be a central division algebra of degree $d$ over a
field $k$.  If $K/k$ is finite and
$\gcd(d,[K:k])=1$, then $A\otimes_k K$ is division.
\end{lemma}
\begin{proof}
Write $A\otimes K\simeq M_s(B)$, where $B$ is central
division over $K$ of degree $d/s$.  A simple left module
has $K$-dimension $s(d/s)^2=d^2/s$.  By restricting scalars
it is a vector space over the division ring $A$, so its
$k$-dimension is a multiple of $d^2$.  Hence
$s\mid[K:k]$.  Since also $s\mid d$, coprimality gives
$s=1$.
\end{proof}

This proves that the fixed degree-36 form remains anisotropic
whenever $\gcd(36,[K:\mathbb Q])=1$.  The assertion must not
be justified by choosing one completion of coprime local degree:
that inference is false for composite $d$, whereas the module
argument applies directly.

The same form also remains anisotropic over $\mathbb Q(i)$.
Here is a local proof, independent of the coprime-degree lemma.
At 37 put $\pi=\zeta_{37}-1$.  The residue field of
$L_{37}$ is $\mathbb F_{37}$ and
$\sigma(\pi)/\pi\equiv2\pmod\pi$.
The algebra $\mathcal O_{L_{37}}[u]$, $u^{36}=2$, has residue
algebra
\[
                    \mathbb F_{37}[u]/(u^{36}-2),
\]
which is a field.  Indeed a root $\alpha$ satisfies
$\alpha^{37}=2\alpha$, so its Frobenius orbit has length
exactly 36 because 2 has multiplicative order 36.
For $x=\sum a_j u^j$ define $w(x)=\min_j v_\pi(a_j)$.
After factoring out its lowest power of $\pi$, a nonzero
element has nonzero reduction in this residue field.
Commuting $\pi^s$ through powers of $u$ changes reduction
by the automorphism $u\mapsto2^su$; consequently the lowest
reductions of a product are still a product of nonzero
elements.  Thus $w(xy)=w(x)+w(y)$ and the local algebra
is division.  Finally 37 splits in $\mathbb Q(i)$ because
$6^2\equiv-1\pmod{37}$, so a completion after this base
change is still $\mathbb Q_{37}$.  Division at that completion
proves division over $\mathbb Q(i)$.

Combining the two preferred quadratic systems gives polynomials
with respectively 1315 and 1339 total coordinates and
\emph{degree exactly 72}.  The exact residual inventories put
the input monomial $X_0^2$ with coefficient one in the first
residual; no other residual contains $X_0$.  The first norm
coordinate is the coefficient of the algebra identity, so the
coefficient of $X_0^{72}$ in the resulting polynomial is one.
In particular the top degree cannot cancel, and each resulting
integer polynomial is primitive.

The coefficient estimate becomes
\[
                      \|N\|_1\le36!\,72^{36}<2^{361}.
\]
The exact residual tables have coefficient norm at most
$2^{511}$.  The 1339-coordinate version, with ordinary
integer coefficients already materialized in those residuals,
therefore has coefficient height at most
\[
                              2^{18757}.
\]
For the 1315-coordinate version substitute the fixed integer
$\kappa=3^{2^{28}}$ into residuals affine in $\kappa$.
Using $\log_2\kappa<429496730$ gives coefficient height at
most
\[
                              2^{15461901037}.
\]
The artifact names
\nolinkurl{norm_repair_1315_degree72_20261106_v1.json.gz}
and
\nolinkurl{norm_repair_1339_degree72_20261106_v1.json.gz}
specify the complete residual lists, cyclotomic basis,
matrix entries, and determinant.  They are exact finite
polynomial specifications, not expanded coefficient tables
of degree 72.  The symbol $\zeta_{37}$ is eliminated by
reduction modulo its cyclotomic polynomial and extraction
of the constant coefficient; it is not a quantified variable.

For an arbitrary specified number field the fixed form may
cease to be anisotropic.  Theorem~\ref{thm:division-norm-conjunction-phase3}
provides a field-dependent choice of primes $q,p$ with the
same degree 36 and rational integer coefficients.  Thus
degree 72 and the coordinate counts persist, while the
coefficient bound becomes the explicit prime-dependent bound
in that theorem.

\subsection{The preceding degree-37 form and its comparison bounds}

The choices
\[
                         d=37,\qquad q=149,\qquad p=5
\]
are explicit.  The degree-37 subfield of $\mathbb Q(\zeta_{149})$
is unramified at five, and five is inert there: the Galois group
has prime order 37 and $5^4\equiv29\not\equiv1\pmod{149}$.
The local valuation argument above proves division.

This same rational norm form remains anisotropic over every
number field $K$ for which $37\nmid[K:\mathbb Q]$.  To see
this directly, use
$\sum_{v\mid5}e_vf_v=[K:\mathbb Q]$ to choose a place with
$37\nmid e_vf_v$.  The unramified degree-37 extension remains
a field after this local base change.  Normalizing the valuation
of $K_v$ gives $v(5)=e_v$, and the weights
\[
                       v(a_j)+\frac{e_vj}{37}
\]
again have distinct fractional parts.  The same lowest-term
argument proves that the local cyclic algebra is division.
For fields with degree divisible by 37, the field-dependent
prime choices in the theorem still supply degree 37 and rational
integer coefficients.

For comparison, there are 1322 primitive quadratic residuals in the
1367-coordinate system and 1346 in the 1391-coordinate system.
Both numbers are at most $37^2=1369$.  Thus the norm construction
combines either system into one polynomial of degree exactly 74,
with the same respective total coordinate counts.  Exactness of
the degree follows without expansion: the highest homogeneous
quadratic parts of the residuals are not all zero.  At some
rational specialization this tuple is nonzero, and anisotropy
shows that its norm is nonzero.  It is exactly the homogeneous
degree-74 part of the combined polynomial.

For the explicit $(37,149,5)$ choice,
\[
               \|N\|_1\le37!\,740^{37}<2^{496}.
\]
The exact residual tables give maximum coefficient norm at most
$2^{511}$.  Hence the ordinary integer version with 1391 total
coordinates has coefficient height at most
\[
                              2^{19403}.
\]
For the 1367-coordinate version the residuals are affine in
the fixed integer $\kappa=3^{2^{28}}$.  Since
$\log_2\kappa<(8/5)2^{28}<429496730$, a valid integer
coefficient-height bound is
\[
                              2^{15891398413}.
\]
These bounds concern the norm-compiled polynomials.  Their
degree-74 determinant specifications are separate from the
frozen degree-four sum-of-squares tables.

\subsection{The finite-algebra compiler over every number field}

All algebraic steps in the guarded finite-algebra theorem work
over an arbitrary characteristic-zero field: the finite
separating-parameter bound, simple rational eigenline,
trace-adjugate identities, and reciprocal guard are unchanged.
Only the final conjunction by a sum of squares requires an
ordering.  Over a number field with a real embedding it remains
valid as written.  Over an arbitrary number field replace each
branch's sum of squares by the norm form of
Theorem~\ref{thm:division-norm-conjunction-phase3}.

If $n$ bounds the algebra rank and $c$ bounds the number of
closed residuals, a branch contains at most $n+c$ residuals:
one reciprocal residual, at most $n-1$ high-coefficient
conditions, and at most $c$ reversed traces.  Use one norm form
of degree
\[
                         d=\lceil\sqrt{n+c}\rceil
\]
and pad shorter lists with zero.  This introduces no new
coordinate.  Relative to the recorded sum-of-squares degree
bound, multiply each branch bound by $d/2$ before multiplying
over its finite indices.  This operation preserves the rational
projection over the specified number field of the guarded
finite source itself.

The exact recorded integer bounds give the following convenient
rounded estimates.  All counts include the input.
\[
\begin{array}{c|c|r|c}
\text{source}&\text{total coordinates}&d&\text{degree bound over }K\\\hline
\text{direct finite anchor}&41&16385&2^{9566686}\\
\text{direct semilocal}&43&23171&2^{250}\\
\text{proxy finite anchor}&43&23171&2^{9566693}\\
\text{proxy semilocal}&45&32769&2^{257}
\end{array}
\]
The combining form has integer coefficients chosen using $K$;
the number of coordinates and the displayed degree bound do not
depend on $[K:\mathbb Q]$.  These are field-adapted encodings
of the corresponding source systems.  They do not establish
that those systems define $\mathcal O_K$.


\section{Field behavior of the final quartic transformations}
\label{sec:final-quartic-fields-v2}

The last transformations improve the quartic tables without
new arithmetic input.  They give the preferred translated
1315-coordinate quartic, with coefficient height
$64\kappa$ for $\kappa=3^{2^{28}}$, and the preferred ordinary
integer-coefficient quartic with 1335 coordinates and height
less than $2^{1020}$.  Their field behavior follows precisely
from the type of transformation involved.  The four preceding
coefficient tables and their certificates remain unchanged.

\subsection{The affine translation works over every field of characteristic zero}

Put $S=3^{2^{27}}$, so $S^2=\kappa$.  The private anchor
coordinates $u=X_{55}$, $v=X_{56}$, $w=X_{57}$ in the
1315-coordinate table occur only in the row
\[
              -3A^2+16u^2+96v^2-32w^2+32S^2,
                    \qquad A=X_{52}+X_{53}-X_{54}.
\]
Make the rational affine substitution
\begin{equation}\label{eq:final-anchor-affine-map-v2}
 u=S+\frac{U-3V+3W}{4},\qquad
 v=\frac S2+\frac V4,\qquad
 w=\frac{3S}{2}+\frac W4.
\end{equation}
The new row is exactly
\[
 -3A^2+U^2-6UV+6UW+15V^2-18VW+7W^2+8SU.
\]
The inverse map is
\[
       U=4u+12v-12w+8S,\qquad V=4v-2S,
                                      \qquad W=4w-6S.
\]
Both maps are defined over $\mathbb Q$, with $S$ an ordinary
fixed integer.  Therefore the substitution is a bijection of
witness tuples over every characteristic-zero field.  It
preserves the literal polynomial's input projection, as well
as that of the residual system.  No use of an ordering is
required for this assertion.

For the generic $\mathbb Q(i)$ section of the preceding
1315-coordinate quartic, the old private anchor coordinates
are all zero.  Their inverse images are
\[
                   (X_{55},X_{56},X_{57})=(8S,-2S,-6S).
\]
Keep the Gaussian coordinates at indices 72 and 73, and set
all remaining witnesses to zero.  The restriction is the
same $R(T)+a^2+b^2$ as before, with $\kappa=S^2$.
Choosing $a=(1-R)/2$ and $b=i(1+R)/2$ proves that the
translated literal quartic also accepts every input over
every characteristic-zero field containing $i$.

\subsection{Eliminating a private squaring coordinate preserves the system}

In the preceding 1339-coordinate table, the coordinate
$z=X_{81}$ occurs only in the two rows
\[
                       X_{80}^2-z,
                         \qquad R_0+32z,
\]
where $R_0$ is independent of $z$.  Delete the first row and
$z$, and replace the second row by $R_0+32X_{80}^2$.
This gives 1338 coordinates and 1295 quadratic residuals.
Over every field, solutions of the new residual system lift
uniquely by $z=X_{80}^2$, and every old system solution is
obtained this way.

The literal sums of squares need a separate statement.
Every zero of the new sum lifts to a zero of the old sum
over every field, because the deleted row then vanishes and
the remaining rows agree.  The converse follows over fields
with a real embedding: a zero sum of squares forces each
row to vanish, and hence forces $z=X_{80}^2$.  Over a
general nonreal field, an old zero might instead arise by
cancellation among nonzero rows.  System elimination alone
does not prove that such a zero descends to the new sum.
We make no unrestricted converse assertion for literal sums.

Nevertheless the old generic $\mathbb Q(i)$ section has
$X_{80}=X_{81}=0$, so it does descend.  Its Gaussian
coordinate indices 96 and 97 become 95 and 96 after deletion.
Thus the new 1338-coordinate quartic, too, has input
projection equal to the whole field whenever $i$ belongs
to that field.  The script
\nolinkurl{check_quartic_base_change_20261106_v2.py}
verifies these new generic sections from the full sparse
tables and, independently, by squaring the restricted
residuals.  It writes a separate receipt, leaving the
preceding four certificates unchanged.

\subsection{Inlining three fixed constants}

The final ordinary-coefficient refinement substitutes the fixed
integers $3^{32}$, $3^{64}$, and $3^{128}$ for coordinates
58, 59, and 60 of the 1338-coordinate system, and deletes
their three defining rows.  The remaining system has 1335
coordinates and 1292 quadratic residuals.  These are private
polynomial graph equations, so substitution and reinsertion
give inverse maps on system solutions over every field.
As in the preceding elimination, the new literal sum of
squares always lifts to the old one; its converse follows
over fields with a real embedding, and is not asserted over
all nonreal fields merely from the system equivalence.

The largest newly introduced coefficient in a residual is
$3^{256}$, whose bit length is 406.  The associated new
quartic coefficient $3^{512}$ is less than $2^{812}$.
The full unchanged bound $H<2^{1020}$ therefore persists,
as also verified from the complete expanded coefficient
table.  The generic Gaussian section now uses coordinates
92 and 93, with all other witnesses zero.  Direct restriction
of the new table again gives $R(T)+a^2+b^2$ for an explicitly
recorded constant-plus-$T^4$ polynomial $R$; the same Gaussian
witness formulas prove the whole-field projection over any
characteristic-zero field containing $i$.  This section is
checked directly, rather than assuming that an old
cancellation solution survives fixed-constant substitution.

\subsection{The preferred degree-72 norm specifications}

All three transformed residual systems fit the same degree-36
cyclic division norm: their residual counts are 1272, 1295,
and 1292, respectively, each at most $36^2=1296$.
The resulting norm polynomials have 1315, 1338, and 1335 total
coordinates, respectively.  Their degree is exactly 72 and
their coefficient of $X_0^{72}$ is one, since the first
residual still contains $+X_0^2$ and all other residuals
are independent of the input coordinate.

The new translated residuals have only one coefficient
involving $S$, namely $8S$ at $U$ in the anchor row.  The
coefficient norm of every constant part is less than
$2^{511}<S$.  Hence every residual has coefficient norm
strictly less than
\[
                            9S=3^{2^{27}+2}.
\]
Since $3^5<2^8$,
\[
               \log_2(9S)<\frac85(2^{27}+2)=214748368.
\]
Combining this with $36!72^{36}<2^{361}$ gives the explicit
bound
\[
 H(P_{1315}^{\mathrm{norm,tr}})
    <36!72^{36}\,3^{36(2^{27}+2)}
    <2^{7730941609}.
\]
The ordinary 1338- and 1335-coordinate residuals still have coefficient
norm at most $2^{511}$, giving
\[
                       H(P_{1338}^{\mathrm{norm}}),\quad
                       H(P_{1335}^{\mathrm{norm}})
                                         \le2^{18757}.
\]
The complete specifications are
\nolinkurl{norm_repair_1315_translated_degree72_20261106_v2.json.gz}
\nolinkurl{norm_repair_1338_degree72_20261106_v2.json.gz},
and
\nolinkurl{norm_repair_1335_degree72_20261106_v2.json.gz}.
Each supplies the full residual list, cyclic matrix, and
cyclotomic coefficient-extraction rule.  Neither expands
the full degree-72 coefficient table or introduces a
quantified cyclotomic coordinate.

The fixed norm is faithful over $\mathbb Q(i)$ and over
extensions of degree coprime to 36, by the preceding local
and module arguments.  For any other specified number field,
the field-dependent prime construction supplies a faithful
degree-36 norm with the same coordinate counts and degree
72, with its stated prime-dependent coefficient bound.
Because these transformations preserve the residual
systems over all fields, these norm polynomials preserve
their input projections wherever the combining norms are
faithful.

In particular the unshifted scalar-contact contradictions
still hold for the transformed residual systems and their
faithful norm combinations.  Over a field with a real
embedding they also hold for the literal transformed sums
of squares.  These facts concern the unshifted artifacts;
they do not extend the contact counterexample to the
separately checked shift-2402 constructions.  Neither
transformation changes the unresolved arithmetic question
of defining $\mathcal O_K$ over a general number field.


\section{The remaining optimization problem}
\label{sec:strategy}

The preferred construction in this report uses forty-one total
coordinates, including the input. Its degree is bounded by
$2^{9566673}$, and its arithmetic interpretation remains conditional
on the named source inputs and Sun's simultaneous prime-pattern theorem.
The forty-three-coordinate semilocal alternative has the substantially
smaller degree bound $2^{236}$. For degree four, the supplied counts
are 1315 and 1335, the latter with coefficient height less than
$2^{1020}$. These are constructions and upper bounds; no optimality
claim follows from the comparisons.

\subsection{What a reduction to ten coordinates would have to accomplish}

The forty retained coordinates of the preferred source consist of the
input, eighteen components for nine elliptic abscissas, one common
slope, and twenty parameters in ten pole predicates. The finite-algebra
compiler adds one final witness. It has already removed the ordinate
roots, pole-conic roots, and the arithmetic evaluation coordinates.
The direct Lucas test removes the remaining two retained proxy
coordinates and their third, algebraic proxy. Replacing this compiler
by another primitive-element argument would leave the same retained
base and would not, by itself, decrease the count.

The geometric results above explain why each independent elliptic
point cannot simply be encoded by one rational parameter and a finite
algebraic choice. The specific cyclic point group is dense in its
two-dimensional Weil restriction. The denominator results also exclude
compressing independent parity conditions by a rank-deficient list of
monomials, or replacing the entire integrality condition by finitely
many fixed rational pullbacks of perfect-denominator predicates.
These are restrictions on particular methods, not general lower bounds
of eighteen or forty coordinates.

The most promising remaining arithmetic question is whether several
point choices or pole witnesses can be tied together while preserving
simultaneous rational solvability. A proposed sharing must prove this
simultaneous assertion. Separate existence for each pole conic, or
local solvability at every prime for each one independently, does not
provide a common rational parameter. Another possibility is an integer
relation with a different height argument requiring fewer independently
chosen elliptic labels. The direct-polynomial congruence lemma supplies
a precise interface for testing such relations, including the condition
$k>12\deg F$ for its displayed height estimate.

No construction with two through ten total coordinates has been proved
here. Two total coordinates are impossible, even for arbitrary finite
systems. The unrestricted cases of three through ten coordinates remain
unresolved by this work. For quadratic systems the proved restrictions
are stronger: at least five total coordinates and at least three
independent quadratic residuals are necessary. A successful quadratic
system in five through ten coordinates would immediately give a
quartic in exactly the same coordinates over $\mathbb Q$.

\subsection{Degree and coefficients at a fixed coordinate count}

The fixed elliptic exclusion product dominates the enormous degree of
the forty-one-coordinate construction. A better explicit height
constant, or a different size condition costing no new coordinates,
would be more consequential than a small improvement to its finite
algebra. The semilocal alternative illustrates this trade explicitly.
The remaining finite ranks are $2^{28}$ and $2^{29}$; their determinant
and separating-parameter bounds are conservative, and sparse formulas
could improve degree and coefficient bounds without saving coordinates.

For quartics, the preferred 1315-coordinate version also uses an
invertible affine change of its three private anchor witnesses, giving
exact primitive coefficient height $64\cdot3^{2^{28}}$. Removing the
last constant-chain square and inlining its first three constants gives
the 1335-coordinate alternative at height less than $2^{1020}$.
The direct field-valued Lucas test has reduced both the
number of tests and the lengths of their shared symmetric-product
computations. Further useful computation should operate on the actual
quadratic residual system and preserve an exact reconstruction
certificate. Changes of residual basis and invertible coordinate
changes may also reduce coefficient height. All comparisons must use
primitive integer coefficients after clearing numerical denominators.
Computing a huge fixed integer with extra squaring coordinates remains
a distinct trade, which is why both quartic tables are supplied.

\subsection{The arithmetic hypotheses still requiring independent review}

The most consequential review remains the source manuscript's
restricted $2$-converse, its fixed-field coefficient transfer, and
the complete five-contact height argument with its explicit constant.
The shared-point and direct-Lucas constructions additionally use
Sun's Theorem~4.1. The stated theorem and its order-compatibility
consequences were checked; the complete analytic proof of its
prime-pattern estimates has not been independently reproduced here.
The distinction between the unconditional Artin-symbol theorem of
Frei--Tinkov\'a and their conditional primitive-root theorem matters
\cite{FreiTinkova2025}. Sun's further large-prime sieve and its use
of polynomial-prime transference \cite{TaoZiegler2018} remain external
inputs, not conclusions of the polynomial checkers.

This extra input has a specific role. A common comparison index must
satisfy simultaneous congruences at different selected primes. Their
compatibility requires
\[
                    \gcd(d_{p,-},d_{q,+})=4
                         \qquad\text{for all }p,q.
\]
An odd common divisor would make the residues zero and four
incompatible. One cannot justify the common index by reusing only the
original private-witness existence statements. The conditional proofs
in this report retain the stronger simultaneous hypothesis explicitly.

\subsection{Number fields and the limits of a uniform formula}

The explicit $\mathbb Q(i)$ sections and the rejected integral contact
inputs show that the literal newer polynomials do not extend unchanged
to rings of integers. Translating the input square removes the finite
contact obstruction; a field-adapted division norm repairs conjunction
without adding variables. These are useful repairs, but neither restores
the fixed cyclic elliptic-group argument. The totally real quartic
example gives a concrete point outside that group satisfying every
finite anchor exclusion.

The general constructibility theorem further shows that no one fixed
finite existential formula can define the rings of integers in all
number-field extensions. The field-dependent existence question remains
open here. Under the claimed rational definition, an existential
definition of the embedded rational field would permit the integral
basis construction in Theorem~\ref{thm:nf-equivalence}. Undecidability
or MRDP over the ring of integers does not supply that rational-field
isolation. The next arithmetic work should address this missing
assertion or develop a different source adapted to each field.


\section{A conditional geometric obstruction}
\label{sec:mazur}

The numerical lower bounds and the conic obstruction above are
unconditional. There is also a useful
conjectural obstruction of a different kind. Mazur's finite-component
conjecture states that the closure of the rational points of a variety
over $\Q$ in its real points has finitely many connected components
\cite[Conjecture 3]{Mazur1995}. It would rule out a Diophantine
definition of $\Z$ with any degree or number of variables. It must
not be used as an established theorem against the proposed source.

\begin{proposition}
If a polynomial defines $\Z$ in $\Q$, the real closure of its rational
zero set has infinitely many connected components.
\end{proposition}
\begin{proof}
Let $V$ be its affine zero set and let $C$ be the closure of $V(\Q)$
inside $V(\mathbb R)$. The input projection is continuous. Its value
on every convergent sequence of rational points is a limit of integers,
and hence an integer; therefore its image on $C$ is contained in $\Z$.
Every integer occurs already on $V(\Q)$. A connected component of
$C$ has connected image in the discrete set $\Z$, hence has only one
input value. Infinitely many components are necessary.
\end{proof}

Mazur's formulation uses reduced schemes of finite type over $\Q$,
so it applies directly to the reduced affine zero set used here
\cite[Section 1 and Conjecture 3]{Mazur1995}. This topological obstruction
clarifies the strength of the proposed theorem. It does not add a
premise to our exact algebraic transformations.


\appendix
\section{Permissible precision reductions}

The following elementary argument isolates exactly what the growth
comparison needs. It does not establish the arithmetic estimates from
which it starts.

\begin{proposition}\label{prop:precision-thresholds}
Let $a,B,n,\Delta\in\mathbb Z$, with $B\ge a\ge1$, let
$\beta\in\mathbb Q$, and let $H\ge1$ be real.
Write $h(\beta)=\log\max\{|b|,c\}$ when $\beta=b/c$ is in lowest
terms with $c>0$. If $\beta\ne n$, let $N$ be the absolute value
of the numerator of $\beta-n$ in lowest terms.
Suppose an arithmetic certificate supplies the following estimates
for some $\kappa>48$,
$\lambda>12$, and $\mu>24$:
\begin{enumerate}
\item if $\Delta\ne0$, then
$B^2-6<7H|\Delta|^{24/\kappa}$ and $|\Delta|<5B^4$;
\item if $\Delta=0$, then $\log B>a-3$;
\item if $\beta\ne n$, then
$B^2<7HN^{24/\lambda}+6$ and $N\le2e^{h(\beta)}B$;
\item $h(\beta)\le H a^{24/\mu}$.
\end{enumerate}
Put
\[
 C_\kappa=7\,5^{24/\kappa}+6,\qquad
 C_\lambda=7\,2^{24/\lambda}+6,\qquad
 A_\lambda=\frac{\lambda}{2\lambda-24},\qquad
 B_\lambda=\frac{12}{\lambda-12}.
\]
If
\[
 a>\max\left\{6,
 (C_\kappa H)^{\kappa/(2\kappa-96)},
 4A_\lambda\log(C_\lambda H),
 (4B_\lambda H)^{\mu/(\mu-24)}\right\},
\]
then $\Delta=0$ and $\beta=n$.
\end{proposition}
\begin{proof}
For $\Delta\ne0$, absorb the constant six using $B,H\ge1$ to obtain
$B^{2-96/\kappa}<C_\kappa H$, contradicting $B\ge a$.
Thus $\Delta=0$. If $\beta\ne n$, put $r=24/\lambda<2$.
The third hypothesis and $B,H\ge1$ imply
\[
 B^{2-r}<C_\lambda H e^{r h(\beta)}.
\]
Consequently
\[
 \log B<A_\lambda\log(C_\lambda H)+B_\lambda h(\beta)<a/2<a-3,
\]
contrary to the second hypothesis.
\end{proof}

The thresholds $48,12,24$ are thresholds of this particular comparison
of exponents. This proof does not show that the corresponding boundary
values are impossible by a different argument.

\begin{corollary}\label{cor:scalar32}
The choices
\[
 (\kappa,\lambda,\mu)=(52,13,32),\qquad a>(21H)^{13/2}
\]
satisfy the conclusion of Proposition~\ref{prop:precision-thresholds}.
If $\log_2H<4783227$, the existing type of anchor with twenty-six
successive squarings, giving $a\ge3^{2^{25}-1}$, is sufficient.
\end{corollary}
\begin{proof}
Here $C_\kappa<21$, $C_\lambda<32$,
$A_\lambda=13/2$, $B_\lambda=12$, and $24/\mu=3/4$.
The assumed lower bound for $a$ implies $a>(48H)^4$ because
$21^{13/2}>48^4$ and $H\ge1$; hence $h(\beta)<a/48$.
It also implies $a>26\log(32H)$: use $\log(32H)<32H$ and
$21^{13/2}>832$. Thus the last contradiction reads
$\log B<(13/2)\log(32H)+12h(\beta)<a/2<a-3$.
For the numerical anchor,
\[
 \log_2\big((21H)^{13/2}\big)
 <\tfrac{13}{2}(5+4783227)=31091008,
\]
whereas
\[
 \log_2\big(3^{2^{25}-1}\big)
 >\tfrac32(2^{25}-1)=50331646.5.
\]
\end{proof}

For the source system's scalar batch the order $48$ is used solely to
obtain $M_5(\beta)^{48}\mid s$, hence
$h(\beta)\le H\sqrt{|s|}$. Replacing that order by $32$ gives instead
$h(\beta)\le H|s|^{3/4}$, which Corollary~\ref{cor:scalar32} shows
is sufficient. Completeness still allows all the required divisibilities:
the index $s$ is chosen in an arbitrarily divisible nonzero subgroup.
The batch's recursive power construction uses eight multiplications
for the required powers $r^2,r^{96}$, and six for $r^2,r^{64}$.
Thus this change removes two source power definitions. Its exact effect
on the compiled polynomial must be measured after the full pipeline.

For comparison, orders $49$ and $52$ require the same thirteen terms
in the logarithm truncation used by the source, but the source's shared
power recursion is less economical for $(49,13)$: it uses twelve
multiplications for $r^2,r^{26},r^{98}$ and nine for
$r^2,r^{26},r^{104}$. A smaller numerical precision does not automatically
mean fewer variables. Joint addition-chain optimization is the correct
comparison.


The source-preserving released-count packages retain the original precisions.
The $(52,13)$ experiment on the original source did not improve its
coordinate count, because its power chains were less economical.
The new shared-point quartics of Section~\ref{sec:new-quartic}
use precision $52$ and scalar order $32$, together with stronger
size conditions and a different evaluation scheme. Their certified
counts are $1367$ and $1391$, as proved there.


\section{Audit of the arithmetic source and the remaining dependencies}
\label{sec:arithmetic-source-audit}

\subsection{Scope and conclusion}

The source examined here is \texttt{quartic\_12504.tex} at commit
\nolinkurl{28e87b1b62af0effd954605af7f70b2699dd5672}, with SHA--256
\begin{center}\footnotesize\ttfamily
38d458ba4178083766c5ae56a894da595fe27170c1ef07f6b49750fc5edf7c37.
\end{center}
The line numbers below refer to this fixed file.  Labels are also given
because subsequent edits can change line numbers.  We examined the
arithmetic interpretation of the released reciprocal comparison source,
its principal local arguments, its final height comparisons, and selected
substantial inputs to completeness.  We did not independently verify the
entire new restricted converse argument or every application of the
arithmetic-geometric height theorems.

\emph{No rational counterexample to the source system, and no invalid
step in the elementary local and numerical deductions listed below, was
found in this audit.}  This statement is substantially weaker than an
independent proof of the asserted Diophantine definition.  The exact
polynomial transformations elsewhere in this paper preserve the projected
rational zero set.  Identifying that set with $\mathbb Z$ still requires
the arithmetic dependencies isolated here.

Two dependencies deserve particular emphasis.  The pole predicate's
one-sided obstruction has an elementary valuation proof, but its
completeness uses a new restricted $2$-converse.  The asserted
five-contact height theorem has a lengthy geometric and modularity proof;
a calculation of its displayed constant is not a proof of that theorem.
Neither dependency is discharged by checking the polynomial compiler.

\subsection{The local obstruction and norm argument}

For $b\in\mathbb Q$, put $B=17b$ and let $\Psi(b)$ assert the existence
of $t,U,e,p,q,s,h\in\mathbb Q$ satisfying
\[
 U(1+bt^2)=b(1-bt^2),\qquad p^2=B(1-e)(1+3e),\qquad
 h=e(B-U^2),\qquad Bq^2=(s^2-Bh)(s^2+3Bh).
\]
The relevant source statements are
\texttt{hyp:o8-pole} and \nolinkurl{lem:original-seven-parity-soundness}
(lines 7532--7597).

\paragraph{Checked one-sided statement.}
If $b\ne0$ and $\Psi(b)$, then $v_\ell(b)$ is not negative and odd
for any prime $\ell\notin\{2,3,17\}$.  Indeed, suppose
$k=v_\ell(b)<0$ is odd.  The first equation rules out $1+bt^2=0$.
The valuations of $1-bt^2$ and $1+bt^2$ agree, so $v_\ell(U)=k$;
this also holds when $t=0$.  Consequently
$v_\ell(B-U^2)=2k$.  The second equation rules out $e=0$ and forces
$v_\ell(e)=0$: a nonzero valuation would make its right side have
odd valuation.  Thus $h\ne0$ and $v_\ell(Bh)=3k$ is odd.
For any $s$, the right side of the final equation is nonzero and has
even valuation; the two terms in each factor have unequal valuation
when $s\ne0$, and the assertion is immediate when $s=0$.
The left side is either zero or has odd valuation, a contradiction.
This proof includes $e=1$, $e=-1/3$, and $p=0$.  No omitted inverse
condition on a witness is needed.

\paragraph{Checked norm propagation.}
The binary operation used by the compiler is
$N_r(a,b)=a^2+rb^2$.  If $R=v_\ell(r)$ is odd, terms with both inputs
nonzero have distinct valuation parity.  Hence
\[
 v_\ell N_r(a,b)=\min\{2v_\ell(a),R+2v_\ell(b)\}.
\]
In particular a tree containing the constant leaf $1$ cannot have zero
root at such a prime.  If the constant is on the all-left path and
$R<0$, the root $W$ has $v_\ell(W)\le0$; then $rW^2$ has negative
odd valuation, contradicting the pole obstruction.  Thus $R>0$.
The same obstruction then implies
$R+2v_\ell(W)>0$.  Descending the preceding minimum formula shows
that a leaf $X/r^e$ has valuation greater than $-R$.  Therefore
\[
 v_\ell(X)>(e-1)R\ge e-1,
 \qquad\text{and hence}\qquad v_\ell(X)\ge e.
\]
This verifies the strict inequality and integer rounding used in the
mixed-precision batches; it also explains why a global condition
$W\ne0$ is unnecessary for soundness.  See lines 7789--7804 and
8080--8120 of the source.

The finite squareclass correction in lines 7755--7804 is separate from
this soundness argument.  Its extra factors have support in
$\{2,3,5,7,11,13,17\}$, already excluded from the relevant contacts.
Completeness chooses the root to be a unit at primes dividing the raw
integer or the correction.  This condition is needed: multiplying an
arbitrary rational square by an integer does not by itself guarantee a
square reduced denominator.

\subsection{The reciprocal logarithm, common point, and final estimates}

\paragraph{Checked formal approximation.}
For $E:y^2=x^3-Dx$, $D=11-6\sqrt2$, set $f=1/x$, $z=-x/y$,
$a_0=1$, $a_1=-2/5$, and
$a_j=(4j-5)a_{j-1}/(4j+1)$.  The relations
\[
 f=z^2(1-Df^2),\qquad
 4w(1-w)G'(w)+(1+w)G(w)=1-w,
 \quad G(w)=\sum_{j\ge0}a_jw^j,
\]
give $\log_E(z)=zG(Df^2)$.  At a good split prime $p>76176$,
with $r=v_p(z)\ge1$, one has $v_p(f)=2r$.
For $e\in\{16,64\}$, $m=e/4$, and
$L_e=z\sum_{j=0}^{m-1}a_j(Df^2)^j$, the source bounds the entire
tail, including indices whose denominators are divisible by $p$:
\[
 v_p(\log_E(z)-L_e)\ge(4m+1)r,
 \qquad L_e/\log_E(z)\in1+p^e\mathbb Z_p.
\]
The factorial denominator bound used for the tail is sufficient.
For $4j+1\ge p$ one has $j\ge19044$ and
$v_p(a_j)\ge-(4j+1)/(p-1)\ge-j$, which gives the required bound
and convergence.  Truncating before the first denominator divisible by
$p$ alone would not justify the infinite-tail estimate; the source does
not make that shortcut.  See \texttt{lem:quartic-reciprocal-log},
lines 8753--8816.

The separate global nonvanishing argument is also valid for the
truncations used here.  In $\mathbb Q(\sqrt2)$ the valuation extending
$v_2$ has value group $\frac12\mathbb Z$, $v(D)=0$, and $v(a_j)=1$
for $j\ge1$.  If $v(f)\ge0$, the constant term of $G_J(Df^2)$
is uniquely least; if $v(f)<0$ and $J\ge2$, its last term is uniquely
least.  Thus the logarithm inverses introduce no exceptional rational
point.  See \nolinkurl{lem:quartic-reciprocal-nonzero}, lines 8818--8842.

\paragraph{Checked compatibility of the common auxiliary point.}
In the completeness construction the private congruences are
\[
 h_p\equiv0\pmod{d_{p,1}},\quad
 h_p\equiv4\pmod{d_{p,2}},\quad
 h_p\equiv4\pmod{p^{64+v_p(j)}},\qquad
 j=\operatorname{lcm}_{p}d_{p,2}.
\]
The conditions $\gcd(d_{p,1},d_{p,2})=4$ and
$p\nmid d_{p,1}d_{p,2}$ suffice for these congruences at each
individual $p$.  The construction does not need pairwise coprimality
between group orders belonging to different primes.  A selected prime
may divide $j$ because of another group order; the term $v_p(j)$
already accounts for this possibility.  Avoiding the finitely many
forbidden values of $h_p$ preserves the point-chart guards.
The same verification does not authorize replacing all $h_p$ by a single
shared $h$: that would require compatibility between congruences with
moduli $d_{p,1}$ and $d_{q,2}$ for different primes $p,q$, and the source
provides no such cross-prime hypothesis.  Sharing this witness is therefore
a blocked optimization unless an additional prime-pattern theorem or a
different construction is supplied.  This checks the stated use of the
ordinary-prime pattern, rather than independently reproving the full
prime-pattern theorem; see
\texttt{subsec:o8-common-J-proof}, lines 8394--8468, and lines
8964--8984.

\paragraph{Checked final numerical implication.}
Assume the stated height and radical-divisibility inputs, and put
$a=|s|$, $B=\max(|n|,|s|,|t|,|z|)$, $H=H_*$.
The anchor gives $a>(20H)^2$.  If the integer Lucas residual
\[
 \Delta=t^2-s(s-2)tz+(s-2)^2z^2-1
\]
is nonzero, the source obtains
$B^2-6<7H|\Delta|^{3/8}<14HB^{3/2}$.
Dividing by $B^{3/2}$, with $B\ge1$, yields
$\sqrt B<14H+6\le20H$, a contradiction.  Thus $\Delta=0$.
The elementary Lucas-unit argument then gives $\log B>a-3$.

If $\beta\ne n$, let $N$ be the absolute numerator of
$\beta-n$ in lowest terms.  The precision-$16$ argument gives
\[
 B^2<7HN^{3/2}+6,\qquad N\le2e^{h(\beta)}B.
\]
Consequently
$\log B<2\log(32H)+3h(\beta)$.
The scalar estimate $h(\beta)\le H\sqrt a<a/20$ and
$\log(32H)<a/20$ make the right side less than $a/4<a-3$.
There is no loss of the exponent when the beta precision changes
from $64$ to $16$.  Finally $\beta=n$ and
$X^2=\beta-2\in\mathbb Z$ imply $X\in\mathbb Z$ by coprime
numerator and denominator.  See \nolinkurl{thm:quartic-source-arithmetic},
lines 8867--8985.  This is a conditional verification of the final
implication, not a verification of its height hypotheses.

\subsection{The two substantial arithmetic dependencies}

\paragraph{Restricted converse and pole completeness.}
The exact converse needed is
\[
 \operatorname{corank}_{\mathbb Z_2}
  \operatorname{Sel}_{2^\infty}(E_\ell/\mathbb Q)=1
 \quad\Longrightarrow\quad
 \operatorname{ord}_{s=1}L(E_\ell,s)=1,
 \qquad E_\ell:y^2=x(x-\ell)(x+3\ell),
\]
for positive squarefree $\ell\equiv11\pmod{12}$.
It is \texttt{hyp:RC} at lines 1016--1028.  The manuscript supplies
its own long argument in Sections \texttt{sec:rc},
\texttt{sec:arithmetic}, and \texttt{sec:limits} (lines 1771--5755),
with appendices at lines 9589--13699.  It is not an immediate
consequence of ordinary two-descent or of the published forward
analytic-rank-zero-and-one theorems.

The distinction is essential.  A two-Selmer group of dimension three,
in the presence of full rational two-torsion, implies
$r+\dim_{\mathbb F_2}\Sha[2]=1$.  The Cassels--Tate pairing permits
one to infer full $2^\infty$-Selmer corank one, as in
\texttt{lem:corank}; it does not by itself remove a divisible
$\mathbb Q_2/\mathbb Z_2$ summand of $\Sha$ or prove analytic rank one.
The new restricted converse is used precisely at this remaining step.

We checked the algebraic shape of the simultaneous realization and
its reduction to this converse, but have not certified all its local
descent tables or the new converse.  The most important further audit
within the latter is the passage from finite Heegner classes and local
conditions to the uniform valuation bound and the fixed-field
coefficient transfer:
\nolinkurl{ex:limits-uniform-lower-bound} (lines 4920--5217),
\nolinkurl{ex:limits-fixed-field-clearing} (lines 5498--5657), and
\nolinkurl{ex:limits-coefficient-transfer} (lines 5658--5754).
These statements assert uniformity over arbitrarily large active
supports and are part of the new arithmetic argument.  The present
report neither proves nor disproves them.

\paragraph{Five-contact height theorem and its exact constant.}
The required statement is \texttt{hyp:o8-height}
(lines 6666--7520):
\[
 h(a)\le H_*M_5(a)^{24},\qquad \log_2H_*<4783227,
\]
where the contacts are $\infty,0,1,2401,-343/32$ and the product
$M_5(a)$ omits primes at most $76176$.
The displayed bound on $H_*$ is needed for the released anchor;
an unspecified larger constant cannot simply be substituted.
The source derives it using the discriminant-$210$ Shimura curve,
a degree-at-most-$576$ normal fiber field, a comparison with stable
Faltings height, and separate CM and non-CM estimates.

Two plausible objections were checked against primary sources and
should not be presented as gaps.  Yuan's Theorem 1.1 and Section 2.1
of \emph{Explicit Kodaira--Spencer map over Shimura Curves}
\cite[Theorem 1.1 and Section 2.1]{Yuan2024} do provide the stated maximal-stack
regularity and semistability and the isometry whose image is
$d_B\omega^{\otimes2}$.  Pan's Theorem 1.0.4 of
\emph{The Fontaine--Mazur conjecture in the residually reducible case}
\cite[Theorem 1.0.4]{Pan2022} has no extra residual-image restriction
at $p=5$, provided the representation is irreducible, odd, finitely
ramified, potentially semistable, and regular.  Thus residual
reducibility alone does not invalidate the source's use of Pan.

What remains to be checked is the complete chain of hypotheses in the
source's application: the canonical normalizer action and constant
fiber torsor, the integral height comparison with its normalization,
the scalar cocycle splitting and ramification bounds, and the
modular-Jacobian transfer preserving the displayed constant.
The principal non-CM passage is
\texttt{sec:h5-level-repair}, lines 6904--7345; the CM estimate follows
at lines 7349--7500.  These are substantive mathematical arguments,
not consequences of the finite numerical certificate.  Source comments
marking a repair as awaiting review record a review status; they are
not, by themselves, a proof of an error.

\subsection{Consequences for the optimization claims}

The unconditional conclusions established by an exact transformation
certificate have the form
\[
 \{X:\exists\mathbf y\ F(X,\mathbf y)=0\}
 =\{X:\exists\mathbf z\ F_{\rm new}(X,\mathbf z)=0\}.
\]
They remain meaningful independently of the source's main arithmetic
claim.  Until the dependencies above have been independently verified,
the further conclusion that either set equals $\mathbb Z$ must be
stated as conditional on the arithmetic theorem of the source.
A counterexample to an individual intermediate lemma would require
tracing its use and possible repair before concluding that the released
polynomial fails; no such counterexample is supplied by this audit.


\section{Reproducibility and the scope of the computations}
\label{sec:reproducibility}

The proposed files are confined to \nolinkurl{research/optimization/}
in a local checkout of the user's fork at the revision specified in
Section~\ref{sec:scope}. The released polynomial, source templates, and
manuscript are unchanged. No commit, push, or pull request was made.
The complete paper is the standalone file
\nolinkurl{diophantine_optimization.tex}; it needs no external inputs or
BibTeX database. Its bibliography gives the primary references used here.

\subsection{Exact objects and independent checks}

The directories \nolinkurl{quartic-anchor-packed/} and
\nolinkurl{quartic-anchor-packed-large/} contain the 12,493- and
12,465-variable quartics, respectively. Each supplies the complete
integer coefficient table, quadratic residuals, variable map,
transformation certificate, and an independent checker. The
12,449-variable package \nolinkurl{quartic-12449/} stores its full
coefficient table over $\Z[M]$ with the fixed specialization
$M=3^{2^{26}}$. Its checker performs all substitutions and all
sum-of-squares expansions exactly, and proves that the monomial count
and coefficient height survive the specialization.

The directories \nolinkurl{polynomial-853/} and
\nolinkurl{polynomial-854/} contain the complete polynomial circuits,
finite parameter sets, construction reports, and verification receipts.
The predecessor \nolinkurl{polynomial-855/} is included as a checked
dependency. Every operation in these circuits is integer addition,
multiplication, or a nonnegative integer power. The two finite parameters
are replaced by each of their specified integer values and the resulting
branch polynomials are multiplied. There is no rational division or
unquantified variable in the final polynomial specification.

The 855-variable structural checker verifies the complete circuit
syntax, denominator coverage, residual and branch inventories, and
degree and coefficient-norm bounds. Its six reproducible finite-field
comparisons of all 297 rational residuals are explicitly \emph{sampled
checks}, not proofs of identities. The source correspondence is justified
by the algebraic proofs in this paper, the independent review
\nolinkurl{compact-audit-855.md}, and inspection of the pinned templates.
The separate square-root checker verifies twelve identities by exact
integer polynomial arithmetic. The 854- and 853-variable transformation
checkers import no producer code; they prove their anchor and primitive-
element identities exactly. In particular the latter checks every one
of the 295 substitutions by expression-tree induction on
$H\mapsto Q^kH(P/Q)$ and reruns the exact pole-denominator check.

The number-field section is independently verified in $\Q(i)[X]$
against both the complete released coefficient table and all released
residuals. The finite-fiber inventory checker verifies the 557 retained
coordinates and the finite-choice count. The proof of geometric
finiteness is mathematical. The explicit finite-algebra compiler and its
independent checker, described in Section~\ref{sec:explicit-finite-algebra},
now supply the indexed 558-variable polynomial, including all exceptional
cases. The enormous final matrices and products are not expanded.

All primary verification scripts use the Python standard library. From
the root of the pinned repository one may run
\begin{verbatim}
python3 -E -B research/optimization/verify_all.py
\end{verbatim}
The script records the commands and their outcomes in
\nolinkurl{verification/all-checks.json}; detailed logs and individual
receipts remain separate. The checkers are reproducible finite
calculations, not formal verification of the source arithmetic theorem.

\subsection{Identifying the principal polynomial files}

The following SHA-256 digests identify the exact objects used. The
12,449-variable digest is for the uncompressed JSON coefficient table;
the other quartic digests are for the uncompressed text tables.
\begin{description}
\item[Released 12,504-variable quartic:]
\mbox{}\\\nolinkurl{7f5aa3117af82a0b323f0a4e74a18599204f5540408106d3ce99612b5ed830ca}
\item[12,493-variable quartic:]
\mbox{}\\\nolinkurl{06c2eeb37871341588f232d72eee72a49b8a76a2c904dc8465518d280f3378a9}
\item[12,465-variable quartic:]
\mbox{}\\\nolinkurl{de8e8080587ef83e19653e3545a1c15805526ac847547c66dc6cc21dd1bc466c}
\item[12,449-variable quartic:]
\mbox{}\\\nolinkurl{f8813916debae6bd20559e1130ae9786eebf48503582e37da32f1c4b018985d0}
\item[853-variable circuit:]
\mbox{}\\\nolinkurl{67992338a1d904d623d9b1b076f28b9da6b84e44736d34e90784f75f305e5d63}
\item[854-variable circuit:]
\mbox{}\\\nolinkurl{aef6888809817e37e372471fd780f83d0241e58f04ddeafd03fd6b46cfa39d62}
\end{description}

A complete file manifest accompanies the review package. The TeX source
has passed static checks of braces, environments, cross-references, and
bibliography keys. A working LaTeX format was unavailable in the present
environment, so no successful typesetting or visual-layout verification
is claimed. This limitation concerns the document's rendering, not the
exact polynomial checks described above.


\section{Reproducing the constructions with Lucas proxies}
\label{sec:phase2-reproducibility}

The preferred new constructions are stored under
\nolinkurl{research/phase2/}. The complete paper is the standalone file
\nolinkurl{research/optimization/diophantine_optimization.tex}.
All variable counts include the distinguished input. The principal
specifications and their independent verification receipts are:
\begin{center}
\small
\begin{tabular}{r p{82mm}}
\toprule
Total variables & Package relative to \texttt{research/phase2/}\\
\midrule
43 & \nolinkurl{finite_fibre/polynomial-43-conic/}\\
45 & \nolinkurl{finite_fibre/polynomial-45-conic-semilocal/}\\
558 & \nolinkurl{finite_fibre/polynomial-558/}\\
1367 & \nolinkurl{quartic/quartic-best/}\\
1391 & \nolinkurl{quartic/quartic-small-coefficients/}\\
\bottomrule
\end{tabular}
\end{center}
The first three contain a
\nolinkurl{polynomial-specification.json}; the two quartics contain a
\nolinkurl{polynomial.json.gz}. Every package includes its construction
report and a \nolinkurl{verification.json} describing the exact scope
of its independent check. The complete SHA-256 manifest is generated
with the review archive. File names and hashes identify mathematical
objects; they are not additional mathematical hypotheses.

From the repository root, the single command
\begin{verbatim}
python3 -E -B research/phase2/verify_phase2.py
\end{verbatim}
runs eleven checks and records their outputs in
\nolinkurl{research/phase2/verification/all-checks.json} and its
\nolinkurl{logs/} directory. It checks the frozen source DAGs,
the scalar normalization constants, the conic parametrization and
duplication identities, the precision inequalities, the fixed-point
arithmetic, the three preferred finite-algebra specifications, and
both complete quartic tables. The checkers are separate programs from
the producers. Exact integer and rational arithmetic is used for the
new algebraic checks.

The finite-algebra interpreter also executes every determinant,
resultant, coefficient reversal, and outer finite case on eight small
examples, including nilpotents, excluded components, collisions,
rank drops, irrational residue fields, and zero quadratic roots. It
checks 177 factors and supplies a fully expanded degree-sixteen
example. This is a check of the construction on the specified examples.
The general theorem, and not extrapolation from those computations,
justifies the large instances.

To rebuild the quartics, use the producer with
\nolinkurl{--share-height --precision 52 --scalar-order 32}.
The optional \nolinkurl{--small-coefficients} selects the
1391-variable alternative. The package README files give the complete
commands for rebuilding each source and finite-algebra specification,
including all finite-chart and scalar-normalization options.
The original-source comparison checks remain independently runnable
using \nolinkurl{research/optimization/verify_all.py}.

These receipts establish the stated finite identities, inventories,
and bounds. They do not claim to verify the external prime-pattern
theorem, the source's restricted pole converse, or its full explicit
height theorem. Those arithmetic dependencies are stated in
Section~\ref{sec:scope} and in the source-audit appendix.
The huge indexed products and multiplication matrices are specified
exactly but are not materialized. Only the quartics have complete
expanded coefficient tables. No global minimum is claimed.

The assembly program checks balanced environments and braces, unique
labels, defined cross-references, and defined bibliography keys. A
working LaTeX format was not available in the preparation environment;
no successful typesetting or rendered-layout verification is claimed.
The standalone source can be compiled twice with a normal
\texttt{pdflatex} installation. All additions are prepared for review
on the local branch \nolinkurl{research/polynomial-optimization}.
No commit, push, or pull request has been made.


\section{Reproducing the direct constructions and the field counterexamples}
\label{sec:p3-reproduction}

All files in this section are in the new
\nolinkurl{research/phase3/} directory. The previous complete paper,
the preceding source files, and the released quartic remain unchanged.
The standalone paper is
\nolinkurl{diophantine_optimization_20261106_v3.tex}.
It contains its own bibliography and does not require external
\LaTeX{} inputs. Static checks verify balanced environments and braces,
unique labels and bibliography keys, and defined references and
citations. No rendered-PDF verification is claimed.

\subsection{Indexed formulas and quartic tables}

The four complete coefficient sources are
\nolinkurl{compression/source-41-direct.json.gz},
\nolinkurl{compression/source-43-direct-semilocal.json.gz},
\nolinkurl{compression/source-41-shift2402.json.gz}, and
\nolinkurl{compression/source-43-semilocal-shift2402.json.gz}.
They are produced by \nolinkurl{compression/build_direct_source.py}.
The options \texttt{--anchor finite} and
\texttt{--anchor semilocal} select the size condition;
\texttt{--beta-shift 2402} selects the translated input.

The generic compiler and checker from the preceding construction
are reused unchanged. From the repository root, one example is
\begin{verbatim}
python3 -E -B research/phase2/finite_fibre/compile_shared_source.py \
  research/phase3/compression/source-41-direct.json.gz \
  --output research/phase3/finite_fibre/polynomial-41
python3 -E -B research/phase2/finite_fibre/verify_shared_specification.py \
  research/phase3/compression/source-41-direct.json.gz \
  research/phase3/finite_fibre/polynomial-41/polynomial-specification.json
\end{verbatim}
The other three sources have separate matching specification packages
and passing receipts. The verification concerns integer coefficient
conversion, triangular monicity, guards, all exceptional quadratic
cases, the finite index sets, and degree and coefficient majorants.
It is not a verification of the external arithmetic hypotheses.

The direct quartic producer is
\nolinkurl{quartic/build_direct_quartic.py}. Reproduction uses
\begin{verbatim}
python3 -E -B research/phase3/quartic/build_direct_quartic.py \
  --share-height --precision 52 --scalar-order 32 \
  --output research/phase3/quartic/quartic-direct
python3 -E -B research/phase2/quartic/verify_shared_quartic.py \
  research/phase3/quartic/quartic-direct/polynomial.json.gz
\end{verbatim}
Add \texttt{--small-coefficients} and use the corresponding
small-coefficient output directory for the 1339-variable polynomial.
The checker imports no producer. It checks every gate, the reversible
row operations and eliminations, and the complete polynomial obtained
by independent exact squaring and convolution of the final residuals.
No floating-point or randomized comparison is used.

The preferred tables additionally use the certified transformations in
\nolinkurl{quartic/transform_direct_quartics.py}. The corresponding
independent checker is
\nolinkurl{quartic/verify_final_transforms_independent.py}.
It checks every affine and constant substitution, the inverse coordinate
maps or unique lifts, and both the complete substituted polynomial and
the independently squared residual list. It imports no producer. The
fresh transformed packages preserve their frozen parents for comparison.
The distinction between encoded coefficient entries and coordinate
monomials after specializing a fixed integer is checked separately by
\nolinkurl{compression/check_fixed_coefficient_support.py}.

\subsection{Independent checks for the field assertions}

The following scripts use exact arithmetic and have separate receipts:
\begin{description}
\item[Quartic base change.]
\nolinkurl{number_fields/check_quartic_base_change_20261106_v1.py}
checks the generic $\mathbb Q(i)$ section against both the complete
tables and the independently squared restricted residuals, for all
four quartics.
\item[Rejected contacts.]
\nolinkurl{number_fields/check_contact_exclusions_20261106_v1.py}
checks the three exact input relations in the four unshifted indexed
sources and four quadratic systems. The shifted sources are expressly
outside this counterexample's scope.
\item[The totally real extension.]
\nolinkurl{number_fields/check_real_anchor_20261106_v1.py}
checks the field, point, and doubling calculations in the displayed
degree-eight compositum.
\item[Division norms.]
\nolinkurl{number_fields/build_division_norm_specs_20261106_v1.py}
produces all four determinant specifications. The independently
written
\nolinkurl{number_fields/verify_division_norm_specs_independent_20261106_v1.py}
imports neither this producer nor an earlier checker. It compares
every embedded residual with pinned source bytes, reconstructs each
matrix-entry index, verifies formal Galois conjugacy, checks the
explicit inert primes, and proves the stated exact degrees from the
unit coefficient of the highest pure input power. It also recomputes
the coefficient bounds and the unrounded finite-algebra degree
majorants before their displayed powers-of-two rounding.
\end{description}

The accompanying independent proof review explains the local
division arguments, including the degree-36 norm over
$\mathbb Q(i)$, and the coprime-extension module argument. The norm
determinants are exact finite specifications, not expanded tables.
The counterexamples and norm equivalences are unconditional algebraic
results. They do not certify that the changed source defines any
extension field's ring of integers.

The final affine and constant-chain transformations have separate
number-field certificates and degree-72 specifications with dated
suffix \nolinkurl{20261106_v2}. These are checked independently in
addition to the preceding specifications. In particular, elimination of
a defining residual preserves common zero sets over every field, but
does not by itself preserve every cancellation among nonzero squares
over a nonreal field. The explicit Gaussian sections are therefore
verified again for each new literal quartic.

\subsection{Rebuild and integrity records}

\nolinkurl{assemble_phase3.py} builds this paper from the preceding
complete report and the new mathematical supplements, and checks that
the preceding report's bytes are unchanged. The verification index
records the passing receipts and hashes of the checked artifacts.
\nolinkurl{make_checkpoint.py} packages the complete research tree,
the pinned source manuscript, and the unchanged released baseline.
It hashes the same bytes it writes and verifies the resulting ZIP
archive. Its manifest makes the checkpoint independently inspectable.
No commit, push, or pull-request submission is part of these scripts.


\section{Verification of the new source specification}
\label{sec:direct-source-review}

The computational checks in this section concern the explicitly
recorded source equations and their translation.  They do not certify
the external arithmetic hypotheses used to prove the integer input
projection.  That distinction is essential: an exact compiler can
correctly translate an arithmetic system whose completeness or soundness
has not been independently established.

The program \texttt{check\_direct\_source.py} compares the abstract
syntax trees of the frozen producer with proxies and the direct-Lucas
producer.  All eight auxiliary methods are identical.  After removing
the two retained proxies, their finite quadratic root, and their three
comparison tests, the remaining construction statements are identical
apart from the direct polynomial test, its corresponding counts and
metadata, and the optional shift in $\beta$.  The previously preferred
sharing, scalar normalization, conic parametrization and fixed scale
options are now the defaults; their mathematical formulas are unchanged.

The new direct test is also checked by a separate sparse polynomial
implementation over $\mathbb Q$.  Substitute
\[
 u=u_0+u_1\sqrt2,\quad v=v_0+v_1\sqrt2,\quad
 w=w_0+w_1\sqrt2.
\]
Evaluation of the producer's actual expression agrees exactly with
the independently expanded rational and $\sqrt2$ components of
\[
                  v^2-u(u-2)vw+(u-2)^2w^2-1.
\]
These components have respectively $23$ and $19$ nonzero monomials
and total degree four.  The complete coefficient dictionaries are
in \texttt{direct-source-check.json}.

The normalization is verified by exact Laurent-polynomial identities.
For each precision $m\in\{1,16,64\}$ and
$R=A^{-5}F\delta^{-6}$, the coefficient multiplying the square of
the field test is
\[
 \frac{R}{\delta^{2m+2}R^{2m}}
          =A^{10m-5}\delta^{10m-8}F^{1-2m}.
\]
For $W=w_0+w_1\sqrt2$, the guarded reciprocal trace is
$2w_0/(w_0^2-2w_1^2)$, so its pole input is exactly
\[
 \frac{R}{(\operatorname{Tr}_{K/\mathbb Q}(W^{-1}))^2}
       =\frac{F(w_0^2-2w_1^2)^2}{4A^5\delta^6w_0^2}.
\]
The checker reads these expressions from the producer and compares
their expanded dictionaries, rather than relying only on sample
numerical evaluations.

All four frozen source files reproduce byte for byte.  Their integer
coefficient DAGs are acyclic, every monic relation depends only on
retained coordinates and earlier generators, and all counts and finite
parameter lists agree with the displayed source.  There are exactly
sixteen comparison tests and no remaining Lucas proxy variable,
equation or exceptional-root case.  The verified counts are
\[
\begin{array}{c|c|c|c|c}
\text{anchor}&C&\text{retained coordinates}&
   \text{finite relation degrees}&\text{coefficient DAG nodes}\\\hline
\text{finite}&2&40&4^9\,2^{10}&16039\\
\text{semilocal}&2&42&4^9\,2^{11}&16085\\
\text{finite}&2402&40&4^9\,2^{10}&16040\\
\text{semilocal}&2402&42&4^9\,2^{11}&16086
\end{array}
\]
Here $C$ means $\beta=X^2+C$.  The products in the fourth column
describe the nine quartic and ten or eleven quadratic finite relations;
they are not degrees of the final polynomial.  The independent
finite-algebra compiler supplies the final degree and height bounds.

The source check is reproducible with
\[
 \nolinkurl{python research/phase3/compression/check_direct_source.py}.
\]
The separate program \texttt{check\_elliptic\_density\_data.py}
verifies the exact curve and point calculations used in
Corollary~\ref{cor:point-modules-dimension}.  Its scope is those
fixed algebraic calculations; the geometric density assertion is
proved in the text.


\begin{thebibliography}{99}

\bibitem{Repository2026}
\emph{A polynomial definition of the integers in the rationals},
repository \texttt{elliotglazer/diophantine-z-q}, revision
\nolinkurl{28e87b1b62af0effd954605af7f70b2699dd5672}, accessed
6 November 2026. Manuscript \texttt{quartic\_12504.tex}; package
\texttt{quartics/quartic-12504}.
\url{https://github.com/elliotglazer/diophantine-z-q}.
Working fork: \url{https://github.com/JoseBrox/diophantine-z-q_fork}.

\bibitem{Matiyasevich1970}
Yu. V. Matiyasevich, \emph{The Diophantineness of enumerable sets},
Doklady Akademii Nauk SSSR \textbf{191} (1970), no. 2, 279--282;
English translation, Soviet Mathematics Doklady \textbf{11} (1970),
354--357. \url{https://www.mathnet.ru/eng/dan35274}.

\bibitem{MatiyasevichRobinson1975}
Yu. Matiyasevich and J. Robinson, \emph{Reduction of an arbitrary
Diophantine equation to one in 13 unknowns}, Acta Arithmetica
\textbf{27} (1975), 521--553.
\url{https://doi.org/10.4064/aa-27-1-521-553}.

\bibitem{Jones1982}
J. P. Jones, \emph{Universal Diophantine equation}, Journal of Symbolic
Logic \textbf{47} (1982), no. 3, 549--571.
\url{https://doi.org/10.2307/2273588}.

\bibitem{Daans2024}
N. Daans, \emph{Universally defining $\Z$ in $\Q$ with 10 quantifiers},
Journal of the London Mathematical Society \textbf{109} (2024), no. 2,
e12864. \url{https://doi.org/10.1112/jlms.12864}.
Primary preprint: \url{https://arxiv.org/abs/2301.02107}.

\bibitem{DDF2025}
N. Daans, P. Dittmann, and A. Fehm, \emph{Existential rank and
essential dimension of diophantine sets}, arXiv:2102.06941v5 (2025),
in particular Theorem 1.4, Corollary 4.21, and Remark 8.23.
\url{https://arxiv.org/abs/2102.06941v5}.

\bibitem{ZhangSun2022}
G.-R. Zhang and Z.-W. Sun, \emph{$\Q\setminus\Z$ is Diophantine
over $\Q$ with 32 unknowns}, Bulletin of the Polish Academy of Sciences,
Mathematics \textbf{70} (2022), no. 2, 93--106.
\url{https://arxiv.org/abs/2104.02520}.

\bibitem{Sun2026Rational86}
Z.-W. Sun, \emph{86 unknowns for Hilbert's tenth problem over the
rationals}, arXiv:2610.20691v2 (2026).
\url{https://arxiv.org/html/2610.20691v2}.
The new shared-point construction additionally assumes Theorem 4.1,
not merely the algebraic devices of Sections 10.3--10.6. Its use is
explicitly separated from source-preserving algebraic transformations.

\bibitem{FreiTinkova2025}
C. Frei and M. Tinkov\'a, with an appendix by C. Frei, J. K\"onig,
and M. Tinkov\'a, \emph{Linear constellations in primes with arithmetic
restrictions}, arXiv:2408.17441v3 (2025), Theorem 1.4.
\url{https://arxiv.org/html/2408.17441v3}.
The Artin-symbol theorem cited here is unconditional; the separate
primitive-root theorem has a Riemann-hypothesis assumption.

\bibitem{TaoZiegler2018}
T. Tao and T. Ziegler, \emph{Polynomial patterns in the primes},
Forum of Mathematics, Pi \textbf{6} (2018), e1.
\url{https://arxiv.org/abs/1603.07817}.

\bibitem{PatersonStockmeyer1973}
M. S. Paterson and L. J. Stockmeyer, \emph{On the number of nonscalar
multiplications necessary to evaluate polynomials}, SIAM Journal on
Computing \textbf{2} (1973), no. 1, 60--66.
\url{https://doi.org/10.1137/0202007}.

\bibitem{Sun2021}
Z.-W. Sun, \emph{Further results on Hilbert's Tenth Problem},
Science China Mathematics \textbf{64} (2021), no. 2, 281--306;
in particular Theorem 1.1(i) and Equation (1.1).
\url{https://arxiv.org/abs/1704.03504v7}.

\bibitem{Koenigsmann2016}
J. Koenigsmann, \emph{Defining $\Z$ in $\Q$}, Annals of Mathematics
\textbf{183} (2016), no. 1, 73--93.
\url{https://doi.org/10.4007/annals.2016.183.1.2}.
Primary preprint: \url{https://arxiv.org/abs/1011.3424}.

\bibitem{Voight2021}
J. Voight, \emph{Quaternion algebras}, Graduate Texts in Mathematics
\textbf{288}, Springer, 2021; in particular Proposition 14.6.7.
\url{https://doi.org/10.1007/978-3-030-56694-4}.

\bibitem{lbKollar2002}
J. Koll\'ar, \emph{Unirationality of cubic hypersurfaces}, Journal
of the Institute of Mathematics of Jussieu \textbf{1} (2002), no. 3,
467--476, Theorem 2.
\url{https://doi.org/10.1017/S1474748002000117}.
\url{https://arxiv.org/abs/math/0005146}.

\bibitem{lbSchinzel2000}
A. Schinzel, \emph{Polynomials with special regard to reducibility},
Encyclopedia of Mathematics and its Applications \textbf{77},
Cambridge University Press, 2000, Chapter 4, Section 4.
\url{https://doi.org/10.1017/CBO9780511542916}.

\bibitem{lbMuller2001}
P. M\"uller, \emph{Finiteness results for Hilbert's irreducibility
theorem}, arXiv:math/0109071, 2001; introduction and Section 2.1.
\url{https://arxiv.org/abs/math/0109071}.

\bibitem{lbLang1983}
S. Lang, \emph{Fundamentals of Diophantine geometry}, Springer, 1983,
Chapter 8, pp. 188--224.
\url{https://doi.org/10.1007/978-1-4757-1810-2}.

\bibitem{p2lbEGHP2006}
D. Eisenbud, M. Green, K. Hulek, and S. Popescu,
\emph{Small schemes and varieties of minimal degree}, American Journal
of Mathematics \textbf{128} (2006), no. 6, 1363--1389.
\url{https://doi.org/10.1353/ajm.2006.0043}.
\url{https://eisenbud.github.io/papers/pdfs/2006-002.pdf}.

\bibitem{p2lbHarrisTschinkel2000}
J. Harris and Y. Tschinkel, \emph{Rational points on quartics},
Duke Mathematical Journal \textbf{104} (2000), no. 3, 477--500.
\url{https://arxiv.org/abs/math/9809015}.
\url{https://cims.nyu.edu/~tschinke/papers/yuri/00quart/quart.pdf}.

\bibitem{p2lbBergEtAl2024}
J. Berg, C. Pagano, B. Poonen, M. Stoll, N. Triantafillou, B. Viray,
and I. Vogt, \emph{Brauer--Manin obstructions requiring arbitrarily many
Brauer classes}, Bulletin of the London Mathematical Society
\textbf{56} (2024), no. 5, 1587--1604; Theorem 5.1 and Section 6.1.
\url{https://doi.org/10.1112/blms.12999}.
\url{https://math.mit.edu/~poonen/papers/many_brauer.pdf}.

\bibitem{p2lbSerre1981}
J.-P. Serre, \emph{Quelques applications du th\'eor\`eme de densit\'e
de Chebotarev}, Publications Math\'ematiques de l'IH\'ES \textbf{54}
(1981), 123--201; Section 2.1, Theorem 1.
\url{https://numdam.org/item/PMIHES_1981__54__123_0.pdf}.

\bibitem{p2lbVoight2019}
J. Voight, \emph{Identifying the matrix ring: algorithms for quaternion
algebras and quadratic forms}, author's revised text, 21 May 2019,
Proposition 5.5; arXiv:1004.0994.
\url{https://jvoight.github.io/articles/quatalgs-051919.pdf}.

\bibitem{KoymansPagano2026}
P. Koymans and C. Pagano, \emph{Hilbert's tenth problem via additive
combinatorics}, Journal of the American Mathematical Society,
published electronically 10 August 2026.
\url{https://doi.org/10.1090/jams/1081}.
\url{https://arxiv.org/html/2412.01768v3}.

\bibitem{AlpogeBhargavaHoShnidman2026}
L. Alp\"oge, M. Bhargava, W. Ho, and A. Shnidman,
\emph{Rank stability in quadratic extensions and Hilbert's tenth
problem for the ring of integers of a number field}, Inventiones
Mathematicae \textbf{243} (2026), 1129--1139.
\url{https://doi.org/10.1007/s00222-025-01392-3}.
The cited Corollary 1.2 is that of the published version; the arXiv
version 1 states a weaker Diophantine-model conclusion.

\bibitem{Mazur1995}
B. Mazur, \emph{Speculations about the topology of rational points:
an up-date}, Ast\'erisque \textbf{228} (1995), 165--181;
Conjecture 3, p. 165.
\url{https://numdam.org/item/AST_1995__228__165_0/}.

\bibitem{Pan2022}
L. Pan, \emph{The Fontaine--Mazur conjecture in the residually
reducible case}, Journal of the American Mathematical Society
\textbf{35} (2022), 1031--1169.
\url{https://doi.org/10.1090/jams/991}.
\url{https://arxiv.org/abs/1901.07166}.

\bibitem{Yuan2024}
X. Yuan, \emph{Explicit Kodaira--Spencer map over Shimura curves},
Acta Mathematica Sinica, Chinese Series \textbf{67} (2024), no. 2,
227--249. \url{https://doi.org/10.12386/A20220154}.
\url{https://arxiv.org/abs/2205.11334}.

% New bibliography item only. The existing keys lbKollar2002,
% lbLang1983 and Voight2021 are reused by the phase-three supplement.
\bibitem{p3sysMerkurjevParimala2025}
A.~Merkurjev and R.~Parimala,
\emph{Quadratic Forms Beyond Arithmetic},
Notices of the American Mathematical Society \textbf{72} (2025), no.~7,
711--718, \href{https://doi.org/10.1090/noti3192}{doi:10.1090/noti3192}.
Author manuscript:
\url{https://www.math.ucla.edu/~merkurev/papers/quadratic8.pdf}, Section~5.1.


\bibitem{P3Poonen2010}
B.~Poonen, \emph{Multivariable polynomial injections on rational numbers},
Acta Arithmetica \textbf{145} (2010), no.~2, 123--127.
\href{https://doi.org/10.4064/aa145-2-2}{doi:10.4064/aa145-2-2}.
\url{https://math.mit.edu/~poonen/papers/QQ.pdf}.

\bibitem{P3Pasten2020}
H.~Pasten, \emph{Bivariate polynomial injections and elliptic curves},
Selecta Mathematica, New Series \textbf{26} (2020), article 22.
\href{https://doi.org/10.1007/s00029-020-0548-x}{doi:10.1007/s00029-020-0548-x}.
\url{https://arxiv.org/abs/1811.02416}.

\bibitem{P3Hirvensalo2026}
M.~Hirvensalo, \emph{On injective multivariate polynomials over rational
numbers}, Theoretical Computer Science \textbf{1082} (2026), 116158.
\href{https://doi.org/10.1016/j.tcs.2026.116158}{doi:10.1016/j.tcs.2026.116158}.
Theorem~2 has the restricted domain specified in the text.

\bibitem{P3MilneAbelian}
J.~S.~Milne, \emph{Abelian Varieties}, in G.~Cornell and J.~H.~Silverman
(eds.), Arithmetic Geometry, Springer, New York, 1986, 103--150;
corrected author version, 2 January 2022, Sections 8 and 12.
\url{https://www.jmilne.org/math/xnotes/AVs.pdf}.


% Entries for inclusion inside an existing thebibliography environment.
\bibitem{MilneCFT2020}
J.~S.~Milne, \emph{Class Field Theory}, version 4.03,
6 August 2020, Chapter IV, Sections 1--4.
\url{https://www.jmilne.org/math/CourseNotes/CFT.pdf}.

\bibitem{KedlayaCFT}
K.~S.~Kedlaya, \emph{Notes on Class Field Theory},
Section 2.4 (Chebotarev's density theorem) and Section 7.6
(Brauer groups, especially Theorem 7.6.12), accessed
6 November 2026.
\url{https://kskedlaya.org/cft/}.

\bibitem{StacksChevalley}
The Stacks Project Authors, \emph{The Stacks Project},
Tag 00FE, Theorem 10.29.10 (Chevalley's theorem), 2026.
\url{https://stacks.math.columbia.edu/tag/00FE}.

\bibitem{Basu2014}
S.~Basu, \emph{Algorithms in Real Algebraic Geometry:
A Survey}, arXiv:1409.1534, 2014, Section 1.1 and Theorem 2.1.
\url{https://arxiv.org/pdf/1409.1534}.

\bibitem{Plaumann2023}
D.~Plaumann, \emph{Real Algebraic Geometry},
IHP tutorial lectures, 12 October 2023,
Corollary 4.5 and Example 4.6.
\url{https://wwwold.mathematik.tu-dortmund.de/~dplauman/notes/rag-ihp.pdf}.


\end{thebibliography}
\end{document}
